\PassOptionsToPackage{usenames,dvipsnames,svgnames,table}{xcolor}
%\documentclass[a4paper,usenglish]{lipics}
\documentclass[a4paper]{lipics}
\usepackage{hyperref}
\usepackage{multicol}
\usepackage{tikz,tikz-cd}
\usepackage{quiver}
\usepackage[inline]{enumitem}
\usepackage[round]{natbib}
\bibliographystyle{plainnat}

\let\shortcite=\cite
\renewcommand{\slash}{\mathrel/}
\newcommand{\tri}{\filledmedtriangleright}
\renewcommand{\paragraph}[1]{\medskip\par\noindent\textbf{#1.} }
\renewcommand{\slash}{\mathrel/}
\newcommand{\val}[1]{\llbracket #1\rrbracket}

\usetikzlibrary{cd}

\sloppy

\usepackage{pifont}

\usepackage[dvipsnames]{xcolor}

\title{Modular construction of free algebras for distributive combinations of theories}

% Please add yourself if you make a change
\author[1]{Ohad Kammar}
\author[2]{Amaury Mazoyer}

\affil[1]{
  Laboratory for Foundations of
  Computer Science                  \\
  School of Informatics             \\
  University of Edinburgh           \\
  Scotland                          \\
  \url{ohad.kammar@ed.ac.uk}
}

\affil[2]{
  Département d'Informatique      \\
  ENS de Lyon                     \\
  France                          \\
  \url{amaury.mazoyer@ens-lyon.fr}
}

\authorrunning{O.~Kammar}
\Copyright{O.~Kammar}
\subjclass{F.3.1 Specifying and Verifying and Reasoning about Programs}
\keywords{}
%\serieslogo{lipics-logo-bw}
\volumeinfo{}{1}{}{1}{1}{1}
\EventShortName{Proof}
% \EventShortName{Working note}
%\DOI{10.4230/LIPIcs.xxx.yyy.p}
\def\copyrightline{}

\let\nsubset\undefined
\usepackage{amsthm,proof,MnSymbol}
\usepackage{oktheorem,okgeneralmath,okcategory}
\usepackage{bbold}
\usepackage{epiolmec}
\usepackage{oksemantics}
\usepackage{mathpartir}
\usepackage{todonotes}
\usepackage{natbib}
\usepackage{fp}
\usepackage{mathdots}
%next two for no-depth underlining
\usepackage{contour}
\usepackage[normalem]{ulem}
\newcommand{\OK}[1]{\todo[inline,author=OK]{#1}}
\newcommand{\ok}[1]{\todo[size=\tiny]{OK: #1}{}}


%%%%%%%%%%% Cross references %%%%%%%%%%
\newcommand\seclabel[1]{\label{sec:#1}}
\newcommand\secref[1]{Sec.~\ref{sec:#1}}
\newcommand\applabel[1]{\label{app:#1}}
\newcommand\appref[1]{Sec.~\ref{app:#1}}
\newcommand\subseclabel[1]{\label{subsec:#1}}
\newcommand\subsecref[1]{\S~\ref{subsec:#1}}
\newcommand\subsubseclabel[1]{\label{subsubsec:#1}}
\newcommand\subsubsecref[1]{\S~\ref{subsubsec:#1}}
\newcommand\figlabel[1]{\label{fig:#1}}
\newcommand\figref[1]{Fig.~\ref{fig:#1}}

\newcommand\figfile[2]{%
    \begin{figure}%
      \begin{center}%
        \input{#1}%
        \caption{#2}%
        \figlabel{#1}%
      \end{center}%
    \end{figure}%
}
%%%%%%%%%%%%%%%%%%%%%%%%%%%%%%%%%%%%%%


%%%%%%%%%%% Add possessive citations, e.g. "Leroy's (1992)", taken
\newcommand{\citepos}[1]{\citeauthor{#1}'s~\citeyearpar{#1}}
\WithSuffix\newcommand\citepos*[2]{\citeauthor{#1}'s #2~\citeyearpar{#1}}
\WithSuffix\newcommand\citepos-[2]{\citeauthor{#1}#2~\citeyearpar{#1}}
%%%%%%%%%%%%%%%% End of possessive citation %%%%%%%%%%%%%%%%%%%%%
\makeatletter
\newcommand\ibidt[2][]{%
  \citeauthor{#2}~\ibid[#1]{#2}
}
\newcommand\ibid[2][]{%
  \citetext{%
    \hyper@natlinkstart{#2}%
    \textit{ibid.}%
    \hyper@natlinkend%
    \ifx\ibid@empty#1\ibid@empty%
    \else%
    , #1%
    \fi%
  }%
}

\newcommand\citloct[1]{%
  \citeauthor{#1}~\citloc{#1}%
}
\newcommand\citloc[1]{%
  \citetext{%
    \hyper@natlinkstart{#1}%
    \textit{cit.~loc.}%
    \hyper@natlinkend%
  }%
}
\newcommand\citlocpos[1]{%
  \citeauthor{#1}'s~\citloc{#1}%
}
\WithSuffix\newcommand\citlocpos*[2]{%
  \citeauthor{#1}'s #2~\citloc{#1}%
}
\makeatother


\newcommand\diagram[1]{\begin{center}
    \diagram*{#1}
\end{center}}
\WithSuffix\newcommand\diagram*[1]{\begin{tabular}[c]{@{}c@{}}\includegraphics{diagram-#1.mps}\end{tabular}}


%%%%%%%%%%%%%%%%%%%%%%%%%%%%%%%%%%%%%%
\usepackage{tcolorbox}
\newtcbox{\myfigbox}{%
  on line, arc=1pt, outer arc=1pt,colframe=black,
  colback=white,
  boxsep=0pt,
  left=1pt,right=1pt,top=1pt,bottom=1pt,%
  boxrule=.5pt,bottomrule=.5pt,toprule=.5pt}

%%%%%%%%%%%%%%%%%%%%%%%%%%%%%%%%%%%%%%
\makeatletter
\newcommand\newmvar[3][]{%
  \newcommand#2[1][1]{#1{\ifcase ##1\or #3\or\else@ctrerr\fi}}%
}

%% introduce a variation of a defined meta-variable
\newcommand\mvarvar[3]{%
  \WithSuffix\newcommand#1#2[1][0]{%
    \FPeval{\tempindex}{clip(#3+##1)}%
    #1[\tempindex]%
  }
}

\makeatother

\let\C\undefined
\let\a\undefined
\newmvar{\M}{M\or N}
\newmvar{\V}{V\or W \or U}
\newmvar{\a}{a\or b\or c}
\newmvar{\x}{x\or y\or z}
\newmvar{\X}{X\or Y\or Z}
\newmvar{\Aname}{A\or B\or C}
\newmvar{\Cname}{C\or D\or A\or B\or C}
\newmvar{\Fname}{F\or G\or H\or K\or L}
\newmvar{\Uname}{U\or V\or W\or X\or Y}
\let\v\undefined
\newmvar{\v}{\hat{v}\or \hat{w}}
\let\l\undefined
\newmvar{\l}{\hat{\ell}\or \hat{k}}

\newcommand\C[1][1]{\mathcal{\Cname[#1]}}
\newcommand\A[1][1]{\mathcal{\Aname[#1]}}
%%%%%%%%%%%%%%%%%%%%%%%%%%%%%%%%%%%
%%%%%%%%%%% monads %%%%%%%%%%%%%
\newcommand\newop[2]{\newcommand{#1}{\mathop{\mathrm{#2}}\nolimits}}
\newcommand{\bindsymbol}{\mathrel{\small\rrangle}\joinrel\mathrel{\small=}}
\newcommand{\klseq}{\ggg}
\newcommand{\bind}{\mathrel{\bindsymbol}}
\newop{\return}{return}
\newop{\strength}{strength}
\newop{\shadow}{shadow}

\newcommand\spair[2]{\langle #1, #2\rangle}
\newcommand\sunit{\langle\rangle}
\makeatletter
\let\llangle\@undefined
\let\rrangle\@undefined
\DeclareMathDelimiter{\llangle}{\mathopen}%
                     {okMnSymbolLargeSymbols}{'164}{okMnSymbolLargeSymbols}{'164}
\DeclareMathDelimiter{\rrangle}{\mathclose}%
                     {okMnSymbolLargeSymbols}{'171}{okMnSymbolLargeSymbols}{'171}
\makeatother
\newcommand\ppair[2]{\left\llangle#1,#2\right\rrangle}
\newcommand\ptimes{\otimes}
\def\pexplen{0mu}
\def\pexpmidlen{12mu}
\newcommand\pexp{
  \mathrel{%
    \overrightharpoon{%
      {\mspace{-\pexplen}\smash{\mathord{\relbar\mspace{-\pexpmidlen}\relbar}}\vphantom{.}}\mspace{-\pexplen}}%
  }%
}
\usepackage{mathtools}
\usepackage{tikz-cd}
\usetikzlibrary{decorations.pathmorphing}

% code from Gonzalo
\newcommand\xrsquigarrow[1]{%
    \mathrel{%
        \begin{tikzpicture}[%
            baseline={(current bounding box.south)}
            ]
        \node[%
            ,inner sep=.44ex
            ,align=center
            ] (tmp) {$\scriptstyle #1$};
        \path[%
            ,draw,<-
            ,decorate,decoration={%
            ,snake,aspect=0
            ,amplitude=1pt
            ,segment length=1.5mm,pre length=3.5pt
                }
            ]
        (tmp.south east) -- (tmp.south west);
        \end{tikzpicture}
        }
    }
\newcommand\xlsquigarrow[1]{%
    \mathrel{%
        \begin{tikzpicture}[%
            ,baseline={(current bounding box.south)}
            ]
        \node[%
            ,inner sep=.44ex
            ,align=center
            ] (tmp) {$\scriptstyle #1$};
        \path[%
            ,draw,<-
            ,decorate,decoration={%
                ,zigzag
                ,amplitude=0.7pt
                ,segment length=1.2mm,pre length=3.5pt
                }
            ]
        (tmp.south west) -- (tmp.south east);
        \end{tikzpicture}
        }
    }

%%%%%%%%%%%%%%%%%%%%%%%%%%%%%%%%%%%%%%%%%%
\usepackage{xcolor}
\usepackage{listings}
\colorlet{code}{DarkBlue}
\colorlet{types}{Magenta}
%%%%%%%%%%%%%%%%% Syntax %%%%%%%%%%%%%%%%%
\newcommand\gdefinedby{\mathrel{::=}}
\newcommand\gor{\lvert}

\newcommand\opname[1][op]{\mathrm{#1}}
\newcommand\op[1][op]{\mathrm{#1}\,}
\newcommand\typeof[1][\,]{\mathrm{typeof}#1}
\newcommand\arity[1][\,]{\mathrm{arity}#1}
\newcommand\OpType[2]{#2\seq{#1}}

\newcommand\Times{\mathbin{\boldsymbol{*}}}
\newcommand\CTimes{\mathbin{\boldsymbol{\times}}}

\usepackage{scalerel,stackengine}
\stackMath
\newcommand\reallywidehat[1]{%
\savestack{\tmpbox}{\stretchto{%
  \scaleto{%
    \scalerel*[\widthof{\ensuremath{#1}}]{\kern-.6pt\bigwedge\kern-.6pt}%
    {\rule[-\textheight/2]{1ex}{\textheight}}%WIDTH-LIMITED BIG WEDGE
  }{\textheight}%
}{0.5ex}}%
\stackon[1pt]{#1}{\tmpbox}%
}

\newcommand\sem[1]{\scottbrackets{#1}}

\newenvironment{catdef}{%
\begingroup%
\newcommand\objects[1]{\item[Objects $##1$:]}%
\newcommand\morphisms[1]{\item[{Morphisms $##1$:}]}%
\begin{description}}{\end{description}\endgroup}

\newcommand\eval{\mathop{\mathrm{eval}}\nolimits}
\newcommand\curry{\mathop{\mathrm{curry}}\nolimits}

\newcommand\cotuple\coseq
\newcommand\deneq{\mathrel{%
  \raisebox{.5ex}{%
  \scalebox{0.4}{\EOx}%
  }
  }}
\newcommand\vto[3]{%
  \begin{smallmatrix}%
        \hphantom{#2}#1     \\%
        #2\mathord\downarrow\\%
        \hphantom{#2}#3       %
  \end{smallmatrix}%
}
\newcommand\vfrom[3]{%
  \begin{smallmatrix}%
        \hphantom{#2}#1     \\%
        #2\mathord\uparrow\\%
        \hphantom{#2}#3       %
  \end{smallmatrix}%
}

\makeatletter
\newcommand\NewFontImport[6]{
\newcommand#1{\begingroup\mathchoice%
  {\hbox{\fontsize{\tf@size}{\tf@size}\usefont{#2}{#3}{#4}{#5}\char#6}}%
  {\hbox{\fontsize{\tf@size}{\tf@size}\usefont{#2}{#3}{#4}{#5}\char#6}}%
  {\hbox{\fontsize{\sf@size}{\sf@size}\usefont{#2}{#3}{#4}{#5}\char#6}}%
  {\hbox{\fontsize{\ssf@size}{\ssf@size}\usefont{#2}{#3}{#4}{#5}\char#6}}%
  \endgroup%
}}
\makeatother

\DeclareFontFamily{U}{bbold}{}
\DeclareFontShape{U}{bbold}{m}{n}{<->bbold10}{}
\DeclareFontSubstitution{U}{bbold}{m}{n}
\NewFontImport\SinSym{U}{bbold}{m}{n}{'017}
\newcommand\cat[1]{#1_{\text{cat}}}
\newcommand\involve[2][]{\overline{#2}_{#1}}
\WithSuffix\newcommand\involve*[1]{\overline{#1}}
\newcommand\iinvolve[1]{\involve*{\involve{#1}}}
\newcommand\iiinvolve[1]{\involve*{\iinvolve{#1}}}
\newcommand\invlaw[1][]{\iota^{#1}}
\newcommand\Monoid{\Category{Monoid}}
\newcommand\Model[2]{\Category{Model}(#1,#2)}
\newcommand\Rig{\Category{Rig}}
\newcommand\fun{\mathop{\uplambda}}
\newcommand\carrier[1]{\underline{#1}}
\newcommand\pres{\mathbb P}
\newcommand\unit{\eta}
\newcommand\counit{\epsilon}
\newcommand\SC{\mathop{\mathrm{SC}}}
\newcommand\aconj[1][]{j_{#1}}
\newcommand\obj[1]{#1_{\text{obj}}}
\newcommand\ICAT{\Category{ICAT}}
\newcommand\Frexlib{\textsc{Frex}}
\newcommand\icoseq[1]{\overline[#1\overline]}

\newcommand\To\Rightarrow
\newcommand\vcompose\circ
\newcommand\hcompose{\mathbin\bullet}
\newcommand\newquiver[2]{%
  \expandafter\newcommand\csname quiver:#1\endcsname{%
    \protect{#2}
  }%
}
\newcommand\quiver[1]{\expandafter\csname quiver:#1\endcsname}
\newquiver{alpha f g}{
% https://q.uiver.app/?q=WzAsMixbMCwwLCJhIl0sWzIsMCwiYiJdLFswLDEsImYiLDAseyJjdXJ2ZSI6LTJ9XSxbMCwxLCJnIiwyLHsiY3VydmUiOjJ9XSxbMiwzLCJcXGFscGhhIiwwLHsic2hvcnRlbiI6eyJzb3VyY2UiOjIwLCJ0YXJnZXQiOjIwfX1dXQ==
\begin{tikzcd}[ampersand replacement=\&]%
	a \&\& b
 	\arrow[""{name=0, anchor=center, inner sep=0}, "f", curve={height=-12pt}, from=1-1, to=1-3]
 	\arrow[""{name=1, anchor=center, inner sep=0}, "g"', curve={height=12pt}, from=1-1, to=1-3]
 	\arrow["\alpha", shorten <=3pt, shorten >=3pt, Rightarrow, from=0, to=1]
 \end{tikzcd}
}

\newquiver{identity 2 cell}{
  % https://q.uiver.app/?q=WzAsMyxbMCwwLCJhIl0sWzIsMCwiYiJdLFsxLDAsIj0iXSxbMCwxLCJmIiwwLHsiY3VydmUiOi0yfV0sWzAsMSwiZiIsMix7ImN1cnZlIjoyfV1d
\begin{tikzcd}[ampersand replacement=\&]
	a \& {=} \& b
	\arrow["f", curve={height=-12pt}, from=1-1, to=1-3]
	\arrow["f"', curve={height=12pt}, from=1-1, to=1-3]
\end{tikzcd}
}

\newquiver{vertical composition type}{
% https://q.uiver.app/?q=WzAsMixbMCwwLCJhIl0sWzIsMCwiYiJdLFswLDEsImYiLDAseyJjdXJ2ZSI6LTR9XSxbMCwxLCJnIiwyLHsibGFiZWxfcG9zaXRpb24iOjIwfV0sWzAsMSwiaCIsMix7ImN1cnZlIjo0fV0sWzIsMywiXFxhbHBoYSIsMCx7InNob3J0ZW4iOnsic291cmNlIjoyMCwidGFyZ2V0IjoyMH19XSxbMyw0LCJcXGJldGEiLDAseyJzaG9ydGVuIjp7InNvdXJjZSI6MjAsInRhcmdldCI6MjB9fV1d
\begin{tikzcd}[ampersand replacement=\&]
	a \&\& b
	\arrow[""{name=0, anchor=center, inner sep=0}, "f", curve={height=-24pt}, from=1-1, to=1-3]
	\arrow[""{name=1, anchor=center, inner sep=0}, "g"'{pos=0.2}, from=1-1, to=1-3]
	\arrow[""{name=2, anchor=center, inner sep=0}, "h"', curve={height=24pt}, from=1-1, to=1-3]
	\arrow["\alpha", shorten <=3pt, shorten >=3pt, Rightarrow, from=0, to=1]
	\arrow["\beta", shorten <=3pt, shorten >=3pt, Rightarrow, from=1, to=2]
\end{tikzcd}
}

\newquiver{vertical composition}{
  % https://q.uiver.app/?q=WzAsMixbMCwwLCJhIl0sWzIsMCwiYiJdLFswLDEsImYiLDAseyJjdXJ2ZSI6LTJ9XSxbMCwxLCJoIiwyLHsiY3VydmUiOjJ9XSxbMiwzLCJcXGJldGEgXFxjaXJjIFxcYWxwaGEiLDAseyJzaG9ydGVuIjp7InNvdXJjZSI6MjAsInRhcmdldCI6MjB9fV1d
\begin{tikzcd}[ampersand replacement=\&]
	a \&\& b
	\arrow[""{name=0, anchor=center, inner sep=0}, "f", curve={height=-12pt}, from=1-1, to=1-3]
	\arrow[""{name=1, anchor=center, inner sep=0}, "h"', curve={height=12pt}, from=1-1, to=1-3]
	\arrow["{\beta \circ \alpha}", shorten <=3pt, shorten >=3pt, Rightarrow, from=0, to=1]
\end{tikzcd}
}

\newquiver{horizontal composition type}{
% https://q.uiver.app/?q=WzAsMyxbMCwwLCJhIl0sWzIsMCwiYiJdLFs0LDAsImMiXSxbMCwxLCJmXzEiLDAseyJjdXJ2ZSI6LTJ9XSxbMCwxLCJnXzEiLDIseyJjdXJ2ZSI6Mn1dLFsxLDIsImZfMiIsMCx7ImN1cnZlIjotMn1dLFsxLDIsImdfMiIsMix7ImN1cnZlIjoyfV0sWzMsNCwiXFxhbHBoYSIsMCx7InNob3J0ZW4iOnsic291cmNlIjoyMCwidGFyZ2V0IjoyMH19XSxbNSw2LCJcXGJldGEiLDAseyJzaG9ydGVuIjp7InNvdXJjZSI6MjAsInRhcmdldCI6MjB9fV1d
\begin{tikzcd}[ampersand replacement=\&]
	a \&\& b \&\& c
	\arrow[""{name=0, anchor=center, inner sep=0}, "{f_1}", curve={height=-12pt}, from=1-1, to=1-3]
	\arrow[""{name=1, anchor=center, inner sep=0}, "{g_1}"', curve={height=12pt}, from=1-1, to=1-3]
	\arrow[""{name=2, anchor=center, inner sep=0}, "{f_2}", curve={height=-12pt}, from=1-3, to=1-5]
	\arrow[""{name=3, anchor=center, inner sep=0}, "{g_2}"', curve={height=12pt}, from=1-3, to=1-5]
	\arrow["\alpha", shorten <=3pt, shorten >=3pt, Rightarrow, from=0, to=1]
	\arrow["\beta", shorten <=3pt, shorten >=3pt, Rightarrow, from=2, to=3]
\end{tikzcd}
}

\newquiver{horizontal composition}{
% https://q.uiver.app/?q=WzAsMixbMCwwLCJhIl0sWzIsMCwiYyJdLFswLDEsImZfMlxcY2lyYyBmXzEiLDAseyJjdXJ2ZSI6LTJ9XSxbMCwxLCJnXzIgXFxjaXJjIGdfMSIsMix7ImN1cnZlIjoyfV0sWzIsMywiXFxiZXRhIFxcaGNvbXBvc2UgXFxhbHBoYSIsMCx7InNob3J0ZW4iOnsic291cmNlIjoyMCwidGFyZ2V0IjoyMH19XV0=
\begin{tikzcd}[ampersand replacement=\&,row sep=small]
	a \&\& c
	\arrow[""{name=0, anchor=center, inner sep=0}, "{f_2\circ f_1}", curve={height=-12pt}, from=1-1, to=1-3]
	\arrow[""{name=1, anchor=center, inner sep=0}, "{g_2 \circ g_1}"', curve={height=12pt}, from=1-1, to=1-3]
	\arrow["{\beta \hcompose \alpha}", shorten <=3pt, shorten >=3pt, Rightarrow, from=0, to=1]
\end{tikzcd}
}

\newquiver{interchange law vertical}{
% https://q.uiver.app/?q=WzAsNixbMCwwLCJhIl0sWzEsMCwiYiJdLFswLDEsImEiXSxbMSwxLCJiIl0sWzIsMCwiYyJdLFsyLDEsImMiXSxbMSwzLCJcXGNpcmMiLDEseyJzdHlsZSI6eyJib2R5Ijp7Im5hbWUiOiJub25lIn0sImhlYWQiOnsibmFtZSI6Im5vbmUifX19XSxbMCwxXSxbMSw0XSxbMiwzXSxbMyw1XSxbMCwxLCIiLDEseyJjdXJ2ZSI6LTN9XSxbMSw0LCIiLDEseyJjdXJ2ZSI6LTN9XSxbMiwzLCIiLDEseyJjdXJ2ZSI6M31dLFszLDUsIiIsMSx7ImN1cnZlIjozfV0sWzExLDcsIiIsMSx7InNob3J0ZW4iOnsic291cmNlIjoyMCwidGFyZ2V0IjoyMH19XSxbMTIsOCwiIiwxLHsic2hvcnRlbiI6eyJzb3VyY2UiOjIwLCJ0YXJnZXQiOjIwfX1dLFsxMCwxNCwiIiwxLHsic2hvcnRlbiI6eyJzb3VyY2UiOjIwLCJ0YXJnZXQiOjIwfX1dLFs5LDEzLCIiLDEseyJzaG9ydGVuIjp7InNvdXJjZSI6MjAsInRhcmdldCI6MjB9fV1d
\begin{tikzcd}[ampersand replacement=\&]
	a \& b \& c \\
	a \& b \& c
	\arrow["\circ"{description}, draw=none, from=1-2, to=2-2]
	\arrow[""{name=0, anchor=center, inner sep=0}, from=1-1, to=1-2]
	\arrow[""{name=1, anchor=center, inner sep=0}, from=1-2, to=1-3]
	\arrow[""{name=2, anchor=center, inner sep=0}, from=2-1, to=2-2]
	\arrow[""{name=3, anchor=center, inner sep=0}, from=2-2, to=2-3]
	\arrow[""{name=4, anchor=center, inner sep=0}, curve={height=-18pt}, from=1-1, to=1-2]
	\arrow[""{name=5, anchor=center, inner sep=0}, curve={height=-18pt}, from=1-2, to=1-3]
	\arrow[""{name=6, anchor=center, inner sep=0}, curve={height=18pt}, from=2-1, to=2-2]
	\arrow[""{name=7, anchor=center, inner sep=0}, curve={height=18pt}, from=2-2, to=2-3]
	\arrow[shorten <=2pt, shorten >=2pt, Rightarrow, from=4, to=0]
	\arrow[shorten <=2pt, shorten >=2pt, Rightarrow, from=5, to=1]
	\arrow[shorten <=2pt, shorten >=2pt, Rightarrow, from=3, to=7]
	\arrow[shorten <=2pt, shorten >=2pt, Rightarrow, from=2, to=6]
\end{tikzcd}}

\newquiver{interchange law horizontal}{
  % https://q.uiver.app/?q=WzAsNCxbMCwwLCJhIl0sWzEsMCwiYiJdLFszLDAsImMiXSxbMiwwLCJiIl0sWzAsMV0sWzAsMSwiIiwxLHsiY3VydmUiOi0zfV0sWzMsMiwiIiwxLHsiY3VydmUiOi0zfV0sWzMsMl0sWzAsMSwiIiwxLHsiY3VydmUiOjN9XSxbMywyLCIiLDEseyJjdXJ2ZSI6M31dLFsxLDMsIlxcaGNvbXBvc2UiLDEseyJzdHlsZSI6eyJib2R5Ijp7Im5hbWUiOiJub25lIn0sImhlYWQiOnsibmFtZSI6Im5vbmUifX19XSxbNSw0LCIiLDEseyJzaG9ydGVuIjp7InNvdXJjZSI6MjAsInRhcmdldCI6MjB9fV0sWzYsNywiIiwxLHsic2hvcnRlbiI6eyJzb3VyY2UiOjIwLCJ0YXJnZXQiOjIwfX1dLFs0LDgsIiIsMSx7InNob3J0ZW4iOnsic291cmNlIjoyMCwidGFyZ2V0IjoyMH19XSxbNyw5LCIiLDEseyJzaG9ydGVuIjp7InNvdXJjZSI6MjAsInRhcmdldCI6MjB9fV1d
\begin{tikzcd}[ampersand replacement=\&]
	a \& b \& b \& c
	\arrow[""{name=0, anchor=center, inner sep=0}, from=1-1, to=1-2]
	\arrow[""{name=1, anchor=center, inner sep=0}, curve={height=-18pt}, from=1-1, to=1-2]
	\arrow[""{name=2, anchor=center, inner sep=0}, curve={height=-18pt}, from=1-3, to=1-4]
	\arrow[""{name=3, anchor=center, inner sep=0}, from=1-3, to=1-4]
	\arrow[""{name=4, anchor=center, inner sep=0}, curve={height=18pt}, from=1-1, to=1-2]
	\arrow[""{name=5, anchor=center, inner sep=0}, curve={height=18pt}, from=1-3, to=1-4]
	\arrow["\hcompose"{description}, draw=none, from=1-2, to=1-3]
	\arrow[shorten <=2pt, shorten >=2pt, Rightarrow, from=1, to=0]
	\arrow[shorten <=2pt, shorten >=2pt, Rightarrow, from=2, to=3]
	\arrow[shorten <=2pt, shorten >=2pt, Rightarrow, from=0, to=4]
	\arrow[shorten <=2pt, shorten >=2pt, Rightarrow, from=3, to=5]
\end{tikzcd}
}

\newquiver{interchange law}{
% https://q.uiver.app/?q=WzAsMyxbMCwwLCJhIl0sWzEsMCwiYiJdLFsyLDAsImMiXSxbMCwxXSxbMCwxLCIiLDEseyJjdXJ2ZSI6LTN9XSxbMCwxLCIiLDEseyJjdXJ2ZSI6M31dLFsxLDIsIiIsMSx7ImN1cnZlIjotM31dLFsxLDIsIiIsMSx7ImN1cnZlIjozfV0sWzEsMl0sWzQsMywiIiwxLHsic2hvcnRlbiI6eyJzb3VyY2UiOjIwLCJ0YXJnZXQiOjIwfX1dLFszLDUsIiIsMSx7InNob3J0ZW4iOnsic291cmNlIjoyMCwidGFyZ2V0IjoyMH19XSxbNiw4LCIiLDEseyJzaG9ydGVuIjp7InNvdXJjZSI6MjAsInRhcmdldCI6MjB9fV0sWzgsNywiIiwxLHsic2hvcnRlbiI6eyJzb3VyY2UiOjIwLCJ0YXJnZXQiOjIwfX1dXQ==
\begin{tikzcd}[ampersand replacement=\&]
	a \& b \& c
	\arrow[""{name=0, anchor=center, inner sep=0}, from=1-1, to=1-2]
	\arrow[""{name=1, anchor=center, inner sep=0}, curve={height=-18pt}, from=1-1, to=1-2]
	\arrow[""{name=2, anchor=center, inner sep=0}, curve={height=18pt}, from=1-1, to=1-2]
	\arrow[""{name=3, anchor=center, inner sep=0}, curve={height=-18pt}, from=1-2, to=1-3]
	\arrow[""{name=4, anchor=center, inner sep=0}, curve={height=18pt}, from=1-2, to=1-3]
	\arrow[""{name=5, anchor=center, inner sep=0}, from=1-2, to=1-3]
	\arrow[shorten <=2pt, shorten >=2pt, Rightarrow, from=1, to=0]
	\arrow[shorten <=2pt, shorten >=2pt, Rightarrow, from=0, to=2]
	\arrow[shorten <=2pt, shorten >=2pt, Rightarrow, from=3, to=5]
	\arrow[shorten <=2pt, shorten >=2pt, Rightarrow, from=5, to=4]
\end{tikzcd}
}

\newquiver{cat horizontal composition type}{
  % https://q.uiver.app/?q=WzAsMyxbMCwwLCJcXENbM10iXSxbMSwwLCJcXENbNF0iXSxbMiwwLCJcXENbMV0iXSxbMCwxLCJGXzEiLDEseyJjdXJ2ZSI6LTJ9XSxbMSwyLCJGXzIiLDEseyJjdXJ2ZSI6LTJ9XSxbMCwxLCJHXzEiLDEseyJjdXJ2ZSI6Mn1dLFsxLDIsIkdfMiIsMSx7ImN1cnZlIjoyfV0sWzMsNSwiXFxhbHBoYSIsMCx7InNob3J0ZW4iOnsic291cmNlIjoyMCwidGFyZ2V0IjoyMH19XSxbNCw2LCJcXGJldGEiLDAseyJzaG9ydGVuIjp7InNvdXJjZSI6MjAsInRhcmdldCI6MjB9fV1d
\begin{tikzcd}[ampersand replacement=\&]
	{\C[3]} \& {\C[4]} \& {\C[1]}
	\arrow[""{name=0, anchor=center, inner sep=0}, "{F_1}"{description}, curve={height=-12pt}, from=1-1, to=1-2]
	\arrow[""{name=1, anchor=center, inner sep=0}, "{F_2}"{description}, curve={height=-12pt}, from=1-2, to=1-3]
	\arrow[""{name=2, anchor=center, inner sep=0}, "{G_1}"{description}, curve={height=12pt}, from=1-1, to=1-2]
	\arrow[""{name=3, anchor=center, inner sep=0}, "{G_2}"{description}, curve={height=12pt}, from=1-2, to=1-3]
	\arrow["\alpha", shorten <=3pt, shorten >=3pt, Rightarrow, from=0, to=2]
	\arrow["\beta", shorten <=3pt, shorten >=3pt, Rightarrow, from=1, to=3]
\end{tikzcd}
}

\newquiver{cat horizontal composition commutative}{
% https://q.uiver.app/?q=WzAsNCxbMCwwLCJGXzJGXzFhIl0sWzEsMCwiR18yRl8xYSJdLFswLDEsIkZfMkdfMWEiXSxbMSwxLCJHXzJHXzFhIl0sWzAsMSwiXFxiZXRhX3tGXzFhfSJdLFswLDIsIkZfMlxcYWxwaGFfYSIsMl0sWzEsMywiR18yXFxhbHBoYV9hIl0sWzIsMywiXFxiZXRhX3tHXzFhfSIsMl0sWzAsMywiKFxcYmV0YVxcaGNvbXBvc2VcXGFscGhhKV9hIiwxXV0=
\begin{tikzcd}[ampersand replacement=\&]
	{F_2F_1a} \& {G_2F_1a} \\
	{F_2G_1a} \& {G_2G_1a}
	\arrow["{\beta_{F_1a}}", from=1-1, to=1-2]
	\arrow["{F_2\alpha_a}"', from=1-1, to=2-1]
	\arrow["{G_2\alpha_a}", from=1-2, to=2-2]
	\arrow["{\beta_{G_1a}}"', from=2-1, to=2-2]
	\arrow["{(\beta\hcompose\alpha)_a}"{description}, from=1-1, to=2-2]
\end{tikzcd}
}

\newquiver{adjoint triangle one}{
% https://q.uiver.app/?q=WzAsNCxbMCwwLCJhIl0sWzEsMCwiYiJdLFsyLDAsImEiXSxbMywwLCJiIl0sWzAsMSwiZiIsMl0sWzEsMiwidSIsMl0sWzIsMywiZiIsMl0sWzEsMywiIiwyLHsiY3VydmUiOjQsImxldmVsIjoyLCJzdHlsZSI6eyJoZWFkIjp7Im5hbWUiOiJub25lIn19fV0sWzAsMiwiIiwyLHsiY3VydmUiOi00LCJsZXZlbCI6Miwic3R5bGUiOnsiaGVhZCI6eyJuYW1lIjoibm9uZSJ9fX1dLFsyLDcsIlxcY291bml0IiwyLHsic2hvcnRlbiI6eyJ0YXJnZXQiOjIwfX1dLFs4LDEsIlxcdW5pdCIsMix7InNob3J0ZW4iOnsic291cmNlIjoyMH19XV0=
\begin{tikzcd}[ampersand replacement=\&,row sep=2.25em]
	a \& b \& a \& b
	\arrow["f"', from=1-1, to=1-2]
	\arrow["u"', from=1-2, to=1-3]
	\arrow["f"', from=1-3, to=1-4]
	\arrow[""{name=0, anchor=center, inner sep=0}, curve={height=24pt}, Rightarrow, no head, from=1-2, to=1-4]
	\arrow[""{name=1, anchor=center, inner sep=0}, curve={height=-24pt}, Rightarrow, no head, from=1-1, to=1-3]
	\arrow["\counit"', shorten >=2pt, Rightarrow, from=1-3, to=0]
	\arrow["\unit"', shorten <=2pt, Rightarrow, from=1, to=1-2]
\end{tikzcd}
}

\newquiver{adjoint triangle two}{%
% https://q.uiver.app/?q=WzAsNCxbMSwwLCJhIl0sWzIsMCwiYiJdLFszLDAsImEiXSxbMCwwLCJiIl0sWzAsMSwiZiIsMl0sWzEsMiwidSIsMl0sWzAsMiwiIiwyLHsiY3VydmUiOi00LCJsZXZlbCI6Miwic3R5bGUiOnsiaGVhZCI6eyJuYW1lIjoibm9uZSJ9fX1dLFszLDAsInUiLDJdLFsxLDMsIiIsMix7ImN1cnZlIjotNCwibGV2ZWwiOjIsInN0eWxlIjp7ImhlYWQiOnsibmFtZSI6Im5vbmUifX19XSxbNiwxLCJcXHVuaXQiLDIseyJzaG9ydGVuIjp7InNvdXJjZSI6MjB9fV0sWzAsOCwiXFxjb3VuaXQiLDIseyJsYWJlbF9wb3NpdGlvbiI6MjAsInNob3J0ZW4iOnsidGFyZ2V0IjoyMH19XV0=
\begin{tikzcd}[ampersand replacement=\&,row sep=2.25em]
	b \& a \& b \& a
	\arrow["f"', from=1-2, to=1-3]
	\arrow["u"', from=1-3, to=1-4]
	\arrow[""{name=0, anchor=center, inner sep=0}, curve={height=-24pt}, Rightarrow, no head, from=1-2, to=1-4]
	\arrow["u"', from=1-1, to=1-2]
	\arrow[""{name=1, anchor=center, inner sep=0}, curve={height=-24pt}, Rightarrow, no head, from=1-3, to=1-1]
	\arrow["\unit"', shorten <=2pt, Rightarrow, from=0, to=1-3]
	\arrow["\counit"'{pos=0.2}, shorten >=2pt, Rightarrow, from=1-2, to=1]
\end{tikzcd}}

\newquiver{narrow alpha f g}{
% https://q.uiver.app/?q=WzAsMixbMCwwLCJhIl0sWzIsMCwiYiJdLFswLDEsImYiLDAseyJjdXJ2ZSI6LTJ9XSxbMCwxLCJnIiwyLHsiY3VydmUiOjJ9XSxbMiwzLCJcXGFscGhhIiwwLHsic2hvcnRlbiI6eyJzb3VyY2UiOjIwLCJ0YXJnZXQiOjIwfX1dXQ==
\begin{tikzcd}[ampersand replacement=\&,column sep=tiny]
	a \&\& b
	\arrow[""{name=0, anchor=center, inner sep=0}, "f", curve={height=-12pt}, from=1-1, to=1-3]
	\arrow[""{name=1, anchor=center, inner sep=0}, "g"', curve={height=12pt}, from=1-1, to=1-3]
	\arrow["\alpha", shorten <=3pt, shorten >=3pt, Rightarrow, from=0, to=1]
\end{tikzcd}
}

\newquiver{F(alpha f g)}{
% https://q.uiver.app/?q=WzAsMixbMCwwLCJcXEZGXzBhIl0sWzIsMCwiXFxGRl8wYiJdLFswLDEsIlxcRkZfMWYiLDAseyJjdXJ2ZSI6LTJ9XSxbMCwxLCJcXEZGXzFnIiwyLHsiY3VydmUiOjJ9XSxbMiwzLCJcXEZGXzJcXGFscGhhIiwwLHsic2hvcnRlbiI6eyJzb3VyY2UiOjIwLCJ0YXJnZXQiOjIwfX1dXQ==
\begin{tikzcd}[ampersand replacement=\&,column sep=small]
	{\FF_0a} \&\& {\FF_0b}
	\arrow[""{name=0, anchor=center, inner sep=0}, "{\FF_1f}", curve={height=-12pt}, from=1-1, to=1-3]
	\arrow[""{name=1, anchor=center, inner sep=0}, "{\FF_1g}"', curve={height=12pt}, from=1-1, to=1-3]
	\arrow["{\FF_2\alpha}", shorten <=3pt, shorten >=3pt, Rightarrow, from=0, to=1]
\end{tikzcd}}

\newquiver{F(adjoint triangle one)}{
% https://q.uiver.app/?q=WzAsNCxbMCwwLCJcXEZGIGEiXSxbMSwwLCJcXEZGIGIiXSxbMiwwLCJcXEZGIGEiXSxbMywwLCJcXEZGIGIiXSxbMCwxLCJcXEZGIGYiLDJdLFsxLDIsIlxcRkYgdSIsMl0sWzIsMywiXFxGRiBmIiwyXSxbMSwzLCIiLDIseyJjdXJ2ZSI6NCwibGV2ZWwiOjIsInN0eWxlIjp7ImhlYWQiOnsibmFtZSI6Im5vbmUifX19XSxbMCwyLCIiLDIseyJjdXJ2ZSI6LTQsImxldmVsIjoyLCJzdHlsZSI6eyJoZWFkIjp7Im5hbWUiOiJub25lIn19fV0sWzIsNywiXFxGRiBcXGNvdW5pdCIsMix7InNob3J0ZW4iOnsidGFyZ2V0IjoyMH19XSxbOCwxLCJcXEZGXFx1bml0IiwyLHsibGFiZWxfcG9zaXRpb24iOjMwLCJzaG9ydGVuIjp7InNvdXJjZSI6MjB9fV1d
\begin{tikzcd}[ampersand replacement=\&]
	{\FF a} \& {\FF b} \& {\FF a} \& {\FF b}
	\arrow["{\FF f}"', from=1-1, to=1-2]
	\arrow["{\FF u}"', from=1-2, to=1-3]
	\arrow["{\FF f}"', from=1-3, to=1-4]
	\arrow[""{name=0, anchor=center, inner sep=0}, curve={height=24pt}, Rightarrow, no head, from=1-2, to=1-4]
	\arrow[""{name=1, anchor=center, inner sep=0}, curve={height=-24pt}, Rightarrow, no head, from=1-1, to=1-3]
	\arrow["{\FF \counit}"', shorten >=2pt, Rightarrow, from=1-3, to=0]
	\arrow["\FF\unit"'{pos=0.3}, shorten <=2pt, Rightarrow, from=1, to=1-2]
\end{tikzcd}}

\newquiver{F(adjoint triangle two)}{
% https://q.uiver.app/?q=WzAsNCxbMSwwLCJcXEZGIGEiXSxbMiwwLCJcXEZGIGIiXSxbMywwLCJcXEZGIGEiXSxbMCwwLCJcXEZGIGIiXSxbMCwxLCJcXEZGIGYiLDJdLFsxLDIsIlxcRkYgdSIsMl0sWzAsMiwiIiwyLHsiY3VydmUiOi00LCJsZXZlbCI6Miwic3R5bGUiOnsiaGVhZCI6eyJuYW1lIjoibm9uZSJ9fX1dLFszLDAsIlxcRkYgdSIsMl0sWzEsMywiIiwyLHsiY3VydmUiOi00LCJsZXZlbCI6Miwic3R5bGUiOnsiaGVhZCI6eyJuYW1lIjoibm9uZSJ9fX1dLFs2LDEsIlxcRkYgXFx1bml0IiwyLHsibGFiZWxfcG9zaXRpb24iOjMwLCJzaG9ydGVuIjp7InNvdXJjZSI6MjB9fV0sWzAsOCwiXFxGRiBcXGNvdW5pdCIsMix7ImxhYmVsX3Bvc2l0aW9uIjoyMCwic2hvcnRlbiI6eyJ0YXJnZXQiOjIwfX1dXQ==
\begin{tikzcd}[ampersand replacement=\&]
	{\FF b} \& {\FF a} \& {\FF b} \& {\FF a}
	\arrow["{\FF f}"', from=1-2, to=1-3]
	\arrow["{\FF u}"', from=1-3, to=1-4]
	\arrow[""{name=0, anchor=center, inner sep=0}, curve={height=-24pt}, Rightarrow, no head, from=1-2, to=1-4]
	\arrow["{\FF u}"', from=1-1, to=1-2]
	\arrow[""{name=1, anchor=center, inner sep=0}, curve={height=-24pt}, Rightarrow, no head, from=1-3, to=1-1]
	\arrow["{\FF \unit}"'{pos=0.3}, shorten <=2pt, Rightarrow, from=0, to=1-3]
	\arrow["{\FF \counit}"'{pos=0.2}, shorten >=2pt, Rightarrow, from=1-2, to=1]
\end{tikzcd}
}

\newquiver{multi-morphism}{
  % https://q.uiver.app/?q=WzAsNSxbMCwwLCJhXzEiXSxbMCwxLCJcXHZkb3RzIl0sWzAsMiwiYV9uIl0sWzEsMSwiZiJdLFsyLDEsImIiXSxbMCwzLCIiLDIseyJzdHlsZSI6eyJoZWFkIjp7Im5hbWUiOiJub25lIn19fV0sWzIsMywiIiwwLHsic3R5bGUiOnsiaGVhZCI6eyJuYW1lIjoibm9uZSJ9fX1dLFszLDRdLFs1LDYsIiIsMix7ImN1cnZlIjoyLCJzaG9ydGVuIjp7InNvdXJjZSI6MjAsInRhcmdldCI6MjB9LCJsZXZlbCI6MX1dXQ==
\begin{tikzcd}[ampersand replacement=\&,sep=small]
	{a_1} \\
	\vdots \& f \& b \\
	{a_n}
	\arrow[""{name=0, anchor=center, inner sep=0}, no head, from=1-1, to=2-2]
	\arrow[""{name=1, anchor=center, inner sep=0}, no head, from=3-1, to=2-2]
	\arrow[from=2-2, to=2-3]
	\arrow[curve={height=12pt}, shorten <=5pt, shorten >=5pt, from=0, to=1]
\end{tikzcd}
}

\newquiver{multi composition}{
% https://q.uiver.app/?q=WzAsMTMsWzIsMSwiYl8xIl0sWzIsNSwiYl9uIl0sWzMsMywiZyJdLFs0LDMsImMiXSxbMSwxLCJmXzEiXSxbMSw1LCJmX24iXSxbMCwwLCJhX3sxMX0iXSxbMCwyLCJhX3sxbV8xfSJdLFswLDQsImFfe24xfSJdLFswLDYsImFfe25tX259Il0sWzAsMSwiXFx2ZG90cyJdLFswLDUsIlxcdmRvdHMiXSxbMSwzLCJcXHZkb3RzIl0sWzAsMiwiIiwyLHsic3R5bGUiOnsiaGVhZCI6eyJuYW1lIjoibm9uZSJ9fX1dLFsxLDIsIiIsMCx7InN0eWxlIjp7ImhlYWQiOnsibmFtZSI6Im5vbmUifX19XSxbMiwzXSxbNCwwXSxbNiw0LCIiLDIseyJzdHlsZSI6eyJoZWFkIjp7Im5hbWUiOiJub25lIn19fV0sWzcsNCwiIiwyLHsic3R5bGUiOnsiaGVhZCI6eyJuYW1lIjoibm9uZSJ9fX1dLFs1LDFdLFs4LDUsIiIsMCx7InN0eWxlIjp7ImhlYWQiOnsibmFtZSI6Im5vbmUifX19XSxbOSw1LCIiLDAseyJzdHlsZSI6eyJoZWFkIjp7Im5hbWUiOiJub25lIn19fV0sWzEzLDE0LCIiLDIseyJjdXJ2ZSI6Miwic2hvcnRlbiI6eyJzb3VyY2UiOjIwLCJ0YXJnZXQiOjIwfSwibGV2ZWwiOjF9XSxbMTcsMTgsIiIsMix7ImN1cnZlIjoyLCJzaG9ydGVuIjp7InNvdXJjZSI6MjAsInRhcmdldCI6MjB9LCJsZXZlbCI6MX1dLFsyMCwyMSwiIiwwLHsiY3VydmUiOjIsInNob3J0ZW4iOnsic291cmNlIjoyMCwidGFyZ2V0IjoyMH0sImxldmVsIjoxfV1d
\begin{tikzcd}[ampersand replacement=\&,sep=small]
	{a_{11}} \\
	\vdots \& {f_1} \& {b_1} \\
	{a_{1m_1}} \\
	\& \vdots \&\& g \& c \\
	{a_{n1}} \\
	\vdots \& {f_n} \& {b_n} \\
	{a_{nm_n}}
	\arrow[""{name=0, anchor=center, inner sep=0}, no head, from=2-3, to=4-4]
	\arrow[""{name=1, anchor=center, inner sep=0}, no head, from=6-3, to=4-4]
	\arrow[from=4-4, to=4-5]
	\arrow[from=2-2, to=2-3]
	\arrow[""{name=2, anchor=center, inner sep=0}, no head, from=1-1, to=2-2]
	\arrow[""{name=3, anchor=center, inner sep=0}, no head, from=3-1, to=2-2]
	\arrow[from=6-2, to=6-3]
	\arrow[""{name=4, anchor=center, inner sep=0}, no head, from=5-1, to=6-2]
	\arrow[""{name=5, anchor=center, inner sep=0}, no head, from=7-1, to=6-2]
	\arrow[curve={height=12pt}, shorten <=9pt, shorten >=9pt, from=0, to=1]
	\arrow[curve={height=12pt}, shorten <=6pt, shorten >=6pt, from=2, to=3]
	\arrow[curve={height=12pt}, shorten <=6pt, shorten >=6pt, from=4, to=5]
\end{tikzcd}
}

\newquiver{multi composition associativity}{
% https://q.uiver.app/?q=WzAsMjYsWzIsMSwiYl97MTF9Il0sWzIsNSwiYl97MW5fMX0iXSxbMywzLCJnXzEiXSxbNCwzLCJjXzEiXSxbMSwxLCJmX3sxMX0iXSxbMSw1LCJmX3sxbl8xfSJdLFswLDAsImFfezExMX0iXSxbMCwyLCJhX3sxMW1fMX0iXSxbMCw0LCJhX3sxbl8xMX0iXSxbMCw2LCJhX3sxbl8xbV97bl8xfX0iXSxbMCwxLCJcXHZkb3RzIl0sWzAsNSwiXFx2ZG90cyJdLFsxLDMsIlxcdmRvdHMiXSxbMyw3LCJcXHZkb3RzIl0sWzAsOCwiYV97azExfSJdLFswLDEwLCJhX3trMW1fMX0iXSxbMSw5LCJmX3trMX0iXSxbMiw5LCJiX3trMX0iXSxbMywxMSwiZ19rIl0sWzAsMTIsImFfe2tuX2sxfSJdLFswLDE0LCJhX3trbl9rbV97bl9rfX0iXSxbMSwxMywiZl97a25fa30iXSxbMiwxMywiYl97a25fa30iXSxbNCwxMSwiY19rIl0sWzUsNywiaCJdLFs2LDcsImQiXSxbMCwyLCIiLDIseyJzdHlsZSI6eyJoZWFkIjp7Im5hbWUiOiJub25lIn19fV0sWzEsMiwiIiwwLHsic3R5bGUiOnsiaGVhZCI6eyJuYW1lIjoibm9uZSJ9fX1dLFsyLDNdLFs0LDBdLFs2LDQsIiIsMix7InN0eWxlIjp7ImhlYWQiOnsibmFtZSI6Im5vbmUifX19XSxbNyw0LCIiLDIseyJzdHlsZSI6eyJoZWFkIjp7Im5hbWUiOiJub25lIn19fV0sWzUsMV0sWzgsNSwiIiwwLHsic3R5bGUiOnsiaGVhZCI6eyJuYW1lIjoibm9uZSJ9fX1dLFs5LDUsIiIsMCx7InN0eWxlIjp7ImhlYWQiOnsibmFtZSI6Im5vbmUifX19XSxbMTQsMTUsIlxcdmRvdHMiLDEseyJzdHlsZSI6eyJib2R5Ijp7Im5hbWUiOiJub25lIn0sImhlYWQiOnsibmFtZSI6Im5vbmUifX19XSxbMTQsMTYsIiIsMSx7InN0eWxlIjp7ImhlYWQiOnsibmFtZSI6Im5vbmUifX19XSxbMTUsMTYsIiIsMSx7InN0eWxlIjp7ImhlYWQiOnsibmFtZSI6Im5vbmUifX19XSxbMTYsMTddLFsxNywxOCwiIiwxLHsic3R5bGUiOnsiaGVhZCI6eyJuYW1lIjoibm9uZSJ9fX1dLFsxOSwyMCwiXFx2ZG90cyIsMSx7InN0eWxlIjp7ImJvZHkiOnsibmFtZSI6Im5vbmUifSwiaGVhZCI6eyJuYW1lIjoibm9uZSJ9fX1dLFsxOSwyMSwiIiwxLHsic3R5bGUiOnsiaGVhZCI6eyJuYW1lIjoibm9uZSJ9fX1dLFsyMCwyMSwiIiwxLHsic3R5bGUiOnsiaGVhZCI6eyJuYW1lIjoibm9uZSJ9fX1dLFsyMSwyMl0sWzE4LDIzXSxbMjIsMTgsIiIsMSx7InN0eWxlIjp7ImhlYWQiOnsibmFtZSI6Im5vbmUifX19XSxbMjQsMjVdLFszLDI0LCIiLDEseyJzdHlsZSI6eyJoZWFkIjp7Im5hbWUiOiJub25lIn19fV0sWzIzLDI0LCIiLDEseyJzdHlsZSI6eyJoZWFkIjp7Im5hbWUiOiJub25lIn19fV0sWzE2LDIxLCJcXHZkb3RzIiwxLHsic3R5bGUiOnsiYm9keSI6eyJuYW1lIjoibm9uZSJ9LCJoZWFkIjp7Im5hbWUiOiJub25lIn19fV0sWzI2LDI3LCIiLDIseyJjdXJ2ZSI6Miwic2hvcnRlbiI6eyJzb3VyY2UiOjIwLCJ0YXJnZXQiOjIwfSwibGV2ZWwiOjF9XSxbMzAsMzEsIiIsMix7ImN1cnZlIjoyLCJzaG9ydGVuIjp7InNvdXJjZSI6MjAsInRhcmdldCI6MjB9LCJsZXZlbCI6MX1dLFszMywzNCwiIiwwLHsiY3VydmUiOjIsInNob3J0ZW4iOnsic291cmNlIjoyMCwidGFyZ2V0IjoyMH0sImxldmVsIjoxfV0sWzM2LDM3LCIiLDEseyJjdXJ2ZSI6Miwic2hvcnRlbiI6eyJzb3VyY2UiOjIwLCJ0YXJnZXQiOjIwfSwibGV2ZWwiOjF9XSxbNDEsNDIsIiIsMSx7ImN1cnZlIjoyLCJzaG9ydGVuIjp7InNvdXJjZSI6MjAsInRhcmdldCI6MjB9LCJsZXZlbCI6MX1dLFs0Nyw0OCwiIiwxLHsiY3VydmUiOjUsInNob3J0ZW4iOnsic291cmNlIjoyMCwidGFyZ2V0IjoyMH0sImxldmVsIjoxfV0sWzM5LDQ1LCIiLDEseyJjdXJ2ZSI6Miwic2hvcnRlbiI6eyJzb3VyY2UiOjIwLCJ0YXJnZXQiOjIwfSwibGV2ZWwiOjF9XV0=
\begin{tikzcd}[ampersand replacement=\&,sep=small]
	{a_{111}} \\
	\vdots \& {f_{11}} \& {b_{11}} \\
	{a_{11m_1}} \\
	\& \vdots \&\& {g_1} \& {c_1} \\
	{a_{1n_11}} \\
	\vdots \& {f_{1n_1}} \& {b_{1n_1}} \\
	{a_{1n_1m_{n_1}}} \\
	\&\&\& \vdots \&\& h \& d \\
	{a_{k11}} \\
	\& {f_{k1}} \& {b_{k1}} \\
	{a_{k1m_1}} \\
	\&\&\& {g_k} \& {c_k} \\
	{a_{kn_k1}} \\
	\& {f_{kn_k}} \& {b_{kn_k}} \\
	{a_{kn_km_{n_k}}}
	\arrow[""{name=0, anchor=center, inner sep=0}, no head, from=2-3, to=4-4]
	\arrow[""{name=1, anchor=center, inner sep=0}, no head, from=6-3, to=4-4]
	\arrow[from=4-4, to=4-5]
	\arrow[from=2-2, to=2-3]
	\arrow[""{name=2, anchor=center, inner sep=0}, no head, from=1-1, to=2-2]
	\arrow[""{name=3, anchor=center, inner sep=0}, no head, from=3-1, to=2-2]
	\arrow[from=6-2, to=6-3]
	\arrow[""{name=4, anchor=center, inner sep=0}, no head, from=5-1, to=6-2]
	\arrow[""{name=5, anchor=center, inner sep=0}, no head, from=7-1, to=6-2]
	\arrow["\vdots"{description}, draw=none, from=9-1, to=11-1]
	\arrow[""{name=6, anchor=center, inner sep=0}, no head, from=9-1, to=10-2]
	\arrow[""{name=7, anchor=center, inner sep=0}, no head, from=11-1, to=10-2]
	\arrow[from=10-2, to=10-3]
	\arrow[""{name=8, anchor=center, inner sep=0}, no head, from=10-3, to=12-4]
	\arrow["\vdots"{description}, draw=none, from=13-1, to=15-1]
	\arrow[""{name=9, anchor=center, inner sep=0}, no head, from=13-1, to=14-2]
	\arrow[""{name=10, anchor=center, inner sep=0}, no head, from=15-1, to=14-2]
	\arrow[from=14-2, to=14-3]
	\arrow[from=12-4, to=12-5]
	\arrow[""{name=11, anchor=center, inner sep=0}, no head, from=14-3, to=12-4]
	\arrow[from=8-6, to=8-7]
	\arrow[""{name=12, anchor=center, inner sep=0}, no head, from=4-5, to=8-6]
	\arrow[""{name=13, anchor=center, inner sep=0}, no head, from=12-5, to=8-6]
	\arrow["\vdots"{description}, draw=none, from=10-2, to=14-2]
	\arrow[curve={height=12pt}, shorten <=9pt, shorten >=9pt, from=0, to=1]
	\arrow[curve={height=12pt}, shorten <=5pt, shorten >=5pt, from=2, to=3]
	\arrow[curve={height=12pt}, shorten <=5pt, shorten >=5pt, from=4, to=5]
	\arrow[curve={height=12pt}, shorten <=5pt, shorten >=5pt, from=6, to=7]
	\arrow[curve={height=12pt}, shorten <=5pt, shorten >=5pt, from=9, to=10]
	\arrow[curve={height=30pt}, shorten <=19pt, shorten >=19pt, from=12, to=13]
	\arrow[curve={height=12pt}, shorten <=9pt, shorten >=9pt, from=8, to=11]
\end{tikzcd}
}

\newquiver{multi-natural transformation type}{
% https://q.uiver.app/?q=WzAsMixbMCwwLCJcXFVVWzFdIl0sWzEsMCwiXFxVVVsyXSJdLFswLDEsIkYiLDAseyJjdXJ2ZSI6LTJ9XSxbMCwxLCJHIiwyLHsiY3VydmUiOjJ9XSxbMiwzLCJcXGFscGhhIiwwLHsic2hvcnRlbiI6eyJzb3VyY2UiOjIwLCJ0YXJnZXQiOjIwfX1dXQ==
\begin{tikzcd}[ampersand replacement=\&]
	{\UU[1]} \& {\UU[2]}
	\arrow[""{name=0, anchor=center, inner sep=0}, "F", curve={height=-12pt}, from=1-1, to=1-2]
	\arrow[""{name=1, anchor=center, inner sep=0}, "G"', curve={height=12pt}, from=1-1, to=1-2]
	\arrow["\alpha", shorten <=3pt, shorten >=3pt, Rightarrow, from=0, to=1]
\end{tikzcd}
}

\newquiver{multi-naturality}{
% https://q.uiver.app/?q=WzAsOCxbMSwxLCJGYV8xIl0sWzAsMCwiRmFfbiJdLFsyLDEsIkdhXzEiXSxbMywwLCJHYV9uIl0sWzAsMiwiRmYiXSxbMywyLCJHZiJdLFswLDMsIkZiIl0sWzMsMywiR2IiXSxbMSw0LCIiLDAseyJzdHlsZSI6eyJoZWFkIjp7Im5hbWUiOiJub25lIn19fV0sWzAsNCwiIiwyLHsic3R5bGUiOnsiaGVhZCI6eyJuYW1lIjoibm9uZSJ9fX1dLFsyLDUsIiIsMix7InN0eWxlIjp7ImhlYWQiOnsibmFtZSI6Im5vbmUifX19XSxbMyw1LCIiLDAseyJzdHlsZSI6eyJoZWFkIjp7Im5hbWUiOiJub25lIn19fV0sWzYsNywiXFxhbHBoYV9iIiwyXSxbNCw2XSxbNSw3XSxbMCwyLCJcXGFscGhhX3thXzF9IiwyXSxbMSwwLCJcXGRkb3RzIiwxLHsic3R5bGUiOnsiYm9keSI6eyJuYW1lIjoibm9uZSJ9LCJoZWFkIjp7Im5hbWUiOiJub25lIn19fV0sWzMsMiwiXFxpZGRvdHMiLDEseyJzdHlsZSI6eyJib2R5Ijp7Im5hbWUiOiJub25lIn0sImhlYWQiOnsibmFtZSI6Im5vbmUifX19XSxbMSwzLCJcXGFscGhhX3thX259Il0sWzksOCwiIiwyLHsiY3VydmUiOjEsInNob3J0ZW4iOnsic291cmNlIjoyMCwidGFyZ2V0IjoyMH0sImxldmVsIjoxfV0sWzEwLDExLCIiLDIseyJjdXJ2ZSI6LTEsInNob3J0ZW4iOnsic291cmNlIjoyMCwidGFyZ2V0IjoyMH0sImxldmVsIjoxfV0sWzE1LDEyLCI9IiwxLHsic2hvcnRlbiI6eyJzb3VyY2UiOjIwLCJ0YXJnZXQiOjIwfSwic3R5bGUiOnsiYm9keSI6eyJuYW1lIjoibm9uZSJ9LCJoZWFkIjp7Im5hbWUiOiJub25lIn19fV1d
\begin{tikzcd}[ampersand replacement=\&]
	{Fa_n} \&\&\& {Ga_n} \\
	\& {Fa_1} \& {Ga_1} \\
	Ff \&\&\& Gf \\
	Fb \&\&\& Gb
	\arrow[""{name=0, anchor=center, inner sep=0}, no head, from=1-1, to=3-1]
	\arrow[""{name=1, anchor=center, inner sep=0}, no head, from=2-2, to=3-1]
	\arrow[""{name=2, anchor=center, inner sep=0}, no head, from=2-3, to=3-4]
	\arrow[""{name=3, anchor=center, inner sep=0}, no head, from=1-4, to=3-4]
	\arrow[""{name=4, anchor=center, inner sep=0}, "{\alpha_b}"', from=4-1, to=4-4]
	\arrow[from=3-1, to=4-1]
	\arrow[from=3-4, to=4-4]
	\arrow[""{name=5, anchor=center, inner sep=0}, "{\alpha_{a_1}}"', from=2-2, to=2-3]
	\arrow["\ddots"{description}, draw=none, from=1-1, to=2-2]
	\arrow["\iddots"{description}, draw=none, from=1-4, to=2-3]
	\arrow["{\alpha_{a_n}}", from=1-1, to=1-4]
	\arrow[curve={height=6pt}, shorten <=4pt, shorten >=4pt, from=1, to=0]
	\arrow[curve={height=-6pt}, shorten <=4pt, shorten >=4pt, from=2, to=3]
	\arrow["{=}"{description}, draw=none, from=5, to=4]
\end{tikzcd}
}

\newquiver{parallel multi-natural}{
% https://q.uiver.app/?q=WzAsMixbMCwwLCJcXFVVWzFdIl0sWzEsMCwiXFxVVVsyXSJdLFswLDEsIkYiLDAseyJjdXJ2ZSI6LTJ9XSxbMCwxLCJHIiwyLHsiY3VydmUiOjJ9XSxbMiwzLCJcXGFscGhhIiwyLHsib2Zmc2V0IjoxLCJzaG9ydGVuIjp7InNvdXJjZSI6MjAsInRhcmdldCI6MjB9fV0sWzIsMywiXFxiZXRhIiwwLHsib2Zmc2V0IjotMSwic2hvcnRlbiI6eyJzb3VyY2UiOjIwLCJ0YXJnZXQiOjIwfX1dXQ==
\begin{tikzcd}[ampersand replacement=\&]
	{\UU[1]} \& {\UU[2]}
	\arrow[""{name=0, anchor=center, inner sep=0}, "F", curve={height=-12pt}, from=1-1, to=1-2]
	\arrow[""{name=1, anchor=center, inner sep=0}, "G"', curve={height=12pt}, from=1-1, to=1-2]
	\arrow["\alpha"', shift right=1, shorten <=3pt, shorten >=3pt, Rightarrow, from=0, to=1]
	\arrow["\beta", shift left=1, shorten <=3pt, shorten >=3pt, Rightarrow, from=0, to=1]
\end{tikzcd}
}
\newquiver{2-cell action of Mod(O,-)}{
% https://q.uiver.app/?q=WzAsMixbMCwwLCJcXE1vZChcXE9PLFxcVVVbMV0pIl0sWzIsMCwiXFxNb2QoXFxPTyxcXFVVWzJdKSJdLFswLDEsIlxcTW9kKFxcT08sRikiLDAseyJjdXJ2ZSI6LTJ9XSxbMCwxLCJcXE1vZChcXE9PLEcpIiwyLHsiY3VydmUiOjJ9XSxbMiwzLCJcXE1vZChcXE9PLFxcYWxwaGEpIiwwLHsib2Zmc2V0Ijo1LCJzaG9ydGVuIjp7InNvdXJjZSI6MjAsInRhcmdldCI6MjB9fV1d
\begin{tikzcd}[ampersand replacement=\&]
	{\Mod(\OO,\UU[1])} \&\& {\Mod(\OO,\UU[2])}
	\arrow[""{name=0, anchor=center, inner sep=0}, "{\Mod(\OO,F)}", curve={height=-12pt}, from=1-1, to=1-3]
	\arrow[""{name=1, anchor=center, inner sep=0}, "{\Mod(\OO,G)}"', curve={height=12pt}, from=1-1, to=1-3]
	\arrow["{\Mod(\OO,\alpha)}", shift right=5, shorten <=3pt, shorten >=3pt, Rightarrow, from=0, to=1]
\end{tikzcd}}


\newquiver{O-homo as multi-natural}{
  % https://q.uiver.app/?q=WzAsMixbMCwwLCJcXE9PIl0sWzEsMCwiXFxVVSJdLFswLDEsIlxcQVsxXSIsMCx7ImN1cnZlIjotMn1dLFswLDEsIlxcQVsyXSIsMix7ImN1cnZlIjoyfV0sWzIsMywiaCIsMCx7InNob3J0ZW4iOnsic291cmNlIjoyMCwidGFyZ2V0IjoyMH19XV0=
\begin{tikzcd}[ampersand replacement=\&]
	\OO \& \UU
	\arrow[""{name=0, anchor=center, inner sep=0}, "{\A[1]}", curve={height=-12pt}, from=1-1, to=1-2]
	\arrow[""{name=1, anchor=center, inner sep=0}, "{\A[2]}"', curve={height=12pt}, from=1-1, to=1-2]
	\arrow["h", shorten <=3pt, shorten >=3pt, Rightarrow, from=0, to=1]
\end{tikzcd}}

\newquiver{O-homo image under 2-functor}{
% https://q.uiver.app/?q=WzAsMyxbMCwwLCJcXE9PIl0sWzEsMCwiXFxVVSJdLFsyLDAsIlxcVVVbMl0iXSxbMCwxLCJcXEFbMV0iLDAseyJjdXJ2ZSI6LTJ9XSxbMCwxLCJcXEFbMl0iLDIseyJjdXJ2ZSI6Mn1dLFsxLDIsIkYiLDAseyJjdXJ2ZSI6LTJ9XSxbMSwyLCJGIiwyLHsiY3VydmUiOjJ9XSxbMyw0LCJoIiwwLHsic2hvcnRlbiI6eyJzb3VyY2UiOjIwLCJ0YXJnZXQiOjIwfX1dLFs1LDYsIj0iLDEseyJzaG9ydGVuIjp7InNvdXJjZSI6MjAsInRhcmdldCI6MjB9LCJzdHlsZSI6eyJib2R5Ijp7Im5hbWUiOiJub25lIn0sImhlYWQiOnsibmFtZSI6Im5vbmUifX19XV0=
\begin{tikzcd}[ampersand replacement=\&]
	\OO \& \UU \& {\UU[2]}
	\arrow[""{name=0, anchor=center, inner sep=0}, "{\A[1]}", curve={height=-12pt}, from=1-1, to=1-2]
	\arrow[""{name=1, anchor=center, inner sep=0}, "{\A[2]}"', curve={height=12pt}, from=1-1, to=1-2]
	\arrow[""{name=2, anchor=center, inner sep=0}, "F", curve={height=-12pt}, from=1-2, to=1-3]
	\arrow[""{name=3, anchor=center, inner sep=0}, "F"', curve={height=12pt}, from=1-2, to=1-3]
	\arrow["h", shorten <=3pt, shorten >=3pt, Rightarrow, from=0, to=1]
	\arrow["{=}"{description}, draw=none, from=2, to=3]
\end{tikzcd}}

\newquiver{O-homo 2-functoriality:composition}{
% https://q.uiver.app/?q=WzAsMyxbMCwwLCJcXE9PIl0sWzEsMCwiXFxVVSJdLFsyLDAsIlxcVVVbMl0iXSxbMCwxLCJcXEFbMV0iLDAseyJjdXJ2ZSI6LTV9XSxbMCwxLCJcXEFbMl0iLDAseyJsYWJlbF9wb3NpdGlvbiI6MjB9XSxbMSwyLCJGIiwwLHsiY3VydmUiOi0yfV0sWzEsMiwiRiIsMix7ImN1cnZlIjoyfV0sWzAsMSwiXFxBWzNdIiwyLHsiY3VydmUiOjV9XSxbMSwyXSxbMyw0LCJoXzEiLDAseyJzaG9ydGVuIjp7InNvdXJjZSI6MjAsInRhcmdldCI6MjB9fV0sWzQsNywiaF8yIiwwLHsic2hvcnRlbiI6eyJzb3VyY2UiOjIwLCJ0YXJnZXQiOjIwfX1dLFs1LDgsIj0iLDEseyJzaG9ydGVuIjp7InNvdXJjZSI6MjAsInRhcmdldCI6MjB9LCJzdHlsZSI6eyJib2R5Ijp7Im5hbWUiOiJub25lIn0sImhlYWQiOnsibmFtZSI6Im5vbmUifX19XSxbNiw4LCI9IiwxLHsic2hvcnRlbiI6eyJzb3VyY2UiOjIwLCJ0YXJnZXQiOjIwfSwic3R5bGUiOnsiYm9keSI6eyJuYW1lIjoibm9uZSJ9LCJoZWFkIjp7Im5hbWUiOiJub25lIn19fV1d
\begin{tikzcd}[ampersand replacement=\&]
	\OO \& \UU \& {\UU[2]}
	\arrow[""{name=0, anchor=center, inner sep=0}, "{\A[1]}", curve={height=-30pt}, from=1-1, to=1-2]
	\arrow[""{name=1, anchor=center, inner sep=0}, "{\A[2]}"{pos=0.2}, from=1-1, to=1-2]
	\arrow[""{name=2, anchor=center, inner sep=0}, "F", curve={height=-12pt}, from=1-2, to=1-3]
	\arrow[""{name=3, anchor=center, inner sep=0}, "F"', curve={height=12pt}, from=1-2, to=1-3]
	\arrow[""{name=4, anchor=center, inner sep=0}, "{\A[3]}"', curve={height=30pt}, from=1-1, to=1-2]
	\arrow[""{name=5, anchor=center, inner sep=0}, from=1-2, to=1-3]
	\arrow["{h_1}", shorten <=4pt, shorten >=4pt, Rightarrow, from=0, to=1]
	\arrow["{h_2}", shorten <=4pt, shorten >=4pt, Rightarrow, from=1, to=4]
	\arrow["{=}"{description}, draw=none, from=2, to=5]
	\arrow["{=}"{description}, draw=none, from=3, to=5]
\end{tikzcd}
}

\newquiver{O-homo 2-functoriality:identity}{
  % https://q.uiver.app/?q=WzAsMyxbMCwwLCJcXE9PIl0sWzEsMCwiXFxVVSJdLFsyLDAsIlxcVVVbMl0iXSxbMCwxLCJcXEFbMV0iLDAseyJjdXJ2ZSI6LTJ9XSxbMCwxLCJcXEFbMV0iLDIseyJjdXJ2ZSI6Mn1dLFsxLDIsIkYiLDAseyJjdXJ2ZSI6LTJ9XSxbMSwyLCJGIiwyLHsiY3VydmUiOjJ9XSxbMyw0LCI9IiwxLHsic2hvcnRlbiI6eyJzb3VyY2UiOjIwLCJ0YXJnZXQiOjIwfSwic3R5bGUiOnsiYm9keSI6eyJuYW1lIjoibm9uZSJ9LCJoZWFkIjp7Im5hbWUiOiJub25lIn19fV0sWzUsNiwiPSIsMSx7InNob3J0ZW4iOnsic291cmNlIjoyMCwidGFyZ2V0IjoyMH0sInN0eWxlIjp7ImJvZHkiOnsibmFtZSI6Im5vbmUifSwiaGVhZCI6eyJuYW1lIjoibm9uZSJ9fX1dXQ==
\begin{tikzcd}[ampersand replacement=\&]
	\OO \& \UU \& {\UU[2]}
	\arrow[""{name=0, anchor=center, inner sep=0}, "{\A[1]}", curve={height=-12pt}, from=1-1, to=1-2]
	\arrow[""{name=1, anchor=center, inner sep=0}, "{\A[1]}"', curve={height=12pt}, from=1-1, to=1-2]
	\arrow[""{name=2, anchor=center, inner sep=0}, "F", curve={height=-12pt}, from=1-2, to=1-3]
	\arrow[""{name=3, anchor=center, inner sep=0}, "F"', curve={height=12pt}, from=1-2, to=1-3]
	\arrow["{=}"{description}, draw=none, from=0, to=1]
	\arrow["{=}"{description}, draw=none, from=2, to=3]
\end{tikzcd}}

\newquiver{Mod(OO,alpha)_A as a whiskering}{
% https://q.uiver.app/?q=WzAsMyxbMCwwLCJcXE9PIl0sWzEsMCwiXFxVVSJdLFsyLDAsIlxcVVVbMl0iXSxbMCwxLCJcXEFbMV0iLDAseyJjdXJ2ZSI6LTJ9XSxbMCwxLCJcXEFbMV0iLDIseyJjdXJ2ZSI6Mn1dLFsxLDIsIkYiLDAseyJjdXJ2ZSI6LTJ9XSxbMSwyLCJHIiwyLHsiY3VydmUiOjJ9XSxbMyw0LCI9IiwxLHsic2hvcnRlbiI6eyJzb3VyY2UiOjIwLCJ0YXJnZXQiOjIwfSwic3R5bGUiOnsiYm9keSI6eyJuYW1lIjoibm9uZSJ9LCJoZWFkIjp7Im5hbWUiOiJub25lIn19fV0sWzUsNiwiXFxhbHBoYSIsMCx7InNob3J0ZW4iOnsic291cmNlIjoyMCwidGFyZ2V0IjoyMH19XV0=
\begin{tikzcd}[ampersand replacement=\&]
	\OO \& \UU \& {\UU[2]}
	\arrow[""{name=0, anchor=center, inner sep=0}, "{\A[1]}", curve={height=-12pt}, from=1-1, to=1-2]
	\arrow[""{name=1, anchor=center, inner sep=0}, "{\A[1]}"', curve={height=12pt}, from=1-1, to=1-2]
	\arrow[""{name=2, anchor=center, inner sep=0}, "F", curve={height=-12pt}, from=1-2, to=1-3]
	\arrow[""{name=3, anchor=center, inner sep=0}, "G"', curve={height=12pt}, from=1-2, to=1-3]
	\arrow["{=}"{description}, draw=none, from=0, to=1]
	\arrow["\alpha", shorten <=3pt, shorten >=3pt, Rightarrow, from=2, to=3]
\end{tikzcd}}

\newquiver{Mod(OO,alpha) naturality as horizontal composition}{
% https://q.uiver.app/?q=WzAsMyxbMCwwLCJcXE9PIl0sWzEsMCwiXFxVVSJdLFsyLDAsIlxcVVVbMl0iXSxbMCwxLCJcXEFbMV0iLDAseyJjdXJ2ZSI6LTJ9XSxbMCwxLCJcXEFbMV0iLDIseyJjdXJ2ZSI6Mn1dLFsxLDIsIkYiLDAseyJjdXJ2ZSI6LTJ9XSxbMSwyLCJHIiwyLHsiY3VydmUiOjJ9XSxbMyw0LCJoIiwwLHsic2hvcnRlbiI6eyJzb3VyY2UiOjIwLCJ0YXJnZXQiOjIwfX1dLFs1LDYsIlxcYWxwaGEiLDAseyJzaG9ydGVuIjp7InNvdXJjZSI6MjAsInRhcmdldCI6MjB9fV1d
\begin{tikzcd}[ampersand replacement=\&]
	\OO \& \UU \& {\UU[2]}
	\arrow[""{name=0, anchor=center, inner sep=0}, "{\A[1]}", curve={height=-12pt}, from=1-1, to=1-2]
	\arrow[""{name=1, anchor=center, inner sep=0}, "{\A[1]}"', curve={height=12pt}, from=1-1, to=1-2]
	\arrow[""{name=2, anchor=center, inner sep=0}, "F", curve={height=-12pt}, from=1-2, to=1-3]
	\arrow[""{name=3, anchor=center, inner sep=0}, "G"', curve={height=12pt}, from=1-2, to=1-3]
	\arrow["h", shorten <=3pt, shorten >=3pt, Rightarrow, from=0, to=1]
	\arrow["\alpha", shorten <=3pt, shorten >=3pt, Rightarrow, from=2, to=3]
\end{tikzcd}}

\newquiver{Mod(OO,-) preservers vertical composition}{% https://q.uiver.app/?q=WzAsMyxbMCwwLCJcXE9PIl0sWzEsMCwiXFxVVSJdLFsyLDAsIlxcVVVbMl0iXSxbMCwxLCJcXEFbMV0iLDAseyJjdXJ2ZSI6LTJ9XSxbMCwxLCJcXEFbMV0iLDIseyJjdXJ2ZSI6Mn1dLFsxLDIsIkYiLDAseyJjdXJ2ZSI6LTN9XSxbMSwyLCJIIiwyLHsiY3VydmUiOjN9XSxbMSwyLCJHIiwwLHsibGFiZWxfcG9zaXRpb24iOjIwfV0sWzMsNCwiPSIsMSx7InNob3J0ZW4iOnsic291cmNlIjoyMCwidGFyZ2V0IjoyMH0sInN0eWxlIjp7ImJvZHkiOnsibmFtZSI6Im5vbmUifSwiaGVhZCI6eyJuYW1lIjoibm9uZSJ9fX1dLFs1LDcsIlxcYWxwaGEiLDAseyJzaG9ydGVuIjp7InNvdXJjZSI6MjAsInRhcmdldCI6MjB9fV0sWzcsNiwiXFxiZXRhIiwwLHsic2hvcnRlbiI6eyJzb3VyY2UiOjIwLCJ0YXJnZXQiOjIwfX1dXQ==
\begin{tikzcd}[ampersand replacement=\&]
	\OO \& \UU \& {\UU[2]}
	\arrow[""{name=0, anchor=center, inner sep=0}, "{\A[1]}", curve={height=-12pt}, from=1-1, to=1-2]
	\arrow[""{name=1, anchor=center, inner sep=0}, "{\A[1]}"', curve={height=12pt}, from=1-1, to=1-2]
	\arrow[""{name=2, anchor=center, inner sep=0}, "F", curve={height=-18pt}, from=1-2, to=1-3]
	\arrow[""{name=3, anchor=center, inner sep=0}, "H"', curve={height=18pt}, from=1-2, to=1-3]
	\arrow[""{name=4, anchor=center, inner sep=0}, "G"{pos=0.2}, from=1-2, to=1-3]
	\arrow["{=}"{description}, draw=none, from=0, to=1]
	\arrow["\alpha", shorten <=2pt, shorten >=2pt, Rightarrow, from=2, to=4]
	\arrow["\beta", shorten <=2pt, shorten >=2pt, Rightarrow, from=4, to=3]
\end{tikzcd}}

\newquiver{Mod(OO,-) preservers horizontal composition}{% https://q.uiver.app/?q=WzAsNCxbMCwwLCJcXE9PIl0sWzEsMCwiXFxVVSJdLFsyLDAsIlxcVVVbMl0iXSxbMywwLCJcXFVVWzNdIl0sWzAsMSwiXFxBWzFdIiwwLHsiY3VydmUiOi0yfV0sWzAsMSwiXFxBWzFdIiwyLHsiY3VydmUiOjJ9XSxbMSwyLCJGXzEiLDAseyJjdXJ2ZSI6LTJ9XSxbMSwyLCJHXzEiLDIseyJjdXJ2ZSI6Mn1dLFsyLDMsIkZfMiIsMCx7ImN1cnZlIjotMn1dLFsyLDMsIkdfMiIsMix7ImN1cnZlIjoyfV0sWzQsNSwiPSIsMSx7InNob3J0ZW4iOnsic291cmNlIjoyMCwidGFyZ2V0IjoyMH0sInN0eWxlIjp7ImJvZHkiOnsibmFtZSI6Im5vbmUifSwiaGVhZCI6eyJuYW1lIjoibm9uZSJ9fX1dLFs2LDcsIlxcYWxwaGEiLDAseyJzaG9ydGVuIjp7InNvdXJjZSI6MjAsInRhcmdldCI6MjB9fV0sWzgsOSwiXFxiZXRhIiwwLHsic2hvcnRlbiI6eyJzb3VyY2UiOjIwLCJ0YXJnZXQiOjIwfX1dXQ==
\begin{tikzcd}[ampersand replacement=\&]
	\OO \& \UU \& {\UU[2]} \& {\UU[3]}
	\arrow[""{name=0, anchor=center, inner sep=0}, "{\A[1]}", curve={height=-12pt}, from=1-1, to=1-2]
	\arrow[""{name=1, anchor=center, inner sep=0}, "{\A[1]}"', curve={height=12pt}, from=1-1, to=1-2]
	\arrow[""{name=2, anchor=center, inner sep=0}, "{F_1}", curve={height=-12pt}, from=1-2, to=1-3]
	\arrow[""{name=3, anchor=center, inner sep=0}, "{G_1}"', curve={height=12pt}, from=1-2, to=1-3]
	\arrow[""{name=4, anchor=center, inner sep=0}, "{F_2}", curve={height=-12pt}, from=1-3, to=1-4]
	\arrow[""{name=5, anchor=center, inner sep=0}, "{G_2}"', curve={height=12pt}, from=1-3, to=1-4]
	\arrow["{=}"{description}, draw=none, from=0, to=1]
	\arrow["\alpha", shorten <=3pt, shorten >=3pt, Rightarrow, from=2, to=3]
	\arrow["\beta", shorten <=3pt, shorten >=3pt, Rightarrow, from=4, to=5]
\end{tikzcd}}

\newquiver{Mod(OO,-) preservers identities}{% https://q.uiver.app/?q=WzAsMyxbMCwwLCJcXE9PIl0sWzEsMCwiXFxVVSJdLFsyLDAsIlxcVVVbMl0iXSxbMCwxLCJcXEFbMV0iLDAseyJjdXJ2ZSI6LTJ9XSxbMCwxLCJcXEFbMV0iLDIseyJjdXJ2ZSI6Mn1dLFsxLDIsIkYiLDAseyJjdXJ2ZSI6LTJ9XSxbMSwyLCJHIiwyLHsiY3VydmUiOjJ9XSxbMyw0LCI9IiwxLHsic2hvcnRlbiI6eyJzb3VyY2UiOjIwLCJ0YXJnZXQiOjIwfSwic3R5bGUiOnsiYm9keSI6eyJuYW1lIjoibm9uZSJ9LCJoZWFkIjp7Im5hbWUiOiJub25lIn19fV0sWzUsNiwiPSIsMSx7InNob3J0ZW4iOnsic291cmNlIjoyMCwidGFyZ2V0IjoyMH0sInN0eWxlIjp7ImJvZHkiOnsibmFtZSI6Im5vbmUifSwiaGVhZCI6eyJuYW1lIjoibm9uZSJ9fX1dXQ==
\begin{tikzcd}[ampersand replacement=\&]
	\OO \& \UU \& {\UU[2]}
	\arrow[""{name=0, anchor=center, inner sep=0}, "{\A[1]}", curve={height=-12pt}, from=1-1, to=1-2]
	\arrow[""{name=1, anchor=center, inner sep=0}, "{\A[1]}"', curve={height=12pt}, from=1-1, to=1-2]
	\arrow[""{name=2, anchor=center, inner sep=0}, "F", curve={height=-12pt}, from=1-2, to=1-3]
	\arrow[""{name=3, anchor=center, inner sep=0}, "G"', curve={height=12pt}, from=1-2, to=1-3]
	\arrow["{=}"{description}, draw=none, from=0, to=1]
	\arrow["{=}"{description}, draw=none, from=2, to=3]
\end{tikzcd}}

\newquiver{nat trans between pp functors}{
  % https://q.uiver.app/?q=WzAsMixbMCwwLCJcXENbMV0iXSxbMSwwLCJcXENbMl0iXSxbMCwxLCJIIiwwLHsiY3VydmUiOi0yfV0sWzAsMSwiSyIsMix7ImN1cnZlIjoyfV0sWzIsMywiXFxhbHBoYSIsMCx7InNob3J0ZW4iOnsic291cmNlIjoyMCwidGFyZ2V0IjoyMH19XV0=
\begin{tikzcd}[ampersand replacement=\&]
	{\C[1]} \& {\C[2]}
	\arrow[""{name=0, anchor=center, inner sep=0}, "H", curve={height=-12pt}, from=1-1, to=1-2]
	\arrow[""{name=1, anchor=center, inner sep=0}, "K"', curve={height=12pt}, from=1-1, to=1-2]
	\arrow["\alpha", shorten <=3pt, shorten >=3pt, Rightarrow, from=0, to=1]
\end{tikzcd}
}

\newquiver{Mod(TT, alpha) type}{
% https://q.uiver.app/?q=WzAsMixbMCwwLCJcXE11bHRpKFxcVFQsXFxDWzFdKSJdLFszLDAsIlxcTXVsdGkoXFxUVCxcXENbMl0pIl0sWzAsMSwiXFxNdWx0aShcXFRULEgpIiwwLHsiY3VydmUiOi0yfV0sWzAsMSwiXFxNdWx0aShcXFRULEspIiwyLHsiY3VydmUiOjJ9XSxbMiwzLCJcXE11bHRpKFxcVFQsXFxhbHBoYSkiLDAseyJvZmZzZXQiOjUsInNob3J0ZW4iOnsic291cmNlIjoyMCwidGFyZ2V0IjoyMH19XV0=
\begin{tikzcd}[ampersand replacement=\&]
	{\Multi(\TT,\C[1])} \&\&\& {\Multi(\TT,\C[2])}
	\arrow[""{name=0, anchor=center, inner sep=0}, "{\Multi(\TT,H)}", curve={height=-12pt}, from=1-1, to=1-4]
	\arrow[""{name=1, anchor=center, inner sep=0}, "{\Multi(\TT,K)}"', curve={height=12pt}, from=1-1, to=1-4]
	\arrow["{\Multi(\TT,\alpha)}", shift right=5, shorten <=3pt, shorten >=3pt, Rightarrow, from=0, to=1]
\end{tikzcd}}

\newquiver{strict natural 2-transformation: 1-cells}{
% https://q.uiver.app/?q=WzAsNCxbMCwwLCJcXEZGWzFdXzBhIl0sWzEsMCwiXFxGRlsxXV8wYiJdLFswLDEsIlxcRkZbMl1fMGEiXSxbMSwxLCJcXEZGWzJdXzBiIl0sWzAsMSwiXFxGRlsxXV8xZiJdLFsyLDMsIlxcRkZbMl1fMWYiLDJdLFswLDIsIlxcWGlfYSIsMl0sWzEsMywiXFxYaV9iIl0sWzIsMSwiPSIsMSx7InN0eWxlIjp7ImJvZHkiOnsibmFtZSI6Im5vbmUifSwiaGVhZCI6eyJuYW1lIjoibm9uZSJ9fX1dXQ==
\begin{tikzcd}[ampersand replacement=\&]
	{\FF[1]_0a} \& {\FF[1]_0b} \\
	{\FF[2]_0a} \& {\FF[2]_0b}
	\arrow["{\FF[1]_1f}", from=1-1, to=1-2]
	\arrow["{\FF[2]_1f}"', from=2-1, to=2-2]
	\arrow["{\Xi_a}"', from=1-1, to=2-1]
	\arrow["{\Xi_b}", from=1-2, to=2-2]
	\arrow["{=}"{description}, draw=none, from=2-1, to=1-2]
\end{tikzcd}
}

\newquiver{strict natural 2-transformation: 2-cells}{
% https://q.uiver.app/?q=WzAsOCxbMCwxLCJcXEZGIGEiXSxbMSwyLCJcXEZGIGIiXSxbMiwxLCJcXEZGWzJdIGIiXSxbMSwwLCJcXEZGWzJdYSJdLFszLDEsIlxcRkYgYSJdLFs0LDIsIlxcRkYgYiJdLFs1LDEsIlxcRkZbMl0gYiJdLFs0LDAsIlxcRkZbMl0gYSJdLFswLDEsIlxcRkYgZyIsMl0sWzEsMiwiXFxYaV9iIiwyXSxbMCwzLCJcXFhpX2EiXSxbMywyLCJcXEZGWzJdZiIsMCx7ImN1cnZlIjotMn1dLFszLDIsIlxcRkZbMl1nIiwyLHsiY3VydmUiOjJ9XSxbMSwzLCI9IiwwLHsibGFiZWxfcG9zaXRpb24iOjMwLCJvZmZzZXQiOjMsInN0eWxlIjp7ImJvZHkiOnsibmFtZSI6Im5vbmUifSwiaGVhZCI6eyJuYW1lIjoibm9uZSJ9fX1dLFs0LDcsIlxcWGlfYSJdLFs0LDUsIlxcRkYgZiIsMCx7ImN1cnZlIjotMn1dLFs1LDYsIlxcWGlfYiIsMl0sWzcsNiwiXFxGRlsyXSBmIl0sWzQsNSwiXFxGRiBnIiwyLHsiY3VydmUiOjJ9XSxbNSw3LCI9IiwwLHsibGFiZWxfcG9zaXRpb24iOjcwLCJvZmZzZXQiOjIsInN0eWxlIjp7ImJvZHkiOnsibmFtZSI6Im5vbmUifSwiaGVhZCI6eyJuYW1lIjoibm9uZSJ9fX1dLFsyLDQsIj0iLDEseyJzdHlsZSI6eyJib2R5Ijp7Im5hbWUiOiJub25lIn0sImhlYWQiOnsibmFtZSI6Im5vbmUifX19XSxbMTEsMTIsIlxcRkZbMl1cXGFscGhhIiwwLHsic2hvcnRlbiI6eyJzb3VyY2UiOjIwLCJ0YXJnZXQiOjIwfX1dLFsxNSwxOCwiXFxGRlxcYWxwaGEiLDAseyJzaG9ydGVuIjp7InNvdXJjZSI6MjAsInRhcmdldCI6MjB9fV1d
\begin{tikzcd}[ampersand replacement=\&]
	\& {\FF[2]a} \&\&\& {\FF[2] a} \\
	{\FF a} \&\& {\FF[2] b} \& {\FF a} \&\& {\FF[2] b} \\
	\& {\FF b} \&\&\& {\FF b}
	\arrow["{\FF g}"', from=2-1, to=3-2]
	\arrow["{\Xi_b}"', from=3-2, to=2-3]
	\arrow["{\Xi_a}", from=2-1, to=1-2]
	\arrow[""{name=0, anchor=center, inner sep=0}, "{\FF[2]f}", curve={height=-12pt}, from=1-2, to=2-3]
	\arrow[""{name=1, anchor=center, inner sep=0}, "{\FF[2]g}"', curve={height=12pt}, from=1-2, to=2-3]
	\arrow["{=}"{pos=0.3}, shift right=3, draw=none, from=3-2, to=1-2]
	\arrow["{\Xi_a}", from=2-4, to=1-5]
	\arrow[""{name=2, anchor=center, inner sep=0}, "{\FF f}", curve={height=-12pt}, from=2-4, to=3-5]
	\arrow["{\Xi_b}"', from=3-5, to=2-6]
	\arrow["{\FF[2] f}", from=1-5, to=2-6]
	\arrow[""{name=3, anchor=center, inner sep=0}, "{\FF g}"', curve={height=12pt}, from=2-4, to=3-5]
	\arrow["{=}"{pos=0.7}, shift right=2, draw=none, from=3-5, to=1-5]
	\arrow["{=}"{description}, draw=none, from=2-3, to=2-4]
	\arrow["{\FF[2]\alpha}", shorten <=4pt, shorten >=4pt, Rightarrow, from=0, to=1]
	\arrow["\FF\alpha", shorten <=4pt, shorten >=4pt, Rightarrow, from=2, to=3]
\end{tikzcd}}

\newquiver{natural 2-transformation type}{
% https://q.uiver.app/?q=WzAsMixbMCwwLCJcXENDWzFdIl0sWzEsMCwiXFxDQ1syXSJdLFswLDEsIlxcRkZbMV0iLDAseyJjdXJ2ZSI6LTJ9XSxbMCwxLCJcXEZGWzJdIiwyLHsiY3VydmUiOjJ9XSxbMiwzLCJcXFhpIiwwLHsic2hvcnRlbiI6eyJzb3VyY2UiOjIwLCJ0YXJnZXQiOjIwfX1dXQ==
\begin{tikzcd}[ampersand replacement=\&]
	{\CC[1]} \& {\CC[2]}
	\arrow[""{name=0, anchor=center, inner sep=0}, "{\FF[1]}", curve={height=-12pt}, from=1-1, to=1-2]
	\arrow[""{name=1, anchor=center, inner sep=0}, "{\FF[2]}"', curve={height=12pt}, from=1-1, to=1-2]
	\arrow["\Xi", shorten <=3pt, shorten >=3pt, Rightarrow, from=0, to=1]
\end{tikzcd}}

\newquiver{forgetful functor 2-transformation: type}{
% https://q.uiver.app/?q=WzAsMixbMCwwLCJcXHBwQ0FUIl0sWzIsMCwiXFxNdWx0aUNBVCJdLFswLDEsIlxcTXVsdGkoXFxUVCwgLSkiLDAseyJjdXJ2ZSI6LTN9XSxbMCwxLCJcXHNlcS0iLDIseyJjdXJ2ZSI6M31dLFsyLDMsIlxcY2Fycmllci0iLDAseyJzaG9ydGVuIjp7InNvdXJjZSI6MjAsInRhcmdldCI6MjB9fV1d
\begin{tikzcd}[ampersand replacement=\&]
	\ppCAT \&\& \MultiCAT
	\arrow[""{name=0, anchor=center, inner sep=0}, "{\Multi(\TT, -)}", curve={height=-18pt}, from=1-1, to=1-3]
	\arrow[""{name=1, anchor=center, inner sep=0}, "{\seq-}"', curve={height=18pt}, from=1-1, to=1-3]
	\arrow["{\carrier-}", shorten <=5pt, shorten >=5pt, Rightarrow, from=0, to=1]
\end{tikzcd}}

\newquiver{seq alpha type}{
  % https://q.uiver.app/?q=WzAsMixbMCwwLCJcXHBhaXJ7XFxDWzFdfVxccHJvZCJdLFsyLDAsIlxccGFpcntcXENbMl19XFxwcm9kIl0sWzAsMSwiXFxzZXEgSCIsMCx7ImN1cnZlIjotMn1dLFswLDEsIlxcc2VxIEsiLDIseyJjdXJ2ZSI6Mn1dLFsyLDMsIlxcc2VxXFxhbHBoYSIsMCx7InNob3J0ZW4iOnsic291cmNlIjoyMCwidGFyZ2V0IjoyMH19XV0=
\begin{tikzcd}[ampersand replacement=\&]
	{\pair{\C[1]}\prod} \&\& {\pair{\C[2]}\prod}
	\arrow[""{name=0, anchor=center, inner sep=0}, "{\seq H}", curve={height=-12pt}, from=1-1, to=1-3]
	\arrow[""{name=1, anchor=center, inner sep=0}, "{\seq K}"', curve={height=12pt}, from=1-1, to=1-3]
	\arrow["\seq\alpha", shorten <=3pt, shorten >=3pt, Rightarrow, from=0, to=1]
\end{tikzcd}}

\newquiver{pp vertical composition}{
  % https://q.uiver.app/?q=WzAsMixbMCwwLCJcXEMiXSxbMiwwLCJcXENbMl0iXSxbMCwxLCJIIiwwLHsiY3VydmUiOi0zfV0sWzAsMSwiTCIsMix7ImN1cnZlIjozfV0sWzAsMSwiSyIsMV0sWzIsNCwiXFxhbHBoYSIsMCx7InNob3J0ZW4iOnsic291cmNlIjoyMCwidGFyZ2V0IjoyMH19XSxbNCwzLCJcXGJldGEiLDAseyJzaG9ydGVuIjp7InNvdXJjZSI6MjAsInRhcmdldCI6MjB9fV1d
\begin{tikzcd}[ampersand replacement=\&]
	\C \&\& {\C[2]}
	\arrow[""{name=0, anchor=center, inner sep=0}, "H", curve={height=-18pt}, from=1-1, to=1-3]
	\arrow[""{name=1, anchor=center, inner sep=0}, "L"', curve={height=18pt}, from=1-1, to=1-3]
	\arrow[""{name=2, anchor=center, inner sep=0}, "K"{description}, from=1-1, to=1-3]
	\arrow["\alpha", shorten <=2pt, shorten >=2pt, Rightarrow, from=0, to=2]
	\arrow["\beta", shorten <=2pt, shorten >=2pt, Rightarrow, from=2, to=1]
\end{tikzcd}
}

\newquiver{pp horizontal composition}{
  % https://q.uiver.app/?q=WzAsMyxbMCwwLCJcXENbM10iXSxbMSwwLCJcXENbNF0iXSxbMiwwLCJcXEMiXSxbMCwxLCJIXzEiLDAseyJjdXJ2ZSI6LTJ9XSxbMCwxLCJLXzEiLDIseyJjdXJ2ZSI6Mn1dLFsxLDIsIkhfMiIsMCx7ImN1cnZlIjotMn1dLFsxLDIsIktfMiIsMix7ImN1cnZlIjoyfV0sWzUsNiwiXFxiZXRhIiwwLHsic2hvcnRlbiI6eyJzb3VyY2UiOjIwLCJ0YXJnZXQiOjIwfX1dLFszLDQsIlxcYWxwaGEiLDAseyJzaG9ydGVuIjp7InNvdXJjZSI6MjAsInRhcmdldCI6MjB9fV1d
\begin{tikzcd}[ampersand replacement=\&]
	{\C[3]} \& {\C[4]} \& \C
	\arrow[""{name=0, anchor=center, inner sep=0}, "{H_1}", curve={height=-12pt}, from=1-1, to=1-2]
	\arrow[""{name=1, anchor=center, inner sep=0}, "{K_1}"', curve={height=12pt}, from=1-1, to=1-2]
	\arrow[""{name=2, anchor=center, inner sep=0}, "{H_2}", curve={height=-12pt}, from=1-2, to=1-3]
	\arrow[""{name=3, anchor=center, inner sep=0}, "{K_2}"', curve={height=12pt}, from=1-2, to=1-3]
	\arrow["\beta", shorten <=3pt, shorten >=3pt, Rightarrow, from=2, to=3]
	\arrow["\alpha", shorten <=3pt, shorten >=3pt, Rightarrow, from=0, to=1]
\end{tikzcd}
}

\newquiver{carrier 2-naturality}{
% https://q.uiver.app/?q=WzAsMjEsWzAsMSwiXFxNdWx0aShcXFRULFxcQykiXSxbMSwyLCJcXE11bHRpKFxcVFQsXFxDWzJdKSJdLFsyLDEsIlxccGFpcntcXENbMl19XFxwcm9kIl0sWzEsMCwiXFxwYWlyXFxDXFxwcm9kIl0sWzEsNSwiXFxNdWx0aShcXFRULFxcQykiXSxbMiw2LCJcXE11bHRpKFxcVFQsXFxDWzJdKSJdLFszLDUsIlxccGFpcntcXENbMl19XFxwcm9kIl0sWzIsNCwiXFxwYWlyXFxDXFxwcm9kIl0sWzAsMF0sWzAsMl0sWzIsMl0sWzIsMF0sWzIsMywiSFxcY2FycmllclxcQSJdLFsxLDMsIlxccHJvZFxcc2V0e0hcXGNhcnJpZXJcXEF9Il0sWzMsMywiS1xcY2FycmllclxcQSJdLFsxLDRdLFsxLDZdLFszLDZdLFszLDRdLFswLDNdLFswLDVdLFswLDEsIlxcTXVsdGkoXFxUVCwgSykiLDJdLFsxLDIsIlxcY2Fycmllci1fe1xcQ1syXX0iLDJdLFswLDMsIlxcWGlfYSJdLFszLDIsIlxcc2Vxe0h9IiwwLHsiY3VydmUiOi0yfV0sWzMsMiwiXFxzZXEgSyIsMix7ImN1cnZlIjoyfV0sWzEsMywiPSIsMCx7InN0eWxlIjp7ImJvZHkiOnsibmFtZSI6Im5vbmUifSwiaGVhZCI6eyJuYW1lIjoibm9uZSJ9fX1dLFs0LDcsIlxcY2Fycmllci1fe1xcQ30iXSxbNCw1LCJcXE11bHRpKFxcVFQsSCkiLDAseyJjdXJ2ZSI6LTN9XSxbNSw2LCJcXGNhcnJpZXItX3tcXENbMl19IiwyXSxbNyw2LCJcXHNlcXtIfSJdLFs0LDUsIlxcTXVsdGkoXFxUVCxLKSIsMix7ImN1cnZlIjozfV0sWzUsNywiPSIsMCx7ImxhYmVsX3Bvc2l0aW9uIjo3MCwib2Zmc2V0IjoyLCJzdHlsZSI6eyJib2R5Ijp7Im5hbWUiOiJub25lIn0sImhlYWQiOnsibmFtZSI6Im5vbmUifX19XSxbOCw5LCIiLDAseyJvZmZzZXQiOjUsImN1cnZlIjozLCJzdHlsZSI6eyJoZWFkIjp7Im5hbWUiOiJub25lIn19fV0sWzExLDEwLCJcXEEiLDAseyJsYWJlbF9wb3NpdGlvbiI6MTAwLCJjdXJ2ZSI6LTMsInN0eWxlIjp7ImhlYWQiOnsibmFtZSI6Im5vbmUifX19XSxbMTMsMTIsIlxcaXNvbW9ycGhpYyJdLFsxMiwxNCwiXFxhbHBoYV97XFxjYXJyaWVyIFxcQX0iXSxbMTUsMTYsIiIsMCx7Im9mZnNldCI6NSwiY3VydmUiOjUsInN0eWxlIjp7ImhlYWQiOnsibmFtZSI6Im5vbmUifX19XSxbMTgsMTcsIlxcQSIsMCx7ImxhYmVsX3Bvc2l0aW9uIjoxMDAsImN1cnZlIjotNSwic3R5bGUiOnsiaGVhZCI6eyJuYW1lIjoibm9uZSJ9fX1dLFsxMywxOSwiIiwwLHsib2Zmc2V0IjotMSwic2hvcnRlbiI6eyJzb3VyY2UiOjMwLCJ0YXJnZXQiOjMwfSwic3R5bGUiOnsiaGVhZCI6eyJuYW1lIjoibm9uZSJ9fX1dLFsxMywxOSwiIiwwLHsib2Zmc2V0IjoxLCJzaG9ydGVuIjp7InNvdXJjZSI6MzAsInRhcmdldCI6MzB9LCJzdHlsZSI6eyJoZWFkIjp7Im5hbWUiOiJub25lIn19fV0sWzI4LDMxLCJcXE11bHRpKFxcVFQsXFxhbHBoYSkiLDAseyJvZmZzZXQiOjUsInNob3J0ZW4iOnsic291cmNlIjoyMCwidGFyZ2V0IjoyMH19XSxbMjQsMjUsIlxcc2VxXFxhbHBoYSIsMCx7InNob3J0ZW4iOnsic291cmNlIjoyMCwidGFyZ2V0IjoyMH19XSxbMjAsMzcsIiIsMCx7Im9mZnNldCI6MSwic2hvcnRlbiI6eyJzb3VyY2UiOjMwLCJ0YXJnZXQiOjMwfSwibGV2ZWwiOjEsInN0eWxlIjp7ImhlYWQiOnsibmFtZSI6Im5vbmUifX19XSxbMjAsMzcsIiIsMCx7Im9mZnNldCI6LTEsInNob3J0ZW4iOnsic291cmNlIjozMCwidGFyZ2V0IjozMH0sImxldmVsIjoxLCJzdHlsZSI6eyJoZWFkIjp7Im5hbWUiOiJub25lIn19fV1d
\begin{tikzcd}[ampersand replacement=\&]
	{} \& \pair\C\prod \& {} \\
	{\Multi(\TT,\C)} \&\& {\pair{\C[2]}\prod} \\
	{} \& {\Multi(\TT,\C[2])} \& {} \\
	{} \& {\prod\set{H\carrier\A}} \& H\carrier\A \& K\carrier\A \\
	\& {} \& \pair\C\prod \& {} \\
	{} \& {\Multi(\TT,\C)} \&\& {\pair{\C[2]}\prod} \\
	\& {} \& {\Multi(\TT,\C[2])} \& {}
	\arrow["{\Multi(\TT, K)}"', from=2-1, to=3-2]
	\arrow["{\carrier-_{\C[2]}}"', from=3-2, to=2-3]
	\arrow["{\Xi_a}", from=2-1, to=1-2]
	\arrow[""{name=0, anchor=center, inner sep=0}, "{\seq{H}}", curve={height=-12pt}, from=1-2, to=2-3]
	\arrow[""{name=1, anchor=center, inner sep=0}, "{\seq K}"', curve={height=12pt}, from=1-2, to=2-3]
	\arrow["{=}", draw=none, from=3-2, to=1-2]
	\arrow["{\carrier-_{\C}}", from=6-2, to=5-3]
	\arrow[""{name=2, anchor=center, inner sep=0}, "{\Multi(\TT,H)}", curve={height=-18pt}, from=6-2, to=7-3]
	\arrow["{\carrier-_{\C[2]}}"', from=7-3, to=6-4]
	\arrow["{\seq{H}}", from=5-3, to=6-4]
	\arrow[""{name=3, anchor=center, inner sep=0}, "{\Multi(\TT,K)}"', curve={height=18pt}, from=6-2, to=7-3]
	\arrow["{=}"{pos=0.7}, shift right=2, draw=none, from=7-3, to=5-3]
	\arrow[shift right=5, curve={height=18pt}, no head, from=1-1, to=3-1]
	\arrow["\A"{pos=1}, curve={height=-18pt}, no head, from=1-3, to=3-3]
	\arrow["\isomorphic", from=4-2, to=4-3]
	\arrow["{\alpha_{\carrier \A}}", from=4-3, to=4-4]
	\arrow[""{name=4, anchor=center, inner sep=0}, shift right=5, curve={height=30pt}, no head, from=5-2, to=7-2]
	\arrow["\A"{pos=1}, curve={height=-30pt}, no head, from=5-4, to=7-4]
	\arrow[shift left=1, shorten <=10pt, shorten >=10pt, no head, from=4-2, to=4-1]
	\arrow[shift right=1, shorten <=10pt, shorten >=10pt, no head, from=4-2, to=4-1]
	\arrow["{\Multi(\TT,\alpha)}", shift right=5, shorten <=7pt, shorten >=7pt, Rightarrow, from=2, to=3]
	\arrow["\seq\alpha", shorten <=4pt, shorten >=4pt, Rightarrow, from=0, to=1]
	\arrow[shift right=1, shorten <=9pt, shorten >=9pt, no head, from=6-1, to=4]
	\arrow[shift left=1, shorten <=9pt, shorten >=9pt, no head, from=6-1, to=4]
\end{tikzcd}}

\newquiver{mult 2-functoriality vertical composition: a}{
  % https://q.uiver.app/?q=WzAsMyxbMCwwLCJGYSJdLFsyLDAsIkZiIl0sWzMsMCwiR2IiXSxbMCwxLCJGZiIsMCx7ImN1cnZlIjotNH1dLFswLDEsIkZoIiwyLHsiY3VydmUiOjR9XSxbMSwyLCJcXFhpX2IiXSxbMyw0LCJGKFxcYmV0YVxcY2lyY1xcYWxwaGEpIiwwLHsib2Zmc2V0IjozLCJzaG9ydGVuIjp7InNvdXJjZSI6MjAsInRhcmdldCI6MjB9fV1d
\begin{tikzcd}[ampersand replacement=\&]
	Fa \&\& Fb \& Gb
	\arrow[""{name=0, anchor=center, inner sep=0}, "Ff", curve={height=-24pt}, from=1-1, to=1-3]
	\arrow[""{name=1, anchor=center, inner sep=0}, "Fh"', curve={height=24pt}, from=1-1, to=1-3]
	\arrow["{\Xi_b}", from=1-3, to=1-4]
	\arrow["{F(\beta\circ\alpha)}", shift right=3, shorten <=6pt, shorten >=6pt, Rightarrow, from=0, to=1]
\end{tikzcd}
}

\newquiver{mult 2-functoriality vertical composition: b}{
% https://q.uiver.app/?q=WzAsNSxbMCwzLCJGYSJdLFsyLDMsIkZiIl0sWzMsMywiR2IiXSxbMiwxLCJHYSJdLFsyLDAsIkZhIl0sWzAsMSwiRmYiLDAseyJjdXJ2ZSI6LTR9XSxbMCwxLCJGaCIsMix7ImN1cnZlIjo0fV0sWzEsMiwiXFxYaV9iIl0sWzAsMywiXFxYaV9hIiwxLHsiY3VydmUiOi0yfV0sWzMsMiwiR2YiLDEseyJjdXJ2ZSI6LTJ9XSxbMCw0LCJGZiIsMSx7ImN1cnZlIjotM31dLFs0LDIsIlxcWGlfYiIsMSx7ImN1cnZlIjotM31dLFs0LDMsIj0iLDEseyJzdHlsZSI6eyJib2R5Ijp7Im5hbWUiOiJub25lIn0sImhlYWQiOnsibmFtZSI6Im5vbmUifX19XSxbMywxLCI9IiwxLHsibGV2ZWwiOjIsInN0eWxlIjp7ImJvZHkiOnsibmFtZSI6Im5vbmUifSwiaGVhZCI6eyJuYW1lIjoibm9uZSJ9fX1dLFs1LDYsIkYoXFxiZXRhXFxjaXJjXFxhbHBoYSkiLDAseyJvZmZzZXQiOjMsInNob3J0ZW4iOnsic291cmNlIjoyMCwidGFyZ2V0IjoyMH19XV0=
\begin{tikzcd}[ampersand replacement=\&]
	\&\& Fa \\
	\&\& Ga \\
	\\
	Fa \&\& Fb \& Gb
	\arrow[""{name=0, anchor=center, inner sep=0}, "Ff", curve={height=-24pt}, from=4-1, to=4-3]
	\arrow[""{name=1, anchor=center, inner sep=0}, "Fh"', curve={height=24pt}, from=4-1, to=4-3]
	\arrow["{\Xi_b}", from=4-3, to=4-4]
	\arrow["{\Xi_a}"{description}, curve={height=-12pt}, from=4-1, to=2-3]
	\arrow["Gf"{description}, curve={height=-12pt}, from=2-3, to=4-4]
	\arrow["Ff"{description}, curve={height=-18pt}, from=4-1, to=1-3]
	\arrow["{\Xi_b}"{description}, curve={height=-18pt}, from=1-3, to=4-4]
	\arrow["{=}"{description}, draw=none, from=1-3, to=2-3]
	\arrow["{=}"{description}, draw=none, from=2-3, to=4-3]
	\arrow["{F(\beta\circ\alpha)}", shift right=3, shorten <=6pt, shorten >=6pt, Rightarrow, from=0, to=1]
\end{tikzcd}
}

\newquiver{mult 2-functoriality vertical composition: c}{
% https://q.uiver.app/?q=WzAsNSxbMCwzLCJGYSJdLFsxLDMsIkZiIl0sWzIsMywiR2IiXSxbMSwxLCJHYSJdLFsxLDAsIkZhIl0sWzAsMSwiRmgiLDJdLFsxLDIsIlxcWGlfYiIsMl0sWzAsMywiXFxYaV9hIiwxLHsiY3VydmUiOi0yfV0sWzMsMiwiR2YiLDEseyJjdXJ2ZSI6LTJ9XSxbMCw0LCJGZiIsMSx7ImN1cnZlIjotM31dLFs0LDIsIlxcWGlfYiIsMSx7ImN1cnZlIjotM31dLFs0LDMsIj0iLDEseyJzdHlsZSI6eyJib2R5Ijp7Im5hbWUiOiJub25lIn0sImhlYWQiOnsibmFtZSI6Im5vbmUifX19XSxbMywyLCJHZyIsMix7ImN1cnZlIjozfV0sWzgsMTIsIkcoXFxiZXRhIFxcY2lyYyBcXGFscGhhKSIsMCx7InNob3J0ZW4iOnsic291cmNlIjoyMCwidGFyZ2V0IjoyMH19XV0=
\begin{tikzcd}[ampersand replacement=\&]
	\& Fa \\
	\& Ga \\
	\\
	Fa \& Fb \& Gb
	\arrow["Fh"', from=4-1, to=4-2]
	\arrow["{\Xi_b}"', from=4-2, to=4-3]
	\arrow["{\Xi_a}"{description}, curve={height=-12pt}, from=4-1, to=2-2]
	\arrow[""{name=0, anchor=center, inner sep=0}, "Gf"{description}, curve={height=-12pt}, from=2-2, to=4-3]
	\arrow["Ff"{description}, curve={height=-18pt}, from=4-1, to=1-2]
	\arrow["{\Xi_b}"{description}, curve={height=-18pt}, from=1-2, to=4-3]
	\arrow["{=}"{description}, draw=none, from=1-2, to=2-2]
	\arrow[""{name=1, anchor=center, inner sep=0}, "Gg"', curve={height=18pt}, from=2-2, to=4-3]
	\arrow["{G(\beta \circ \alpha)}", shorten <=5pt, shorten >=5pt, Rightarrow, from=0, to=1]
\end{tikzcd}}

\newquiver{mult 2-functoriality vertical composition: d}{
% https://q.uiver.app/?q=WzAsNSxbMCwzLCJGYSJdLFsxLDMsIkZiIl0sWzIsMywiR2IiXSxbMSwxLCJHYSJdLFsxLDAsIkZhIl0sWzAsMSwiRmgiLDJdLFsxLDIsIlxcWGlfYiIsMl0sWzAsMywiXFxYaV9hIiwxLHsiY3VydmUiOi0yfV0sWzMsMiwiR2YiLDAseyJsYWJlbF9wb3NpdGlvbiI6MTAsImN1cnZlIjotMn1dLFswLDQsIkZmIiwxLHsiY3VydmUiOi0zfV0sWzQsMiwiXFxYaV9iIiwxLHsiY3VydmUiOi0zfV0sWzQsMywiPSIsMSx7InN0eWxlIjp7ImJvZHkiOnsibmFtZSI6Im5vbmUifSwiaGVhZCI6eyJuYW1lIjoibm9uZSJ9fX1dLFszLDIsIkdnIiwyLHsiY3VydmUiOjN9XSxbMywyXSxbOCwxMywiR1xcYWxwaGEiLDAseyJzaG9ydGVuIjp7InNvdXJjZSI6MjAsInRhcmdldCI6MjB9fV0sWzEzLDEyLCJHXFxiZXRhIiwwLHsic2hvcnRlbiI6eyJzb3VyY2UiOjIwLCJ0YXJnZXQiOjIwfX1dXQ==
\begin{tikzcd}[ampersand replacement=\&]
	\& Fa \\
	\& Ga \\
	\\
	Fa \& Fb \& Gb
	\arrow["Fh"', from=4-1, to=4-2]
	\arrow["{\Xi_b}"', from=4-2, to=4-3]
	\arrow["{\Xi_a}"{description}, curve={height=-12pt}, from=4-1, to=2-2]
	\arrow[""{name=0, anchor=center, inner sep=0}, "Gf"{pos=0.1}, curve={height=-12pt}, from=2-2, to=4-3]
	\arrow["Ff"{description}, curve={height=-18pt}, from=4-1, to=1-2]
	\arrow["{\Xi_b}"{description}, curve={height=-18pt}, from=1-2, to=4-3]
	\arrow["{=}"{description}, draw=none, from=1-2, to=2-2]
	\arrow[""{name=1, anchor=center, inner sep=0}, "Gg"', curve={height=18pt}, from=2-2, to=4-3]
	\arrow[""{name=2, anchor=center, inner sep=0}, from=2-2, to=4-3]
	\arrow["G\alpha", shorten <=2pt, shorten >=2pt, Rightarrow, from=0, to=2]
	\arrow["G\beta", shorten <=3pt, shorten >=3pt, Rightarrow, from=2, to=1]
\end{tikzcd}
}

\newquiver{mult 2-functoriality vertical composition: e}{
  % https://q.uiver.app/?q=WzAsNSxbMCwzLCJGYSJdLFsxLDMsIkZiIl0sWzIsMywiR2IiXSxbMSwxLCJHYSJdLFsxLDAsIkZhIl0sWzAsMSwiRmciLDIseyJjdXJ2ZSI6Mn1dLFsxLDIsIlxcWGlfYiIsMl0sWzAsMywiXFxYaV9hIiwxLHsiY3VydmUiOi0yfV0sWzMsMiwiR2YiLDAseyJsYWJlbF9wb3NpdGlvbiI6MTAsImN1cnZlIjotMn1dLFswLDQsIkZmIiwxLHsiY3VydmUiOi0zfV0sWzQsMiwiXFxYaV9iIiwxLHsiY3VydmUiOi0zfV0sWzQsMywiPSIsMSx7InN0eWxlIjp7ImJvZHkiOnsibmFtZSI6Im5vbmUifSwiaGVhZCI6eyJuYW1lIjoibm9uZSJ9fX1dLFszLDJdLFswLDEsIkZoIiwwLHsiY3VydmUiOi0yfV0sWzgsMTIsIkdcXGFscGhhIiwwLHsic2hvcnRlbiI6eyJzb3VyY2UiOjIwLCJ0YXJnZXQiOjIwfX1dLFsxMyw1LCJGXFxiZXRhIiwwLHsic2hvcnRlbiI6eyJzb3VyY2UiOjIwLCJ0YXJnZXQiOjIwfX1dLFszLDEzLCI9IiwwLHsic2hvcnRlbiI6eyJ0YXJnZXQiOjIwfSwic3R5bGUiOnsiYm9keSI6eyJuYW1lIjoibm9uZSJ9LCJoZWFkIjp7Im5hbWUiOiJub25lIn19fV1d
\begin{tikzcd}[ampersand replacement=\&]
	\& Fa \\
	\& Ga \\
	\\
	Fa \& Fb \& Gb
	\arrow[""{name=0, anchor=center, inner sep=0}, "Fg"', curve={height=12pt}, from=4-1, to=4-2]
	\arrow["{\Xi_b}"', from=4-2, to=4-3]
	\arrow["{\Xi_a}"{description}, curve={height=-12pt}, from=4-1, to=2-2]
	\arrow[""{name=1, anchor=center, inner sep=0}, "Gf"{pos=0.1}, curve={height=-12pt}, from=2-2, to=4-3]
	\arrow["Ff"{description}, curve={height=-18pt}, from=4-1, to=1-2]
	\arrow["{\Xi_b}"{description}, curve={height=-18pt}, from=1-2, to=4-3]
	\arrow["{=}"{description}, draw=none, from=1-2, to=2-2]
	\arrow[""{name=2, anchor=center, inner sep=0}, from=2-2, to=4-3]
	\arrow[""{name=3, anchor=center, inner sep=0}, "Fh", curve={height=-12pt}, from=4-1, to=4-2]
	\arrow["G\alpha", shorten <=2pt, shorten >=2pt, Rightarrow, from=1, to=2]
	\arrow["F\beta", shorten <=3pt, shorten >=3pt, Rightarrow, from=3, to=0]
	\arrow["{=}", draw=none, from=2-2, to=3]
\end{tikzcd}
}

\newquiver{mult 2-functoriality vertical composition: f}{
% https://q.uiver.app/?q=WzAsNSxbMCwzLCJGYSJdLFsyLDMsIkZiIl0sWzMsMywiR2IiXSxbMiwxLCJHYSJdLFsyLDAsIkZhIl0sWzAsMSwiRmgiLDIseyJjdXJ2ZSI6Mn1dLFsxLDIsIlxcWGlfYiIsMl0sWzAsMywiXFxYaV9hIiwxLHsiY3VydmUiOi0yfV0sWzMsMiwiR2YiLDAseyJsYWJlbF9wb3NpdGlvbiI6MTAsImN1cnZlIjotMn1dLFswLDQsIkZmIiwxLHsiY3VydmUiOi0zfV0sWzQsMiwiXFxYaV9iIiwxLHsiY3VydmUiOi0zfV0sWzQsMywiPSIsMSx7InN0eWxlIjp7ImJvZHkiOnsibmFtZSI6Im5vbmUifSwiaGVhZCI6eyJuYW1lIjoibm9uZSJ9fX1dLFswLDEsIiIsMCx7ImxhYmVsX3Bvc2l0aW9uIjozMH1dLFswLDEsIkZmIiwwLHsiY3VydmUiOi0yfV0sWzMsMSwiPSIsMSx7InN0eWxlIjp7ImJvZHkiOnsibmFtZSI6Im5vbmUifSwiaGVhZCI6eyJuYW1lIjoibm9uZSJ9fX1dLFsxMiw1LCJGXFxiZXRhIiwwLHsic2hvcnRlbiI6eyJzb3VyY2UiOjIwLCJ0YXJnZXQiOjIwfX1dLFsxMywxMiwiRlxcYWxwaGEiLDAseyJzaG9ydGVuIjp7InNvdXJjZSI6MjAsInRhcmdldCI6MjB9fV1d
\begin{tikzcd}[ampersand replacement=\&]
	\&\& Fa \\
	\&\& Ga \\
	\\
	Fa \&\& Fb \& Gb
	\arrow[""{name=0, anchor=center, inner sep=0}, "Fh"', curve={height=12pt}, from=4-1, to=4-3]
	\arrow["{\Xi_b}"', from=4-3, to=4-4]
	\arrow["{\Xi_a}"{description}, curve={height=-12pt}, from=4-1, to=2-3]
	\arrow["Gf"{pos=0.1}, curve={height=-12pt}, from=2-3, to=4-4]
	\arrow["Ff"{description}, curve={height=-18pt}, from=4-1, to=1-3]
	\arrow["{\Xi_b}"{description}, curve={height=-18pt}, from=1-3, to=4-4]
	\arrow["{=}"{description}, draw=none, from=1-3, to=2-3]
	\arrow[""{name=1, anchor=center, inner sep=0}, from=4-1, to=4-3]
	\arrow[""{name=2, anchor=center, inner sep=0}, "Ff", curve={height=-12pt}, from=4-1, to=4-3]
	\arrow["{=}"{description}, draw=none, from=2-3, to=4-3]
	\arrow["F\beta", shorten <=2pt, shorten >=2pt, Rightarrow, from=1, to=0]
	\arrow["F\alpha", shorten <=2pt, shorten >=2pt, Rightarrow, from=2, to=1]
\end{tikzcd}
}

\newquiver{mult 2-functoriality vertical composition: g}{
% https://q.uiver.app/?q=WzAsMyxbMCwwLCJGYSJdLFsyLDAsIkZiIl0sWzMsMCwiR2IiXSxbMCwxLCJGaCIsMix7ImN1cnZlIjoyfV0sWzEsMiwiXFxYaV9iIiwyXSxbMCwxLCIiLDAseyJsYWJlbF9wb3NpdGlvbiI6MzB9XSxbMCwxLCJGZiIsMCx7ImN1cnZlIjotMn1dLFs1LDMsIkZcXGJldGEiLDAseyJzaG9ydGVuIjp7InNvdXJjZSI6MjAsInRhcmdldCI6MjB9fV0sWzYsNSwiRlxcYWxwaGEiLDAseyJzaG9ydGVuIjp7InNvdXJjZSI6MjAsInRhcmdldCI6MjB9fV1d
\begin{tikzcd}[ampersand replacement=\&]
	Fa \&\& Fb \& Gb
	\arrow[""{name=0, anchor=center, inner sep=0}, "Fh"', curve={height=12pt}, from=1-1, to=1-3]
	\arrow["{\Xi_b}"', from=1-3, to=1-4]
	\arrow[""{name=1, anchor=center, inner sep=0}, from=1-1, to=1-3]
	\arrow[""{name=2, anchor=center, inner sep=0}, "Ff", curve={height=-12pt}, from=1-1, to=1-3]
	\arrow["F\beta", shorten <=2pt, shorten >=2pt, Rightarrow, from=1, to=0]
	\arrow["F\alpha", shorten <=2pt, shorten >=2pt, Rightarrow, from=2, to=1]
\end{tikzcd}
}

\newquiver{mult 2-functoriality horizontal composition: a}{
  % https://q.uiver.app/?q=WzAsNSxbMCwxLCJGYSJdLFszLDEsIkdjIl0sWzIsMSwiRmMiXSxbMSwwLCJGYiJdLFsxLDEsIkZiIl0sWzIsMSwiXFxYaV9jIl0sWzAsMywiRmZfMSIsMCx7ImN1cnZlIjotMn1dLFszLDIsIkZnXzEiLDAseyJjdXJ2ZSI6LTJ9XSxbMCw0LCJGZ18xIiwyXSxbNCwyLCJGZ18yIiwyXSxbMyw0LCJGKFxcYmV0YVxcYnVsbGV0XFxhbHBoYSkiLDAseyJzaG9ydGVuIjp7InNvdXJjZSI6MjAsInRhcmdldCI6MjB9LCJsZXZlbCI6Mn1dXQ==
\begin{tikzcd}[ampersand replacement=\&]
	\& Fb \\
	Fa \& Fb \& Fc \& Gc
	\arrow["{\Xi_c}", from=2-3, to=2-4]
	\arrow["{Ff_1}", curve={height=-12pt}, from=2-1, to=1-2]
	\arrow["{Fg_1}", curve={height=-12pt}, from=1-2, to=2-3]
	\arrow["{Fg_1}"', from=2-1, to=2-2]
	\arrow["{Fg_2}"', from=2-2, to=2-3]
	\arrow["{F(\beta\bullet\alpha)}", shorten <=2pt, shorten >=2pt, Rightarrow, from=1-2, to=2-2]
\end{tikzcd}}

\newquiver{mult 2-functoriality horizontal composition: b}{
  % https://q.uiver.app/?q=WzAsOSxbMSwzLCJGYSJdLFs0LDMsIkdjIl0sWzMsMywiRmMiXSxbMiwyLCJGYiJdLFsyLDMsIkZiIl0sWzAsMCwiRmIiXSxbMSwxLCJHYSJdLFszLDEsIkdiIl0sWzQsMCwiRmMiXSxbMiwxLCJcXFhpX2MiXSxbMCwzLCJGZl8xIiwxLHsiY3VydmUiOi0yfV0sWzMsMiwiRmdfMSIsMSx7ImN1cnZlIjotMn1dLFswLDQsIkZnXzEiLDJdLFs0LDIsIkZnXzIiLDJdLFszLDQsIkYoXFxiZXRhXFxidWxsZXRcXGFscGhhKSIsMCx7InNob3J0ZW4iOnsic291cmNlIjoyMCwidGFyZ2V0IjoyMH0sImxldmVsIjoyfV0sWzAsNSwiRmZfMSIsMV0sWzUsOCwiRmdfMSIsMV0sWzgsMSwiXFxYaV9jIiwxXSxbMCw2LCJcXFhpX2EiLDFdLFs2LDcsIkdmXzEiLDFdLFs3LDEsIkdmXzIiLDFdLFs3LDExLCI9IiwxLHsic2hvcnRlbiI6eyJ0YXJnZXQiOjIwfSwic3R5bGUiOnsiYm9keSI6eyJuYW1lIjoibm9uZSJ9LCJoZWFkIjp7Im5hbWUiOiJub25lIn19fV0sWzE2LDE5LCI9IiwxLHsic2hvcnRlbiI6eyJzb3VyY2UiOjIwLCJ0YXJnZXQiOjIwfSwic3R5bGUiOnsiYm9keSI6eyJuYW1lIjoibm9uZSJ9LCJoZWFkIjp7Im5hbWUiOiJub25lIn19fV1d
\begin{tikzcd}[ampersand replacement=\&]
	Fb \&\&\&\& Fc \\
	\& Ga \&\& Gb \\
	\&\& Fb \\
	\& Fa \& Fb \& Fc \& Gc
	\arrow["{\Xi_c}", from=4-4, to=4-5]
	\arrow["{Ff_1}"{description}, curve={height=-12pt}, from=4-2, to=3-3]
	\arrow[""{name=0, anchor=center, inner sep=0}, "{Fg_1}"{description}, curve={height=-12pt}, from=3-3, to=4-4]
	\arrow["{Fg_1}"', from=4-2, to=4-3]
	\arrow["{Fg_2}"', from=4-3, to=4-4]
	\arrow["{F(\beta\bullet\alpha)}", shorten <=2pt, shorten >=2pt, Rightarrow, from=3-3, to=4-3]
	\arrow["{Ff_1}"{description}, from=4-2, to=1-1]
	\arrow[""{name=1, anchor=center, inner sep=0}, "{Fg_1}"{description}, from=1-1, to=1-5]
	\arrow["{\Xi_c}"{description}, from=1-5, to=4-5]
	\arrow["{\Xi_a}"{description}, from=4-2, to=2-2]
	\arrow[""{name=2, anchor=center, inner sep=0}, "{Gf_1}"{description}, from=2-2, to=2-4]
	\arrow["{Gf_2}"{description}, from=2-4, to=4-5]
	\arrow["{=}"{description}, draw=none, from=2-4, to=0]
	\arrow["{=}"{description}, draw=none, from=1, to=2]
\end{tikzcd}
}
\newquiver{mult 2-functoriality horizontal composition: c}{
% https://q.uiver.app/?q=WzAsOSxbMSwzLCJGYSJdLFs0LDMsIkdjIl0sWzMsMywiRmMiXSxbMiwzLCJGYiJdLFswLDAsIkZiIl0sWzEsMSwiR2EiXSxbMywxLCJHYiJdLFs0LDAsIkZjIl0sWzIsMiwiR2IiXSxbMiwxLCJcXFhpX2MiLDJdLFswLDMsIkZnXzEiLDJdLFszLDIsIkZnXzIiLDJdLFswLDQsIkZmXzEiLDFdLFs0LDcsIkZnXzEiLDFdLFs3LDEsIlxcWGlfYyIsMV0sWzAsNSwiXFxYaV9hIiwxXSxbNSw2LCJHZl8xIiwxXSxbNiwxLCJHZl8yIiwxXSxbNSw4LCJHZ18xIiwxXSxbOCwxLCJHZ18yIiwxLHsiY3VydmUiOjF9XSxbNiw4LCIiLDEseyJzdHlsZSI6eyJib2R5Ijp7Im5hbWUiOiJub25lIn0sImhlYWQiOnsibmFtZSI6Im5vbmUifX19XSxbNiw4LCJHKFxcYmV0YVxcYnVsbGV0XFxhbHBoYSkiLDAseyJsZXZlbCI6Mn1dLFs4LDMsIj0iLDAseyJzdHlsZSI6eyJib2R5Ijp7Im5hbWUiOiJub25lIn0sImhlYWQiOnsibmFtZSI6Im5vbmUifX19XSxbMTMsMTYsIj0iLDEseyJzaG9ydGVuIjp7InNvdXJjZSI6MjAsInRhcmdldCI6MjB9LCJzdHlsZSI6eyJib2R5Ijp7Im5hbWUiOiJub25lIn0sImhlYWQiOnsibmFtZSI6Im5vbmUifX19XV0=
\begin{tikzcd}[ampersand replacement=\&]
	Fb \&\&\&\& Fc \\
	\& Ga \&\& Gb \\
	\&\& Gb \\
	\& Fa \& Fb \& Fc \& Gc
	\arrow["{\Xi_c}"', from=4-4, to=4-5]
	\arrow["{Fg_1}"', from=4-2, to=4-3]
	\arrow["{Fg_2}"', from=4-3, to=4-4]
	\arrow["{Ff_1}"{description}, from=4-2, to=1-1]
	\arrow[""{name=0, anchor=center, inner sep=0}, "{Fg_1}"{description}, from=1-1, to=1-5]
	\arrow["{\Xi_c}"{description}, from=1-5, to=4-5]
	\arrow["{\Xi_a}"{description}, from=4-2, to=2-2]
	\arrow[""{name=1, anchor=center, inner sep=0}, "{Gf_1}"{description}, from=2-2, to=2-4]
	\arrow["{Gf_2}"{description}, from=2-4, to=4-5]
	\arrow["{Gg_1}"{description}, from=2-2, to=3-3]
	\arrow["{Gg_2}"{description}, curve={height=6pt}, from=3-3, to=4-5]
	\arrow[draw=none, from=2-4, to=3-3]
	\arrow["{G(\beta\bullet\alpha)}", Rightarrow, from=2-4, to=3-3]
	\arrow["{=}", draw=none, from=3-3, to=4-3]
	\arrow["{=}"{description}, draw=none, from=0, to=1]
\end{tikzcd}
}
\newquiver{mult 2-functoriality horizontal composition: d}{
% https://q.uiver.app/?q=WzAsOCxbMSwzLCJGYSJdLFs0LDMsIkdjIl0sWzMsMywiRmMiXSxbMiwzLCJGYiJdLFswLDAsIkZiIl0sWzEsMSwiR2EiXSxbMywxLCJHYiJdLFs0LDAsIkZjIl0sWzIsMSwiXFxYaV9jIiwyXSxbMCwzLCJGZ18xIiwyXSxbMywyLCJGZ18yIiwyXSxbMCw0LCJGZl8xIiwxXSxbNCw3LCJGZ18xIiwxXSxbNywxLCJcXFhpX2MiLDFdLFswLDUsIlxcWGlfYSIsMV0sWzUsNiwiR2ZfMSIsMSx7ImN1cnZlIjotMX1dLFs2LDEsIkdmXzIiLDEseyJsYWJlbF9wb3NpdGlvbiI6NDAsImN1cnZlIjotMX1dLFs1LDYsIkdnXzEiLDEseyJjdXJ2ZSI6Mn1dLFs2LDEsIkdnXzIiLDEseyJjdXJ2ZSI6Mn1dLFs2LDMsIj0iLDAseyJzdHlsZSI6eyJib2R5Ijp7Im5hbWUiOiJub25lIn0sImhlYWQiOnsibmFtZSI6Im5vbmUifX19XSxbMTIsMTUsIj0iLDEseyJzaG9ydGVuIjp7InNvdXJjZSI6MjAsInRhcmdldCI6MjB9LCJzdHlsZSI6eyJib2R5Ijp7Im5hbWUiOiJub25lIn0sImhlYWQiOnsibmFtZSI6Im5vbmUifX19XSxbMTUsMTcsIkdcXGFscGhhIiwwLHsic2hvcnRlbiI6eyJzb3VyY2UiOjIwLCJ0YXJnZXQiOjIwfX1dLFsxNiwxOCwiR1xcYmV0YSIsMCx7InNob3J0ZW4iOnsic291cmNlIjoyMCwidGFyZ2V0IjoyMH19XV0=
\begin{tikzcd}[ampersand replacement=\&]
	Fb \&\&\&\& Fc \\
	\& Ga \&\& Gb \\
	\\
	\& Fa \& Fb \& Fc \& Gc
	\arrow["{\Xi_c}"', from=4-4, to=4-5]
	\arrow["{Fg_1}"', from=4-2, to=4-3]
	\arrow["{Fg_2}"', from=4-3, to=4-4]
	\arrow["{Ff_1}"{description}, from=4-2, to=1-1]
	\arrow[""{name=0, anchor=center, inner sep=0}, "{Fg_1}"{description}, from=1-1, to=1-5]
	\arrow["{\Xi_c}"{description}, from=1-5, to=4-5]
	\arrow["{\Xi_a}"{description}, from=4-2, to=2-2]
	\arrow[""{name=1, anchor=center, inner sep=0}, "{Gf_1}"{description}, curve={height=-6pt}, from=2-2, to=2-4]
	\arrow[""{name=2, anchor=center, inner sep=0}, "{Gf_2}"{description, pos=0.4}, curve={height=-6pt}, from=2-4, to=4-5]
	\arrow[""{name=3, anchor=center, inner sep=0}, "{Gg_1}"{description}, curve={height=12pt}, from=2-2, to=2-4]
	\arrow[""{name=4, anchor=center, inner sep=0}, "{Gg_2}"{description}, curve={height=12pt}, from=2-4, to=4-5]
	\arrow["{=}", draw=none, from=2-4, to=4-3]
	\arrow["{=}"{description}, draw=none, from=0, to=1]
	\arrow["G\alpha", shorten <=2pt, shorten >=2pt, Rightarrow, from=1, to=3]
	\arrow["G\beta", shorten <=3pt, shorten >=3pt, Rightarrow, from=2, to=4]
\end{tikzcd}
}
\newquiver{mult 2-functoriality horizontal composition: e}{
% https://q.uiver.app/?q=WzAsOCxbMSwzLCJGYSJdLFs0LDMsIkdjIl0sWzMsMywiRmMiXSxbMiwzLCJGYiJdLFswLDAsIkZiIl0sWzEsMSwiR2EiXSxbMywxLCJHYiJdLFs0LDAsIkZjIl0sWzIsMSwiXFxYaV9jIiwyXSxbMCwzLCJGZ18xIiwyXSxbMywyLCJGZ18yIiwyXSxbMCw0LCJGZl8xIiwxXSxbNCw3LCJGZ18xIiwxXSxbNywxLCJcXFhpX2MiLDFdLFswLDUsIlxcWGlfYSIsMV0sWzUsNiwiR2ZfMSIsMSx7ImN1cnZlIjotMX1dLFs2LDEsIkdmXzIiLDEseyJsYWJlbF9wb3NpdGlvbiI6NDAsImN1cnZlIjotMX1dLFs2LDEsIkdnXzIiLDEseyJjdXJ2ZSI6Mn1dLFs1LDYsIkdnXzEiLDEseyJjdXJ2ZSI6Mn1dLFszLDYsIlxcWGlfYiJdLFsxNiwxNywiR1xcYmV0YSIsMCx7InNob3J0ZW4iOnsic291cmNlIjoyMCwidGFyZ2V0IjoyMH19XSxbMTUsMTgsIkdcXGFscGhhIiwwLHsic2hvcnRlbiI6eyJzb3VyY2UiOjIwLCJ0YXJnZXQiOjIwfX1dLFsxMiwxNSwiPSIsMSx7InNob3J0ZW4iOnsic291cmNlIjoyMCwidGFyZ2V0IjoyMH0sInN0eWxlIjp7ImJvZHkiOnsibmFtZSI6Im5vbmUifSwiaGVhZCI6eyJuYW1lIjoibm9uZSJ9fX1dLFsxNCwxOSwiPSIsMSx7InNob3J0ZW4iOnsic291cmNlIjoyMCwidGFyZ2V0IjoyMH0sInN0eWxlIjp7ImJvZHkiOnsibmFtZSI6Im5vbmUifSwiaGVhZCI6eyJuYW1lIjoibm9uZSJ9fX1dLFsxOSwyLCI9IiwwLHsic2hvcnRlbiI6eyJzb3VyY2UiOjIwfSwic3R5bGUiOnsiYm9keSI6eyJuYW1lIjoibm9uZSJ9LCJoZWFkIjp7Im5hbWUiOiJub25lIn19fV1d
\begin{tikzcd}[ampersand replacement=\&]
	Fb \&\&\&\& Fc \\
	\& Ga \&\& Gb \\
	\\
	\& Fa \& Fb \& Fc \& Gc
	\arrow["{\Xi_c}"', from=4-4, to=4-5]
	\arrow["{Fg_1}"', from=4-2, to=4-3]
	\arrow["{Fg_2}"', from=4-3, to=4-4]
	\arrow["{Ff_1}"{description}, from=4-2, to=1-1]
	\arrow[""{name=0, anchor=center, inner sep=0}, "{Fg_1}"{description}, from=1-1, to=1-5]
	\arrow["{\Xi_c}"{description}, from=1-5, to=4-5]
	\arrow[""{name=1, anchor=center, inner sep=0}, "{\Xi_a}"{description}, from=4-2, to=2-2]
	\arrow[""{name=2, anchor=center, inner sep=0}, "{Gf_1}"{description}, curve={height=-6pt}, from=2-2, to=2-4]
	\arrow[""{name=3, anchor=center, inner sep=0}, "{Gf_2}"{description, pos=0.4}, curve={height=-6pt}, from=2-4, to=4-5]
	\arrow[""{name=4, anchor=center, inner sep=0}, "{Gg_2}"{description}, curve={height=12pt}, from=2-4, to=4-5]
	\arrow[""{name=5, anchor=center, inner sep=0}, "{Gg_1}"{description}, curve={height=12pt}, from=2-2, to=2-4]
	\arrow[""{name=6, anchor=center, inner sep=0}, "{\Xi_b}", from=4-3, to=2-4]
	\arrow["G\beta", shorten <=3pt, shorten >=3pt, Rightarrow, from=3, to=4]
	\arrow["G\alpha", shorten <=2pt, shorten >=2pt, Rightarrow, from=2, to=5]
	\arrow["{=}"{description}, draw=none, from=0, to=2]
	\arrow["{=}"{description}, draw=none, from=1, to=6]
	\arrow["{=}", draw=none, from=6, to=4-4]
\end{tikzcd}
}
\newquiver{mult 2-functoriality horizontal composition: f}{
% https://q.uiver.app/?q=WzAsOCxbMSwzLCJGYSJdLFs0LDMsIkdjIl0sWzMsMywiRmMiXSxbMiwzLCJGYiJdLFswLDAsIkZiIl0sWzEsMSwiR2EiXSxbMywxLCJHYiJdLFs0LDAsIkZjIl0sWzIsMSwiXFxYaV9jIiwyXSxbMCwzLCJGZ18xIiwyLHsiY3VydmUiOjJ9XSxbMywyLCJGZ18yIiwyXSxbMCw0LCJGZl8xIiwxXSxbNCw3LCJGZ18xIiwxXSxbNywxLCJcXFhpX2MiLDFdLFswLDUsIlxcWGlfYSIsMV0sWzUsNiwiR2ZfMSIsMV0sWzYsMSwiR2ZfMiIsMSx7ImxhYmVsX3Bvc2l0aW9uIjo0MCwiY3VydmUiOi0xfV0sWzYsMSwiR2dfMiIsMSx7ImN1cnZlIjoyfV0sWzMsNiwiXFxYaV9iIl0sWzAsMywiRmZfMSIsMCx7ImN1cnZlIjotMn1dLFsxNiwxNywiR1xcYmV0YSIsMCx7InNob3J0ZW4iOnsic291cmNlIjoyMCwidGFyZ2V0IjoyMH19XSxbMTIsMTUsIj0iLDEseyJzaG9ydGVuIjp7InNvdXJjZSI6MjAsInRhcmdldCI6MjB9LCJzdHlsZSI6eyJib2R5Ijp7Im5hbWUiOiJub25lIn0sImhlYWQiOnsibmFtZSI6Im5vbmUifX19XSxbMTQsMTgsIj0iLDEseyJzaG9ydGVuIjp7InNvdXJjZSI6MjAsInRhcmdldCI6MjB9LCJzdHlsZSI6eyJib2R5Ijp7Im5hbWUiOiJub25lIn0sImhlYWQiOnsibmFtZSI6Im5vbmUifX19XSxbMTgsMiwiPSIsMCx7InNob3J0ZW4iOnsic291cmNlIjoyMH0sInN0eWxlIjp7ImJvZHkiOnsibmFtZSI6Im5vbmUifSwiaGVhZCI6eyJuYW1lIjoibm9uZSJ9fX1dLFsxOSw5LCJGXFxhbHBoYSIsMCx7Im9mZnNldCI6Miwic2hvcnRlbiI6eyJzb3VyY2UiOjIwLCJ0YXJnZXQiOjIwfX1dXQ==
\begin{tikzcd}[ampersand replacement=\&]
	Fb \&\&\&\& Fc \\
	\& Ga \&\& Gb \\
	\\
	\& Fa \& Fb \& Fc \& Gc
	\arrow["{\Xi_c}"', from=4-4, to=4-5]
	\arrow[""{name=0, anchor=center, inner sep=0}, "{Fg_1}"', curve={height=12pt}, from=4-2, to=4-3]
	\arrow["{Fg_2}"', from=4-3, to=4-4]
	\arrow["{Ff_1}"{description}, from=4-2, to=1-1]
	\arrow[""{name=1, anchor=center, inner sep=0}, "{Fg_1}"{description}, from=1-1, to=1-5]
	\arrow["{\Xi_c}"{description}, from=1-5, to=4-5]
	\arrow[""{name=2, anchor=center, inner sep=0}, "{\Xi_a}"{description}, from=4-2, to=2-2]
	\arrow[""{name=3, anchor=center, inner sep=0}, "{Gf_1}"{description}, from=2-2, to=2-4]
	\arrow[""{name=4, anchor=center, inner sep=0}, "{Gf_2}"{description, pos=0.4}, curve={height=-6pt}, from=2-4, to=4-5]
	\arrow[""{name=5, anchor=center, inner sep=0}, "{Gg_2}"{description}, curve={height=12pt}, from=2-4, to=4-5]
	\arrow[""{name=6, anchor=center, inner sep=0}, "{\Xi_b}", from=4-3, to=2-4]
	\arrow[""{name=7, anchor=center, inner sep=0}, "{Ff_1}", curve={height=-12pt}, from=4-2, to=4-3]
	\arrow["G\beta", shorten <=3pt, shorten >=3pt, Rightarrow, from=4, to=5]
	\arrow["{=}"{description}, draw=none, from=1, to=3]
	\arrow["{=}"{description}, draw=none, from=2, to=6]
	\arrow["{=}", draw=none, from=6, to=4-4]
	\arrow["F\alpha", shift right=2, shorten <=3pt, shorten >=3pt, Rightarrow, from=7, to=0]
\end{tikzcd}
}
\newquiver{mult 2-functoriality horizontal composition: g}{
% https://q.uiver.app/?q=WzAsOCxbMSwzLCJGYSJdLFs0LDMsIkdjIl0sWzMsMywiRmMiXSxbMiwzLCJGYiJdLFswLDAsIkZiIl0sWzEsMSwiR2EiXSxbMywxLCJHYiJdLFs0LDAsIkZjIl0sWzIsMSwiXFxYaV9jIiwyXSxbMCwzLCJGZ18xIiwyLHsiY3VydmUiOjJ9XSxbMywyLCJGZ18yIiwyLHsiY3VydmUiOjJ9XSxbMCw0LCJGZl8xIiwxXSxbNCw3LCJGZ18xIiwxXSxbNywxLCJcXFhpX2MiLDFdLFswLDUsIlxcWGlfYSIsMV0sWzUsNiwiR2ZfMSIsMV0sWzYsMSwiR2ZfMiIsMSx7ImxhYmVsX3Bvc2l0aW9uIjo0MH1dLFszLDYsIlxcWGlfYiJdLFswLDMsIkZmXzEiLDAseyJjdXJ2ZSI6LTJ9XSxbMywyLCJGZl8yIiwwLHsiY3VydmUiOi0yfV0sWzEyLDE1LCI9IiwxLHsic2hvcnRlbiI6eyJzb3VyY2UiOjIwLCJ0YXJnZXQiOjIwfSwic3R5bGUiOnsiYm9keSI6eyJuYW1lIjoibm9uZSJ9LCJoZWFkIjp7Im5hbWUiOiJub25lIn19fV0sWzE0LDE3LCI9IiwxLHsic2hvcnRlbiI6eyJzb3VyY2UiOjIwLCJ0YXJnZXQiOjIwfSwic3R5bGUiOnsiYm9keSI6eyJuYW1lIjoibm9uZSJ9LCJoZWFkIjp7Im5hbWUiOiJub25lIn19fV0sWzE3LDIsIj0iLDAseyJzaG9ydGVuIjp7InNvdXJjZSI6MjB9LCJzdHlsZSI6eyJib2R5Ijp7Im5hbWUiOiJub25lIn0sImhlYWQiOnsibmFtZSI6Im5vbmUifX19XSxbMTgsOSwiRlxcYWxwaGEiLDAseyJvZmZzZXQiOjIsInNob3J0ZW4iOnsic291cmNlIjoyMCwidGFyZ2V0IjoyMH19XSxbMTksMTAsIkZcXGJldGEiLDAseyJvZmZzZXQiOjEsInNob3J0ZW4iOnsic291cmNlIjoyMCwidGFyZ2V0IjoyMH19XV0=
\begin{tikzcd}[ampersand replacement=\&]
	Fb \&\&\&\& Fc \\
	\& Ga \&\& Gb \\
	\\
	\& Fa \& Fb \& Fc \& Gc
	\arrow["{\Xi_c}"', from=4-4, to=4-5]
	\arrow[""{name=0, anchor=center, inner sep=0}, "{Fg_1}"', curve={height=12pt}, from=4-2, to=4-3]
	\arrow[""{name=1, anchor=center, inner sep=0}, "{Fg_2}"', curve={height=12pt}, from=4-3, to=4-4]
	\arrow["{Ff_1}"{description}, from=4-2, to=1-1]
	\arrow[""{name=2, anchor=center, inner sep=0}, "{Fg_1}"{description}, from=1-1, to=1-5]
	\arrow["{\Xi_c}"{description}, from=1-5, to=4-5]
	\arrow[""{name=3, anchor=center, inner sep=0}, "{\Xi_a}"{description}, from=4-2, to=2-2]
	\arrow[""{name=4, anchor=center, inner sep=0}, "{Gf_1}"{description}, from=2-2, to=2-4]
	\arrow["{Gf_2}"{description, pos=0.4}, from=2-4, to=4-5]
	\arrow[""{name=5, anchor=center, inner sep=0}, "{\Xi_b}", from=4-3, to=2-4]
	\arrow[""{name=6, anchor=center, inner sep=0}, "{Ff_1}", curve={height=-12pt}, from=4-2, to=4-3]
	\arrow[""{name=7, anchor=center, inner sep=0}, "{Ff_2}", curve={height=-12pt}, from=4-3, to=4-4]
	\arrow["{=}"{description}, draw=none, from=2, to=4]
	\arrow["{=}"{description}, draw=none, from=3, to=5]
	\arrow["{=}", draw=none, from=5, to=4-4]
	\arrow["F\alpha", shift right=2, shorten <=3pt, shorten >=3pt, Rightarrow, from=6, to=0]
	\arrow["F\beta", shift right=1, shorten <=3pt, shorten >=3pt, Rightarrow, from=7, to=1]
\end{tikzcd}
}
\newquiver{mult 2-functoriality horizontal composition: h}{
% https://q.uiver.app/?q=WzAsNCxbMCwwLCJGYSJdLFsxLDAsIkZiIl0sWzMsMCwiR2MiXSxbMiwwLCJGYyJdLFswLDEsIkZmXzEiLDAseyJjdXJ2ZSI6LTJ9XSxbMCwxLCJGZ18xIiwyLHsiY3VydmUiOjJ9XSxbMywyLCJcXFhpX2MiXSxbMSwzLCJGZ18xIiwwLHsiY3VydmUiOi0yfV0sWzEsMywiRmdfMiIsMix7ImN1cnZlIjoyfV0sWzQsNSwiRlxcYWxwaGEiLDAseyJvZmZzZXQiOjEsInNob3J0ZW4iOnsic291cmNlIjoyMCwidGFyZ2V0IjoyMH19XSxbNyw4LCJGXFxiZXRhIiwwLHsic2hvcnRlbiI6eyJzb3VyY2UiOjIwLCJ0YXJnZXQiOjIwfX1dXQ==
\begin{tikzcd}[ampersand replacement=\&]
	Fa \& Fb \& Fc \& Gc
	\arrow[""{name=0, anchor=center, inner sep=0}, "{Ff_1}", curve={height=-12pt}, from=1-1, to=1-2]
	\arrow[""{name=1, anchor=center, inner sep=0}, "{Fg_1}"', curve={height=12pt}, from=1-1, to=1-2]
	\arrow["{\Xi_c}", from=1-3, to=1-4]
	\arrow[""{name=2, anchor=center, inner sep=0}, "{Fg_1}", curve={height=-12pt}, from=1-2, to=1-3]
	\arrow[""{name=3, anchor=center, inner sep=0}, "{Fg_2}"', curve={height=12pt}, from=1-2, to=1-3]
	\arrow["F\alpha", shift right=1, shorten <=3pt, shorten >=3pt, Rightarrow, from=0, to=1]
	\arrow["F\beta", shorten <=3pt, shorten >=3pt, Rightarrow, from=2, to=3]
\end{tikzcd}
}

\newquiver{Cop morphism}{
% https://q.uiver.app/?q=WzAsMTIsWzAsMSwiXFxhdmVjXzEiXSxbMCwzLCJcXGF2ZWNfbiJdLFsxLDEsImZfMSJdLFsyLDEsIlxccGFydGl0X3sxe1xcYWJze1xcYnZlY18xfX19Il0sWzIsMCwiXFxwYXJ0aXRfezExfSJdLFsxLDMsImZfbiJdLFsyLDMsIlxccGFydGl0X3tuMX0iXSxbMiw0LCJcXHBhcnRpdF97bntcXGFic3tcXGJ2ZWNfbn19fSJdLFszLDAsIlxcYnZlY197XFxwYWlyMTF9Il0sWzMsMSwiXFxidmVjX3tcXHBhaXIxe1xcYWJze1xcYnZlY18xfX19Il0sWzMsMywiXFxidmVjX3tcXHBhaXIgbjF9Il0sWzMsNCwiXFxidmVjX3tcXHBhaXIgbntcXGFic3tcXGJ2ZWNfbn19fSJdLFswLDEsIlxcdmRvdHMiLDEseyJzdHlsZSI6eyJib2R5Ijp7Im5hbWUiOiJub25lIn0sImhlYWQiOnsibmFtZSI6Im5vbmUifX19XSxbMiwwXSxbNCwyLCIiLDEseyJzdHlsZSI6eyJoZWFkIjp7Im5hbWUiOiJub25lIn19fV0sWzMsMiwiIiwxLHsic3R5bGUiOnsiaGVhZCI6eyJuYW1lIjoibm9uZSJ9fX1dLFs0LDMsIlxcdmRvdHMiLDEseyJzdHlsZSI6eyJib2R5Ijp7Im5hbWUiOiJub25lIn0sImhlYWQiOnsibmFtZSI6Im5vbmUifX19XSxbNiw3LCJcXHZkb3RzIiwxLHsic3R5bGUiOnsiYm9keSI6eyJuYW1lIjoibm9uZSJ9LCJoZWFkIjp7Im5hbWUiOiJub25lIn19fV0sWzYsNSwiIiwxLHsic3R5bGUiOnsiaGVhZCI6eyJuYW1lIjoibm9uZSJ9fX1dLFs3LDUsIiIsMSx7InN0eWxlIjp7ImhlYWQiOnsibmFtZSI6Im5vbmUifX19XSxbNSwxXSxbNCw4LCI9IiwxLHsic3R5bGUiOnsiYm9keSI6eyJuYW1lIjoibm9uZSJ9LCJoZWFkIjp7Im5hbWUiOiJub25lIn19fV0sWzMsOSwiPSIsMSx7InN0eWxlIjp7ImJvZHkiOnsibmFtZSI6Im5vbmUifSwiaGVhZCI6eyJuYW1lIjoibm9uZSJ9fX1dLFsxMCw2LCI9IiwxLHsic3R5bGUiOnsiYm9keSI6eyJuYW1lIjoibm9uZSJ9LCJoZWFkIjp7Im5hbWUiOiJub25lIn19fV0sWzExLDcsIj0iLDEseyJzdHlsZSI6eyJib2R5Ijp7Im5hbWUiOiJub25lIn0sImhlYWQiOnsibmFtZSI6Im5vbmUifX19XSxbMTQsMTUsIiIsMSx7ImN1cnZlIjotMSwic2hvcnRlbiI6eyJzb3VyY2UiOjIwLCJ0YXJnZXQiOjIwfSwibGV2ZWwiOjF9XSxbMTgsMTksIiIsMSx7ImN1cnZlIjotMSwic2hvcnRlbiI6eyJzb3VyY2UiOjIwLCJ0YXJnZXQiOjIwfSwibGV2ZWwiOjF9XV0=
\begin{tikzcd}[ampersand replacement=\&]
	\&\& {\partit_{11}} \& {\bvec_{\pair11}} \\
	{\avec_1} \& {f_1} \& {\partit_{1{\abs{\bvec_1}}}} \& {\bvec_{\pair1{\abs{\bvec_1}}}} \\
	\\
	{\avec_n} \& {f_n} \& {\partit_{n1}} \& {\bvec_{\pair n1}} \\
	\&\& {\partit_{n{\abs{\bvec_n}}}} \& {\bvec_{\pair n{\abs{\bvec_n}}}}
	\arrow["\vdots"{description}, draw=none, from=2-1, to=4-1]
	\arrow[from=2-2, to=2-1]
	\arrow[""{name=0, anchor=center, inner sep=0}, no head, from=1-3, to=2-2]
	\arrow[""{name=1, anchor=center, inner sep=0}, no head, from=2-3, to=2-2]
	\arrow["\vdots"{description}, draw=none, from=1-3, to=2-3]
	\arrow["\vdots"{description}, draw=none, from=4-3, to=5-3]
	\arrow[""{name=2, anchor=center, inner sep=0}, no head, from=4-3, to=4-2]
	\arrow[""{name=3, anchor=center, inner sep=0}, no head, from=5-3, to=4-2]
	\arrow[from=4-2, to=4-1]
	\arrow["{=}"{description}, draw=none, from=1-3, to=1-4]
	\arrow["{=}"{description}, draw=none, from=2-3, to=2-4]
	\arrow["{=}"{description}, draw=none, from=4-4, to=4-3]
	\arrow["{=}"{description}, draw=none, from=5-4, to=5-3]
	\arrow[curve={height=-6pt}, shorten <=3pt, shorten >=3pt, from=0, to=1]
	\arrow[curve={height=-6pt}, shorten <=3pt, shorten >=3pt, from=2, to=3]
\end{tikzcd}}

\newquiver{Cop composition}{
% https://q.uiver.app/?q=WzAsMjAsWzAsMiwiXFxhdmVjXzEiXSxbMCw1LCJcXGF2ZWNfbiJdLFsxLDIsImZfMSJdLFsyLDIsIlxcYnZlY197XFxwYWlyMXtuXzF9fSJdLFsyLDEsIlxcYnZlY197XFxwYWlyIDExfSJdLFsxLDUsImZfbiJdLFsyLDUsIlxcYnZlY197XFxwYWlyIG4xfSJdLFsyLDYsIlxcYnZlY197XFxwYWlyIG57bV9ufX0iXSxbMywxLCJnX3tcXHBhaXIxMX0iXSxbMywyLCJnX3tcXHBhaXIxe21fMX19Il0sWzQsMCwiXFxjdmVjX3tcXHRyaXBsZTExMX0iXSxbNCwxLCJcXGN2ZWNfe1xcdHJpcGxlMTF7a197MTF9fX0iXSxbNCwyLCJcXGN2ZWNfe1xcdHJpcGxlMXtuXzF9MX0iXSxbNCwzLCJcXGN2ZWNfe3tcXHRyaXBsZTF7bl8xfXtrX3sxbl8xfX19fSJdLFszLDUsImdfe1xccGFpciBuMX0iXSxbMyw2LCJnX3tcXHBhaXIgbnttX259fSJdLFs0LDQsIlxcY3ZlY197XFx0cmlwbGUgbjExfSJdLFs0LDUsIlxcY3ZlY197XFx0cmlwbGUgbjF7a197bjF9fX0iXSxbNCw2LCJcXGN2ZWNfe1xcdHJpcGxlIG57bV9ufTF9Il0sWzQsNywiXFxjdmVjX3tcXHRyaXBsZSBue21fbn17a197bm1fbn19fSJdLFswLDEsIlxcdmRvdHMiLDEseyJzdHlsZSI6eyJib2R5Ijp7Im5hbWUiOiJub25lIn0sImhlYWQiOnsibmFtZSI6Im5vbmUifX19XSxbMiwwXSxbNCwyLCIiLDEseyJzdHlsZSI6eyJoZWFkIjp7Im5hbWUiOiJub25lIn19fV0sWzMsMiwiIiwxLHsic3R5bGUiOnsiaGVhZCI6eyJuYW1lIjoibm9uZSJ9fX1dLFs0LDMsIlxcdmRvdHMiLDEseyJzdHlsZSI6eyJib2R5Ijp7Im5hbWUiOiJub25lIn0sImhlYWQiOnsibmFtZSI6Im5vbmUifX19XSxbNiw3LCJcXHZkb3RzIiwxLHsic3R5bGUiOnsiYm9keSI6eyJuYW1lIjoibm9uZSJ9LCJoZWFkIjp7Im5hbWUiOiJub25lIn19fV0sWzYsNSwiIiwxLHsic3R5bGUiOnsiaGVhZCI6eyJuYW1lIjoibm9uZSJ9fX1dLFs3LDUsIiIsMSx7InN0eWxlIjp7ImhlYWQiOnsibmFtZSI6Im5vbmUifX19XSxbNSwxXSxbOCw5LCJcXHZkb3RzIiwxLHsic3R5bGUiOnsiYm9keSI6eyJuYW1lIjoibm9uZSJ9LCJoZWFkIjp7Im5hbWUiOiJub25lIn19fV0sWzgsNF0sWzksM10sWzExLDgsIiIsMSx7InN0eWxlIjp7ImhlYWQiOnsibmFtZSI6Im5vbmUifX19XSxbMTAsOCwiIiwxLHsic3R5bGUiOnsiaGVhZCI6eyJuYW1lIjoibm9uZSJ9fX1dLFsxMCwxMSwiXFx2ZG90cyIsMSx7InN0eWxlIjp7ImJvZHkiOnsibmFtZSI6Im5vbmUifSwiaGVhZCI6eyJuYW1lIjoibm9uZSJ9fX1dLFsxMiw5LCIiLDEseyJzdHlsZSI6eyJoZWFkIjp7Im5hbWUiOiJub25lIn19fV0sWzEzLDksIiIsMSx7InN0eWxlIjp7ImhlYWQiOnsibmFtZSI6Im5vbmUifX19XSxbMTQsNl0sWzE1LDddLFsxNiwxNCwiIiwxLHsic3R5bGUiOnsiaGVhZCI6eyJuYW1lIjoibm9uZSJ9fX1dLFsxNywxNCwiIiwxLHsic3R5bGUiOnsiaGVhZCI6eyJuYW1lIjoibm9uZSJ9fX1dLFsxOCwxNSwiIiwxLHsic3R5bGUiOnsiaGVhZCI6eyJuYW1lIjoibm9uZSJ9fX1dLFsxOSwxNSwiIiwxLHsic3R5bGUiOnsiaGVhZCI6eyJuYW1lIjoibm9uZSJ9fX1dLFsxMiwxMywiXFx2ZG90cyIsMSx7InN0eWxlIjp7ImJvZHkiOnsibmFtZSI6Im5vbmUifSwiaGVhZCI6eyJuYW1lIjoibm9uZSJ9fX1dLFsxNiwxNywiXFx2ZG90cyIsMSx7InN0eWxlIjp7ImJvZHkiOnsibmFtZSI6Im5vbmUifSwiaGVhZCI6eyJuYW1lIjoibm9uZSJ9fX1dLFsxOCwxOSwiXFx2ZG90cyIsMSx7InN0eWxlIjp7ImJvZHkiOnsibmFtZSI6Im5vbmUifSwiaGVhZCI6eyJuYW1lIjoibm9uZSJ9fX1dLFsyMiwyMywiIiwxLHsiY3VydmUiOi0xLCJzaG9ydGVuIjp7InNvdXJjZSI6MjAsInRhcmdldCI6MjB9LCJsZXZlbCI6MX1dLFsyNiwyNywiIiwxLHsiY3VydmUiOi0xLCJzaG9ydGVuIjp7InNvdXJjZSI6MjAsInRhcmdldCI6MjB9LCJsZXZlbCI6MX1dLFszMywzMiwiIiwxLHsiY3VydmUiOi0xLCJzaG9ydGVuIjp7InNvdXJjZSI6MjAsInRhcmdldCI6MjB9LCJsZXZlbCI6MX1dLFszNSwzNiwiIiwxLHsiY3VydmUiOi0xLCJzaG9ydGVuIjp7InNvdXJjZSI6MjAsInRhcmdldCI6MjB9LCJsZXZlbCI6MX1dLFszOSw0MCwiIiwxLHsiY3VydmUiOi0xLCJzaG9ydGVuIjp7InNvdXJjZSI6MjAsInRhcmdldCI6MjB9LCJsZXZlbCI6MX1dLFs0MSw0MiwiIiwxLHsiY3VydmUiOi0xLCJzaG9ydGVuIjp7InNvdXJjZSI6MjAsInRhcmdldCI6MjB9LCJsZXZlbCI6MX1dXQ==
\begin{tikzcd}[ampersand replacement=\&]
	\&\&\&\& {\cvec_{\triple111}} \\
	\&\& {\bvec_{\pair 11}} \& {g_{\pair11}} \& {\cvec_{\triple11{k_{11}}}} \\
	{\avec_1} \& {f_1} \& {\bvec_{\pair1{n_1}}} \& {g_{\pair1{m_1}}} \& {\cvec_{\triple1{n_1}1}} \\
	\&\&\&\& {\cvec_{{\triple1{n_1}{k_{1n_1}}}}} \\
	\&\&\&\& {\cvec_{\triple n11}} \\
	{\avec_n} \& {f_n} \& {\bvec_{\pair n1}} \& {g_{\pair n1}} \& {\cvec_{\triple n1{k_{n1}}}} \\
	\&\& {\bvec_{\pair n{m_n}}} \& {g_{\pair n{m_n}}} \& {\cvec_{\triple n{m_n}1}} \\
	\&\&\&\& {\cvec_{\triple n{m_n}{k_{nm_n}}}}
	\arrow["\vdots"{description}, draw=none, from=3-1, to=6-1]
	\arrow[from=3-2, to=3-1]
	\arrow[""{name=0, anchor=center, inner sep=0}, no head, from=2-3, to=3-2]
	\arrow[""{name=1, anchor=center, inner sep=0}, no head, from=3-3, to=3-2]
	\arrow["\vdots"{description}, draw=none, from=2-3, to=3-3]
	\arrow["\vdots"{description}, draw=none, from=6-3, to=7-3]
	\arrow[""{name=2, anchor=center, inner sep=0}, no head, from=6-3, to=6-2]
	\arrow[""{name=3, anchor=center, inner sep=0}, no head, from=7-3, to=6-2]
	\arrow[from=6-2, to=6-1]
	\arrow["\vdots"{description}, draw=none, from=2-4, to=3-4]
	\arrow[from=2-4, to=2-3]
	\arrow[from=3-4, to=3-3]
	\arrow[""{name=4, anchor=center, inner sep=0}, no head, from=2-5, to=2-4]
	\arrow[""{name=5, anchor=center, inner sep=0}, no head, from=1-5, to=2-4]
	\arrow["\vdots"{description}, draw=none, from=1-5, to=2-5]
	\arrow[""{name=6, anchor=center, inner sep=0}, no head, from=3-5, to=3-4]
	\arrow[""{name=7, anchor=center, inner sep=0}, no head, from=4-5, to=3-4]
	\arrow[from=6-4, to=6-3]
	\arrow[from=7-4, to=7-3]
	\arrow[""{name=8, anchor=center, inner sep=0}, no head, from=5-5, to=6-4]
	\arrow[""{name=9, anchor=center, inner sep=0}, no head, from=6-5, to=6-4]
	\arrow[""{name=10, anchor=center, inner sep=0}, no head, from=7-5, to=7-4]
	\arrow[""{name=11, anchor=center, inner sep=0}, no head, from=8-5, to=7-4]
	\arrow["\vdots"{description}, draw=none, from=3-5, to=4-5]
	\arrow["\vdots"{description}, draw=none, from=5-5, to=6-5]
	\arrow["\vdots"{description}, draw=none, from=7-5, to=8-5]
	\arrow[curve={height=-6pt}, shorten <=3pt, shorten >=3pt, from=0, to=1]
	\arrow[curve={height=-6pt}, shorten <=3pt, shorten >=3pt, from=2, to=3]
	\arrow[curve={height=-6pt}, shorten <=3pt, shorten >=3pt, from=5, to=4]
	\arrow[curve={height=-6pt}, shorten <=3pt, shorten >=3pt, from=6, to=7]
	\arrow[curve={height=-6pt}, shorten <=3pt, shorten >=3pt, from=8, to=9]
	\arrow[curve={height=-6pt}, shorten <=3pt, shorten >=3pt, from=10, to=11]
\end{tikzcd}
}
\newquiver{psh adjoint equivalence}{
% https://q.uiver.app/?q=WzAsNCxbMSwwLCJcXFBzaFxcVVUiXSxbMiwwLCJcXFBzaCdcXFVVIl0sWzMsMCwiXFxQc2hcXFVVIl0sWzAsMCwiXFxQc2gnXFxVVSJdLFswLDEsIkUnIl0sWzMsMCwiRSJdLFsxLDIsIkUiXSxbMywxLCIiLDAseyJjdXJ2ZSI6LTQsImxldmVsIjoyLCJzdHlsZSI6eyJoZWFkIjp7Im5hbWUiOiJub25lIn19fV0sWzAsMiwiIiwwLHsiY3VydmUiOjQsImxldmVsIjoyLCJzdHlsZSI6eyJoZWFkIjp7Im5hbWUiOiJub25lIn19fV0sWzcsMCwiXFx1bml0IiwwLHsibGFiZWxfcG9zaXRpb24iOjYwLCJzaG9ydGVuIjp7InNvdXJjZSI6MjB9fV0sWzEsOCwiXFxjb3VuaXQiLDAseyJsYWJlbF9wb3NpdGlvbiI6MzAsInNob3J0ZW4iOnsidGFyZ2V0IjoyMH19XV0=
\begin{tikzcd}[ampersand replacement=\&]
	{\Psh'\UU} \& \Psh\UU \& {\Psh'\UU} \& \Psh\UU
	\arrow["{E'}", from=1-2, to=1-3]
	\arrow["E", from=1-1, to=1-2]
	\arrow["E", from=1-3, to=1-4]
	\arrow[""{name=0, anchor=center, inner sep=0}, curve={height=-24pt}, Rightarrow, no head, from=1-1, to=1-3]
	\arrow[""{name=1, anchor=center, inner sep=0}, curve={height=24pt}, Rightarrow, no head, from=1-2, to=1-4]
	\arrow["\unit"{pos=0.6}, shorten <=2pt, Rightarrow, from=0, to=1-2]
	\arrow["\counit"{pos=0.3}, shorten >=2pt, Rightarrow, from=1-3, to=1]
\end{tikzcd}
}

\newquiver{psh 2-cell transformer}{
 % https://q.uiver.app/?q=WzAsMyxbMCwxLCJcXFVVIl0sWzEsMCwiXFxQc2gnXFxVVSJdLFsyLDEsIlxcUHNoXFxVVSJdLFswLDEsIlxceW9uZWRhJyJdLFswLDIsIlxceW9uZWRhIiwyXSxbMSwyLCJFIl0sWzEsNCwiXFx0YXUiLDAseyJsYWJlbF9wb3NpdGlvbiI6NDAsInNob3J0ZW4iOnsic291cmNlIjoyMCwidGFyZ2V0IjoyMH19XV0=
\begin{tikzcd}[ampersand replacement=\&]
	\& {\Psh'\UU} \\
	\UU \&\& \Psh\UU
	\arrow["{\yoneda'}", from=2-1, to=1-2]
	\arrow[""{name=0, anchor=center, inner sep=0}, "\yoneda"', from=2-1, to=2-3]
	\arrow["E", from=1-2, to=2-3]
	\arrow["\tau"{pos=0.4}, shorten <=3pt, shorten >=3pt, Rightarrow, from=1-2, to=0]
\end{tikzcd}
}

\newquiver{psh 2-cell transformer'}{
% https://q.uiver.app/?q=WzAsMyxbMCwxLCJcXFVVIl0sWzEsMCwiXFxQc2hcXFVVIl0sWzIsMSwiXFxQc2gnXFxVVSJdLFswLDEsIlxceW9uZWRhIl0sWzAsMiwiXFx5b25lZGEnIiwyXSxbMSwyLCJFJyJdLFsxLDQsIlxcdGF1JyIsMCx7ImxhYmVsX3Bvc2l0aW9uIjo0MCwic2hvcnRlbiI6eyJzb3VyY2UiOjIwLCJ0YXJnZXQiOjIwfX1dXQ==
\begin{tikzcd}[ampersand replacement=\&]
	\& \Psh\UU \\
	\UU \&\& {\Psh'\UU}
	\arrow["\yoneda", from=2-1, to=1-2]
	\arrow[""{name=0, anchor=center, inner sep=0}, "{\yoneda'}"', from=2-1, to=2-3]
	\arrow["{E'}", from=1-2, to=2-3]
	\arrow["{\tau'}"{pos=0.4}, shorten <=3pt, shorten >=3pt, Rightarrow, from=1-2, to=0]
\end{tikzcd}
}

\newquiver{psh transformer condition}{
% https://q.uiver.app/?q=WzAsNyxbMCwyLCJcXFVVIl0sWzEsMCwiXFxQc2gnXFxVVSJdLFsxLDEsIlxcUHNoXFxVVSJdLFsyLDIsIlxcUHNoJ1xcVVUiXSxbMywyLCJcXFVVIl0sWzQsMCwiXFxQc2gnXFxVVSJdLFs1LDIsIlxcUHNoJ1xcVVUiXSxbMCwxLCJcXHlvbmVkYSciLDAseyJjdXJ2ZSI6LTJ9XSxbMSwyLCJFIl0sWzIsMywiRSciLDJdLFswLDMsIlxceW9uZWRhJyIsMix7ImN1cnZlIjozfV0sWzAsMiwiXFx5b25lZGEiLDJdLFs0LDUsIlxceW9uZWRhJyIsMCx7ImN1cnZlIjotM31dLFs0LDYsIlxceW9uZWRhJyIsMix7ImN1cnZlIjozfV0sWzUsNiwiIiwxLHsiY3VydmUiOi0zLCJsZXZlbCI6Miwic3R5bGUiOnsiaGVhZCI6eyJuYW1lIjoibm9uZSJ9fX1dLFsxLDMsIiIsMix7ImN1cnZlIjotMiwibGV2ZWwiOjJ9XSxbMSwxMSwiXFx0YXUiLDIseyJjdXJ2ZSI6MSwic2hvcnRlbiI6eyJzb3VyY2UiOjMwLCJ0YXJnZXQiOjIwfX1dLFsyLDEwLCJcXHRhdSciLDAseyJzaG9ydGVuIjp7InNvdXJjZSI6NDAsInRhcmdldCI6MzB9fV0sWzE1LDIsIlxcdW5pdCIsMCx7ImN1cnZlIjotMSwic2hvcnRlbiI6eyJzb3VyY2UiOjQwfX1dLFs1LDEzLCI9IiwwLHsiY3VydmUiOjEsInNob3J0ZW4iOnsic291cmNlIjoyMCwidGFyZ2V0IjoyMH0sInN0eWxlIjp7ImJvZHkiOnsibmFtZSI6Im5vbmUifSwiaGVhZCI6eyJuYW1lIjoibm9uZSJ9fX1dLFsxMiwxNSwiPSIsMCx7InNob3J0ZW4iOnsic291cmNlIjoyMCwidGFyZ2V0IjoyMH0sInN0eWxlIjp7ImJvZHkiOnsibmFtZSI6Im5vbmUifSwiaGVhZCI6eyJuYW1lIjoibm9uZSJ9fX1dXQ==
\begin{tikzcd}[ampersand replacement=\&]
	\& {\Psh'\UU} \&\&\& {\Psh'\UU} \\
	\& \Psh\UU \\
	\UU \&\& {\Psh'\UU} \& \UU \&\& {\Psh'\UU}
	\arrow["{\yoneda'}", curve={height=-12pt}, from=3-1, to=1-2]
	\arrow["E", from=1-2, to=2-2]
	\arrow["{E'}"', from=2-2, to=3-3]
	\arrow[""{name=0, anchor=center, inner sep=0}, "{\yoneda'}"', curve={height=18pt}, from=3-1, to=3-3]
	\arrow[""{name=1, anchor=center, inner sep=0}, "\yoneda"', from=3-1, to=2-2]
	\arrow[""{name=2, anchor=center, inner sep=0}, "{\yoneda'}", curve={height=-18pt}, from=3-4, to=1-5]
	\arrow[""{name=3, anchor=center, inner sep=0}, "{\yoneda'}"', curve={height=18pt}, from=3-4, to=3-6]
	\arrow[curve={height=-18pt}, Rightarrow, no head, from=1-5, to=3-6]
	\arrow[""{name=4, anchor=center, inner sep=0}, curve={height=-12pt}, Rightarrow, from=1-2, to=3-3]
	\arrow["\tau"', curve={height=6pt}, shorten <=9pt, shorten >=6pt, Rightarrow, from=1-2, to=1]
	\arrow["{\tau'}", shorten <=11pt, shorten >=8pt, Rightarrow, from=2-2, to=0]
	\arrow["\unit", curve={height=-6pt}, shorten <=8pt, Rightarrow, from=4, to=2-2]
	\arrow["{=}", curve={height=6pt}, draw=none, from=1-5, to=3]
	\arrow["{=}", draw=none, from=2, to=4]
\end{tikzcd}}

\newquiver{psh transformer condition'}{
% https://q.uiver.app/?q=WzAsOCxbMCwyLCJcXFVVIl0sWzEsMCwiXFxQc2hcXFVVIl0sWzIsMSwiXFxQc2gnXFxVVSJdLFsyLDIsIlxcUHNoXFxVVSJdLFszLDIsIlxcVVUiXSxbNCwwLCJcXFBzaFxcVVUiXSxbNSwyLCJcXFBzaFxcVVUiXSxbNSwxLCJcXFBzaCdcXFVVIl0sWzAsMSwiXFx5b25lZGEiLDAseyJjdXJ2ZSI6LTJ9XSxbMSwyLCJFJyIsMCx7ImN1cnZlIjotMn1dLFsyLDMsIkUiLDAseyJjdXJ2ZSI6LTF9XSxbMCwzLCJcXHlvbmVkYSIsMix7ImN1cnZlIjozfV0sWzAsMiwiXFx5b25lZGEnIiwyXSxbNCw1LCJcXHlvbmVkYSIsMCx7ImN1cnZlIjotM31dLFs1LDcsIkUnIiwwLHsiY3VydmUiOi0xfV0sWzcsNiwiRSIsMCx7ImN1cnZlIjotMX1dLFs1LDYsIiIsMSx7ImN1cnZlIjoyLCJsZXZlbCI6Miwic3R5bGUiOnsiaGVhZCI6eyJuYW1lIjoibm9uZSJ9fX1dLFs0LDYsIlxceW9uZWRhIiwyLHsiY3VydmUiOjN9XSxbMSwxMiwiXFx0YXUnIiwyLHsiY3VydmUiOjEsInNob3J0ZW4iOnsic291cmNlIjozMCwidGFyZ2V0IjoyMH19XSxbMiwxMSwiXFx0YXUiLDAseyJzaG9ydGVuIjp7InNvdXJjZSI6NDAsInRhcmdldCI6MzB9fV0sWzcsMTYsIlxcY291bml0IiwyLHsiY3VydmUiOjEsInNob3J0ZW4iOnsidGFyZ2V0IjoyMH19XSxbMTMsMiwiPSIsMCx7ImxhYmVsX3Bvc2l0aW9uIjo4MCwib2Zmc2V0IjotMiwic2hvcnRlbiI6eyJzb3VyY2UiOjIwfSwic3R5bGUiOnsiYm9keSI6eyJuYW1lIjoibm9uZSJ9LCJoZWFkIjp7Im5hbWUiOiJub25lIn19fV0sWzEzLDE3LCI9IiwwLHsiY3VydmUiOi0xLCJzaG9ydGVuIjp7InNvdXJjZSI6MjAsInRhcmdldCI6MjB9LCJzdHlsZSI6eyJib2R5Ijp7Im5hbWUiOiJub25lIn0sImhlYWQiOnsibmFtZSI6Im5vbmUifX19XV0=
\begin{tikzcd}[ampersand replacement=\&]
	\& \Psh\UU \&\&\& \Psh\UU \\
	\&\& {\Psh'\UU} \&\&\& {\Psh'\UU} \\
	\UU \&\& \Psh\UU \& \UU \&\& \Psh\UU
	\arrow["\yoneda", curve={height=-12pt}, from=3-1, to=1-2]
	\arrow["{E'}", curve={height=-12pt}, from=1-2, to=2-3]
	\arrow["E", curve={height=-6pt}, from=2-3, to=3-3]
	\arrow[""{name=0, anchor=center, inner sep=0}, "\yoneda"', curve={height=18pt}, from=3-1, to=3-3]
	\arrow[""{name=1, anchor=center, inner sep=0}, "{\yoneda'}"', from=3-1, to=2-3]
	\arrow[""{name=2, anchor=center, inner sep=0}, "\yoneda", curve={height=-18pt}, from=3-4, to=1-5]
	\arrow["{E'}", curve={height=-6pt}, from=1-5, to=2-6]
	\arrow["E", curve={height=-6pt}, from=2-6, to=3-6]
	\arrow[""{name=3, anchor=center, inner sep=0}, curve={height=12pt}, Rightarrow, no head, from=1-5, to=3-6]
	\arrow[""{name=4, anchor=center, inner sep=0}, "\yoneda"', curve={height=18pt}, from=3-4, to=3-6]
	\arrow["{\tau'}"', curve={height=6pt}, shorten <=9pt, shorten >=6pt, Rightarrow, from=1-2, to=1]
	\arrow["\tau", shorten <=18pt, shorten >=13pt, Rightarrow, from=2-3, to=0]
	\arrow["\counit"', curve={height=6pt}, shorten >=4pt, Rightarrow, from=2-6, to=3]
	\arrow["{=}"{pos=0.8}, shift left=2, draw=none, from=2, to=2-3]
	\arrow["{=}", curve={height=-6pt}, draw=none, from=2, to=4]
\end{tikzcd}
}

\newquiver{yoneda structure structure}{
% https://q.uiver.app/?q=WzAsMyxbMCwwLCJcXFVVIl0sWzIsMCwiXFxVVVsyXSJdLFsyLDEsIlxcUHNoXFxVVSJdLFswLDEsIkYiXSxbMCwyLCJcXHlvbmVkYSIsMl0sWzEsMiwiXFxVVVsyXShGOy0pIl0sWzQsMSwiXFxtYXAgRiIsMCx7ImxhYmVsX3Bvc2l0aW9uIjo2MCwib2Zmc2V0Ijo1LCJjdXJ2ZSI6LTEsInNob3J0ZW4iOnsic291cmNlIjo0MCwidGFyZ2V0IjoxMH19XV0=
\begin{tikzcd}[ampersand replacement=\&]
	\UU \&\& {\UU[2]} \\
	\&\& \Psh\UU
	\arrow["F", from=1-1, to=1-3]
	\arrow[""{name=0, anchor=center, inner sep=0}, "\yoneda"', from=1-1, to=2-3]
	\arrow["{\UU[2](F;-)}", from=1-3, to=2-3]
	\arrow["{\map F}"{pos=0.6}, shift right=5, curve={height=-6pt}, shorten <=11pt, shorten >=3pt, Rightarrow, from=0, to=1-3]
\end{tikzcd}
}

\newquiver{axiom 1 spelled out}{
% https://q.uiver.app/?q=WzAsNSxbNSwwLCJcXFVVWzJdIl0sWzUsMSwiXFxQc2hcXFVVIl0sWzAsMCwiXFxVVSJdLFsxLDAsIlxcVVVbMl0iXSxbMSwxLCJcXFBzaFxcVVUiXSxbMCwxLCJcXFVVWzJdKEY7LSkiLDIseyJjdXJ2ZSI6M31dLFswLDEsIkciLDAseyJjdXJ2ZSI6LTN9XSxbMiwzLCJGIl0sWzMsNCwiRyIsMCx7ImN1cnZlIjotMn1dLFsyLDQsIlxceW9uZWRhX3tcXFVVfSIsMl0sWzUsNiwiXFxiZXRhIiwwLHsic2hvcnRlbiI6eyJzb3VyY2UiOjIwLCJ0YXJnZXQiOjIwfX1dLFs5LDMsIlxcYWxwaGEiLDAseyJsYWJlbF9wb3NpdGlvbiI6NDAsIm9mZnNldCI6LTEsImN1cnZlIjoxfV0sWzUsOCwiXFxzdWJzdGFja3tcXHRleHR7cGFzdGluZyB9XFxcXFxcbWFwIEZ9IiwwLHsic2hvcnRlbiI6eyJzb3VyY2UiOjMwLCJ0YXJnZXQiOjMwfSwibGV2ZWwiOjEsInN0eWxlIjp7InRhaWwiOnsibmFtZSI6ImFycm93aGVhZCJ9fX1dXQ==
\begin{tikzcd}[ampersand replacement=\&]
	\UU \& {\UU[2]} \&\&\&\& {\UU[2]} \\
	\& \Psh\UU \&\&\&\& \Psh\UU
	\arrow[""{name=0, anchor=center, inner sep=0}, "{\UU[2](F;-)}"', curve={height=18pt}, from=1-6, to=2-6]
	\arrow[""{name=1, anchor=center, inner sep=0}, "G", curve={height=-18pt}, from=1-6, to=2-6]
	\arrow["F", from=1-1, to=1-2]
	\arrow[""{name=2, anchor=center, inner sep=0}, "G", curve={height=-12pt}, from=1-2, to=2-2]
	\arrow[""{name=3, anchor=center, inner sep=0}, "{\yoneda_{\UU}}"', from=1-1, to=2-2]
	\arrow["\beta", shorten <=7pt, shorten >=7pt, Rightarrow, from=0, to=1]
	\arrow["\alpha"{pos=0.4}, shift left=1, curve={height=6pt}, Rightarrow, from=3, to=1-2]
	\arrow["{\substack{\text{pasting }\\\map F}}", shorten <=34pt, shorten >=34pt, tail reversed, from=0, to=2]
\end{tikzcd}}

\newquiver{axiom-2 lifting}{
% https://q.uiver.app/?q=WzAsNSxbMCwwLCJcXFVVWzNdIl0sWzEsMCwiXFxVVSJdLFsyLDAsIlxcVVVbMl0iXSxbMiwyLCJcXFBzaFxcVVUiXSxbMSwxLCJcXFVVIl0sWzQsMywiXFx5b25lZGFfe1xcVVV9IiwyLHsibGFiZWxfcG9zaXRpb24iOjIwfV0sWzAsNCwiRyIsMl0sWzEsMiwiRiJdLFswLDEsIkciXSxbMSw0LCIiLDAseyJsZXZlbCI6Miwic3R5bGUiOnsiaGVhZCI6eyJuYW1lIjoibm9uZSJ9fX1dLFsyLDMsIlxcVVVbMl0oRjstKSJdLFs0LDIsIlxcbWFwIEYiLDAseyJsYWJlbF9wb3NpdGlvbiI6NzAsIm9mZnNldCI6Mywic2hvcnRlbiI6eyJzb3VyY2UiOjIwLCJ0YXJnZXQiOjIwfSwibGV2ZWwiOjJ9XSxbNiwxLCI9IiwxLHsibGFiZWxfcG9zaXRpb24iOjcwLCJzaG9ydGVuIjp7InNvdXJjZSI6MjB9LCJzdHlsZSI6eyJib2R5Ijp7Im5hbWUiOiJub25lIn0sImhlYWQiOnsibmFtZSI6Im5vbmUifX19XV0=
\begin{tikzcd}[ampersand replacement=\&]
	{\UU[3]} \& \UU \& {\UU[2]} \\
	\& \UU \\
	\&\& \Psh\UU
	\arrow["{\yoneda_{\UU}}"'{pos=0.2}, from=2-2, to=3-3]
	\arrow[""{name=0, anchor=center, inner sep=0}, "G"', from=1-1, to=2-2]
	\arrow["F", from=1-2, to=1-3]
	\arrow["G", from=1-1, to=1-2]
	\arrow[Rightarrow, no head, from=1-2, to=2-2]
	\arrow["{\UU[2](F;-)}", from=1-3, to=3-3]
	\arrow["{\map F}"{pos=0.7}, shift right=3, shorten <=5pt, shorten >=5pt, Rightarrow, from=2-2, to=1-3]
	\arrow["{=}"{description, pos=0.7}, draw=none, from=0, to=1-2]
\end{tikzcd}
}

\newquiver{axiom-2 bijection}{
% https://q.uiver.app/?q=WzAsNyxbMCwwLCJcXFVVWzNdIl0sWzIsMCwiXFxVVVsyXSJdLFsyLDIsIlxcUHNoXFxVVSJdLFsxLDEsIlxcVVUiXSxbNCwwLCJcXFVVWzNdIl0sWzYsMCwiXFxVVVsyXSJdLFs1LDEsIlxcVVUiXSxbMywyLCJcXHlvbmVkYV97XFxVVX0iLDIseyJsYWJlbF9wb3NpdGlvbiI6MjB9XSxbMCwzLCJHIiwyXSxbMSwyLCJcXFVVWzJdKEY7LSkiXSxbMCwxLCJIIl0sWzMsMSwiXFxhbHBoYSIsMCx7ImxldmVsIjoyfV0sWzQsNSwiSCJdLFs0LDYsIkciLDIseyJjdXJ2ZSI6MX1dLFs2LDUsIkYiLDIseyJjdXJ2ZSI6MX1dLFs2LDEyLCJcXFlvbmVkYV97XFxVVSxGLEgsR1xcc2VxLX1cXGFscGhhIiwwLHsibGFiZWxfcG9zaXRpb24iOjcwLCJvZmZzZXQiOjUsImN1cnZlIjoxLCJzaG9ydGVuIjp7InNvdXJjZSI6MjAsInRhcmdldCI6MjB9fV0sWzksMTMsIlxcbG9uZ2xlZnRyaWdodGFycm93IiwwLHsibGFiZWxfcG9zaXRpb24iOjYwLCJzaG9ydGVuIjp7InNvdXJjZSI6MjAsInRhcmdldCI6MjB9LCJzdHlsZSI6eyJib2R5Ijp7Im5hbWUiOiJub25lIn0sImhlYWQiOnsibmFtZSI6Im5vbmUifX19XV0=
\begin{tikzcd}[ampersand replacement=\&]
	{\UU[3]} \&\& {\UU[2]} \&\& {\UU[3]} \&\& {\UU[2]} \\
	\& \UU \&\&\&\& \UU \\
	\&\& \Psh\UU
	\arrow["{\yoneda_{\UU}}"'{pos=0.2}, from=2-2, to=3-3]
	\arrow["G"', from=1-1, to=2-2]
	\arrow[""{name=0, anchor=center, inner sep=0}, "{\UU[2](F;-)}", from=1-3, to=3-3]
	\arrow["H", from=1-1, to=1-3]
	\arrow["\alpha", Rightarrow, from=2-2, to=1-3]
	\arrow[""{name=1, anchor=center, inner sep=0}, "H", from=1-5, to=1-7]
	\arrow[""{name=2, anchor=center, inner sep=0}, "G"', curve={height=6pt}, from=1-5, to=2-6]
	\arrow["F"', curve={height=6pt}, from=2-6, to=1-7]
	\arrow["{\Yoneda_{\UU,F,H,G\seq-}\alpha}"{pos=0.7}, shift right=5, curve={height=6pt}, shorten <=3pt, shorten >=3pt, Rightarrow, from=2-6, to=1]
	\arrow["\longleftrightarrow"{pos=0.6}, draw=none, from=0, to=2]
\end{tikzcd}
}

\newquiver{axiom-2 proof: uniqueness data}{
% https://q.uiver.app/?q=WzAsNyxbMCwwLCJcXFVVWzNdIl0sWzIsMCwiXFxVVVsyXSJdLFsyLDIsIlxcUHNoXFxVVSJdLFsxLDEsIlxcVVUiXSxbNCwwLCJcXFVVWzNdIl0sWzYsMCwiXFxVVVsyXSJdLFs1LDEsIlxcVVUiXSxbMywyLCJcXHlvbmVkYV97XFxVVX0iLDIseyJsYWJlbF9wb3NpdGlvbiI6MjB9XSxbMCwzLCJHIiwyXSxbMSwyLCJcXFVVWzJdKEY7LSkiXSxbMCwxLCJIIl0sWzMsMSwiXFxhbHBoYSIsMCx7ImxldmVsIjoyfV0sWzQsNSwiSCJdLFs0LDYsIkciLDIseyJjdXJ2ZSI6MX1dLFs2LDUsIkYiLDIseyJjdXJ2ZSI6MX1dLFs2LDEyLCJcXGJldGEiLDAseyJsYWJlbF9wb3NpdGlvbiI6NDAsInNob3J0ZW4iOnsic291cmNlIjoxMCwidGFyZ2V0IjozMH19XV0=
\begin{tikzcd}[ampersand replacement=\&]
	{\UU[3]} \&\& {\UU[2]} \&\& {\UU[3]} \&\& {\UU[2]} \\
	\& \UU \&\&\&\& \UU \\
	\&\& \Psh\UU
	\arrow["{\yoneda_{\UU}}"'{pos=0.2}, from=2-2, to=3-3]
	\arrow["G"', from=1-1, to=2-2]
	\arrow["{\UU[2](F;-)}", from=1-3, to=3-3]
	\arrow["H", from=1-1, to=1-3]
	\arrow["\alpha", Rightarrow, from=2-2, to=1-3]
	\arrow[""{name=0, anchor=center, inner sep=0}, "H", from=1-5, to=1-7]
	\arrow["G"', curve={height=6pt}, from=1-5, to=2-6]
	\arrow["F"', curve={height=6pt}, from=2-6, to=1-7]
	\arrow["\beta"{pos=0.4}, shorten <=2pt, shorten >=5pt, Rightarrow, from=2-6, to=0]
\end{tikzcd}
}
\newquiver{axiom-2 proof: uniqueness assumption}{
  % https://q.uiver.app/?q=WzAsOCxbMCwwLCJcXFVVWzNdIl0sWzIsMCwiXFxVVVsyXSJdLFsyLDIsIlxcUHNoXFxVVSJdLFsxLDEsIlxcVVUiXSxbNCwwLCJcXFVVWzNdIl0sWzYsMCwiXFxVVVsyXSJdLFs2LDIsIlxcUHNoXFxVVSJdLFs1LDEsIlxcVVUiXSxbMywyLCJcXHlvbmVkYV97XFxVVX0iLDIseyJsYWJlbF9wb3NpdGlvbiI6MjB9XSxbMCwzLCJHIiwyXSxbMSwyLCJcXFVVWzJdKEY7LSkiXSxbNCw3LCJHIiwyXSxbNyw2LCJcXHlvbmVkYV97XFxVVX0iLDJdLFs0LDUsIkgiLDAseyJjdXJ2ZSI6LTJ9XSxbNSw2LCJcXFVVWzJdKEY7LSkiXSxbNyw1LCJcXGFscGhhIiwwLHsic2hvcnRlbiI6eyJzb3VyY2UiOjIwLCJ0YXJnZXQiOjIwfSwibGV2ZWwiOjJ9XSxbMCwxLCJIIiwwLHsiY3VydmUiOi0zfV0sWzEsNCwiPSIsMix7InN0eWxlIjp7ImJvZHkiOnsibmFtZSI6Im5vbmUifSwiaGVhZCI6eyJuYW1lIjoibm9uZSJ9fX1dLFszLDEsIkYiLDAseyJjdXJ2ZSI6LTF9XSxbMywxNiwiXFxiZXRhIiwwLHsib2Zmc2V0IjotMywic2hvcnRlbiI6eyJzb3VyY2UiOjMwLCJ0YXJnZXQiOjMwfX1dLFs4LDEsIlxcbWFwIEYiLDAseyJvZmZzZXQiOi0zLCJjdXJ2ZSI6Miwic2hvcnRlbiI6eyJzb3VyY2UiOjQwLCJ0YXJnZXQiOjQwfX1dXQ==
\begin{tikzcd}[ampersand replacement=\&]
	{\UU[3]} \&\& {\UU[2]} \&\& {\UU[3]} \&\& {\UU[2]} \\
	\& \UU \&\&\&\& \UU \\
	\&\& \Psh\UU \&\&\&\& \Psh\UU
	\arrow[""{name=0, anchor=center, inner sep=0}, "{\yoneda_{\UU}}"'{pos=0.2}, from=2-2, to=3-3]
	\arrow["G"', from=1-1, to=2-2]
	\arrow["{\UU[2](F;-)}", from=1-3, to=3-3]
	\arrow["G"', from=1-5, to=2-6]
	\arrow["{\yoneda_{\UU}}"', from=2-6, to=3-7]
	\arrow["H", curve={height=-12pt}, from=1-5, to=1-7]
	\arrow["{\UU[2](F;-)}", from=1-7, to=3-7]
	\arrow["\alpha", shorten <=5pt, shorten >=5pt, Rightarrow, from=2-6, to=1-7]
	\arrow[""{name=1, anchor=center, inner sep=0}, "H", curve={height=-18pt}, from=1-1, to=1-3]
	\arrow["{=}"', draw=none, from=1-3, to=1-5]
	\arrow["F", curve={height=-6pt}, from=2-2, to=1-3]
	\arrow["\beta", shift left=3, shorten <=8pt, shorten >=8pt, Rightarrow, from=2-2, to=1]
	\arrow["{\map F}", shift left=3, curve={height=12pt}, shorten <=14pt, shorten >=14pt, Rightarrow, from=0, to=1-3]
\end{tikzcd}
}

\newquiver{internal yoneda lemma}{
% https://q.uiver.app/?q=WzAsNSxbMCwwLCJcXFVVIl0sWzEsMSwiXFxQc2hcXFVVIl0sWzEsMCwiXFxQc2hcXFVVIl0sWzMsMCwiXFxQc2hcXFVVIl0sWzMsMSwiXFxQc2hcXFVVIl0sWzAsMiwiXFx5b25lZGEiXSxbMiwxLCJcXElkIl0sWzAsMSwiXFx5b25lZGEiLDJdLFszLDQsIlxcUHNoXFxVVShcXHlvbmVkYTstKSIsMix7ImN1cnZlIjoyfV0sWzMsNCwiXFxJZCIsMCx7ImN1cnZlIjotMn1dLFs3LDIsIj0iLDEseyJzaG9ydGVuIjp7InNvdXJjZSI6MjB9LCJzdHlsZSI6eyJib2R5Ijp7Im5hbWUiOiJub25lIn0sImhlYWQiOnsibmFtZSI6Im5vbmUifX19XSxbOCw5LCJcXFlvbmVkYSIsMCx7InNob3J0ZW4iOnsic291cmNlIjoyMCwidGFyZ2V0IjoyMH19XV0=
\begin{tikzcd}[ampersand replacement=\&]
	\UU \& \Psh\UU \&\& \Psh\UU \\
	\& \Psh\UU \&\& \Psh\UU
	\arrow["\yoneda", from=1-1, to=1-2]
	\arrow["\Id", from=1-2, to=2-2]
	\arrow[""{name=0, anchor=center, inner sep=0}, "\yoneda"', from=1-1, to=2-2]
	\arrow[""{name=1, anchor=center, inner sep=0}, "{\Psh\UU(\yoneda;-)}"', curve={height=12pt}, from=1-4, to=2-4]
	\arrow[""{name=2, anchor=center, inner sep=0}, "\Id", curve={height=-12pt}, from=1-4, to=2-4]
	\arrow["{=}"{description}, draw=none, from=0, to=1-2]
	\arrow["\Yoneda", shorten <=5pt, shorten >=5pt, Rightarrow, from=1, to=2]
\end{tikzcd}}

\newquiver{hom-profunctor morphism input type}{
  % https://q.uiver.app/?q=WzAsMixbMCwwLCJcXFVVWzFdIl0sWzEsMCwiXFxVVVsyXSJdLFswLDEsIkYiLDAseyJjdXJ2ZSI6LTJ9XSxbMCwxLCJHIiwyLHsiY3VydmUiOjJ9XSxbMywyLCJcXGFscGhhIiwyLHsic2hvcnRlbiI6eyJzb3VyY2UiOjIwLCJ0YXJnZXQiOjIwfX1dXQ==
\begin{tikzcd}[ampersand replacement=\&]
	{\UU[1]} \& {\UU[2]}
	\arrow[""{name=0, anchor=center, inner sep=0}, "F", curve={height=-12pt}, from=1-1, to=1-2]
	\arrow[""{name=1, anchor=center, inner sep=0}, "G"', curve={height=12pt}, from=1-1, to=1-2]
	\arrow["\alpha"', shorten <=3pt, shorten >=3pt, Rightarrow, from=1, to=0]
\end{tikzcd}
}

\newquiver{hom-profunctor morphism output type}{
  % https://q.uiver.app/?q=WzAsMixbMCwwLCJcXFVVWzJdIl0sWzIsMCwiXFxQc2h7XFxVVVsxXX0iXSxbMCwxLCJcXFVVWzJdKEY7XFxJZC0pIiwwLHsiY3VydmUiOi0yfV0sWzAsMSwiXFxVVVsyXShHO1xcSWQtKSIsMix7ImN1cnZlIjoyfV0sWzIsMywiXFxVVVsyXShcXGFscGhhOyBcXElkLSkiLDAseyJvZmZzZXQiOjUsInNob3J0ZW4iOnsic291cmNlIjoyMCwidGFyZ2V0IjoyMH19XV0=
\begin{tikzcd}[ampersand replacement=\&]
	{\UU[2]} \&\& {\Psh{\UU[1]}}
	\arrow[""{name=0, anchor=center, inner sep=0}, "{\UU[2](F;\Id-)}", curve={height=-12pt}, from=1-1, to=1-3]
	\arrow[""{name=1, anchor=center, inner sep=0}, "{\UU[2](G;\Id-)}"', curve={height=12pt}, from=1-1, to=1-3]
	\arrow["{\UU[2](\alpha; \Id-)}", shift right=5, shorten <=3pt, shorten >=3pt, Rightarrow, from=0, to=1]
\end{tikzcd}
}

\newquiver{hom-profunctor special case up}{
% https://q.uiver.app/?q=WzAsNixbMywwLCJcXFVVWzJdIl0sWzMsMiwiXFxQc2h7XFxVVVsxXX0iXSxbMCwwLCJcXFVVWzFdIl0sWzUsMCwiXFxVVVsxXSJdLFs3LDAsIlxcVVVbMl0iXSxbNywyLCJcXFBzaHtcXFVVWzFdfSJdLFswLDEsIlxcVVVbMl0oRjtcXElkLSkiLDIseyJsYWJlbF9wb3NpdGlvbiI6MzAsImN1cnZlIjo0fV0sWzAsMSwiXFxVVVsyXShHO1xcSWQtKSIsMCx7ImN1cnZlIjotNX1dLFsyLDEsIlxceW9uZWRhIiwyLHsiY3VydmUiOjN9XSxbMiwwLCJGIiwwLHsiY3VydmUiOi0yfV0sWzMsNCwiRiIsMCx7ImN1cnZlIjotMn1dLFszLDUsIlxceW9uZWRhIiwyXSxbNCw1LCJcXFVVWzJdKEc7XFxJZC0pIiwwLHsiY3VydmUiOi01fV0sWzMsNCwiRyIsMix7ImN1cnZlIjoyfV0sWzYsNywiXFxVVVsyXShcXGFscGhhOyBcXElkLSkiLDAseyJvZmZzZXQiOjEsInNob3J0ZW4iOnsic291cmNlIjoyMCwidGFyZ2V0IjoyMH19XSxbMTMsMTAsIlxcYWxwaGEiLDIseyJzaG9ydGVuIjp7InNvdXJjZSI6MjAsInRhcmdldCI6MjB9fV0sWzExLDQsIlxcbWFwIEciLDIseyJjdXJ2ZSI6Miwic2hvcnRlbiI6eyJzb3VyY2UiOjMwLCJ0YXJnZXQiOjEwfX1dLFs3LDExLCI9IiwxLHsibGFiZWxfcG9zaXRpb24iOjcwLCJzaG9ydGVuIjp7InNvdXJjZSI6MjAsInRhcmdldCI6MjB9LCJzdHlsZSI6eyJib2R5Ijp7Im5hbWUiOiJub25lIn0sImhlYWQiOnsibmFtZSI6Im5vbmUifX19XSxbOCw5LCJcXG1hcCBGIiwwLHsiY3VydmUiOi0yLCJzaG9ydGVuIjp7InNvdXJjZSI6MjAsInRhcmdldCI6MjB9fV1d
\begin{tikzcd}[ampersand replacement=\&]
	{\UU[1]} \&\&\& {\UU[2]} \&\& {\UU[1]} \&\& {\UU[2]} \\
	\\
	\&\&\& {\Psh{\UU[1]}} \&\&\&\& {\Psh{\UU[1]}}
	\arrow[""{name=0, anchor=center, inner sep=0}, "{\UU[2](F;\Id-)}"'{pos=0.3}, curve={height=24pt}, from=1-4, to=3-4]
	\arrow[""{name=1, anchor=center, inner sep=0}, "{\UU[2](G;\Id-)}", curve={height=-30pt}, from=1-4, to=3-4]
	\arrow[""{name=2, anchor=center, inner sep=0}, "\yoneda"', curve={height=18pt}, from=1-1, to=3-4]
	\arrow[""{name=3, anchor=center, inner sep=0}, "F", curve={height=-12pt}, from=1-1, to=1-4]
	\arrow[""{name=4, anchor=center, inner sep=0}, "F", curve={height=-12pt}, from=1-6, to=1-8]
	\arrow[""{name=5, anchor=center, inner sep=0}, "\yoneda"', from=1-6, to=3-8]
	\arrow["{\UU[2](G;\Id-)}", curve={height=-30pt}, from=1-8, to=3-8]
	\arrow[""{name=6, anchor=center, inner sep=0}, "G"', curve={height=12pt}, from=1-6, to=1-8]
	\arrow["{\UU[2](\alpha; \Id-)}", shift right=1, shorten <=11pt, shorten >=11pt, Rightarrow, from=0, to=1]
	\arrow["\alpha"', shorten <=3pt, shorten >=3pt, Rightarrow, from=6, to=4]
	\arrow["{\map G}"', curve={height=12pt}, shorten <=12pt, shorten >=4pt, Rightarrow, from=5, to=1-8]
	\arrow["{=}"{description, pos=0.7}, draw=none, from=1, to=5]
	\arrow["{\map F}", curve={height=-12pt}, shorten <=9pt, shorten >=9pt, Rightarrow, from=2, to=3]
\end{tikzcd}
}

\newquiver{hom profunctor general input}{
% https://q.uiver.app/?q=WzAsMyxbMCwwLCJcXFVVWzFdIl0sWzEsMCwiXFxVVVsyXSJdLFsyLDAsIlxcVVVbM10iXSxbMCwxLCJGIiwwLHsiY3VydmUiOi0yfV0sWzIsMSwiSCIsMix7ImN1cnZlIjoyfV0sWzAsMSwiRyIsMix7ImN1cnZlIjoyfV0sWzIsMSwiSyIsMCx7ImN1cnZlIjotMn1dLFs0LDYsIlxcZ2FtbWEiLDIseyJzaG9ydGVuIjp7InNvdXJjZSI6MjAsInRhcmdldCI6MjB9fV0sWzUsMywiXFxhbHBoYSIsMix7InNob3J0ZW4iOnsic291cmNlIjoyMCwidGFyZ2V0IjoyMH19XV0=
\begin{tikzcd}[ampersand replacement=\&]
	{\UU[1]} \& {\UU[2]} \& {\UU[3]}
	\arrow[""{name=0, anchor=center, inner sep=0}, "F", curve={height=-12pt}, from=1-1, to=1-2]
	\arrow[""{name=1, anchor=center, inner sep=0}, "H"', curve={height=12pt}, from=1-3, to=1-2]
	\arrow[""{name=2, anchor=center, inner sep=0}, "G"', curve={height=12pt}, from=1-1, to=1-2]
	\arrow[""{name=3, anchor=center, inner sep=0}, "K", curve={height=-12pt}, from=1-3, to=1-2]
	\arrow["\gamma"', shorten <=3pt, shorten >=3pt, Rightarrow, from=1, to=3]
	\arrow["\alpha"', shorten <=3pt, shorten >=3pt, Rightarrow, from=2, to=0]
\end{tikzcd}
}

\newquiver{hom profunctor general def via hcompose}{
% https://q.uiver.app/?q=WzAsNSxbOCwwLCJcXFBzaHtcXFVVWzFdfSJdLFs2LDAsIlxcVVVbMl0iXSxbNCwwLCJcXFVVWzNdIl0sWzAsMCwiXFxVVVszXSJdLFsyLDAsIlxcUHNoe1xcVVVbMV19Il0sWzEsMCwiXFxVVVsyXShGO1xcSWQtKSIsMCx7ImN1cnZlIjotMn1dLFsyLDEsIkgiLDAseyJjdXJ2ZSI6LTJ9XSxbMSwwLCJcXFVVWzJdKEc7XFxJZC0pIiwyLHsiY3VydmUiOjJ9XSxbMiwxLCJLIiwyLHsiY3VydmUiOjJ9XSxbMyw0LCJcXFVVWzJdKEY7SC0pIiwwLHsiY3VydmUiOi0zfV0sWzMsNCwiXFxVVVsyXShHO0stKSIsMix7ImN1cnZlIjozfV0sWzQsMiwiXFxkZWZpbmVkYnkiLDEseyJzdHlsZSI6eyJib2R5Ijp7Im5hbWUiOiJub25lIn0sImhlYWQiOnsibmFtZSI6Im5vbmUifX19XSxbNiw4LCJcXGdhbW1hIiwyLHsic2hvcnRlbiI6eyJzb3VyY2UiOjIwLCJ0YXJnZXQiOjIwfX1dLFs1LDcsIlxcYWxwaGEiLDAseyJzaG9ydGVuIjp7InNvdXJjZSI6MjAsInRhcmdldCI6MjB9fV0sWzksMTAsIlxcVVVbMl0oXFxhbHBoYTtcXGdhbW1hKSIsMCx7Im9mZnNldCI6NSwic2hvcnRlbiI6eyJzb3VyY2UiOjIwLCJ0YXJnZXQiOjIwfX1dXQ==
\begin{tikzcd}[ampersand replacement=\&]
	{\UU[3]} \&\& {\Psh{\UU[1]}} \&\& {\UU[3]} \&\& {\UU[2]} \&\& {\Psh{\UU[1]}}
	\arrow[""{name=0, anchor=center, inner sep=0}, "{\UU[2](F;\Id-)}", curve={height=-12pt}, from=1-7, to=1-9]
	\arrow[""{name=1, anchor=center, inner sep=0}, "H", curve={height=-12pt}, from=1-5, to=1-7]
	\arrow[""{name=2, anchor=center, inner sep=0}, "{\UU[2](G;\Id-)}"', curve={height=12pt}, from=1-7, to=1-9]
	\arrow[""{name=3, anchor=center, inner sep=0}, "K"', curve={height=12pt}, from=1-5, to=1-7]
	\arrow[""{name=4, anchor=center, inner sep=0}, "{\UU[2](F;H-)}", curve={height=-18pt}, from=1-1, to=1-3]
	\arrow[""{name=5, anchor=center, inner sep=0}, "{\UU[2](G;K-)}"', curve={height=18pt}, from=1-1, to=1-3]
	\arrow["\definedby"{description}, draw=none, from=1-3, to=1-5]
	\arrow["\gamma"', shorten <=3pt, shorten >=3pt, Rightarrow, from=1, to=3]
	\arrow["\alpha", shorten <=3pt, shorten >=3pt, Rightarrow, from=0, to=2]
	\arrow["{\UU[2](\alpha;\gamma)}", shift right=5, shorten <=5pt, shorten >=5pt, Rightarrow, from=4, to=5]
\end{tikzcd}
}
\newquiver{hom profunctor general input-output}{
% https://q.uiver.app/?q=WzAsNSxbMCwwLCJcXFVVWzFdIl0sWzEsMCwiXFxVVVsyXSJdLFsyLDAsIlxcVVVbM10iXSxbNCwwLCJcXFVVWzNdIl0sWzYsMCwiXFxQc2h7XFxVVVsxXX0iXSxbMCwxLCJGIiwwLHsiY3VydmUiOi0yfV0sWzIsMSwiSCIsMix7ImN1cnZlIjoyfV0sWzAsMSwiRyIsMix7ImN1cnZlIjoyfV0sWzIsMSwiSyIsMCx7ImN1cnZlIjotMn1dLFszLDQsIlxcVVVbMl0oRjtILSkiLDAseyJjdXJ2ZSI6LTN9XSxbMyw0LCJcXFVVWzJdKEc7Sy0pIiwyLHsiY3VydmUiOjN9XSxbMiwzLCJcXFVVWzJdKC07LSAtKSIsMCx7InNob3J0ZW4iOnsic291cmNlIjoxMCwidGFyZ2V0IjoxMH0sInN0eWxlIjp7InRhaWwiOnsibmFtZSI6Im1hcHMgdG8ifX19XSxbNiw4LCJcXGdhbW1hIiwyLHsic2hvcnRlbiI6eyJzb3VyY2UiOjIwLCJ0YXJnZXQiOjIwfX1dLFs3LDUsIlxcYWxwaGEiLDIseyJzaG9ydGVuIjp7InNvdXJjZSI6MjAsInRhcmdldCI6MjB9fV0sWzksMTAsIlxcVVVbMl0oXFxhbHBoYTtcXGdhbW1hKSIsMCx7Im9mZnNldCI6NSwic2hvcnRlbiI6eyJzb3VyY2UiOjIwLCJ0YXJnZXQiOjIwfX1dXQ==
\begin{tikzcd}[ampersand replacement=\&]
	{\UU[1]} \& {\UU[2]} \& {\UU[3]} \&\& {\UU[3]} \&\& {\Psh{\UU[1]}}
	\arrow[""{name=0, anchor=center, inner sep=0}, "F", curve={height=-12pt}, from=1-1, to=1-2]
	\arrow[""{name=1, anchor=center, inner sep=0}, "H"', curve={height=12pt}, from=1-3, to=1-2]
	\arrow[""{name=2, anchor=center, inner sep=0}, "G"', curve={height=12pt}, from=1-1, to=1-2]
	\arrow[""{name=3, anchor=center, inner sep=0}, "K", curve={height=-12pt}, from=1-3, to=1-2]
	\arrow[""{name=4, anchor=center, inner sep=0}, "{\UU[2](F;H-)}", curve={height=-18pt}, from=1-5, to=1-7]
	\arrow[""{name=5, anchor=center, inner sep=0}, "{\UU[2](G;K-)}"', curve={height=18pt}, from=1-5, to=1-7]
	\arrow["{\UU[2](-;- -)}", shorten <=4pt, shorten >=4pt, maps to, from=1-3, to=1-5]
	\arrow["\gamma"', shorten <=3pt, shorten >=3pt, Rightarrow, from=1, to=3]
	\arrow["\alpha"', shorten <=3pt, shorten >=3pt, Rightarrow, from=2, to=0]
	\arrow["{\UU[2](\alpha;\gamma)}", shift right=5, shorten <=5pt, shorten >=5pt, Rightarrow, from=4, to=5]
\end{tikzcd}
}

\newquiver{multi-natural transformation coop}{
% https://q.uiver.app/?q=WzAsMixbMCwwLCJcXFVVWzFdIl0sWzEsMCwiXFxVVVsyXSJdLFsxLDAsIkYiLDIseyJjdXJ2ZSI6Mn1dLFsxLDAsIkciLDAseyJjdXJ2ZSI6LTJ9XSxbMywyLCJcXGFscGhhIiwyLHsic2hvcnRlbiI6eyJzb3VyY2UiOjIwLCJ0YXJnZXQiOjIwfX1dXQ==
\begin{tikzcd}[ampersand replacement=\&]
	{\UU[1]} \& {\UU[2]}
	\arrow[""{name=0, anchor=center, inner sep=0}, "F"', curve={height=12pt}, from=1-2, to=1-1]
	\arrow[""{name=1, anchor=center, inner sep=0}, "G", curve={height=-12pt}, from=1-2, to=1-1]
	\arrow["\alpha"', shorten <=3pt, shorten >=3pt, Rightarrow, from=1, to=0]
\end{tikzcd}
}

\newquiver{psh 2-cell action type}{
  % https://q.uiver.app/?q=WzAsMixbMCwwLCJcXFBzaHtcXFVVWzFdfSJdLFsxLDAsIlxcUHNoe1xcVVVbMl19Il0sWzAsMSwiXFxQc2ggRiIsMCx7ImN1cnZlIjotMn1dLFswLDEsIlxcUHNoIEciLDIseyJjdXJ2ZSI6Mn1dLFsyLDMsIlxcUHNoXFxhbHBoYSIsMCx7InNob3J0ZW4iOnsic291cmNlIjoyMCwidGFyZ2V0IjoyMH19XV0=
\begin{tikzcd}[ampersand replacement=\&]
	{\Psh{\UU[1]}} \& {\Psh{\UU[2]}}
	\arrow[""{name=0, anchor=center, inner sep=0}, "{\Psh F}", curve={height=-12pt}, from=1-1, to=1-2]
	\arrow[""{name=1, anchor=center, inner sep=0}, "{\Psh G}"', curve={height=12pt}, from=1-1, to=1-2]
	\arrow["\Psh\alpha", shorten <=3pt, shorten >=3pt, Rightarrow, from=0, to=1]
\end{tikzcd}
}

\newquiver{iota type}{
% https://q.uiver.app/?q=WzAsMixbMCwwLCJcXFBzaHtcXFVVWzFdfSJdLFsyLDAsIlxcUHNoe1xcVVVbMV19Il0sWzAsMSwiXFxQc2h7XFxJZFt7XFxVVVsxXX1dfSIsMCx7ImN1cnZlIjotMn1dLFswLDEsIlxcSWRbXFxQc2h7XFxVVVsxXX1dIiwyLHsiY3VydmUiOjJ9XSxbMiwzLCJcXFlvbmVkYV97XFxVVVsxXX0iLDAseyJzaG9ydGVuIjp7InNvdXJjZSI6MjAsInRhcmdldCI6MjB9fV1d
\begin{tikzcd}[ampersand replacement=\&]
	{\Psh{\UU[1]}} \&\& {\Psh{\UU[1]}}
	\arrow[""{name=0, anchor=center, inner sep=0}, "{\Psh{\Id[{\UU[1]}]}}", curve={height=-12pt}, from=1-1, to=1-3]
	\arrow[""{name=1, anchor=center, inner sep=0}, "{\Id[\Psh{\UU[1]}]}"', curve={height=12pt}, from=1-1, to=1-3]
	\arrow["{\Yoneda_{\UU[1]}}", shorten <=3pt, shorten >=3pt, Rightarrow, from=0, to=1]
\end{tikzcd}
}

\newquiver{gamma type}{
% https://q.uiver.app/?q=WzAsMixbMCwwLCJcXFBzaHtcXFVVWzFdfSJdLFsyLDAsIlxcUHNoe1xcVVVbM119Il0sWzAsMSwiXFxQc2h7RiBcXGNvbXBvc2UgR30iLDAseyJjdXJ2ZSI6LTJ9XSxbMCwxLCJcXFBzaCBHIFxcY29tcG9zZSBcXFBzaCBGIiwyLHsiY3VydmUiOjJ9XSxbMiwzLCJcXGdhbW1hX3tGLEd9IiwwLHsic2hvcnRlbiI6eyJzb3VyY2UiOjIwLCJ0YXJnZXQiOjIwfX1dXQ==
\begin{tikzcd}[ampersand replacement=\&]
	{\Psh{\UU[1]}} \&\& {\Psh{\UU[3]}}
	\arrow[""{name=0, anchor=center, inner sep=0}, "{\Psh{F \compose G}}", curve={height=-12pt}, from=1-1, to=1-3]
	\arrow[""{name=1, anchor=center, inner sep=0}, "{\Psh G \compose \Psh F}"', curve={height=12pt}, from=1-1, to=1-3]
	\arrow["{\gamma_{F,G}}", shorten <=3pt, shorten >=3pt, Rightarrow, from=0, to=1]
\end{tikzcd}
}

\newquiver{theta type}{
% https://q.uiver.app/?q=WzAsMyxbMCwwLCJcXFVVWzFdIl0sWzEsMCwiXFxQc2h7XFxVVVsyXX0iXSxbMSwxLCJcXFBzaHtcXFVVWzNdfSJdLFswLDEsIlxcVVUoSCwtKSIsMCx7ImN1cnZlIjotMn1dLFsxLDIsIlxcUHNoIEciXSxbMCwyLCJcXFVVWzFdKEhcXGNvbXBvc2UgRywgLSkiLDIseyJjdXJ2ZSI6Mn1dLFs1LDEsIlxcdGhldGFfe0csSH0iLDAseyJsYWJlbF9wb3NpdGlvbiI6NzAsInNob3J0ZW4iOnsic291cmNlIjoyMH19XV0=
\begin{tikzcd}[ampersand replacement=\&]
	{\UU[1]} \& {\Psh{\UU[2]}} \\
	\& {\Psh{\UU[3]}}
	\arrow["{\UU(H,-)}", curve={height=-12pt}, from=1-1, to=1-2]
	\arrow["{\Psh G}", from=1-2, to=2-2]
	\arrow[""{name=0, anchor=center, inner sep=0}, "{\UU[1](H\compose G, -)}"', curve={height=12pt}, from=1-1, to=2-2]
	\arrow["{\theta_{G,H}}"{pos=0.7}, shorten <=5pt, Rightarrow, from=0, to=1-2]
\end{tikzcd}
}

\newquiver{theta UP}{
% https://q.uiver.app/?q=WzAsOCxbNSwwLCJcXFVVWzFdIl0sWzYsMCwiXFxQc2h7XFxVVVsyXX0iXSxbNiwxLCJcXFBzaHtcXFVVWzNdfSJdLFswLDAsIlxcVVVbM10iXSxbMSwwLCJcXFVVWzJdIl0sWzIsMCwiXFxVVVsxXSJdLFsyLDEsIlxcUHNoe1xcVVVbMl19Il0sWzIsMiwiXFxQc2h7XFxVVVszXX0iXSxbMSwyLCJcXFBzaCBHIl0sWzAsMiwiXFxVVVsxXShIXFxjb21wb3NlIEcsIC0pIiwyLHsiY3VydmUiOjJ9XSxbMCwxLCJcXFVVKEgsLSkiLDAseyJjdXJ2ZSI6LTJ9XSxbMyw0LCJHIl0sWzQsNSwiSCJdLFs0LDYsIlxceW9uZWRhIiwyLHsibGFiZWxfcG9zaXRpb24iOjQwLCJjdXJ2ZSI6Mn1dLFs1LDYsIlxcVVVbMV0oSCwtKSJdLFs2LDcsIlxcUHNoIEciXSxbMyw3LCJcXHlvbmVkYSIsMix7ImN1cnZlIjozfV0sWzksMSwiXFx0aGV0YV97RyxIfSIsMCx7ImxhYmVsX3Bvc2l0aW9uIjo3MCwic2hvcnRlbiI6eyJzb3VyY2UiOjIwfX1dLFsxMyw1LCJcXG1hcCBIIiwyLHsib2Zmc2V0IjotNCwic2hvcnRlbiI6eyJzb3VyY2UiOjIwLCJ0YXJnZXQiOjMwfX1dLFs2LDksIlxcbGVmdHJpZ2h0YXJyb3ciLDEseyJsYWJlbF9wb3NpdGlvbiI6MzAsImxldmVsIjoxLCJzdHlsZSI6eyJib2R5Ijp7Im5hbWUiOiJub25lIn0sImhlYWQiOnsibmFtZSI6Im5vbmUifX19XSxbMTYsMTMsIlxcbWFwIChcXFBzaHtcXFVVWzJdfSgtLEctKSkiLDIseyJsYWJlbF9wb3NpdGlvbiI6MTAsIm9mZnNldCI6LTUsImN1cnZlIjotMSwic2hvcnRlbiI6eyJzb3VyY2UiOjIwLCJ0YXJnZXQiOjMwfX1dXQ==
\begin{tikzcd}[ampersand replacement=\&]
	{\UU[3]} \& {\UU[2]} \& {\UU[1]} \&\&\& {\UU[1]} \& {\Psh{\UU[2]}} \\
	\&\& {\Psh{\UU[2]}} \&\&\&\& {\Psh{\UU[3]}} \\
	\&\& {\Psh{\UU[3]}}
	\arrow["{\Psh G}", from=1-7, to=2-7]
	\arrow[""{name=0, anchor=center, inner sep=0}, "{\UU[1](H\compose G, -)}"', curve={height=12pt}, from=1-6, to=2-7]
	\arrow["{\UU(H,-)}", curve={height=-12pt}, from=1-6, to=1-7]
	\arrow["G", from=1-1, to=1-2]
	\arrow["H", from=1-2, to=1-3]
	\arrow[""{name=1, anchor=center, inner sep=0}, "\yoneda"'{pos=0.4}, curve={height=12pt}, from=1-2, to=2-3]
	\arrow["{\UU[1](H,-)}", from=1-3, to=2-3]
	\arrow["{\Psh G}", from=2-3, to=3-3]
	\arrow[""{name=2, anchor=center, inner sep=0}, "\yoneda"', curve={height=18pt}, from=1-1, to=3-3]
	\arrow["{\theta_{G,H}}"{pos=0.7}, shorten <=5pt, Rightarrow, from=0, to=1-7]
	\arrow["{\map H}"', shift left=4, shorten <=5pt, shorten >=7pt, Rightarrow, from=1, to=1-3]
	\arrow["\leftrightarrow"{description, pos=0.3}, draw=none, from=2-3, to=0]
	\arrow["{\map (\Psh{\UU[2]}(-,G-))}"'{pos=0.1}, shift left=5, curve={height=-6pt}, shorten <=6pt, shorten >=8pt, Rightarrow, from=2, to=1]
\end{tikzcd}
}

\newquiver{axiom 3b proof goal}{
% https://q.uiver.app/?q=WzAsMTAsWzgsMCwiXFxVVVsxXSJdLFs5LDAsIlxcUHNoe1xcVVVbMl19Il0sWzksMSwiXFxQc2h7XFxVVVszXX0iXSxbMCwwLCJcXFVVWzNdIl0sWzEsMCwiXFxVVVsyXSJdLFsyLDAsIlxcVVVbMV0iXSxbMiwxLCJcXFBzaHtcXFVVWzJdfSJdLFsyLDIsIlxcUHNoe1xcVVVbM119Il0sWzUsMiwiXFxVVVszXSJdLFs1LDAsIlxcVVVbMl0iXSxbMSwyLCJcXFBzaCBHIl0sWzAsMiwiXFxVVVsxXShIXFxjb21wb3NlIEcsIC0pIiwyLHsibGFiZWxfcG9zaXRpb24iOjAsImN1cnZlIjoyfV0sWzAsMSwiXFxVVShILC0pIiwwLHsiY3VydmUiOi0yfV0sWzMsNCwiRyJdLFs0LDUsIkgiXSxbNCw2LCJcXHlvbmVkYSIsMix7ImxhYmVsX3Bvc2l0aW9uIjo0MCwiY3VydmUiOjJ9XSxbNSw2LCJcXFVVWzFdKEgsLSkiXSxbNiw3LCJcXFBzaCBHIl0sWzMsNywiXFx5b25lZGEiLDIseyJjdXJ2ZSI6M31dLFs4LDIsIiIsMix7ImN1cnZlIjoxfV0sWzgsOSwiRyIsMCx7ImN1cnZlIjotMX1dLFs5LDAsIkgiXSxbMTEsMSwiXFxpbnZcXFlvbmVkYSIsMCx7ImxhYmVsX3Bvc2l0aW9uIjo3MCwic2hvcnRlbiI6eyJzb3VyY2UiOjIwfX1dLFsxNSw1LCJcXG1hcCBIIiwyLHsib2Zmc2V0IjotNCwic2hvcnRlbiI6eyJzb3VyY2UiOjIwLCJ0YXJnZXQiOjMwfX1dLFs2LDExLCI9IiwxLHsibGFiZWxfcG9zaXRpb24iOjIwLCJsZXZlbCI6MSwic3R5bGUiOnsiYm9keSI6eyJuYW1lIjoibm9uZSJ9LCJoZWFkIjp7Im5hbWUiOiJub25lIn19fV0sWzE4LDE1LCJcXG1hcCAoXFxQc2h7XFxVVVsyXX0oLSxHLSkpIiwyLHsibGFiZWxfcG9zaXRpb24iOjEwLCJvZmZzZXQiOi01LCJjdXJ2ZSI6LTEsInNob3J0ZW4iOnsic291cmNlIjoyMCwidGFyZ2V0IjozMH19XSxbOCwxMSwiXFxtYXAgKEggXFxjb21wb3NlIEcpIiwwLHsiY3VydmUiOi00LCJzaG9ydGVuIjp7InRhcmdldCI6MjB9fV1d
\begin{tikzcd}[ampersand replacement=\&]
	{\UU[3]} \& {\UU[2]} \& {\UU[1]} \&\&\& {\UU[2]} \&\&\& {\UU[1]} \& {\Psh{\UU[2]}} \\
	\&\& {\Psh{\UU[2]}} \&\&\&\&\&\&\& {\Psh{\UU[3]}} \\
	\&\& {\Psh{\UU[3]}} \&\&\& {\UU[3]}
	\arrow["{\Psh G}", from=1-10, to=2-10]
	\arrow[""{name=0, anchor=center, inner sep=0}, "{\UU[1](H\compose G, -)}"'{pos=0}, curve={height=12pt}, from=1-9, to=2-10]
	\arrow["{\UU(H,-)}", curve={height=-12pt}, from=1-9, to=1-10]
	\arrow["G", from=1-1, to=1-2]
	\arrow["H", from=1-2, to=1-3]
	\arrow[""{name=1, anchor=center, inner sep=0}, "\yoneda"'{pos=0.4}, curve={height=12pt}, from=1-2, to=2-3]
	\arrow["{\UU[1](H,-)}", from=1-3, to=2-3]
	\arrow["{\Psh G}", from=2-3, to=3-3]
	\arrow[""{name=2, anchor=center, inner sep=0}, "\yoneda"', curve={height=18pt}, from=1-1, to=3-3]
	\arrow[curve={height=6pt}, from=3-6, to=2-10]
	\arrow["G", curve={height=-6pt}, from=3-6, to=1-6]
	\arrow["H", from=1-6, to=1-9]
	\arrow["\inv\Yoneda"{pos=0.7}, shorten <=5pt, Rightarrow, from=0, to=1-10]
	\arrow["{\map H}"', shift left=4, shorten <=5pt, shorten >=7pt, Rightarrow, from=1, to=1-3]
	\arrow["{=}"{description, pos=0.2}, draw=none, from=2-3, to=0]
	\arrow["{\map (\Psh{\UU[2]}(-,G-))}"'{pos=0.1}, shift left=5, curve={height=-6pt}, shorten <=6pt, shorten >=8pt, Rightarrow, from=2, to=1]
	\arrow["{\map (H \compose G)}", curve={height=-24pt}, shorten >=24pt, Rightarrow, from=3-6, to=0]
\end{tikzcd}
}

\newquiver{axiom 3* assumption}{
  % https://q.uiver.app/?q=WzAsNSxbMCwwLCJcXFVVWzJdIl0sWzAsMSwiXFxQc2hcXFVVIl0sWzIsMCwiXFxVVSJdLFs0LDAsIlxcVVVbMl0iXSxbNCwxLCJcXFBzaFxcVVUiXSxbMCwxLCJcXFVVWzJdKEYsLSkiLDIseyJjdXJ2ZSI6Mn1dLFswLDEsIkgiLDAseyJjdXJ2ZSI6LTJ9XSxbMiwzLCJGIl0sWzIsNCwiXFx5b25lZGEiLDIseyJjdXJ2ZSI6M31dLFszLDQsIkgiLDAseyJjdXJ2ZSI6LTN9XSxbMyw0LCJcXFVVWzJdKEYsLSkiLDIseyJsYWJlbF9wb3NpdGlvbiI6NzAsImN1cnZlIjoyfV0sWzUsNiwiXFxzaWdtYSIsMCx7InNob3J0ZW4iOnsic291cmNlIjoyMCwidGFyZ2V0IjoyMH19XSxbOCwzLCJcXG1hcCBGIiwwLHsibGFiZWxfcG9zaXRpb24iOjQwLCJvZmZzZXQiOjUsImN1cnZlIjotMywic2hvcnRlbiI6eyJzb3VyY2UiOjQwLCJ0YXJnZXQiOjMwfX1dLFsxMCw5LCJcXHNpZ21hIiwyLHsic2hvcnRlbiI6eyJzb3VyY2UiOjIwLCJ0YXJnZXQiOjIwfX1dXQ==
\begin{tikzcd}[ampersand replacement=\&]
	{\UU[2]} \&\& \UU \&\& {\UU[2]} \\
	\Psh\UU \&\&\&\& \Psh\UU
	\arrow[""{name=0, anchor=center, inner sep=0}, "{\UU[2](F,-)}"', curve={height=12pt}, from=1-1, to=2-1]
	\arrow[""{name=1, anchor=center, inner sep=0}, "H", curve={height=-12pt}, from=1-1, to=2-1]
	\arrow["F", from=1-3, to=1-5]
	\arrow[""{name=2, anchor=center, inner sep=0}, "\yoneda"', curve={height=18pt}, from=1-3, to=2-5]
	\arrow[""{name=3, anchor=center, inner sep=0}, "H", curve={height=-18pt}, from=1-5, to=2-5]
	\arrow[""{name=4, anchor=center, inner sep=0}, "{\UU[2](F,-)}"'{pos=0.7}, curve={height=12pt}, from=1-5, to=2-5]
	\arrow["\sigma", shorten <=5pt, shorten >=5pt, Rightarrow, from=0, to=1]
	\arrow["{\map F}"{pos=0.4}, shift right=5, curve={height=-18pt}, shorten <=19pt, shorten >=14pt, Rightarrow, from=2, to=1-5]
	\arrow["\sigma"', shorten <=6pt, shorten >=6pt, Rightarrow, from=4, to=3]
\end{tikzcd}
}

\newquiver{Axiom 3* counterexample goal}{
% https://q.uiver.app/?q=WzAsOCxbMCwwLCJcXFVVWzNdIl0sWzIsMCwiXFxVVVsyXSJdLFswLDIsIlxcVVVbMV0iXSxbMiwyLCJcXFBzaFxcVVUiXSxbNCwwLCJcXFVVWzNdIl0sWzYsMCwiXFxVVVsyXSJdLFs0LDIsIlxcVVVbMV0iXSxbNiwyLCJcXFBzaFxcVVUiXSxbMSwzLCJIIiwwLHsiY3VydmUiOi0yfV0sWzAsMiwiVyIsMix7ImN1cnZlIjoyfV0sWzIsMywiXFx5b25lZGEiLDIseyJjdXJ2ZSI6Mn1dLFswLDEsIkciXSxbMiwxLCJcXGFscGhhIiwwLHsiY3VydmUiOjIsInNob3J0ZW4iOnsic291cmNlIjoyMCwidGFyZ2V0IjoyMH0sImxldmVsIjoyfV0sWzQsNiwiVyIsMl0sWzYsNSwiRiIsMCx7ImN1cnZlIjotMn1dLFs2LDcsIlxceW9uZWRhIiwyLHsiY3VydmUiOjJ9XSxbNSw3LCJIIiwwLHsiY3VydmUiOi0yfV0sWzUsNywiXFxVVVsyXShGLC0pIiwyLHsib2Zmc2V0IjotMiwiY3VydmUiOjN9XSxbNCw1LCJHIl0sWzE1LDE3LCJcXG1hcCBGIiwwLHsib2Zmc2V0IjotNSwiY3VydmUiOjIsInNob3J0ZW4iOnsic291cmNlIjozMCwidGFyZ2V0IjozMH19XSxbMTcsMTYsIlxcc2lnbWEiLDAseyJzaG9ydGVuIjp7InNvdXJjZSI6MjAsInRhcmdldCI6MjB9fV0sWzEzLDE4LCJcXGJldGEiLDAseyJjdXJ2ZSI6LTEsInNob3J0ZW4iOnsic291cmNlIjoyMCwidGFyZ2V0IjoyMH19XV0=
\begin{tikzcd}[ampersand replacement=\&]
	{\UU[3]} \&\& {\UU[2]} \&\& {\UU[3]} \&\& {\UU[2]} \\
	\\
	{\UU[1]} \&\& \Psh\UU \&\& {\UU[1]} \&\& \Psh\UU
	\arrow["H", curve={height=-12pt}, from=1-3, to=3-3]
	\arrow["W"', curve={height=12pt}, from=1-1, to=3-1]
	\arrow["\yoneda"', curve={height=12pt}, from=3-1, to=3-3]
	\arrow["G", from=1-1, to=1-3]
	\arrow["\alpha", curve={height=12pt}, shorten <=13pt, shorten >=13pt, Rightarrow, from=3-1, to=1-3]
	\arrow[""{name=0, anchor=center, inner sep=0}, "W"', from=1-5, to=3-5]
	\arrow["F", curve={height=-12pt}, from=3-5, to=1-7]
	\arrow[""{name=1, anchor=center, inner sep=0}, "\yoneda"', curve={height=12pt}, from=3-5, to=3-7]
	\arrow[""{name=2, anchor=center, inner sep=0}, "H", curve={height=-12pt}, from=1-7, to=3-7]
	\arrow[""{name=3, anchor=center, inner sep=0}, "{\UU[2](F,-)}"', shift left=2, curve={height=18pt}, from=1-7, to=3-7]
	\arrow[""{name=4, anchor=center, inner sep=0}, "G", from=1-5, to=1-7]
	\arrow["{\map F}", shift left=5, curve={height=12pt}, shorten <=12pt, shorten >=12pt, Rightarrow, from=1, to=3]
	\arrow["\sigma", shorten <=5pt, shorten >=5pt, Rightarrow, from=3, to=2]
	\arrow["\beta", curve={height=-6pt}, shorten <=9pt, shorten >=9pt, Rightarrow, from=0, to=4]
\end{tikzcd}
}

\newquiver{adjunction type}{
  % https://q.uiver.app/?q=WzAsMixbMCwwLCJcXFVVWzFdIl0sWzEsMCwiXFxVVVsyXSJdLFswLDEsIkYiLDAseyJjdXJ2ZSI6LTJ9XSxbMSwwLCJVIiwwLHsiY3VydmUiOi0yfV0sWzIsMywiIiwwLHsibGV2ZWwiOjEsInN0eWxlIjp7Im5hbWUiOiJhZGp1bmN0aW9uIn19XV0=
\begin{tikzcd}[ampersand replacement=\&]
	{\UU[1]} \& {\UU[2]}
	\arrow[""{name=0, anchor=center, inner sep=0}, "F", curve={height=-12pt}, from=1-1, to=1-2]
	\arrow[""{name=1, anchor=center, inner sep=0}, "U", curve={height=-12pt}, from=1-2, to=1-1]
	\arrow["\dashv"{anchor=center, rotate=-90}, draw=none, from=0, to=1]
\end{tikzcd}
}

\newquiver{adjunction hom-set bijection}{
% https://q.uiver.app/?q=WzAsMyxbMCwxLCJcXFVVWzJdIl0sWzEsMCwiXFxVVVsxXSJdLFsyLDEsIlxcUHNoe1xcVVVbMV19Il0sWzEsMiwiXFxVVVsxXSgtLC0tKSIsMCx7ImN1cnZlIjotMX1dLFswLDEsIlUiLDAseyJjdXJ2ZSI6LTF9XSxbMCwyLCJcXFVVWzJdKEYsLSkiLDIseyJjdXJ2ZSI6MX1dLFs1LDEsIlxccGhpIiwwLHsibGFiZWxfcG9zaXRpb24iOjMwLCJjdXJ2ZSI6MSwic2hvcnRlbiI6eyJzb3VyY2UiOjIwfX1dLFs1LDEsIlxcaXNvbW9ycGhpYyIsMix7ImN1cnZlIjoxLCJzaG9ydGVuIjp7InNvdXJjZSI6MjB9LCJzdHlsZSI6eyJib2R5Ijp7Im5hbWUiOiJub25lIn0sImhlYWQiOnsibmFtZSI6Im5vbmUifX19XV0=
\begin{tikzcd}[ampersand replacement=\&]
	\& {\UU[1]} \\
	{\UU[2]} \&\& {\Psh{\UU[1]}}
	\arrow["{\UU[1](-,--)}", curve={height=-6pt}, from=1-2, to=2-3]
	\arrow["U", curve={height=-6pt}, from=2-1, to=1-2]
	\arrow[""{name=0, anchor=center, inner sep=0}, "{\UU[2](F,-)}"', curve={height=6pt}, from=2-1, to=2-3]
	\arrow["\phi"{pos=0.3}, curve={height=6pt}, shorten <=4pt, Rightarrow, from=0, to=1-2]
	\arrow["\isomorphic"', curve={height=6pt}, draw=none, from=0, to=1-2]
\end{tikzcd}}

\newquiver{hom-algebra carrier}{
  % https://q.uiver.app/?q=WzAsNCxbMCwwLCJcXGNhcnJpZXJ7QVxcbXVsdGltYXAgQn0iXSxbMSwwLCJcXGNhcnJpZXIgQl57XFxjYXJyaWVyIEF9Il0sWzIsMSwiKFRcXGNhcnJpZXIgQilee1RcXGNhcnJpZXIgQX0iXSxbMywwLCJcXGNhcnJpZXIgQl57VFxcY2FycmllciBBfSJdLFsxLDIsIlxcdW5kZXJsaW5lIFQiLDIseyJjdXJ2ZSI6Mn1dLFsyLDMsIkJcXHNlbS1ee1RcXGNhcnJpZXIgQX0iLDIseyJjdXJ2ZSI6Mn1dLFsxLDMsIlxcY2FycmllciBCXntBXFxzZW0tfSIsMCx7ImN1cnZlIjotNH1dLFswLDEsImUiXV0=
\begin{tikzcd}[ampersand replacement=\&]
	{\carrier{A\multimap B}} \& {\carrier B^{\carrier A}} \&\& {\carrier B^{T\carrier A}} \\
	\&\& {(T\carrier B)^{T\carrier A}}
	\arrow["{\underline T}"', curve={height=12pt}, from=1-2, to=2-3]
	\arrow["{B\sem-^{T\carrier A}}"', curve={height=12pt}, from=2-3, to=1-4]
	\arrow["{\carrier B^{A\sem-}}", curve={height=-24pt}, from=1-2, to=1-4]
	\arrow["e", from=1-1, to=1-2]
\end{tikzcd}
}

\newquiver{A multimap limit preservation diagram}{
  % https://q.uiver.app/?q=WzAsMixbMCwwLCJcXGNhcnJpZXIge0RpfV57XFxjYXJyaWVyIEF9Il0sWzIsMCwiXFxjYXJyaWVyIHtEaX1ee1RcXGNhcnJpZXIgQX0iXSxbMCwxLCJcXGNhcnJpZXIge0RpfV57QVxcc2VtLX0iLDAseyJjdXJ2ZSI6LTJ9XSxbMCwxLCJEaVxcc2VtLV57VFxcY2FycmllciBBfVxcY29tcG9zZSBcXHVuZGVybGluZSBUIiwyLHsiY3VydmUiOjJ9XV0=
\begin{tikzcd}[ampersand replacement=\&]
	{\carrier {Di}^{\carrier A}} \&\& {\carrier {Di}^{T\carrier A}}
	\arrow["{\carrier {Di}^{A\sem-}}", curve={height=-12pt}, from=1-1, to=1-3]
	\arrow["{Di\sem-^{T\carrier A}\compose \underline T}"', curve={height=12pt}, from=1-1, to=1-3]
\end{tikzcd}
}

\newquiver{limit Di first}{
  % https://q.uiver.app/?q=WzAsMixbMCwwLCJcXGNhcnJpZXIge0x9XntcXGNhcnJpZXIgQX0iXSxbMiwwLCJcXGNhcnJpZXIge0x9XntUXFxjYXJyaWVyIEF9Il0sWzAsMSwiXFxjYXJyaWVyIHtMfV57QVxcc2VtLX0iLDAseyJjdXJ2ZSI6LTJ9XSxbMCwxLCJMXFxzZW0tXntUXFxjYXJyaWVyIEF9XFxjb21wb3NlIFxcdW5kZXJsaW5lIFQiLDIseyJjdXJ2ZSI6Mn1dXQ==
\begin{tikzcd}[ampersand replacement=\&]
	{\carrier {L}^{\carrier A}} \&\& {\carrier {L}^{T\carrier A}}
	\arrow["{\carrier {L}^{A\sem-}}", curve={height=-12pt}, from=1-1, to=1-3]
	\arrow["{L\sem-^{T\carrier A}\compose \underline T}"', curve={height=12pt}, from=1-1, to=1-3]
\end{tikzcd}
}

\newcommand\CAT{\Category{CAT}}
\newcommand\ppCAT{\Category{CAT}_{\times}}
\newcommand\LocallySmall[1]{{#1}_{\mathcal L}}
\newcommand\MultiCAT{\Category{MultiCAT}}
\newcommand\AdMultiCAT{\Category{MultiCAT}_{\mathrm{admissible}}}
\newcommand\LargeCatName[1]{\mathsf{#1}}
\newcommand\CC[1][1]{\LargeCatName{\Cname[#1]}}
\newcommand\FF[1][1]{\LargeCatName{\Fname[#1]}}
\newcommand\UU[1][1]{\mathcal{\Uname[#1]}}
\newcommand\Field{\mathbb{F}}
\newcommand\Multi{\mathbf{Multi}}
\newcommand\ordinary[1]{#1_{\mathbf{o}}}
\newcommand\OO{\mathcal O}
\newcommand\TT{\mathcal T}
\newcommand\MM{\mathcal M}
\newcommand\DD{\mathcal A}
\newcommand\Mod{\mathbf{Mod}}
\newcommand\Mon{\mathbf{Mon}}
\newcommand\str{\mathrm{str}}
\newcommand\List{\mathop{\mathsf{List}}\nolimits}
\newcommand\Tree{\mathop{\mathsf{Tree}}\nolimits}
\newcommand\Branch{\mathop{\mathsf{Branch}}\nolimits}
\newcommand\Partit{\mathop{\mathsf{Partition}}\nolimits}
\newcommand\abs[1]{\left\lvert#1\right\rvert}
\newcommand\norm[1]{\left\lVert#1\right\rVert}
\newcommand\floor[1]{\left\lfloor#1\right\rfloor}
\newmvar{\partit}{\pi\or\sigma\or\tau}
\WithSuffix\newcommand\++{\mathbin{\mathord +\!\!\!\mathord +}}
\newcommand\mvec\mathbf
\newmvar{\xvec}{\mathbf{x}}
\newmvar{\Xvec}{\mathbf{X}}
\newmvar{\yvec}{\mathbf{y}}
\newmvar{\zvec}{\mathbf{z}}
\newmvar{\avec}{\mathbf{a}}
\newmvar{\bvec}{\mathbf{b}}
\newmvar{\cvec}{\mathbf{c}}
\newmvar{\dvec}{\mathbf{d}}
\newmvar{\evec}{\mathbf{e}}
\newmvar{\fvec}{\mathbf{f}}
\newmvar{\gvec}{\mathbf{g}}
\newmvar{\Fvec}{\mathbf{F}}
\newmvar{\Gvec}{\mathbf{G}}
\newmvar{\Hvec}{\mathbf{H}}
\newmvar{\Kvec}{\mathbf{K}}
\newmvar{\Pvec}{\mathbf{P}}
\newmvar{\uvec}{\mathbf{u}}
\newmvar{\vvec}{\mathbf{v}}
\newmvar{\svec}{\mathbf{s}}
\newmvar{\xivec}{\boldsymbol{\chi}}
\newcommand\concat{\mathop{\mathchoice%
  {\bigplus\mspace{-17mu}\bigplus}%
  {\bigplus\mspace{-15mu}\bigplus}%
  {\bigplus\mspace{-13mu}\bigplus}%
  {\bigplus\mspace{-11mu}\bigplus}%
  }%
}
\newcommand\Psh\widehat
\WithSuffix\newcommand\Psh'[1]{\Psh{#1}'}
% an interruptible version of seq (supports &)
\newcommand\tabmark{&}
\newcommand\iseq[1]{
  \left\langle\vphantom{\renewcommand\tabmark{}#1}\right.\!\!%
  #1%
  \!\!\left.\vphantom{\renewcommand\tabmark{}#1}\right\rangle%
  }

\WithSuffix\newcommand\seq*[1]{
  \langle#1\rangle}
\renewcommand\yoneda{\boldsymbol{\mathsf{y}}}
\newcommand\Yoneda{\boldsymbol\Upsilon}
\newcommand\map{\mathop{\mathsf{map}}}
\newcommand\coop[1]{#1^{\mathrm{coop}}}
\newcommand\Trans{\mathfrak T}
\newcommand\surtrans{\mathrel{\raisebox{.5ex}{\begin{tikzpicture}
% arrows.meta
    \draw [-{stealth}{stealth}{stealth}] (0,0)--(.5,0);
\end{tikzpicture}}}
}
\newcommand\opsurtrans{\mathrel{\raisebox{.5ex}{\begin{tikzpicture}
% arrows.meta
    \draw [{stealth}{stealth}{stealth}-] (0,0)--(.5,0);
\end{tikzpicture}}}
}


\newcommand\ren\rho
\newcommand\Sig[1][]{\Sigma_{#1}}
\newcommand\types\vdash
\newcommand\dist{\mathbin\triangleright}
\newcommand\Ax[1][]{\mathrm{Ax}_{#1}}
\newcommand\Theory{\Category}
\newcommand\rev{\mathrm{rev}}
\newcommand\perm\pi
\newcommand\Sym[1]{\mathbf S_{#1}}
\newcommand\lp{\left(}
\newcommand\rp{\right)}

\begin{document}
\maketitle
Started: 28 Nov 2022
\begin{abstract}
  A constructive proof of the construction of free algebras in semi-ring varieties.
\end{abstract}

\begin{itemize}
  \item \secref{sec:2-cat} and \secref{sec:multicat} define the 2-categorical and
  multicategorical notions we need, mostly to fix definitions and notations.
  \item \secref{sec:operads} defines the notion of operad and a 2-functor used in the
  construction of free algebras (\propositionref{Mod(O,-) is a 2-functor}).
  \item \secref{sec:multihom} defines the notion of multi-homomorphism and a multi-category
  and multi-functor used in the construction of free algebras
  (\propositionref{Multi(TT,C) multi-category}).
  \item \secref{sec:yoneda} gives a multi-categorical representability theorem for
  adjunctions using the Yoneda structure (\propositionref{representability multi-categorical adjunctions}).
  \item \secref{sec:multihom-extension} proves the existence of a left adjoint to a multi-functor under
  certain conditions (\theoremref{left adjoint iff commutative}).
  \item \secref{sec:distributivity} defines the concept of translation and of distributive
  combination, and proves that the category of models of a distributive combination is a pullback
  (\propositionref{Mod(Trans dist A) is a pullback}).
  \item \secref{sec:fral} constructs the free algebra for the distributive combination
  of two theories and proves that it is indeed free (\lemmaref{A is free}).
\end{itemize}

\section{2-categories}
\label{sec:2-cat}

We'll need basic concepts in higher-dimensional category theory.

A \emph{(strict) 2-category} $\CC$ consists of the following data and
structure:
\begin{itemize}
\item Cellular structure:
\begin{itemize}
  \item A collection (set, class, etc.) $\CC_0$ of $0$-cells $a, b, c,
    \ldots$;
  \item A collection (set, class, etc.) $\CC_1(a,b)$ of $1$-cells
    $f,g,h,\ldots : a \to b$ for every pair of 0-cells $a,b \in \CC_0$;
  \item A collection (set, class, etc.) $\CC_2(f,g)$ of $2$-cells
    $\alpha, \beta, \gamma, \ldots : f \Rightarrow g$ for every pair
    parallel of 1-cells $f,g : a \rightrightarrows b$.
    We'll depict $2$-cells with oriented faces:
    $\quiver{alpha f g}$.

\end{itemize}
\item Identities:
  \begin{itemize}
  \item A $1$-cell $\id[a] : a \to a$, the \emph{identity} $1$-cell,
    for every $0$-cell $a \in \CC_0$.
  \item A $2$-cell $\id[f] : f \To f$, the \emph{identity} $2$-cell,
    for every $1$-cell $f : a \to b$. We depict identities by marking
    the face with an equality symbol: \quiver{identity 2 cell}
  \end{itemize}
\item Composition:
  \begin{itemize}
  \item A \emph{composition} operation on $1$-cells, assigning a
    $1$-cell $g \compose f$ for every pair of composeable $1$-cells $a
    \xto{f} b \xto{g} c$.
  \item A \emph{vertical composition} operation on $2$-cells,
    assigning a $2$-cell $\beta \vcompose \alpha : f \To h$ for every
    pair of composeable $2$-cells $f \xto{\alpha} g \xto{\beta} h : a
    \to b$, which we depict it by pasting the faces along their shared
    $1$-cell:
    \[
    \quiver{vertical composition type} = \quiver{vertical composition}
    \]
  \item A \emph{horizontal composition} operation on $2$-cells,
    assigning a $2$-cell $\beta \hcompose \alpha : g_1\compose f_1 \To
    g_2 \compose f_2$ for every pair of parallel $2$-cells $a \xto{f_1
      \xto{\alpha} g_1} b \xto{f_2 \xto{\beta} g_2} c$. We depict it
    by pasting the faces along the shared vertex:
    \[
      \quiver{horizontal composition type} = \quiver{horizontal composition}
    \]
  \end{itemize}
\end{itemize}
such that:
\begin{itemize}
\item $1$-cell composition is associative and identities are neutral:
  \[
  h \compose (g \compose f) = (h \compose g) \compose f
  \qquad
  \id[b] \compose f = f = f \compose \id[a]
  \]
\item $2$-cell vertical composition is associative and identities are neutral:
  \[
  \gamma \vcompose (\beta \vcompose \alpha) = (\gamma \vcompose \beta) \compose \alpha
  \qquad
  \id[g] \compose \alpha = \alpha = \alpha \compose \id[f]
  \]
\item $2$-cell horizontal composition is associative and identities
  between identities are neutral, and composition of $2$-cell
  identities preserves $1$-cell composition:
  \[
  \gamma \hcompose (\beta \hcompose \alpha) = (\gamma \hcompose \beta) \hcompose \alpha
  \qquad
  \id[{\id[b]}] \hcompose \alpha = \alpha = \alpha \hcompose \id[{\id[a]}]
  \qquad
  \id[g] \hcompose \id[f] = \id[g \compose f]
  \]
\item The interchange law holds:
  \begin{gather*}
  (\beta_2 \vcompose \alpha_2) \hcompose (\beta_1 \vcompose \alpha_1)
  =
  (\beta_2 \hcompose \beta_1) \vcompose (\alpha_2 \vcompose \alpha_1)
  \\
  \quiver{interchange law vertical}
  =
  \quiver{interchange law horizontal}
  \end{gather*}
  This law means we can depict compositions of $2$-cells using
  commutative diagrams, since all decompositions of the diagram into
  pasting expressions are equal $2$-cells. For example, we will depict
  both expressions in the interchange law by this pasting diagram:
  \quiver{interchange law}
\end{itemize}
\begin{example}
  The paradigmatic example is the $2$-category $\CAT$ that has:
  \begin{itemize}
  \item potentially large categories as $0$-cells;
  \item functors between them as $1$-cells;
  \item natural transormations between them as $2$-cells;
  \item pointwise composition of functors as $1$-cell composition;
  \item pointwise composition of natural transformations as horizontal
    composition;
  \item horizontal composition of natural transformations given, for
    $\quiver{cat horizontal composition type}$ by:
    $\begin{aligned}
      (\beta \hcompose \alpha)_a
          &{}= \beta_{G_1a} \compose F_2\alpha_{a}
        \\&{}= G_2\alpha_a \compose \beta_{F_1}a
    \end{aligned}$
    i.e., by the naturality square:
    $\quiver{cat horizontal composition commutative}$.
  \end{itemize}
  The required properties are well-known~\citep{mac-lane:cwm}.
\end{example}

An \emph{adjunction} in a $2$-category $\CC$ consists of:
\begin{itemize}
\item two $1$-cells $a \xto{f} b \xto{u} a$ called the left ($f$) and
  right ($u$) adjoints;
\item two $2$-cells $u \compose f \xto{\unit} \id[a]$ and $\id[b]
  \xto{\counit} u \compose f$ called the unit ($\unit$) and counit
  ($\counit$).
\end{itemize}
such that:
\begin{itemize}
\item $\quiver{adjoint triangle one} = \id[f]$ and $\quiver{adjoint triangle two} = \id[u]$
\end{itemize}

\begin{example}
  An adjunction in $\CAT$ is the ordinary notion of an adjunction.
\end{example}


A \emph{$2$-functor} $\FF : \CC[1] \to \CC[2]$ consists of three mappings:
\begin{itemize}
  \item $\FF_0 : \CC[1]_0 \to \CC[2]_0$ sending
    $0$-cells $a$ in $\CC[1]$ to $0$-cells $\FF_0a$ in $\CC[2]$.
  \item $\FF_1 : \CC[1]_1(a,b) \to \CC[2]_1(\FF_0a,
    \FF_0b)$ sending $1$-cells $f : a \to b$ in $\CC[1]$ to $1$-cells
    $\FF_1f : \FF_0a \to \FF_0b$ in $\CC[2]$.
  \item $\FF_2 : \CC[1]_2(f,g) \to \CC[2]_2(\FF_1f, \FF_1g)$ sending
    $2$-cells $\!\!\quiver{narrow alpha f g}\!\!$ in $\CC[1]$ to $2$-cells
    $\quiver{F(alpha f g)}$ in~$\CC[2]$.
\end{itemize}
such that $\FF$ preserves all identities and compositions.

The point of $2$-functors is that they transport pasting diagrams to
pasting diagrams, while preserving the identities.
Our main usage of $2$-functors is the following lemma that lets us
transport adjunctions:
\begin{lemma}
  Let $\FF : \CC[1] \to \CC[2]$ be a $2$-functor. If $\seq{f : a \to
    b,u,\unit,\counit}$ is an adjunction in $\CC[1]$ then $\seq{\FF f
    : \FF a \to \FF b, \FF u, \FF\unit, \FF\counit}$ is an adjunction
  in $\CC[2]$.
\end{lemma}
\begin{proof}
  Type-checking, the candidate adjunction has the right type. Since
  $2$-functors preserve pasting diagrams and identities:
  \[
  \quiver{F(adjoint triangle one)}
  = \FF\id[f] = \id[\FF f]
  \
  \quiver{F(adjoint triangle two)}
  = \FF\id[u] = \id[\FF u]
  \]
  and so this candidate is in fact an adjunction.
\end{proof}

We will also observe several (strict) natural $2$-transformations.
Let $\FF[1],\FF[2] : \CC[1] \to \CC[2]$ be two $2$-functors between
$2$-categories. A \emph{(strict) natural $2$-transformation}
$\quiver{natural 2-transformation type}$ consists of:
\begin{itemize}
\item A $1$-cell $\Xi_a : \FF[1]_0 a \to \FF[2]_0 a$, for each $0$-cell $a$;
\end{itemize}
such that:
\begin{itemize}
\item For every $1$-cell $f : a \to b$ in $\CC[1]$:
  \[
  \quiver{strict natural 2-transformation: 1-cells}
  \]
\item For every $2$-cell $\quiver{alpha f g}$:
  \[
  \quiver{strict natural 2-transformation: 2-cells}
  \]
\end{itemize}

\section{Multicategories}
\label{sec:multicat}

A \emph{multicategory} $\UU$ consists of:
\begin{itemize}
\item A collection (set, class, etc.) $\Obj \UU$ of \emph{objects}
  $a,b,c,\ldots$.
\item A collection (set, class, etc.) $\UU(a_1, \ldots, a_n; b)$ of
  \emph{multi-morphisms}, for every (possibly empty) list of
  \emph{source} objects $a_1, \ldots, a_n$ and \emph{target} object
  $b$, in which case we write $f : (a_1, \ldots, a_n) \to b$ and
  depict this multi-morphisms like so:
  \[
  \quiver{multi-morphism}
  \] The
  \emph{morphisms} are those multi-morphisms whose source is of length
  $1$.
\item An \emph{identity} morphism $\id[a] : (a) \to a$ for each object $a$.

  We'll depict morphisms by single-source arrows and identity
  morphisms by elongated equality symbols, as in ordinary commutative
  diagrams.

\item The \emph{composition} operation, defined for every
  multi-morphism $g : (b_1, \ldots, b_n) \to c$ and compatible list of
  multi-morphisms $f_i : (a_{i1},\ldots, a_{im_i}) \to b_i$, for each
  $1 \leq i \leq n$, yielding a multi-morphism $g \compose (f_1,
  \ldots, f_n) : (a_{11}, \ldots, a_{1m_1}, \ldots, a_{n1}, \ldots,
  a_{nm_n}) \to c$. We depict the composed morphism by pasting the
  shared objects:
  \[
  \quiver{multi composition}
  \]
\end{itemize}
such that:
\begin{itemize}
\item identities are neutral, i.e., for all $f : (a_1, \ldots, a_n) \to b$:
  \[
  f \compose (\id[a_1], \ldots, \id[a_n]) = f = \id[b] \compose f
  \]
\item composition is associative:
  \begin{multline*}
  h \compose (g_1 \compose (f_{11}, \ldots, f_{1m_1}), \ldots, g_n\compose (f_{n1}, \ldots, f_{nm_n})) \\=
  (h \compose (g_1, \ldots, g_n)) \compose (f_{11}, \ldots, f_{nm_1}, \ldots, f_{n1}, \ldots, f_{nm_n})
  \end{multline*}
  Associativity means that, like for ordinary commutative diagrams, we
  may decompose a multi-morphism diagram in any several equivalent
  ways. For example, we depict the multi morphism of the associativity
  equation like so:
  \[
  \quiver{multi composition associativity}
  \]
\end{itemize}
The paradigmatic example is multi-linear maps:
% \begin{example}
%   Given a field $\Field$, we define the following multicategory
%   $\Multi\Vect_{\Field}$:
%   \begin{itemize}
%   \item Objects are vector spaces $U, V, W, \ldots$ over $\Field$.
%   \item Multi-morphisms $T : (U_1, \ldots, U_n) \to V$ are the
%     \emph{multi-linear} forms/transformations, i.e., the functions $T
%     : U_1 \times \cdots \times U_n \to V$ that are linear in each
%     component:
%     \begin{multline*}
%     T(u_1, \ldots, u_{i-1}, \lambda u_i + \lambda' u'_i, u_{i+1}, \ldots, u_n)
%     \\=
%     \lambda T(u_1, \ldots, u_{i-1}, u_i, u_{i+1}, \ldots, u_n)
%     +
%     \lambda'T(u_1, \ldots, u_{i-1}, u'_i, u_{i+1}, \ldots, u_n)
%     \end{multline*}
%   \end{itemize}
%   For example, the inner product over the field of reals is a
%   bi-linear transformation.

%   We compose multi-linear functions in the evident way:
%   \begin{multline*}
%     g \compose (f_1, \ldots, f_n) (u_{11}, \ldots, u_{1m_1}, \ldots, u_{n1}, \ldots, u_{nm_n})
%     \\:=
%   g(f_1(u_{11}, \ldots, u_{1m_1}), \ldots, f_n(u_{n1}, \ldots, u_{nm_n}))
%   \end{multline*}
%   and the morphisms in $\Multi\Vect_{\Field}$ are the linear
%   transformations, including the identity functions as identities.
% \end{example}

\begin{example}\examplelabel{multi-cat from chosen products}
  Let $\pair{\C}{\prod}$ be a category with specified finite
  products. We equip it with a multicategory structure as follows:
  \begin{itemize}
  \item The objects are the same, $\Obj{\pair\C\prod} := \Obj\C$.
  \item A multi-morphism $f : (a_1, \ldots, a_n) \to b$ is any
    morphism $f : \prod_{i = 1}^n a_i \to b$.
  \item The identities rely on the canonical isomorphism:
    \[
    \id[a] := \prod_{i=1}^1 a \isomorphic a
    \]
  \item Similarly, composition relies on a canonical
    isomorphism. Letting:
    \begin{itemize}
    \item $[n] \definedby \set{1, \ldots, n}$ be the
    enumeration of the index set;
  \item $p := \sum_{i = 1}^n m_i$ the total number of objects in the source of the comosition morphism; and
  \item $\phi : [p] \xto{\isomorphic} \coprod_{i=1}^n[m_i]$ be an
    enumeration of the nested set of indices that's compatible with
    the source list of the composition;
    \end{itemize}
    define:
    \[
    g \compose (f_1, \ldots, f_n) : \prod_{\ell = 1}^{p}a_{\phi \ell}
    \xto{\isomorphic} \prod_{i = 1}^n\prod_{j = 1}^{m_i} a_{i,j}
    \xto{\prod_{i}f_i} \prod_{i=1}^n b_i \xto{g} c
    \]
  \item Tedious verification of the axioms.
  \end{itemize}
\end{example}
Here is a slightly different way to construct a similar multicategory,
but this time without postulating a choice of finite products. Not
sure it's a good way to do it.
\begin{example}\examplelabel{fin prod multi cat}
Let $\C$ be a category that merely has finite products.
\begin{itemize}
\item Objects are the same.
\item A multi-morphism $f : (a_1, \ldots, a_n) \to b$ consists of a
  chosen product $\Dom f$ for $\pi_{i = 1}^na_i$ together with a
  morphism $f : \Dom f \to b$. For identities $\id[a]$ we can then
  take the chosen product $\prod_{i=1}^1$ to be $a$ itself together
  with the identity, and for composition we can take
  $\prod_{i=1}^n\prod_{j=1}^{m_i} a_i$, although here we need to
  employ some choice for the outer product.
\end{itemize}
\end{example}
Generalising the example $\pair\C\prod$:
% \begin{example}
%   Let $\C = \seq{\carrier\C,(\otimes),I,a,l,r}$ be a monoidal
%   category. We equip it with a multi-category structure:
%   \begin{itemize}
%   \item The objects are the same $\Obj\C := \Obj{\carrier\C}$.
%   \item A multi-morphism $f : (a_1, \ldots, a_n) \to b$ is a morphism
%     $f : a_1 \otimes \cdots \otimes a_n \to b$ (say right-nesting the
%     monoidal product).
%   \end{itemize}
%   with the remainder following the same principles as the finite
%   product example.
% \end{example}

A \emph{multi-functor} $F : \UU[1] \to \UU[2]$ consists of two mappings:
\begin{itemize}
\item $F : \Obj{\UU[1]} \to \Obj{\UU[2]}$ sending $\UU[1]$ objects to $\UU[2]$ objects; and
\item $F : \UU[1](a_1, \ldots, a_n; b) \to \UU[2](Fa_1, \ldots, Fa_n;
  Fb)$ sending a multi-morphism $f : (a_1, \ldots, a_n) \to b$ to a
  multi-morphism $Ff : (Fa_1, \ldots, Fa_n) \to Fb$.
\end{itemize}
such that $F$ preserves:
\begin{itemize}
\item identities: $F\id[a] = \id[Fa] : (Fa) \to a$; and
\item composition: $F(g \compose (f_1, \ldots, f_n)) = F g \compose
  (Ff_1, \ldots Ff_n)$.
\end{itemize}
As usual, multi-functors are defined to preserve the structure of
diagrams with multi-morphisms.
\begin{example}
  The identity multi-functor $\Id : \UU \to \UU$ sends every object
  and multi-morphism to themselves. The composition of multi-functors
  is a multi-functor.
\end{example}

A \emph{multi-natural transformation} $\quiver{multi-natural
  transformation type}$ between multi-functors consists of:
\begin{itemize}
\item a morphism $\alpha_a : Fa \to Ga$ for every object $a \in \Obj\UU$
\end{itemize}
such that:
\begin{itemize}
\item
  \[
  \quiver{multi-naturality}
  \]
\end{itemize}

We can define horizontal and vertical composition of multi-natural
transformations analogously to their category theoretic way:
\[
(\beta\vcompose\alpha)_a := \beta_a \compose (\alpha_a) \qquad
(\beta\hcompose\alpha)_a := \beta_{G_1a}\compose (F_2\alpha_a)
                          =G_2\alpha_a\compose (\beta_{F_1a})
\]
and they are multi-natural for analogous reasons. As a consequence we
have a $2$-category $\MultiCAT$ where:
\begin{itemize}
\item $0$-cells are multi-categories;
\item $1$-cells are multi-functors; and
\item $2$-cells are multi-natural transformations
\item with multi-functor composition and identities:
  \[
  \Id[{\CC}] a := a \qquad \Id[{\CC}]f := f
  \]
  vertical, horizontal composition of multi-natural transformations,
  with identities:
  \[
  (\id[{F}])_a := \id[{Fa}] : (a) \to a
  \]

\end{itemize}

For example, we have a $2$-functor structure $\ordinary- : \MultiCAT \to
\CAT$:
\begin{itemize}
\item sending each multicategory $\UU$ to its (\textbf{o}rdinary)
  \emph{category of morphisms}:
  \begin{itemize}
  \item $\Obj {\ordinary\UU} := \Obj \UU$;
  \item $\ordinary\UU(a,b) := \UU(a;b)$, i.e., the morphisms in
    $\ordinary\UU$ are the multi-morphisms with a length-$1$ source.
  \item identities and composition are given by identity morphisms and
    composition in the multi-category.
  \end{itemize}
\item sending each multi-functor $F : \UU[1] \to \UU[2]$ to the
  functor:
  \begin{itemize}
  \item sending $a \in \Obj{\UU[1]}$ to $Fa \in \Obj{\UU[2]}$;
  \item sending $f : (a) \to b$ to $\ordinary Ff := Ff : (Fa) \to Fb$.
  \end{itemize}
\item sending each multi-natural transformation $\quiver{multi-natural
  transformation type}$ to $(\ordinary\alpha)_a := \alpha_a$.
\end{itemize}
\begin{proposition}
  The above assignment is a $2$-functor $\ordinary- : \MultiCAT \to
  \CAT$.

  It is $2$-faithful in the following sense: given two parallel
  multi-natural transformations $\quiver{parallel multi-natural}$, we have
  \(
  \ordinary \alpha = \ordinary \beta \implies \alpha = \beta
  \) and so $\ordinary-$ reflects adjunctions.
\end{proposition}
\begin{proof}
  To show that the $0$-cell mapping is well-defined, we need to show
  we have a category. Composition of morphisms $a \xto f b \xto g c$
  in $\ordinary\UU$ is their composition in the multi-category, i.e.,
  given by \( g \compose f := g \compose (f) \). The multi-category
  axioms then specialise to:
  \[
  \id \compose (f) = f = f \compose (\id)
  \qquad
  h \compose (g \compose (f)) = (h \compose (g)) \compose (f)
  \]
  and so we have a category.

  To show that the $1$-cell mapping is functorial, note that the
  construction type-checks, hence we have a functor structure, and
  then immediately validate functoriality:
  \[
  \ordinary F (\id) = \id
  \qquad
  \ordinary F(g \compose (f)) = F g \compose (F f)
  \]

  For $2$-cells, the resulting transformation $\ordinary\alpha :
  \ordinary F \to \ordinary G$ is natural since the multi-naturality
  condition, when applied to a morphism, recovers the naturality
  condition for the ordinary transformation.

  The $2$-faithfulness follows by observing that the data for the
  multi-natural transformation is preserved by
  $\ordinary-$. Explicitly, if $\ordinary\alpha = \ordinary\beta$,
  then for each $a \in \Obj\UU$, we have $\alpha_a = \beta_a$, and
  that means that $\alpha = \beta$.
\end{proof}

\section{Operads}
\label{sec:operads}

An \emph{operad} $\OO$ consists of:
\begin{itemize}
\item a set $\OO_n$, whose elements we call \emph{$n$-ary operations},
  for each natural $n \in \naturals$;
\item an identity operation $\id \in \OO_1$; and
\item a \emph{composed} operation $g \compose (f_1, \ldots, f_n) \in
  \OO_{\sum_i n_i}$, for every operation composable pair of operation
  $g \in \OO_n$ and sequence of operations $f_i \in \OO_{n_i}$
\end{itemize}
such that composition is associative and the identity is neutral:
\begin{multline*}
h \compose (g_1(f_{11}, \ldots, f_{1m_1}), \ldots, g_n(f_{n1}, \ldots, f_{nm_n}))
\\\qquad=
(h \compose (g_1, \ldots, g_n))\compose(f_{11}, \ldots, f_{1m_1},
\ldots, f_{n1}, \ldots, f_{nm_n})
\qquad
f \compose (\id, \ldots, \id) = f = \id \compose (f)
\end{multline*}
An operad is then just a $1$-object multi-category $\OO$, by
identifying the $n$-ary operations $\OO_n$ with the $n$-ary
multi-morphisms $\OO(\underbrace{\star, \ldots, \star}_{\text{$n$
    times}}; \star)$.

% \begin{example}\examplelabel{monoid operad}
%   Let $\Monoid$ be the following operad:
%   \begin{itemize}
%   \item Every set of operations $\Monoid_n$ has a unique element
%     $(\bigodot_{i = 1}^n)$.
%   \end{itemize}
%   Since there is only one operation possible at each arity, there is
%   only one choice for identities and compositions, and the operad
%   axioms hold.
% \end{example}
The notion of a model may give more intuition about the definition of
operads.  Let $\OO$ be an operad and $\UU$ a multi-category. A
\emph{model} $\A$ for $\OO$ in $\UU$ consists of:
\begin{itemize}
\item an object $\carrier\A \in \Obj\UU$
\item for each operation $f \in \OO_n$, a multi-morphism
  $\A\sem f : (\overbrace{\carrier\A, \ldots, \carrier \A}^{\text{$n$ times}}) \to \carrier\A$
\end{itemize}
such that composition and identities are preserved:
\[
\A\sem{\id} = \id[\carrier\A] \qquad \A\sem{g \compose (f_1, \ldots,
  f_n)} = \A\sem g \compose (\A\sem{f_1}, \ldots, \A\sem{f_n})
\]
Thus a model of an operad $\OO$ in $\UU$ is just a multi-functor $\A :
\OO \to \UU$.

% \begin{example}\examplelabel{monoid operad models}
%   Recall the multicategory $\pair\Set\bigtimes$ from \exampleref{fin
%     prod multi cat}, i.e., the multicategory induced by choosing the
%   finite cartesian product of sets.  A model for the operad $\Monoid$
%   in $\pair\Set\bigtimes$ is then just a monoid. First, given a monoid
%   $\A = \seq{\carrier \A, (*), e}$, define a model structure over
%   $\carrier\A$ by:
%   \begin{itemize}
%   \item $\A\sem{\bigodot_{i=1}^n}(x_1, \ldots, x_n) := x_1 * \cdots * x_n$
%   \end{itemize}
%   with the usual convention about edge cases, i.e.:
%   \begin{itemize}
%   \item $\A\sem{\bigodot_{i=1}^0}() := e$
%   \item $\A\sem{\bigodot_{i=1}^1}(x) := \A\sem{\id}(x) := x$
%   \item $\A\sem{\bigodot_{i=1}^n}(x_1, \ldots, x_n) := x_1 * \cdots
%     * x_n$ when $n > 1$.
%   \end{itemize}
%   Since $\A$ is associative, we have:
%   \begin{align*}
%   \A&\sem{\bigodot_{i=1}^n\compose(\bigodot_{j=1}^{m_1}, \ldots, \bigodot_{j=1}^{m_n})}(x_{11}, \ldots, x_{1m_1}, \ldots, x_{n1}, \ldots, x_{nm_n})
%   \\&{}=
%   \A\sem{\bigodot_{\ell=1}^{\sum_{i=1}^nm_i}}(x_{11}, \ldots, x_{1m_1}, \ldots, x_{n1}, \ldots, x_{nm_n})
%   \\&{}=
%   x_{11} * \cdots * x_{1m_1}*\cdots * x_{n1} *\cdots* x_{nm_n}
%   \\&{}=
%   (x_{11} * \cdots * x_{1m_1})* \cdots * (x_{n1} *\cdots* x_{nm_n})
%   \\&{}=
%   \A\sem{\bigodot_{i=1}^n}
%   (\A\sem{\bigodot_{j=1}^{m_1}}(x_{11}, \ldots, x_{1m_1})
%   ,\ldots,
%   \A\sem{\bigodot_{j=1}^{m_n}}(x_{n1}, \ldots, x_{nm_n})
%   )
%   \\&{}=
%   \A\sem{\bigodot_{i=1}^n} \compose
%   (\A\sem{\bigodot_{j=1}^{m_1}}
%   ,\ldots,
%   \A\sem{\bigodot_{j=1}^{m_n}}
%   )(x_{11}, \ldots, x_{1m_1}, \ldots, x_{n1}, \ldots, x_{nm_n})
%   \end{align*}
%   Similarly, we have:
%   \[
%   \A\sem{\id \compose \bigodot_{i=1}^n}(x_1, \ldots, x_n) =
%   \A\sem{\bigodot_{i=1}^n}(x_1, \ldots, x_n) =
%   \A\sem{\id} \compose \A\sem{\bigodot_{i=1}^n}(x_1, \ldots, x_n)
%   \]
%   and
%   \begin{align*}
%     \A&\sem{\bigodot_{i=1}^n\compose (\id, \ldots, \id)}(x_1, \ldots, x_n)
%     \\&{}=
%     \A\sem{\bigodot_{i=1}^n}(x_1, \ldots, x_n)
%     \\&{}=
%     x_1 *\cdots *x_n = \id (x_1) *\cdots *(x_n)
%     \\&{}=
%   \A\sem{\bigodot_{i=1}^n}\compose (\A\sem\id, \ldots, \A\sem\id)(x_1, \ldots, x_n)
%   \end{align*}
%   and so $\A$ is a model for $\Monoid$.

%   Conversely, given a model $\A$, define:
%   \begin{itemize}
%   \item $e := \A\sem{\bigodot_{i=1}^0}()$
%   \item $x*y := \A\sem{\bigodot_{i=1}^2}(x,y)$
%   \end{itemize}
%   and observe:
%   \begin{itemize}
%   \item Since $\A\sem{\bigodot_{i=1}^1} = \A\sem{\id} = \id$, we
%     have $\A\sem\id x = x$.
%   \item $e*x = \A\sem{\bigodot_{i=1}^2\compose (\bigodot_{i=1}^0,
%     \id)}(x) = \A\sem{\bigodot_{i=1}^1}(x) = x$
%     and similarly $x*e = x$.
%   \item $(x * y) * z = \A\sem{\bigodot_{i=1}^2(\bigodot_{i=1}^2, \id)}(x, y, z) =
%     \A\sem{\bigodot_{i=1}^2(\id, \bigodot_{i=1}^2)}(x, y, z) = x * (y*z)$.
%   \end{itemize}
%   and so $\A$ is a monoid.

%   Both roundtrips are the identity. When $\A$ is a monoid, then the
%   round-trip monoid is given by:
%   \begin{itemize}
%   \item $e' := \A\sem{\bigodot_{i=1}^0}() = e$;
%   \item $x*'y := \A\sem{\bigodot{i=1}^2}(x,y) = x * y$
%   \end{itemize}
%   Conversely, when $\A$ is a model, then the round-trip model gives:
%   \begin{align*}
%     \A'\sem{\bigodot_{i=1}^n}(x_1, \ldots, x_n)
%     &{}= x_1 * \cdots * x_n
%     \\&{}=
%     \A\sem{\bigodot_{i=1}^2\compose
%       (\id, \bigodot_{i=1}^2\compose
%       (\id, \ldots, \bigodot_{i=1}^2\compose(\id, \id)))}(x_1, \ldots, x_n)
%   \\&{}=
%   \A\sem{\bigodot_{i=1}^n}(x_1, \ldots, x_n)
%   \end{align*}
%   And so the two notions coincide.
% \end{example}
% \begin{example}
%   Given an algebraic theory $\TT$, we define an operad $\OO^{\TT}$:
%   \begin{itemize}
%   \item The $n$-ary operations are the $n$-ary terms $\OO^\TT_n := \TT_n$.
%   \item The identity is given by the projection/variable $\id := x$.
%   \item Composition is given by substitution, taking care of the
%     appropriate weakening:
%     \[
%     x_{11}, \ldots, x_{nm_n} \vdash
%     t\compose (s_1, \ldots, s_n)
%     \definedby
%     t\compose \seq{s_1\compose \seq{x_{11}, \ldots, x_{1m_1}},
%         \ldots, s_n\compose \seq{x_{n1}, \ldots, x_{nm_n}}}
%     \]
%   \end{itemize}
%   The $\Monoid$ operad is quite different from the operad for the
%   algebraic theory $\Monoid$ of monoids. The former has only one
%   operation for each arity, whereas the latter has all possible words
%   over $n$-variables:
%   \[
%   \OO^{\Monoid}_n := \set{x_1, \ldots, x_n \vdash x_{i_1} * \cdots * x_{i_k}
%       \suchthat k \in \naturals, 1 \leq i_1, \ldots, i_k \leq n}
%   \]
%   For example, $\Monoid_1$ only has the identity unary operation, but
%   $\OO^{\Monoid}_1$ also has $e, x^2, \ldots, x^n, \ldots$ as unary
%   operations.
% \end{example}
A \emph{homomorphism} $h : \A[1] \to \A[2]$ between $\OO$-models in
$\UU$ consists of:
\begin{itemize}
\item a morphism $\carrier h : (\carrier{\A[1]}) \to \carrier{\A[2]}$
\end{itemize}
such that:
\begin{itemize}
\item for all $f \in \OO_n$, we have:
  \[
  \carrier h\compose(\A[1]\sem f) = \A[2]\sem f(\carrier h, \ldots, \carrier h)
  \]
\end{itemize}
Identities and composition of homomorphisms are homomorphisms, and so
we have a category $\Mod(\OO, \UU)$ of $\OO$-models in $\UU$.
We also have that a homomorphism of algebras is just a multi-natural
transformation between them as multi-fucntors $\quiver{O-homo as multi-natural}$.

% \begin{example}
%   For every algebraic theory $\TT$ and category $\C$ with chosen
%   products $\prod$, $\Mod(\TT, \C) \isomorphic \Mod(\OO^{\TT},
%   \pair\C\prod)$. Indeed, given a $\TT$-algebra $\A$, define an
%   $\OO^{\TT}$-model on the same carrier by sending each term $t$ to
%   its interpretation $\A\sem{t}$. Conversely, given an
%   $\OO^{\TT}$-model $\A$, we already have interpretations for all the
%   $\TT$-terms, and this interpretation validates the equations. A
%   homomorphism of between $\TT$-algebra is then the same thing as a
%   homomorphism between their corresponding $\OO^{\TT}$-models.
%   In particular the models of $\OO^{\Monoid}$ are the monoids in $\C$:
%   \[
%   \Mod(\OO^{\Monoid}, \pair\C\prod)
%   \isomorphic \Mod(\Monoid, C)
%   \isomorphic \Mon(\C)
%   \]
%   We'll show that we also have an isomorphism $\Mod(\Monoid,
%   \pair\C\prod) \isomorphic \Mon(\C)$. Categorifying the arguments in
%   \exampleref{monoid operad models} shows gives a bijection between
%   the objects of the two categories. The notion of a homomorphisms is
%   maintained across this bijection without changing the underlying
%   morphism, and so we have an isomorphism. Therefore, like with
%   algebraic theories, morita equivalence for operads is a coarser
%   notion than adjoint equivalence (in the $2$-category $\MultiCAT$).
% \end{example}

We'll conclude with a $2$-functor $\Mod(\OO, -) : \MultiCAT \to \CAT$.
\begin{itemize}
\item For every multicategory $\UU$, let $\Mod(\OO, \UU)$ be the
  category of $\OO$-models in $\UU$.
\item For every multi-functor $F : \UU[1] \to \UU[2]$, the
  functor $\Mod(\OO, F) : \Mod(\OO, \UU[1]) \to \Mod(\OO, \UU[2])$
  sends:
  \begin{itemize}
  \item each $\OO$-model $\A$ in $\UU[1]$ to the $\OO$-model:
    \begin{itemize}
    \item $\carrier{F\A} \definedby F\carrier\A$
    \item $(F\A)\sem f : (F\carrier\A, \ldots, F\carrier\A)
      \xto{F(\A\sem f)} F\carrier\A$
    \end{itemize}
  \item each $\OO$-homomorphism $h : \A[1] \to \A[2]$ to the
    $\OO$-homomorphism over $F\carrier h : F\A[1] \to F\A[2]$.
  \end{itemize}
\item For every multi-natural transformation $\quiver{multi-natural
  transformation type}$ we have a natural transformation
  $\quiver{2-cell action of Mod(O,-)}$ given, for every model $\A$, by:
  \[
  \Mod(\OO, \alpha)_{\A} : (F\carrier\A) \xto{\alpha_{\carrier\A}} G\carrier\A
  \]
\end{itemize}
\begin{proposition}
\propositionlabel{Mod(O,-) is a 2-functor}
  This data defines a $2$-functor $\Mod(\OO, -) : \MultiCAT \to \CAT$.
\end{proposition}
\begin{proof}
  To see why the action on $1$-cells is well-defined and functorial:
  \begin{itemize}
    \item For any $\OO$-model $\A$ in $\UU$, view the model as a
      multi-functor $\A : \OO \to \UU$, and note that the model $F\A$
      is given by $F\A : \OO \xto{\A} \UU \xto{F} \UU[2]$. Therefore
      it is a model.
    \item For any $\OO$-model homomorphism $h : \A[1] \to \A[2]$ in
      $\UU$, view it as a multi-natural transformation $\quiver{O-homo
      as multi-natural}$ and note that $\Mod(\OO, F)$ sends it to the
      horizontal compositition $\quiver{O-homo image under 2-functor}$.
      Therefore:
      \begin{itemize}
      \item It is a homomorphism (by type-checking); and
      \item $\Mod(\OO, F)$ is functorial, by the interchange law and
        by preservation of composition:
        \[
        \quiver{O-homo 2-functoriality:composition}\qquad
        \quiver{O-homo 2-functoriality:identity}
        \]
      \end{itemize}
  \end{itemize}

  To see why the action on $2$-cells is well-defined and functorial:
  \begin{itemize}
  \item Consider any multi-natural transormation
    $\quiver{multi-natural transformation type}$. Since the
    multicategory $\OO$ has only one object, $\Mod(\OO,\alpha)_{\A} :
    F \compose \A \star \to G\compose\A \star$ would completely
    determine a multi-natural transformation, and it is the
    multi-natural transformation given by whiskering by $\A$:
    \[
    \quiver{Mod(OO,alpha)_A as a whiskering}
    \]
    It is therefore natural, i.e., each component is an $\OO$-homomorphism
    \[
    \Mod(\OO,\alpha)_{\A} : \Mod(\OO, F)\A = F\A \to G\A = \Mod(\OO, G)\A
    \]
  \item To see that $\Mod(\OO,\alpha)$ is natural, take any
    homomorphism $\quiver{O-homo as multi-natural}$. The naturality
    condition then manifests as the two ways to express the horizontal
    composition at its unique component:
    \[
    \quiver{Mod(OO,alpha) naturality as horizontal composition}
    \]
  \item The preservation of vertical and horizontal composition and
    identities of $2$-cells then follows by taking an $\OO$-model,
    then partitioning these pasting diagrams accordingly:
    \[
    \quiver{Mod(OO,-) preservers vertical composition}
    \qquad
    \quiver{Mod(OO,-) preservers horizontal composition}
    \qquad
    \quiver{Mod(OO,-) preservers identities}
    \]
  \end{itemize}
  We therefore have a $2$-functor.
\end{proof}

\section{Multi-homomorphisms}
\label{sec:multihom}

Let $\TT$ be an algebraic theory, and $\pair\CC\prod$ a category with
chosen finite products. For all objects $X,Y$ in $\CC$, and
natural number $k$, define the strength (for the reader
monad $X^k$):
\[
\str : X^k \times Y \xto{\seq[i=1][k]{\projection_i\times \id}}
  \parent{X\times Y}^k
\]
Let $\A_1, \ldots,
\A_n,\A[2] \in \Mod(\TT, \CC)$ be algebras. A function $f :
\prod_{i=1}^n\carrier{\A_i} \to \carrier{\A[2]}$ is a
\emph{multi-$\TT$-homomorphism}~\citep{hyland-power-discrete-lawvere-theories-journal}
when, for every term $t \in \TT_k$ and $1 \leq i_0 \leq n$, the
following diagram commutes: \diagram{01}
%
When $n = 1$, this condition degenerates to the following diagram,
which states that $f : \carrier\A \to \carrier{\A[2]}$ is a
$\TT$-homomorphism: \diagram{02}
%
Therefore, multi-$\TT$-homomorphisms generalise $\TT$-homomorphisms:
\begin{proposition}\propositionlabel{unary multi-TT-homomorphisms}
  Consider the bijection between $1$-ary multi-morphisms $h :
  (\carrier\A) \to \carrier\A$ and $\C$-morphisms $h' : \carrier\A
  \isomorphic \prod\set{\carrier\A} \xto{h} \carrier\A$.  A
  multi-morphism $h$ is a multi-$\TT$-homomorphism iff the
  corresponding $h'$ is a $\TT$-homomorphism.
\end{proposition}


% FIXME: maybe replace with the proof that multi-A-homomorphisms are multi-linear
% \begin{example}
%   Let $\Field$ be a field and $\Vect_{\Field}$ be the algebraic theory
%   of vector spaces over $\Field$. Consider the term $x_1, x_2 \vdash t
%   \definedby \lambda_1x_1 + \lambda_2x_2$ for some scalars $\lambda_1,
%   \lambda_2 \in \Field$. Then the multi-homomorphism condition, when
%   applied to vector spaces $\V[3], \V[1], \V[2]$, function $f :
%   \carrier{\V[3]} \times \carrier{\V[1]} \to \carrier{\V[2]}$ becomes
%   the usual bilinearity condition, i.e., for $i=1$ (and similarly for $i=2$):
%   \[
%   f(\lambda_1u_1 + \lambda_2u_2, v) = \lambda_1f(u_1, v) +
%   \lambda_2f(u_2,v)
%   \]
%   Since: \diagram{03} Therefore, the multi-linear maps, such as the
%   inner product $\pair-- : \Field^n\times\Field^n \to \Field$, are
%   multi-$\Vect_{\Field}$-homomorphisms.
% \end{example}
We will define a multi-category structure $\carrier- : \Multi(\TT,\C) \to
\pair\C\prod$ over the finite-product multi-category:
\begin{itemize}
\item The objects are $\TT$-algebras.
\item The multi-morphisms are the multi-$\TT$-homomorphisms.
\item The multi-functor $\carrier-$ sends an algebra to its carrier
  and a multi-$\TT$-homomorphism to its underlying morphism.
\item The identities are the identity morphisms.
\item Composition is inherited from $\pair\C\prod$.
\end{itemize}
The next few propositions show that this data defines a
multi-category.

\begin{proposition}\propositionlabel{multi-TT-homomorphism composition}
  Let $g : (\A[2]_1, \ldots, \A[2]_n) \to \A[3]$ and $f_i :
  (\A[1]_{i1}, \ldots, \A[1]_{im_i}) \to \A[2]_i$ be
  multi-$\TT$-homomorphisms. Then $g \compose (f_1, \ldots, f_n) :
  (\A[1]_{11}, \ldots, \A[1]_{nm_n}) \to \A[3]$ is a
  multi-$\TT$-homomorphism.
\end{proposition}
\begin{proof}
  Let $p := \sum_{i=1}^nm_i$ and $\phi : [p] \xto{\isomorphic}
  \coprod_{i=1}^n[m_i]$ be the enumeration from the definition of
  composition in $\pair\C\prod$ for this case. Take any $1 \leq \ell_0
  \leq p$, and let $\pair{i_0}{j_0} := \phi\ell_0$ be its
  corresponding pair of indices. Take any $k$-ary term $t \in \TT_k$,
  and calculate as in \figref{multi-hom composition}.
  \begin{figure}
    \diagram{04}
    \caption{Multi-$\TT$-homomorphism property of the composition}\figlabel{multi-hom composition}
  \end{figure}
\end{proof}

\begin{proposition}\propositionlabel{Multi(TT,C) multi-category}
  The data above defines a multi-category $\Multi(\TT, \C)$ together
  with a faithful multi-functor $\carrier- : \Multi(\TT,\C) \to
  \pair\C\prod$: $\carrier{f_1} = \carrier{f_2} \implies f_1 = f_2$
  for every parallel pair of multi-$\TT$-homomorphisms $f_1,f_2 :
  (\A_1, \ldots, \A_n) \to \A[2]$.
\end{proposition}
\begin{proof}
  Identities are multi-$\TT$-homomorphisms by \propositionref{unary
    multi-TT-homomorphisms}, and composition of
  multi-$\TT$-homomorphisms is a multi-$\TT$-homomorphism by
  \propositionref{multi-TT-homomorphism composition}, and so we have a
  multi-category structure.

  Since composition and identities are inherited from $\pair\C\prod$,
  we have the structure of multi-functor. This multi-functor structure
  is faithful since the multi-morphisms in $\Multi(\TT, \C)$ are those
  of $\pair\C\prod$ with additional structure. Since faithful
  multi-functors reflect multi-categories, we conclude.
\end{proof}
% Let $\ppCAT$ be the $2$-category consisting of:
% \begin{itemize}
% \item $0$-cells are $\pair\C\prod$, categories with chosen finite products;
% \item $1$-cells are finite product preserving functors, not
%   necessarily preserving the chosen products;
% \item $2$-cells are the natural transformations between these functors.
% \end{itemize}
% This $2$-category is a sub-category of $\CAT$ that is full on
% $2$-cells but not on $1$-cells. This seems a bit weird, but that's the
% structure we seem to need. Given such a product preserving functor $H
% : \pair \C\prod \to \pair{\C[2]}\prod$, by fiat we have canonical
% isomorphisms $H\prod_{i=1}^nA \isomorphic \prod_{i=1}^nHA$ that
% preserve the product structure.

% We define the data of a $2$-functor $\Multi(\TT, -) : \ppCAT \to
% \MultiCAT$ sending:
% \begin{itemize}
% \item Each category with chosen products ($0$-cell) $\C$ to the
%   multi-category ($0$-cell) $\Multi(\TT, \C)$.
% \item Each finite-product preserving functor ($1$-cell) $H : \C[1] \to
%   \C[2]$ to the following multi-functor $\Multi(\TT, H) : \Multi(\TT,
%   \C[1]) \to \Multi(\TT, \C[2])$, sending:
%   \begin{itemize}
%   \item each $\TT$-algebra $\A$ to the $\TT$-algebra $\Multi(\TT,
%     H):=H\A$, given by:
%     \begin{itemize}
%     \item its carrier $\carrier{H\A}$ is $H\carrier{\A}$; and
%     \item the interpretation of each $t \in \TT_n$ is $H\A\sem t :
%       (H\carrier\A)^n \xto\isomorphic H(\carrier\A^n) \xto{\A\sem t}
%       H\carrier\A$.
%     \end{itemize}
%     More abstractly, an algebra is a product-preserving functor $\A :
%     \TT \to \C$, and $H\A$ is the product preserving functor $H\A :
%     \TT\xto\A \C \xto H \C[2]$.
%   \item each multi-$\TT$-homomorphism $h : (\A_1, \ldots, \A_n) \to
%     \A[2]$ to the multi-$\TT$-homomorphism $\Multi(\TT,H)h : (H\A_1,
%     \ldots, H\A_n) \to H\A_n$, given by:
%     \[
%     \Multi(\TT,H)h : \prod_{i=1}^nH\A_i \xto{\isomorphic}
%     H(\prod_{i=1}^n\A_i) \xto{Hh} H\A[2]
%     \]
%   \end{itemize}
% \item Each natural transformation $\quiver{nat trans between pp
%   functors}$ ($2$-cell) to the multi-natural transformation
%   $\quiver{Mod(TT, alpha) type}$, given, at an algebra $\A\in
%   \Mod(\TT,\C[1])$, by the unary multi-homomorphism corresponding to
%   the following $\TT$-homomorphism (cf.~\propositionref{unary
%     multi-TT-homomorphisms}):
%   \[
%   \Multi(\TT,\alpha)_{\A}' : H\carrier{\A} \xto{\alpha_{\carrier \A}} K\carrier{\A}
%   \]
% \end{itemize}
% It may be worthwhile to spell-out some additional related
% structure. The following forms the data of a $2$-functor $\seq- :
% \ppCAT \to \MultiCAT$, sending:
% \begin{itemize}
% \item $0$-cells: categories with specified products $\pair\C\prod$ to
%   their associated multi-category $\pair\C\prod$
%   (cf.~\exampleref{multi-cat from chosen products}).
% \item $1$-cells: finite-product preserving functors $H : \C \to
%   \C[2]$ to the multi-functor sending:
%   \begin{itemize}
%   \item objects $a \in \Obj{\pair\C\prod} = \Obj\C$ to $\seq Ha \in \Obj{\C[2]}$;
%   \item multi-morphisms $f:(a_1, \ldots, a_n) \to b$ to the multimorphism:
%     \[
%     \seq Hf : \prod_{i=1}^n Ha_i \xto{\isomorphic} H(\prod_{i=1}^na_i) \xto{Hf} Hb
%     \]
%   \end{itemize}
% \item $2$-cells: natural transformations $\quiver{nat trans between pp
%   functors}$ to the multi-natural transformation given, for each $a
%   \in \Obj\C$, by:
%   \[
%   \seq[a]{\alpha}{} : \prod_{i=1}^1Ha \isomorphic Ha \xto{\alpha_a} Ka
%   \]
% \end{itemize}
% This data allows us to talk about the following data for a natural
% $2$-transformaion $\quiver{forgetful functor 2-transformation: type}$
% whose components, for each category $\pair\C\prod$ with chosen
% finite-products, consists of the forgetful multi-functors:
% \begin{itemize}
% \item sending each $\TT$-algebra $\A$ to its carrier $\carrier\A$;
% \item sending each multi-$\TT$-homomorphism $h : (\A_1, \ldots, \A_n)
%   \to \A[2]$ to its underlying morphism $\carrier h : \prod_{i=1}^n
%   \carrier\A_i \to \carrier{\A[2]}$.
% \end{itemize}
% \begin{proposition}
%   This data defines:
%   \begin{itemize}
%   \item Two $2$-functors $\seq-, \Multi(\TT,-) : \ppCAT \to
%     \MultiCAT$, the former, $\seq-$, is $1$-faithful and
%     $2$-fully-faithful.
%   \item A natural $2$-transformation $\quiver{forgetful functor
%     2-transformation: type}$ whose components are faithful
%     multi-functors.
%   \end{itemize}
% \end{proposition}
% \begin{proof}
%   We have much to do. Starting with the well-formedness of the data
%   for $\seq-$, we have shown $\pair\C\prod$ is a multi-category. To
%   see that $\seq H$ is a functor, calculate as in \figref{seq h functoriality}.
%   \begin{figure}
%   \diagram{05}
%   \diagram{06}
%   \caption{Functoriality of $\seq H$}\figlabel{seq h functoriality}
%   \end{figure}

%   Next, we'll use the following observation. Given a natural
%   transformation $\quiver{nat trans between pp functors}$ between
%   finite-product preserving functors:
%   \[
%     \diagram*{07}\tag*{($*$)}
%   \]
%   proved by inverting the canonical isomorphisms and using the
%   universal property of products:
%   \diagram{08}
%   Using this property, we show that $\quiver{seq alpha type}$ is a
%   multi-natural transformation. Take any multi-morphism $f : (a_1,
%   \ldots, a_n) \to b$, and calculate:
%   \diagram{09}
%   Therefore the data for the $2$-functor $\seq-$ is well-defined.

%   Next, consider the data for the $2$-functor $\Mod(\TT,-)$.
%   \propositionref{Multi(TT,C) multi-category} shows it is a
%   multi-category, and so we need to check the well-formedness of the
%   higher dimensional data. For $1$-cells, as we explained more
%   abstractly, the data $\Multi(\TT, H)\A$ forms an algebra.  Given a
%   multi-$\TT$-homomorphism $h : (\A_1, \ldots, \A_n) \to \A[2]$, we
%   need to show that the morphism:
%   \[
%   \Multi(\TT, H)h : (H\A_1, \ldots, H\A_n) \to H\A[2]
%   \]
%   is indeed a multi-$\TT$-homomorphism. Take any $t \in
%   \TT_k$, $i_0 < n$, and calculate as in \figref{multi-TT-homomorphism
%     of Multi(TT, H)h}.

%   \begin{figure}
%     \diagram{10}
%     \caption{Multi-homomorphism proof for $\Multi(\TT,H)h$
%     }\figlabel{multi-TT-homomorphism of Multi(TT, H)h}
%   \end{figure}

%   To show that the structure on $2$-cells is well-defined, we need to
%   show that $\Mod(\TT,\alpha)_{\A}' : H\A \to K\A$ is a
%   $\TT$-homomorphism. Take any $t \in \TT_k$, and calculate:
%   \diagram{11}

%   Checking that $\seq-$ is a $2$-functor is straightforward:
%   \begin{itemize}
%   \item For $1$-cells, take two finite-product preserving functors
%     $\C[3] \xto{H} \C[4] \xto{K} \C[1]$.
%     \begin{itemize}
%     \item For objects, we have $\seq{K \compose H}a := KH a := \seq K
%       \compose \seq H a$.
%     \item For multi-morphisms $f : (a_1, \ldots, a_n) \to b$,
%       calculate:
%       \diagram{12}
%     \end{itemize}
%     It's also straightforward to check that $\seq\Id$ is the identity
%     multi-functor:
%     \diagram{15}
%   \item For $2$-cells:
%     \begin{itemize}
%     \item For the identity natural transformation we have:
%     \[
%     \seq[a]{\id[H]} : \prod\set{Ha} \xto\isomorphic Fa = Fa
%     \]
%     \item For vertical composition, calculate for vertically
%       composable $2$-cells $\quiver{pp vertical composition}$:
%       \diagram{13}
%     \item For horizontal composition, calculate for $\quiver{pp
%       horizontal composition}$:
%     \diagram{14}
%     \end{itemize}
%   \end{itemize}
%   And thus we have a $2$-functor $\seq- : \ppCAT \to \MultiCAT$.

%   Finally, we turn to the natural $2$-transformation $\carrier - :
%   \Multi(\TT,-) \to \seq{-}$.
%   \begin{itemize}
%   \item For the $1$-naturality condition, take a product preserving
%     functor $H : \C[1] \to \C[2]$ in $\ppCAT$ and show the
%     following equality of multi-functors:
%     \diagram{16}
%     by chasing an algebra $\A \in \Multi(\TT, \C[1])$ (top) and a
%     multi-$\TT$-homomorphism $h : (\A_1, \ldots, \A_n)\to \A[2]$
%     (bottom):
%     \begin{center}
%       \diagram*{17}
%       \hfill
%       \diagram*{18}
%     \end{center}
%   \item For the $2$-naturality condition, take a natural
%     transformation $\quiver{nat trans between pp functors}$ and
%     algebra $\A$, and consider the component at $\A$ of the
%     appropriate pastings (drawn badly :D):
%     \[
%     \quiver{carrier 2-naturality}
%     \]
%   \end{itemize}
%   And so we have a natural $2$-transformation $\carrier- :
%   \Multi(\TT,-) \to \seq-$ \emph{provided} we show that $\Multi(\TT,-)
%   : \ppCAT \to \MultiCAT$ is indeed a $2$-functor. We can deduce that
%   from the fact that $\carrier-$ is component-wise injective:
%   \begin{itemize}
%   \item For identity $1$-cells:
%     \begin{itemize}
%     \item On objects, we indeed check that the carrier and algebra
%       structure of $\Multi(\TT, \Id)\A$ are the same as of $\A$.
%     \item On multi-homormophisms, we use the faithfulness of
%       $\carrier-_{\C}$:
%       \[
%       \carrier-_{\C}\compose \Multi(\TT, \Id) h
%       \explain={$1$-cell $2$-naturality}
%       \seq\Id \compose \carrier-_{\C}h
%       \explain*={$\seq-$ is a $2$-functor}
%       \carrier-_{\C}h
%       =
%       \carrier-_{\C}\compose \Id h
%       \]
%       and so by faithfulness, $\Multi(\TT, \Id)h = \Id h$.
%     \end{itemize}
%   \item Similarly, for composable $1$-cells $\C[3] \xto H \C[4] \xto K \C$:
%     \begin{itemize}
%     \item On objects, we indeed check that the algebras are the same.
%     \item On multi-homomorphisms, we use faithfulness:
%       \[
%       \carrier{\Multi(\TT, K \compose H)h}_{\C}
%       \explain={$1$-cell $2$-naturality}
%       \seq{K \compose H} (\carrier{h}_{\C[3]})
%       \explain*={$\seq-$ is a $2$-functor}
%       \seq{K} \compose \seq{H} (\carrier{h}_{\C[3]})
%       =
%       \carrier{\Multi(\TT, K)\compose \Multi(\TT, H)h}_{\C}
%       \]
%     \end{itemize}
%     and so $\Mod(\TT,-)$ is functorial on $1$-cells.
%   \item For $2$-cell identities, the same technique applies:
%     \[
%     \carrier{\Multi(\TT, \id[H])_{\A}}_{\C[2]}
%     \explain={$2$-cell $2$-naturality}
%     \seq[{\carrier\A}]{\id[H]}
%     \explain*={$\seq-$ is a $2$-functor}
%     (\id[\seq{H}])_{\carrier\A}
%     =
%     (\id[\seq{H}{\carrier\A}])
%     \explain={$1$-cell $2$-naturality}
%     \id[\carrier{\Multi(\TT, H)\A}]
%     \explain*={$\carrier-_{\C[2]}$ multi-functoriality}
%     \carrier{\id[\Multi(\TT, H)_{\A}]}
%     \]
%   \item For vertical composition, take a composable pair of $2$-cells,
%     and calulate as follows (reverting to the more succinct notation
%     in the definition of natural $2$-transformations):
%     \begin{gather*}
%       \quiver{mult 2-functoriality vertical composition: a}\\
%     \begin{aligned}
%       &\explain={identities}
%       \quiver{mult 2-functoriality vertical composition: b}
%       \explain={$2$-cell\\$2$-naturality}
%       \quiver{mult 2-functoriality vertical composition: c}
%       \explain={$G := \seq-$ is a\\$2$-functor}
%       \quiver{mult 2-functoriality vertical composition: d}\\
%       &\explain={$2$-cell\\$2$-naturality}
%       \quiver{mult 2-functoriality vertical composition: e}
%       \explain={$2$-cell\\$2$-naturality}
%       \quiver{mult 2-functoriality vertical composition: f}
%       \explain={identities}
%       \quiver{mult 2-functoriality vertical composition: g}
%     \end{aligned}
%     \end{gather*}
%   \item For horizontal composition, we use the same technique:
%     \begin{gather*}
%       \quiver{mult 2-functoriality horizontal composition: a}\\
%       \begin{aligned}
%         &\explain={identities}
%         \quiver{mult 2-functoriality horizontal composition: b}
%         \explain={$2$-cell\\$2$-naturality}
%         \quiver{mult 2-functoriality horizontal composition: c}\\
%         &\explain={$G$ is a\\$2$-functor}
%         \quiver{mult 2-functoriality horizontal composition: d}
%         \explain={identities}
%         \quiver{mult 2-functoriality horizontal composition: e}\\
%         &
%         \explain={$2$-naturality}
%         \quiver{mult 2-functoriality horizontal composition: f}
%         \explain={$2$-naturality}
%         \quiver{mult 2-functoriality horizontal composition: g}\\
%         &\explain={identities}
%         \quiver{mult 2-functoriality horizontal composition: h}
%       \end{aligned}
%     \end{gather*}
%   \end{itemize}
%   Therefore, $\Mod(\TT,-)$ is also a $2$-functor.
% \end{proof}

\section{Yoneda structure}
\label{sec:yoneda}

We'll need universal properties in the $2$-category of
multi-categories. It might be useful to develop a rich theory of
representability inside $\MultiCAT$. We'll try to do it using the
\emph{Yoneda structure} of \cite{street-walters-yoneda-structures}.

\subsection{Admissible multi-functors}
A multi-functor $F : \UU \to \UU[2]$ is \emph{admissible} when, for
all $a_1, \ldots, a_n$ in $\UU$ and $b$ in $\UU[2]$, the multi-homset
$\UU[2](Fa_1, \ldots, Fa_n; b) \in \Set$ is small. A multi-category is
\emph{admissible} when its identity multi-functor is admissible, i.e.,
when it is locally small. The admissible multi-functors are closed
under-precomposition, i.e., for $\UU[1] \xto F \UU[2] \xto G \UU[3]$,
if $G$ is admissible, then $G \compose F$ is also admissible. Indeed,
given $a_i$ in $\UU[1]$ and $c$ in $\UU[3]$, we have:
\[
\UU[3](G\compose F a_1, \ldots, G\compose F a_n; C)
=
\UU[3](G(F a_1), \ldots, G(F a_n); C)
\explain \in {$G$ admissibility\\for $Fa_1, \ldots, Fa_n$}
\Set
\]

\subsection{Iterated lists}
As you can see in from previous sections, manipulating
multi-categories in the abstract makes the indexing heavy. Here we
develop some notation to manipulate sequences of multi-morphisms
ensemble, which is probably related to a construction of a free
monoidal category out of a multi-category.

Recall the structure of the list/free monoid monad, given by:
\begin{itemize}
  \item For each set $X$, the set $\List X$ of $X$-sequences $\xvec :=
    \seq{\xvec_1, \ldots, \xvec_{\abs\xvec}}$, i.e., $\abs\xvec$ is
    the length of the sequence, and its components in order are
    $\xvec_i$ for $1 \leq i \leq n$.

    We will also grade the monad, i.e., $\List_n X \subset \List X$
    are the length $n$ sequences.

  \item For each $x \in X$, the singleton $\seq{x} \in \List_1 X$.
  \item Given a sequence $\xvec \in \List_n X$ and, for each $1 \leq i
    \leq \abs\xvec$, another sequence $\yvec_i \in \List_{m_i}Y$, we
    have the concatenated sequence:
    \[
    \concat_{i=1}^{n}\yvec^i := \seq{\yvec_{11}, \ldots,
      \yvec_{1m_1}, \ldots, \yvec_{n1}, \ldots,
      \yvec_{nm_n}}
    \]
    We will also need iterated lists. These too form a monad-like
    structure, that of \emph{unlabelled, rooted, finitely branching
      trees of uniform depth}:
    \begin{itemize}
    \item For each set $X$, the set $\Tree X \definedby \Union_{k \in
      \naturals} \List^kX$ represents finitely-branching rooted trees
      whose nodes are unlabelled, and whose leaves, all of which have
      the same depth, are taken from $X$.
    \item For each $x \in X$, we identify the unique rooted tree with
      $1$ leaf at depth $1$ labelled by $x$ with the singleton list
      $\seq x \in \List^1X$.

      The empty tree, that only has the root node and no leaves, is
      represented by the empty list.

      Note that trees are allowed to have childless \emph{nodes} at
      non-uniform depths, e.g.: $\seq{\seq{},\seq{\seq{}},
        \seq{\seq{\seq{1}}}}$ has two nodes at depths $1$ and $2$ and
      a leaf at depth $3$.
    \item Given a tree whose leaves are trees, we can substitute the
      inner trees into the outer tree.
    \end{itemize}

    Let $[n] := \set{1, \ldots, n}$ be the canonical set of $n$-many
    indices.  For each depth-$k$ iterated list $\xivec \in \List^k X$,
    let $\floor\xivec \in \List X$ be its fully flattened sequence.

    We denote the set of \emph{branches} of a rooted tree $\xivec$ by
    $\Branch \xivec$, i.e., sequences of indices in the tree that mark
    a path from the root to a leaf. The leaf-depths are uniform, we
    have, for each $\xivec \in \List^kX$:
    \[
    \Branch \xivec = \coprod_{i_1 \in
      [\abs\xivec]}\coprod_{i_2 \in [\abs{\xivec_{i_1}}]}\cdots
    \coprod_{i_{k-1}\in[\abs{\xivec_{i_1,i_2\ldots,i_{k-2}}}]}[\abs{\xivec{i_1,
          \ldots, i_{k-1}}}]
    \]

    Let the \emph{fringe size} of a tree be its number of leaves,
    i.e., the cardinality of its set of branches:
    \begin{align*}
    \norm \xivec \definedby \abs{\floor\xivec}
    &{}\definedby \sum_{i_1 \in [\abs\xivec]} \sum_{i_2 \in [\abs{\xivec_{i_1}}]}\cdots
    \sum_{i_{k-1} \in [\abs{\xivec_{i_1,i_2,\ldots,i_{k-2}}}]}\abs{\xivec_{i_1,i_2,\ldots,i_{k-2},i_{k-1}}}
    \\&{}=\cardinality{\Branch \xivec}
    \end{align*}
    For example, for $\xivec \definedby
    \seq{\seq{\seq{1,2},\seq{3,4}},\seq{\seq{5},\seq{6,7},\seq{8}}}$
    we have $\floor\xivec = \seq{1,2,\ldots,8}$ and:
    \[
    \abs{\xivec} = 2\qquad
    \abs{\xivec_1} = 2 \qquad
    \abs{\xivec_2} = 3 \qquad
    \norm{\xivec} :=
    \sum_{i_1 \in [\abs\xivec]}\sum_{i_2 \in [\abs{\xivec_{i_1}}]}\abs{\xivec_{i_1,i_2}}
    = \parent{2 + 2} + \parent{1 + 2 + 1} = 8
    \]
    We then define the bijection:
    \begin{align*}
    \seq[{\xivec}]{-,-,\ldots,-} &{}: \Branch\xivec
    \xto{\isomorphic}
        [\norm\xivec]
    \\\seq[\xivec]{i_1, \ldots, i_k} &{}\definedby
    \sum_{j_1 < i_1}\norm{\xivec_{j_1}}
    +
    \sum_{j_2 < i_2}\norm{\xivec_{i_1,j_2}}
    +
    \ldots
    +
    \sum_{j_{k-1} < i_{k-1}}\norm{\xivec_{i_1,i_1, \ldots, i_{k-2},j_{k-1}}}
    + i_k
    \end{align*}
    so that $\floor\xivec_{\seq{i_1, \ldots, i_k}} = \xivec_{i_1,
      \ldots, i_k}$ and $\seq{\xivec_{\seq[\xivec][-1]{1}}, \ldots,
      \xivec_{\seq[\xivec][-1]{\norm\xivec}}}$. For example,
    continuing with the previous $\xivec$, we have:
    \begin{align*}
      \seq{1,2,1} &= 0 + 2 + 1 = 3 & \floor\xivec_3 &= 3 = \xivec_{1,2,1}
      \\
      \seq{2,2,1} &= 4 + (1 + 2) + 1 = 8 & \floor\xivec_8 &= 8 = \xivec_{2,2,1}
    \end{align*}
\end{itemize}
Given a list $\avec \in \List A$ over some set $A$, an \emph{ordered
  partition of $\avec$} is a list-of-lists $\pi \in \List^2 A$ such
that $\concat_{i = 1}^{\abs \pi} \pi_i = \avec$. Similarly, an
\emph{iterated ordered partition of $\avec$} is a tree $\pi \in \Tree
A$ such that $\floor \pi = \avec$.


\subsection{The multi-category of presheaves}
Let $\UU$ be a multi-category. We will define a \emph{category}
$\opposite\UU$ so that our presheaves will be functors $F :
\opposite\UU \to \Set$:
\begin{itemize}
\item $\Obj{\opposite\UU} := \List \Obj\UU$ are finite-sequences
  $\avec = \seq{\avec_1, \ldots, \avec_n}$ of $\UU$-objects.
\item Morphisms $f \in \opposite\UU(\avec , \bvec)$, which we will
  write as ``$f : \avec \from \bvec$ in $\UU$'', consist of:
  \begin{itemize}
  \item an ordered partition $\partit_f\definedby \seq{\partit_1,
    \ldots, \partit_n} \in \List^2 \Obj\UU$, $\vec b$, i.e.:
    \(
    \bvec = \concat_{i=1}^n \partit_i
    \).
  \item for each $1 \leq i \leq n$, a multi-morphism $f_i : \avec_i \from
    (\bvec_i)$.  We can depict a morphism like so:
    \[
    \quiver{Cop morphism}
    \]
  \item The identity $\id[\vec a]^{[\opposite\UU]}$ consists of:
    \begin{itemize}
    \item $\partit_{\id[\vec a]} := \seq{\seq{a_1}, \ldots,
      \seq{a_n}}$, the partition into singletons; and
    \item $(\id[\vec a])_i := \id[a_i] : a_i \from (a_i)$.
    \end{itemize}
  \item To compose morphisms $\avec \xfrom f \bvec \xfrom g \cvec$,
    use the multiplication of the $\Tree$ monad to get an iterated
    partition of depth $3$, and then flatten all the depth $1$
    iterated lists into lists.  Then, compose the multi-morphisms to
    get a morphism $f \compose g : \avec \from \cvec$ (note the order
    of arguments in this notation). Letting $\abs\avec \defines n$ and
    $\abs{\partit_{f,i}} \defines m_i$:
    \begin{itemize}
    \item $\partit_{f \compose g} \definedby
      \seq{\seq{\partit_{g_{\pair11}}
                \++ \ldots \++
                \partit_{g_{\pair1{m_1}}}}
          , \ldots
          ,\seq{\partit_{g_{\pair n1}}
                \++ \ldots \++
                \partit_{g_{\pair n{m_n}}}}
          }$
    \item $(f \compose g)_{i} := f_i\compose (g_{\pair i1}, \ldots, g_{\pair i{m_i}})$.
    \end{itemize}
      We can visualise this operation so, letting $k_{i,j} \defines
      \abs{\partit_{g_i},j}$:
      \[
      \quiver{Cop composition}
      \]
  \end{itemize}
\end{itemize}
\begin{lemma}
  This data defines a category $\opposite\UU$.
\end{lemma}
\begin{proof}
  We inherit the properties from the monad laws and the multi-category
  axioms.
\end{proof}

A \emph{presheaf over $\UU$} is a functor $F : \opposite\UU \to \Set$.

\begin{example}\examplelabel{Yoneda mapping}
  Let $\UU$ be a locally small multi-category. For each $a \in
  \Obj\UU$, we have a presheaf $\yoneda[\UU]a : \opposite\UU \to \Set$
  called the presheaf \emph{represented by $a$}:
  \[
  \yoneda_{\UU}\,a\,\bvec \definedby \UU(\bvec; a) \in \Set
  \qquad
  \yoneda_{\UU}\,a\,(\bvec \xfrom{f} \cvec) : g \in \UU(\bvec; a)
  \mapsto g \compose (f_1, \ldots, f_{\abs\bvec}) \in \UU(\cvec; a)
  \]
  The functoriality:
  \[
  \yoneda_{\UU}\,a\,(\bvec \xfrom{f} \cvec \xfrom{g} \dvec)
  =
  \yoneda_{\UU}\,a\,(g)
  \compose
  \yoneda_{\UU}\,a\,(f)
  \qquad
  \yoneda_{\UU}\,a\,(\id[\bvec])
  =
  \id[\yoneda_{\UU}\,a\,\bvec]
  \]
  follows from the multi-category axioms and the definition of $\opposite\UU$.
\end{example}

Let $F_1, \ldots, F_n, G : \opposite\UU \to \Set$ be presheaves.  A
\emph{natural multi-transformation $\alpha : (F_1, \ldots, F_n) \to
  G$} consists of:
\begin{itemize}
\item For every sequence $\avec\in \Obj{\opposite\UU}$ and length-$n$
  ordered partition $\partit$ of $\avec$, a function:
  \[
  \alpha_{\partit} : \prod_{i=1}^nF_i\partit_i \to G\avec
  \]
\end{itemize}
such that:
\begin{itemize}
\item For every $f : \avec \from \bvec$ in $\UU$ and for every
  length-$n$ ordered partition $\phi = \seq[i=1][n]{\phi_i : \partit_i
    \from \partit'_i}$ of the sequence $f$: \diagram{19}
\end{itemize}

\begin{example}\examplelabel{Yoneda functoriality}
  Let $g : (\avec) \to b$ be a multi-morphism in $\UU$. We define a
  natural multi-transformation, where $n \definedby \abs\avec$:
  \begin{align*}
  \yoneda_{\UU}\,g :{}& (\yoneda_{\UU}\, \avec_1, \ldots, \yoneda_{\UU}\,\avec_n) \to
  \yoneda_{\UU}\,b\\
  (\yoneda_{\UU}\,g)_{\partit} :{}&
  \seq{ \partit_1 \xto{u_1} \avec_1
      , \ldots
      , \partit_n \xto{u_n} \avec_n
      } \mapsto
  g\compose (u_1, \ldots, u_n)
  \end{align*}
  To see that it is natural, take any $f : \cvec \from \dvec$ in
  $\UU$, and an $n$-ary ordered partition of it $\phi$, and chase
  elements around the naturality diagram:
  \diagram{20}
  and so $\yoneda g$ is a natural multi-transformation.
\end{example}
The multi-category $\Psh\UU$ of \emph{presheaves over $\UU$} consists
of this data:
\begin{itemize}
\item Objects are presheaves $F : \opposite\UU \to \Set$.
\item Multi-morphisms $\alpha : (F_1, \ldots, F_n) \to G$ are the
  natural multi-transformations.
\item The identity multi-transformation $\id[F] : (F) \to F$ has as
  component at each unary partition $\partit = \seq{\avec}$ of $\avec$
  the canonical isomorphism:
  \[
  (\id[F])_{\seq{\avec}} : \prod\set{F\avec} \xto{\isomorphic} F\avec
  \]
\item To define composition, take composable natural
  multi-transformations
  \[
  \alpha = \seq[i=1][n]{(\Fvec_i) \to \Gvec_i}
  \quad\text{and}\quad
  \beta : (\Gvec) \to \Hvec
  \quad\text{where:}\quad
  \Fvec \in \List^2(\opposite\UU \to \Set),
  \Gvec \in \List^1(\opposite\UU \to \Set)
  \] letting as usual $p := \norm
  \Fvec$, $n \definedby \abs \Gvec$, and $m_i \definedby
  \abs{\Fvec_i}$. Given a $p$-ary ordered partition $\partit$ of some
  $\avec$, define these two auxiliary partitions:
  \begin{itemize}
  \item $\partit_{\pair i*} \definedby \seq[j=1][m_i]{\partit_{\pair ij}}$
  \item $\partit_{\pair i{\++}} \definedby
    \concat_{j=1}^{m_i}\partit_{\pair ij} = \floor{\partit'_i}$
  \item So $\alpha_{i, \partit_{\pair i*}}
    :\prod_{j=1}^{m_i}\Fvec_{\pair ij} \to \Gvec_i\partit_{\pair
    i{\++}}$.
  \end{itemize}
  It might easier to visualise how they relate to each other by
  aligning their fringes:
  \[
  \begin{array}{@{}*{4}{l@{}}c@{}l@{}c@{}l@{}@{}l@{}c@{}l@{}*{2}{l@{}}}
    \partit {}&= \tabmark
                  \iseq{\tabmark         \partit_{\pair 11}\
                       \tabmark,\ldots, \tabmark\partit_{\pair 1{m_1}}
                       \tabmark,\ldots,\tabmark\tabmark\partit_{\pair n1}
                       \tabmark,\ldots, \tabmark\partit_{\pair n{m_n}}\tabmark}
    \\
    \partit_{\pair{*}*} &{}:=\tabmark
                  \iseq{\iseq{\tabmark  \partit_{\pair 11}
                             \tabmark,\ldots, \tabmark\partit_{\pair 1{m_1}}}
                             \tabmark,\ldots,\tabmark
                        \iseq{\tabmark\partit_{\pair n1}
                             \tabmark,\ldots,\tabmark\partit_{\pair n{m_n}}}\tabmark}
    \\
    \partit_{\pair{*}{\++}} &{}:=\tabmark
                  \iseq{\tabmark               \partit_{\pair 11}
                        \tabmark\++ \ldots\++\tabmark\partit_{\pair 1{m_1}}
                        \tabmark,\ldots,\tabmark\tabmark\partit_{\pair n1}
                        \tabmark\++ \ldots\++\tabmark\partit_{\pair n{m_n}}
                  \tabmark}
  \end{array}
  \]
  We define the $\partit$-th component of the composition:
  \[
  (\beta \compose (\alpha_1, \ldots, \alpha_n))_{\partit} :
  \prod_{\ell =1}^p\Fvec_{\ell}\partit_{\ell} \xto{\isomorphic}
  \prod_{i = 1}^n\prod_{j=1}^{m_i} \Fvec_{\pair ij}\partit_{\pair ij}
  \xto{\prod_{i=1}^n\alpha_{i,\partit_{\pair i*}}}
  \prod_{i=1}^n\Gvec_i\partit_{\pair i{\++}}
  \xto{\beta_{\partit_{\pair *{\++}}}}
  H\avec
  \]
\end{itemize}
\begin{proposition}
  These data define a multi-category $\Psh\UU$.
\end{proposition}
\begin{proof}
  First, we need to show the identities and composition are
  well-defined, i.e., natural. For the identity, take any unary
  partition $\phi = \seq{f}$ of a morphism $f : \avec \from \bvec$ in
  $\UU$, and calculate:
  \diagram{21}
  Turning to the well-definedness of composition, continue with the
  same notation and take a $p$-ary partition $\phi = \seq{\phi_{\ell} :
    \partit[1]_{\ell} \from \partit[2]_{\ell}}$ of some morphism $f : \avec \from
  \bvec$ in $\UU$. The crucial step is to repeat the same partitioning
  scheme for $\phi$ as well:
  \begin{itemize}
  \item $\phi_{\pair i*} \definedby \seq[j=1][m_i]{\phi_{\pair ij} : \partit_{\pair ij} \from \partit[2]_{\pair ij}}$;
  \item $\phi_{\pair i{\++}} \definedby \concat_{j=1}^{m_i}\phi_{\pair ij} :
    \partit_{\pair i{\++}} \from \partit[2]_{\pair i{\++}}$ so that $\phi_{\pair i*}$ is a partition of $\phi_{\pair i{\++}}$;
  \end{itemize}
  and we can now calculate the naturality condition of the composition:
  \diagram{22}

  The neutrality of the identities w.r.t.~composition follow
  immediately:
  \begin{center}
  \diagram*{23}
  \diagram*{24}
  \end{center}

  To prove that composition is associative, we'll first dwell on the
  partitioning scheme we used to define composition.

  We can organise the data in a sequence of $d$-many composable
  multi-transformations:
  \[
  \Fvec^{d+1} \xto{\alpha^{d}} \cdots \xto{\alpha^1} \Fvec^1 \xto{\alpha^0} F^0
  \tag{($*$)}
  \]
  as follows. We consider a rooted tree $t$ of uniform leaf-depth $d+1$. We label
  each leaf, say in branch $\seq{i_1, \ldots, i_{d+1}}$ in the tree,
  with a single presheaf $\Fvec^{d+1}_{\seq{i1, \ldots,
      i_{d+1}}}$. Each node of depth $1 \leq s \leq d$ in the tree,
  say on the path $\seq{i_1, \ldots, i_s}$ is also labelled with a
  presheaf $\Fvec^s_{i_1, \ldots, i_s}$. Given these labels, we also
  annotate each node, that has, say, $k_{\seq{i_1, \ldots, i_s}}$
  children, with a natural multi-transformation $\alpha^s_{\seq{i_1,
      \ldots, i_s}} : \parent{\seq[i=1][k_{\seq{i_1, \ldots,
          i_s}}]{\Fvec^{s+1}_{i_1, \ldots, i_{s + 1}}}} \to
  \Fvec^s_{\seq{i_1, \ldots, i_s}}$.

  The bijection between the set of length-$s$ branches and the set of
  nodes and leafs at level $s$ allows us to define the lists
  $\Fvec^s$ and $\alpha^s$ to make the notation in ($*$) precise.
  After composing these multi-transformations in some order, we get a
  multi-transformation:
  \[
  \bigcirc \alpha: \parent{\Fvec^{s+1}} \to \Fvec^0
  \]
  Letting $p \definedby \abs{\Fvec^{s+1}}$, when we take any $\avec$
  and a $p$-ary partition of $\avec$, we can extend the tree $t$ by
  another level (turning all the previous leaves into nodes but
  without the multi-transformation label), and the leaves in the newly
  added level are labelled by the objects in $\avec$. It is this tree,
  its labels, and its branches and paths that we are manipulating with
  our notation.

  Given a tree (with any sort of labelling) of uniform leaf depth $d$,
  we denote by $t_{\seq{i_1, \ldots, i_s, *, \ldots, *}}$ the sub-tree
  rooted at node $\seq{i_1, \ldots, i_s}$. It's perhaps convenient to
  pad the remaining components up to a length-$d$ sequence to spot
  typos. We denote by:
  \[
  t_{\seq*{i_1, \ldots, i_{s},
      *, \ldots, *, \explain*{\++}{pos.~$j_1$},
      *, \ldots, *, \explain*{\++}{pos.~$j_2$},
     *,
     \underset{\substack{~\\[8pt]\ldots}}\ldots, *, \explain*{\++}{pos.~$j_r$},
     *, \ldots, *
  }}
  \]
  the sub-tree rooted at $\seq{i_1, \ldots, i_s}$, whose levels $j_1,
  \ldots, j_r$ were flattened. This notation is consistent with the
  notation we used for composition and in the previous parts of the
  proof, apart from identifying the partition $\pi$ with its depth-$3$
  iterated list. The crux of our proof is a kind of substitution
  lemma/monad-like associativity law that means the various iterations
  on this notation are consistent with each other, for example:
  \[
  (t_{i_1, i_2, *, \++, *, \++, *})_{j_1, \pair{j_2}{j_3}, \++}
  =
  t_{i_1, i_2, j_1, j_2, j_3, \++, \++}
  \]

  Back to proving associativity. Consider $3$ levels of composable
  multi-transformations:
  \[
  \Fvec \xto{\alpha} \Gvec \xto{\beta} \Hvec \xto{\gamma} \Kvec
  \]
  i.e.:
  \begin{itemize}
  \item $\gamma : \parent{\seq[i=1][k_{\seq{}}]{\Hvec_{\seq i}}} \to \Kvec$
  \item $\beta_{\seq i} :
    \parent{\seq[j=1][k_{\seq{i}}]{\Gvec_{\seq{i,j}}}} \to \Hvec_{\seq
      i}$, for all $1 \leq i \leq k_{\seq{}}$.
  \item $\alpha_{\seq{i,j}} :
    \parent{\seq[r=1][k_{\seq{i,j}}]{\Fvec_{\seq{i,j,r}}}} \to
    \Gvec_{\seq{i,j}}$, for all $1 \leq i \leq k_{\seq{}}$ and $1 \leq
    j \leq k_{\seq{i}}$.
  \end{itemize}
  Let:
  \begin{itemize}
  \item $f := \abs\Fvec$
  \item $f_{\seq{i,j}} := \abs{\Fvec_{\triple ij*}}$
  \item $g := \abs\Gvec$
  \item $g_i := \norm{\Gvec_{\seq{i, *}}} = k_{\seq i}$
  \item $h := \abs\Hvec = k_{\seq{}}$
  \end{itemize}
  Consider any $f$-ary partition $\partit$ of some $\avec$, and
  calculate:
  \diagram{25}
  Therefore, composition is associative.
\end{proof}

\begin{lemma}
  When $\UU$ is small, $\Psh\UU$ is locally small.
\end{lemma}
\begin{proof}
  When $\Obj\UU$ is small, so is $\List\Obj\UU$, and so is the set
  $\coprod_{\avec \in \List\Obj\UU}\Partit_n\avec$ of $n$-ary ordered
  partitions over some list of objects. Given a presheaf sequences and
  a presheaf $\Fvec,G : \UU \to \Set$, the collection of natural
  multi-transformation is a subset of the set $\prod_{\avec \in
    \Obj{\opposite\UU}, \partit\in\Partit\avec}
  \Set(\prod_i\Fvec_i\partit_i, G\avec)$.
\end{proof}

\subsection{The Yoneda lemma}
Even if it's not immediately clear from the axiomatisation, there's a
Yoneda lemma sitting at the heart of the Yoneda structure axioms.
Recall from examples \exampleref{Yoneda mapping} and
\exampleref{Yoneda functoriality} that, for every locally small
multi-category $\UU$, we have a multi-functor structure:
\[
\begin{array}{@{}l@{}l@{}}
  \yoneda_{\UU} &{}: \UU \to \Psh\UU \\
  \yoneda_{\UU}\,a &{}\definedby \UU(-,a) \\
  (\yoneda_{\UU}\,\parent{\vto \avec f{a'}})_{\partit} &{}\definedby
    \seq[i=1][\abs\avec]{(\partit_i)
    \xto{u_i} \avec_i} \mapsto f \compose \seq[i=1][\abs\avec]{u_i}
\end{array}
\]

\begin{lemma}
  This data defines a multi-functor $\yoneda_{\UU} : \UU \to \Psh\UU$.
\end{lemma}
\begin{proof}
  For identity multimorphisms $\id[a] : (a) \to a$, there is only one
  $1$-partition of any sequence $\bvec$, and we have:
  \[
  (\yoneda\id)_{\seq[i=1][1]{\bvec}} \seq[i=1][1]{u : (\bvec) \to a}
  \mapsto \id \compose (u) = u = (\id[\yoneda_a])_{\seq\bvec}\seq u
  \]
  and so $\yoneda\id = \id[\yoneda_a]$. Taking composable
  multi-morphisms $g : (\cvec) \to d$ and $f_i : (\bvec_i) \to
  \cvec_i$, $i = 1, \ldots, n := \abs\cvec$, take any $p$-ary
  partition $\partit$ of some $\avec$, where $p := \norm\bvec$, and
  chase an arbitrary element of $\prod_{\ell=1}^p\yoneda
  \bvec_\ell\partit_{\ell}$, where $m_i := \abs{\bvec_i}$:
  \diagram{26} and so $\yoneda(g \compose \seq[i]{f_i}) = \yoneda g
  \compose \seq[i]{\yoneda f_i}$.
\end{proof}

\begin{lemma}\emph{Yoneda for multi-categories.}
  Let $\UU$ be a locally-small multi-category, $\avec$ an $n$-ary
  sequence of objects in $\UU$, and $F \in \Psh\UU$ a presheaf. We
  have a bijection:
  \[
  \Yoneda_{\UU, F, \avec} : \Psh\UU(\yoneda \avec_1, \ldots, \yoneda\avec_n; F)
  \xto\isomorphic F\avec
  \]
  given by:
  \begin{gather*}
  \Yoneda_{\UU, F, \avec}(\seq[i=1][n]{\yoneda \avec_i} \xto\alpha F)
  \definedby
  \alpha_{\seq[i=1][n]{\seq[j=1][1]{\avec_i}}}\seq[i=1][n]{\id[a_i]}
  \in F\avec
  \\
  (\inv\Yoneda_{\UU, F, \avec}x)_{\partit} : \seq[i=1][n]{\partit_i
    \xto{u_i} \avec_i} \mapsto F(u_1, \ldots, u_n)x
  \end{gather*}
\end{lemma}
\goodbreak
\begin{proof}
  As usual, the proof is to just do it.

  To show well-definedness, i.e., that $\inv\Yoneda x : (\yoneda
  \avec_1, \ldots, \yoneda\avec_n) \to F$ is natural, take any
  morphism $\bvec \xfrom f\cvec$ in $\UU$, an $n$-ary partition
  $\seq[i=1][n]{\partit_i \xfrom{\phi_i} \partit'_i}$ of it, let
  $m_i:= \abs{\partit_i}$, and chase an arbitrary element of
  $\prod_{i=1}^n\yoneda \avec_i \partit_i$:
  \diagram{27}
  To show $\Yoneda \compose \inv\Yoneda = \id$, take any $x \in
  F\avec$:
  \[
  \Yoneda(\inv\Yoneda x) = F(\id, \ldots, \id)(x) \explain={functoriality} x
  \]
  For the converse, take any natural multi-transformation $\alpha :
  (\yoneda\avec_1, \ldots, \avec_n) \to F$. We need to show that, for
  every $n$-ary partition $\partit$ of some $\bvec$ and $\partit_i
  \xto{u_i} \avec_i$:
  \[
  \alpha_{\partit}(u_1, \ldots, u_n)
  =
  (\inv\Yoneda (\Yoneda \alpha))_{\partit}(u_1, \ldots, u_n) =
  F(u_1, \ldots, u_n)
  \parent{\alpha_{\seq[i=1][n]{\seq[j=1][1]{\avec_i}}}
    (\id[\avec_1], \ldots, \id[\avec_n])}
  \]
  We do so by chasing the element $\seq[i=1][n]{\id[\avec_i]} \in
  \prod_{i=1}^n\yoneda {\avec_i}{\seq[j=1][1]{\avec_i}}$ around
  $\alpha$'s naturality diagram for the partition
  $\seq[i=1][n]{\seq[j=1][1]{\avec_i} \xfrom{u_i} \partit_i}$ of the
  morphism $u : \avec \xfrom{(u_1, \ldots, u_n)} \bvec$ in $\UU$:
  \diagram{28}
  and so $\inv\Yoneda\compose \Yoneda = \id$.
\end{proof}
We call both this result and the bijection $\Yoneda$ the Yoneda lemma.

\begin{lemma}\emph{Yoneda embedding.}
  The Yoneda multi-functor $\yoneda_{\UU} : \UU \to \Psh\UU$ is
  fully-faithful.
\end{lemma}
\begin{proof}
  Take any $\avec$ and $b$, and consider the Yoneda lemma for the
  representable presheaf $F := \yoneda b$. Take any $f \in \UU(\avec;
  b) = \yoneda b\avec$.  For every partition $\partit$ of every
  $n$-tuple $\cvec$ and $\seq[i=1][n]{u_i : \partit_i \to \avec_i}$,
  we have:
  \[
  (\inv\Yoneda_{\UU, \yoneda b, \avec} f)_{\partit}\seq[i=1][n]{u_i}
  =
  \yoneda b(u_1, \ldots, u_n) f
  =
  f \compose (u_1, \ldots, u_n
  =
  (\yoneda f)_{\partit}\seq[i=1][n]{u_i}
  \]
  and so $\inv\Yoneda_{\UU, \yoneda b, \avec}f = \yoneda_{\UU}f$,
  i.e., $\inv\Yoneda_{\UU, \yoneda b, \avec}$ is the multi-functorial
  action of $\yoneda$ on the multi-homset $\UU(\avec; b)$. Since the
  Yoneda lemma is bijective, this functorial action is bijective,
  i.e., $\yoneda_{\UU}$ is fully-faithful.
\end{proof}
\begin{lemma}\emph{Naturality of the Yoneda Lemma.}
  The Yoneda lemma $\Yoneda : \Psh\UU(\yoneda[\avec], F)
  \xto\isomorphic F\avec$ is natural in the following senses:
  \begin{itemize}
  \item Naturally in $\avec \in \opposite\UU$, i.e.~for all $f : \avec \from \bvec$:
    \[
    Ff (\Yoneda(\yoneda[\avec] \xto\alpha F)) =
    \Yoneda(\yoneda[\bvec] \xto{\yoneda[f]} \yoneda[\avec] \xto\alpha F)
    \]
  \item Multi-naturally in $F \in \Psh\UU$, i.e.~for all $(F_1,
    \ldots, F_n) \xto{\beta} G$ and partition $\partit$ of $\avec$:
    \[
    \beta_{\partit}\seq[i=1][n]{\Yoneda(\yoneda[\partit_i] \xto{\alpha_i} F_i)}
    =
    \Yoneda\parent{\yoneda[\avec] \xto{\beta\compose\seq[i=1][n]{\alpha_i}} G}
    \]
  \end{itemize}
\end{lemma}
We can recover these from the general theory of Yoneda structures, but
we use them to establish the Yoneda structure axioms.
\begin{proof}
  In both cases, employ mindless calculation. For naturality:
  \begin{align*}
    Ff (\Yoneda(\yoneda[\avec] \xto\alpha F))
    &{}=
    Ff \parent{\alpha_{\seq[i=1][n]{\seq[j=1][1]{\avec_i}}}
      \seq[i=1][n]{\id[\avec_i]}}
    \explain={$\alpha$-naturality}
    \alpha_{\seq[i=1][n]{\bvec_{\pair i*}}}
    \seq[i=1][n]{(\yoneda \avec_if_i)\seq[j=1][1]{\id[\avec_i]}}
    \\&{}\explain={identities}
    \alpha_{\seq[i=1][n]{\bvec_{\pair i*}}}
    \seq[i=1][n]{(\yoneda f_i)_{
        \seq[j=1][\abs{\bvec_{\pair i*}}]{\seq[k=1][1]{\bvec_{\pair ij}}}}
      \seq[j=1][\abs{\bvec_{\pair i*}}]{\id[\bvec_{\pair ij}]}}
    \\&{}\explain={$\compose$-def}
    (\alpha\compose\seq[i=1][n]{\yoneda f_i})_{
      \seq[\ell=1][\norm\bvec]{\seq[j=1][1]{\bvec_{\ell}}}}
    \seq[\ell=1][\norm\bvec]{\id[\bvec_\ell]}
    =
    \Yoneda(\alpha \compose\seq[i=1][n]{\yoneda f_i})
  \end{align*}
  as we wanted.

  For multi-naturality:
  \begin{align*}
    \beta_{\partit}\seq[i=1][n]{\Yoneda(\yoneda[\partit_i] \xto{\alpha_i} F_i)}
    &{}=
    \beta_{\partit}\seq[i=1][n]{
      \alpha_{i, \seq[j=1][\abs{\partit_i}]{\seq[k=1][1]{\partit_{\pair ij}}}}
      \seq[j=1][\abs{\partit_{i}}]{\id[\partit_{\pair ij}]}}
    \\&{}\explain={$\compose$-def}
    (\beta\compose\seq[i=1][n]{\alpha_i})_{\partit}
    \seq[\ell=1][\norm \partit]{\id[\partit_{\seq[][-1]{\ell}}]}
    =
    \Yoneda(\beta\compose\seq[i=1][n]{\alpha_i})
  \end{align*}
  as we wanted.
\end{proof}

\begin{lemma}\emph{Special Yoneda reduction.}
  \lemmalabel{special Yoneda reduction}
  For every $\avec \xto{u} b$, we have:
  \(
  \Yoneda_{\UU, \yoneda b, \avec}(\yoneda[\avec] \xto{\yoneda u} \yoneda b)
  = u
  \).
\end{lemma}
\begin{proof}
  Calculate:
  $\Yoneda(\yoneda u) = \yoneda u_{\avec} (\id[\avec_1], \ldots, \id[\avec_n])
  = u \compose (\id, \ldots, \id) = u$.
\end{proof}

\begin{lemma}\lemmalabel{admissible yoneda embedding}
  The Yoneda embedding is, up to equivalence, admissible: for every multi-category $\UU$ there are:
  \begin{itemize}
  \item a multi-category $\Psh'\UU$;
  \item an admissible functor $\yoneda_{\UU}' : \UU \to \Psh'\UU$; and
  \item an adjoint equivalence in $\MultiCAT$, i.e.~an adjunction
    with invertible unit and counit:
    \[
    \quiver{psh adjoint equivalence}
    \]
  \item $2$-cells:
    \begin{center}
      \quiver{psh 2-cell transformer}
      \quiver{psh 2-cell transformer'}
    \end{center}
  \end{itemize}
  such that:
  \begin{center}
    \quiver{psh transformer condition}
    \quiver{psh transformer condition'}
  \end{center}
\end{lemma}
\begin{proof}
  Take any $\avec$ and presheaf $F$. By the Yoneda lemma:
  \[
  \Psh\UU(\yoneda\avec_1, \ldots, \yoneda\avec_n; F) \isomorphic
  F\avec \in \Set
  \]
  is a small set. So now we need to tediously construct a different
  presheaf multi-category $\Psh' \UU$ which is admissible:
  \begin{itemize}
  \item $\Obj{\Psh'\UU} := \Obj{\Psh\UU} + \Obj{\UU}$, i.e., either a
    presheaf or an object, together with a tag telling us which is
    which. We call the left-tagged elements `actual presheaves' and
    the right-tagged elements `formal representables'.  For each $x
    \in \Obj{\Psh'\UU}$, we can associate a presheaf $\sem x$:
    \[
    \sem{\injection_1 F} := F \qquad \sem{\injection_2 a} := \yoneda a
    \]
  \item We define the homs as follows:
    \[
    \Psh'\UU(x_1, \ldots, x_n; y) \definedby \begin{cases}
      \forall i. x_i = \injection_2 a_i: & \sem y(a_1, \ldots, a_n) \\
      \text{otherwise}: & \Psh\UU(\sem{x_1}, \ldots, \sem{x_n}; \sem y)
    \end{cases}
    \]
    We can extend the semantics function into a multi-functor structure
    $\sem- : \Psh'\UU \to \Psh\UU$ using the Yoneda lemma:
    \[
    \sem{f}_{x_1, \ldots, x_n; y} \definedby
    \begin{cases}
      \forall i. x_i = \injection_2 a_i: &
        \inv\Yoneda_{\UU, \sem y, (a_1, \ldots, a_n)}f \\
      \text{otherwise:} & f
    \end{cases}
    \]
  \item The identities are the identity natural multi-transformations
    over actual presheaves, and the identity multi-morphism at $a$
    over the formal representable by $a$. By fiat, the semantics
    functor-structure preserves identities.
  \item To compose multi-morphisms, apply the semantics
    functor-structure to turn them into composable multi-natural
    transformations, compose them in $\Psh\UU$, and then, if needed,
    apply the Yoneda lemma to get the corresponding element of the
    codomain presheaf. By fiat, the semantics functor-structure
    preserves composition.
  \item We therefore also have a functor from the multi-category
    structure $\Psh'\UU$ to $\Psh\UU$.
  \end{itemize}
  The functor $\sem- : \Psh'\UU \to \Psh\UU$ is fully-faithful, and in
  particular faithful, and therefore $\Psh'\UU$ is a multi-category.
  We will take $E := \sem-$, and $E' \definedby \injection_1 : \Psh\UU
  \to \Psh'\UU$.

  We can define the unit and counit, since the component tells us
  whether we need a multi-transformation or merely its
  representation. For the unit we need to split into cases:
  \[
  \unit_{\injection_1F} \definedby \id[F]
  \qquad
  \unit_{\injection_2a} \definedby \id[a] \in \yoneda aa = \sem{\injection_1\sem{\injection_2a}}a
  \]
  For the counit, we know the domain is an actual presheaf, and so use
  the identity natural multi-transformation.

  I expect these are all natural and satisfy the triangle laws.

  The formal Yoneda functor $\yoneda' : \UU \to \Psh'\UU$ sends each
  $a \in \Obj\UU$ to its formal representable and each multi-morphism
  to itself. Since we started with a locally-small category, we have
  by fiat that $\yoneda'$ is admissble:
  \[
  \Psh'\UU(\yoneda'a_1, \ldots, \yoneda'a_n; x) = \sem x(a_1, \ldots,
  a_n) \in \Set
  \]
  The two transformations $\tau : \sem{\yoneda'-} \to \yoneda$ and
  $\tau' : \injection_1\yoneda- \to \injection_1-$ are then
  identities, since the hom definitions at each object take care of
  everything else.

  What makes the two additional laws hold, in the case that we don't
  have an identity, is that the both the unit for the adjunction and
  our case-splitting composition will choose the correct representative.
\end{proof}

\subsection{The Restricted Yoneda multi-functor}
The structure underlying a Yoneda structure is a choice of 4 components:
\begin{itemize}
\item The admissible Yoneda $1$-cells $\yoneda_{\UU} : \UU \to \Psh\UU$,
  together with their presheaf construction. We covered both in the
  previous section, up to equivalence (\lemmaref{admissible
    yoneda embedding}).
\item For every admissible functor $F : \UU \to \UU[2]$, a $1$-cell
  (non-admissible perhaps) $\UU[2](F;-) : \UU[2] \to \Psh\UU$;
\item and a $2$-cell $\map F : \yoneda_{\UU} \to \UU[2](F;F-)$.
\end{itemize}
\begin{center}
$\quiver{yoneda structure structure}$
\end{center}
This section concerns constructing the remainder of this data, and
proving the first axiom of the Yoneda structure.

Let $F : \UU[1] \to \UU[2]$ be an admissible multi-functor. To keep
the notation less ellipsis-cluttered, given a sequence of
objects/multi-morphisms $\xivec$ in $\UU[1]$, we write $F[\xivec]
\definedby \seq[i=1][\abs\xivec]{F\xivec_i}$ for the componentwise
application of $F$.

We define
the structure of the \emph{$F$-restricted Yoneda multi-functor}
$\UU[2](F;-) : \UU[2] \to \Psh\UU$:
\begin{itemize}
  \item It sends each object $x \in \Obj{\UU[2]}$ to the homset functor
    $\UU[2](F[-]; x) : \opposite\UU \to \Set$ sending:
    \begin{itemize}
    \item each sequence $\avec \in \Obj{\opposite\UU}$ to the
      multi-homset $\UU[2](F[\avec]; x)$; and
    \item each morphism $f : \avec \from \bvec$ to the function:
      \[
      \UU[2](F[f]; x) : \parent{(F[\avec]) \xto u x}
      \mapsto (F[\bvec]) \xto{u \compose (F[f])} x
      \]
    \end{itemize}
  \item It sends each morphism $(\xvec) \xto{v} y$, $n \definedby
    \abs\xvec$ to the following natural multi-transformation
    $\UU[2](F; v): \parent{\UU[2](F; \xvec_1), \ldots,
      \UU[2](F;\xvec_n)} \to \UU[2](F;y)$ given, for every $n$-ary
    partition $\partit$ of $\avec$, by:
    \[
    \UU[2](F; v)_{\partit} : \seq[i=1][n]{F[\partit_i]\xto{v_i} \xvec_i}
    \mapsto (F[\avec])\xto{v\compose(v_1, \ldots, v_n)} y
    \]
\end{itemize}

\begin{lemma}
  These data define a multi-functor $\UU[2](F;-) : \UU[2] \to
  \Psh\UU$.
\end{lemma}
\begin{proof}
  The restricted Yoneda functor is well-defined:
  \begin{itemize}
  \item Each hom $\UU[2](F[\avec]; x)$ is a set since $F$ is admissible.
  \item Given a morphism $\avec \xfrom f \bvec$ in $\UU$, $F[\avec]
    \xfrom{F[f]} F[\bvec]$ in $\UU[2]$, and so given $F[\avec] \xto{u}
    x$ the composition $u \compose (F[f]) : F[\bvec] \to x$
    type-checks as a well-defined multi-morphism.
  \item For identity preservation, calculate:
    \[
    \UU[2](F[{\id[\avec]}]; x)(u) \definedby
      u \compose\seq[i=1][\abs\avec]{F \id[ \avec_i]}
    \explain={$F$ multi-functor}
      u \compose\seq[i=1][\abs\avec]{  \id[F\avec_i]}
    \explain={multi-category}
      u
    \]
    and so $\UU[2](F[{\id[\avec]}]; x)  = \id[{\UU[2](F[\avec]; x)}]$.
  \item Take two composable morphisms in $\opposite\UU$, i.e.: $\avec
    \xfrom f \bvec \xfrom g \cvec$ in $\UU$ and calculate for a
    generic $u : (F[\avec]) \to x$:
    \begin{align*}
      \UU[2](F[f \compose g]; x)u
      &{}= u \compose (F[f \compose g])
      \explain={$\compose$-def in $\opposite\UU$}
      u \compose \seq[i=1][\abs f]{F\parent{f_i \compose
        \seq[j=1][\abs {g_{\pair i*}}]{g_{\pair ij}}}}
      \\&{}\explain={$F$ multi-functor}
      u \compose \seq[i=1][\abs f]{Ff_i\compose
        \seq[j=1][\abs {g_{\pair i*}}]{Fg_{\pair ij}}}
        \explain={associativity of composition\\in a multi-category}
      (u \compose \seq[i=1][\abs f]{Ff_i})\compose
        \seq[\ell=1][\abs g]{Fg_\ell}
      \\{}&=
      \UU[2](F[g]; x) \compose \UU[2](F[f]; x) u
    \end{align*}
    and so $\UU[2](F[f \compose g]; x) = \UU[2](F[g]; x) \compose \UU[2](F[f]; x)$.
    Thus, $\UU[2](F;x) : \opposite\UU \to \Set$ is a presheaf.
  \end{itemize}

  To show that the multi-functorial action is well-defined, take
  $(\xvec) \xto v y$ in $\UU[2]$ with $n \definedby \abs\xvec$, and
  show the naturality condition by taking an $n$-ary partition
  $\seq[i=1][n]{\partit_i \xfrom{\phi_i} \partit'_i}$ of some $\avec
  \xfrom f \bvec$ in $\UU$, chasing an element of
  $\prod_{i=1}^n\UU[2](F[\avec];\xvec _i)$:
  \diagram{29}

  To show this action is multi-functorial, take an identity $\id[x]$
  in $\UU[2]$ and the unique unary partition $\seq{\avec}$ of some
  $\avec$:
  \diagram{30}
%
  Next, given composable multi-morphisms $(\yvec) \xto v z$, $n
  \definedby \abs\yvec$ and $\seq[i=1][n]{(\xvec_i) \xto v_i \yvec_i}$
  in $\UU[2]$, letting $m \definedby \norm\xvec$ and $m_i \definedby
  \abs{\xvec_i}$ and taking any $m$-ary partition of some $\avec$,
  calculate: \diagram{31}
  So $\UU[2](F;v\compose\seq[i=1][n]{v_i})=
  \UU[2](F;v)\compose\seq[i=1][n]{\UU[2](F;v_i)}$ and we have a
  functor.
\end{proof}

The multi-natural transformation $\map F : \yoneda_{\UU} \To
\UU[2](F;F-)$ we now define internalises the multi-functorial action of
$F$. Explicitly, the component at $b \in \Obj\UU$ is given by the
natural multi-transformation $(\map F)_{b} : (\yoneda_{\UU}b) \to
\UU[2](F;Fb)$ whose component at any unary partition
$\seq[i=1][1]\avec$ of some sequence $\avec$ is given by:
\[
(\map F)_{b, \seq[i=1][1]\avec} :
  \prod\set{\yoneda b\avec} \ni
  \seq[i=1][1]{(\avec) \xto{u} b} \mapsto
  (F[\avec] \xto{Fu} Fb) \in \UU[2](F[\avec]; Fb)
\]
or, equivalently, using the Yoneda lemma:
\[
(\map F)_{b, \seq[i=1][1]\avec} \seq{u} \definedby \Yoneda_{F;\avec}(\yoneda(F[\avec])\xto{\yoneda (Fu)}\yoneda (Fb)) \in \yoneda (Fb)(F[\avec]) = \UU[2](F[\avec];Fb)
\]

\begin{lemma}
  This data defines a multi-natural transformation:
  $\map F : \yoneda_{\UU} \To \UU[2](F;F-)$.
\end{lemma}
\begin{proof}
  By Yoneda reduction, the two definitions are
  equivalent. We need to show:
  \begin{itemize}
  \item Well-definedness: each component is multi-natural.
  \item Naturality of the components over-all.
  \end{itemize}
  Unfortunately, these don't just follow automatically from our
  formulation using the Yoneda lemma together with its naturality and
  multi-naturality, since we're applying $F$ to $u$. However, the
  formulation using Yoneda might give some intuition why this
  operation is multi-natural, and our later results will allow us to
  recover it internally.

  To show well-definedness, we need to show that each component $(\map
  F)_{b}$ is a natural multi-transformation. Take the unique unary
  partition $\seq[i=1][1] f$ of some morphism $f : \avec \from \avec'$
  in $\UU$, and chase an element of $\prod\set{\yoneda b\avec}$ around the
  naturality square:
  \diagram{32}

  To show multi-naturality, take any $f : \bvec \to b$ in $\UU$, let
  $n \definedby \abs\bvec$ and consider one component of the two sides
  of the multi-naturality diagram in $\Psh\UU$:
  \[
  (\map F)_{b}\compose(\yoneda f)
  =
  \UU[2](F;Ff)\compose\seq[i=1][n]{(\map F)_{\bvec_i}}
  \]
  by taking an $n$-ary partition $\partit$ of some $\avec$ and
  calculating:
  \diagram{33}
  and so $\map F$ is multi-natural.
\end{proof}

\begin{proposition}\emph{Axiom 1.}
  The multi-functor $\UU[2](F;-)$ together with the multi-natural
  transformation $\map F$ form the left Kan extension of
  $\yoneda_{\UU}$ along $F$:
  \begin{center}
    $\quiver{yoneda structure structure}$
  \end{center}
\end{proposition}
\begin{proof}
  Like the Yoneda lemma, we roll up our sleeves and just do
  it. We need to show a bijective correspondence:
  \begin{center}
    $\quiver{axiom 1 spelled out}$
  \end{center}

  Starting with uniqueness, let $\alpha := (\beta \hcompose F)
  \vcompose (\map F)$ be the pasted multi-natural transformation.


  First, unpack the determining structure of $\beta$.  Take any
  component $x \in \Obj\UU[2]$ of the multi-natural transformation
  $\beta : \UU[2](F;-) \To G$, giving a natural multi-transformation
  $\beta_{x} : (\UU[2](F;x)) \to Gx$. Take any component of
  $\beta_{x}$, i.e.~a uniary partition of some $\avec \in
  \Obj{\opposite\UU}$ and we now have:
  \[
  \beta_{x, \seq[k=1][1]{\avec}} :
  \UU[2](F[\avec]; x) \to Gx(\avec)
  \]
  and we need to determine the image of $\beta_{x, \seq[k=1][1]\avec}$
  on some $F[\avec]\xto v x$, and show that it only depends on $\alpha$.

  Note that we have a natural multi-transformation $Gv : (GF[\avec])
  \to Gx$ in $\Psh\UU$ and so we may use the naturality of each
  component of $\alpha$, form the composite and apply the Yoneda
  lemma to it:
  \[
  Gv \compose \seq[i=1][n]{\yoneda \avec_i \xto{\alpha_{\avec_i}} GF\avec_i}
  : (\yoneda[\avec]) \to Gx
  \implies
  \Yoneda_{Gx, \avec}(Gv\compose \seq[i=1][n]{\alpha_{\avec_i}}) \in Gx(\avec)
  \]
  We will show that:
  \[
  \beta_{x, \seq[k=1][1]{\avec}}\seq[k=1][1]v =
    \Yoneda_{Gx, \avec}(Gv\compose \seq[i=1][n]{\alpha_{\avec_i}})
  \]
  and since the latter only involves $\alpha$, and not $\beta$, then
  pasting with $\map F$ is injective. Indeed:
  \begin{align*}
    \Yoneda_{Gx, \avec}(Gv\compose \seq[i=1][n]{\alpha_{\avec_i}})
    &{}=
    (Gv\compose \seq[i=1][n]{\alpha_{\avec_i}})
      _{\seq[i=1][n]{\seq[j=1][1]{\avec_i}}}\seq[i=1][n]{\id[\avec_i]}
    \\&{}=
    (Gv)_{\seq[i=1][n]{\seq[j=1][1]{\avec_i}}}\seq[i=1][n]{\alpha_{\avec_i}(\id[\avec_i])}
    \\&{}\explain={$\alpha = (\beta \hcompose F)\vcompose (\map F)$}
    (Gv)_{\seq[i=1][n]{\seq[j=1][1]{\avec_i}}}\seq[i=1][n]{\beta_{F\avec_i}(F\id[\avec_i])}
    \\&{}\explain={multi-functoriality}
    (Gv)_{\seq[i=1][n]{\seq[j=1][1]{\avec_i}}}\seq[i=1][n]{\beta_{F\avec_i}(\id[F\avec_i])}
    \\&{}\explain={$\compose$-def in $\Psh\UU$}
    (Gv\compose \seq[i=1][n]{\beta_{F\avec_i}})
       _{\seq[i=1][n]{\seq[j=1][1]{\avec_i}}}\seq[i=1][n]{\id[F\avec_i]}
    \\&{}\explain={multi-naturality}
    (\beta_x \compose (\UU[2](F;v)))
      _{\seq[i=1][n]{\seq[j=1][1]{\avec_i}}}\seq[i=1][n]{\id[F\avec_i]}
    \\&{}\explain={$\compose$-def}
    \beta_{x, \seq[k=1][1]{\avec}} \seq[k=1][1]{
      \UU[2](F;v)_{\seq[i=1][n]{\avec_i}}
      \seq[i=1][n]{\id[F\avec_i]}}
    \\&{}\explain={$\UU[2](F;v)$-def}
    \beta_{x, \seq[k=1][1]{\avec}} \seq[k=1][1]{
      v\compose\seq[i=1][n]{\id[F\avec_i]}}
    \explain={multi-category}
    \beta_{x, \seq[k=1][1]{\avec}} \seq[k=1][1]{v}
  \end{align*}
  and so the correspondence is injective.

  To show surjectivity, consider any multi-natural $\alpha$, and set:
  \[
  \beta_{x, \seq[k=1][1]{\avec}}\seq[k=1][1]v \definedby
    \Yoneda_{Gx, \avec}(Gv\compose \seq[i=1][n]{\alpha_{\avec_i}})
  \]
  We need to show:
  \begin{itemize}
  \item $\beta_{x}$ is a natural multi-transformation;
  \item $\beta$ is a multi-natural transformation;
  \item pasting with $\map F$ recovers $\alpha$, i.e.: $\alpha = (\beta
    \hcompose F)\vcompose (\map F)$.
  \end{itemize}

  To show that $\beta_{x} : \UU[2](F;x) \to Gx$ is natural, take any
  $f : \avec \from \bvec$ in $\UU$ and calculate:
  \begin{align*}
  Gx f(\beta_{x, \seq[k=1][1]{\avec}}\seq[k=1][1]v)
  &{}=
  Gx f\parent{\Yoneda(Gv \compose \seq[i=1][\abs\avec]{\alpha_{\avec_i}})}
  \explain={$\Yoneda$-naturality}
  \Yoneda\parent{Gv \compose
    \seq[i=1][\abs\avec]{\alpha_{\avec_i}\compose \yoneda f_i}}
  \\&{}\explain={$\alpha$-multi-naturality}
  \Yoneda\parent{Gv \compose
    \seq[i=1][\abs\avec]{GF f_i\compose
      \seq[j=1][\abs{\bvec_{\pair i*}}]{
        \alpha_{\bvec_{\pair ij}}}}}
  \\&{}\explain={$G$-multi-functor}
    \Yoneda\parent{G(v \compose Ff_i)\compose
      \seq[\ell=1][\abs{\bvec}]{
        \alpha_{\bvec_{\ell}}}}
  =\beta_{x, \seq[k=1][1]{\bvec}}\seq[k=1][1]{\UU[2](F[f];x)v}
  \end{align*}

  To show that $\beta$ is multi-natural, take any $v : (\xvec) \to y$
  in $\UU[2]$ and $v_i : F[\avec_i] \to \xvec_i$ for $1 \leq i \leq n
  \definedby \abs\avec$:
  \begin{align*}
    (Gv \compose \seq[i=1][n]{\beta_{\xvec_i}})
      _{\seq[i=1][n]{\avec_i}}(v_1, \ldots, v_n)
    &{}=
    Gv\seq[i=1][n]{\Yoneda(Gv_i \compose
      \seq[j=1][\abs{\avec_i}]{\alpha_{\avec_{\pair ij}}})}
    \\&\explain={$\Yoneda$-multi-naturality, associativity}
    \Yoneda((Gv\compose\seq[i=1][n]{Gv_i}) \compose
    \seq[\ell=1][\norm{\avec}]{\alpha_{\avec_{\ell}}})
    \\&{}\explain={$G$-multi-functoriality}
    \Yoneda(G(v\compose\seq[i=1][n]{v_i}) \compose
    \seq[\ell=1][\norm{\avec}]{\alpha_{\avec_{\ell}}})
    \\&{}=
    (\beta_{y}\compose \UU[2](F;v))\seq[\ell=1][n]{v_1, \ldots, v_n}
  \end{align*}

  To show that we recover $\alpha$, take any $u : (\avec) \to b$ and
  calculate:
  \begin{align*}
    ((\beta \hcompose F) \vcompose (\map F))_{b, \seq[k=1][1]\avec}(u)
    =
    \beta(Fu)
    =
    \Yoneda(GFu \compose \seq[i=1][\abs\avec]{\alpha_{\avec_i}})
    \explain={$\alpha$-multi-naturality}
    \Yoneda(\alpha_{b}\compose\yoneda u)
    =\alpha_{b, \seq[k=1][1]\avec}u
  \end{align*}
  concluding the proof.
\end{proof}

\subsection{The Relative Yoneda lemma}
\begin{proposition}\emph{Axiom 2.}
  For every multi-functor (admissible or not) $G : \UU[3] \to \UU$,
  the following triangle is a left Kan lifting of $\yoneda_{\UU}$
  through $\UU[2](F;-)$:
  \begin{center}
    $\quiver{axiom-2 lifting}$
  \end{center}
  I.e., for every $H : \UU[3] \to \UU[2]$, pasting
  $(\map F)\hcompose G$ along $\UU[3] \xto G \UU \xto F \UU[2]$ gives a bijection:
  \begin{center}
    $\quiver{axiom-2 bijection}$
  \end{center}
  We call the bijection $\alpha \mapsto \Yoneda_{\UU,
    F,H,G\seq-}\alpha$ the \emph{Yoneda lemma relative to $G$ along
    $F$}.
\end{proposition}
\begin{proof}
  Starting with uniqueness, given the data:
  \begin{center}
    $\quiver{axiom-2 proof: uniqueness data}$
  \end{center}
  such that
  \begin{center}
    $\quiver{axiom-2 proof: uniqueness assumption}$
  \end{center}
  we will show that $\beta$ is uniquely determined by $\alpha$,
  i.e., the relative Yoneda lemma is injective.

  Take any $e \in \Obj{\UU[3]}$, and then: $\beta_e : (FGe) \to He$ is
  a multi-morphism in $\UU[2]$, i.e.~an element $\beta_e \in
  \UU[2](F[(Ge)]; He)$, and so by the Yoneda lemma we have a
  multi-natural transformation:
  \[
  \alpha'_e \definedby
  \inv\Yoneda_{\UU,\UU[2](F;He),\seq{Ge}}\beta_e : \yoneda (Ge) \to
  \UU[2](F[-];He)
  \]
  which is the same type as the component $\alpha_e$. We will show
  that $\alpha'_e = \alpha_e$, and so $\beta_e$ is uniquely determined
  by $\alpha$:
  \[
  \beta_e = \Yoneda_{\UU,\UU[2](F;He),\seq{Ge}}\alpha_e
  \]
  Take any $\avec \in \Obj{\opposite\UU}$ and $u : (\avec) \to Ge$ in
  $\UU$. On the one hand:
  \begin{align*}
    \alpha'_{e, \avec}(u)
    &{}\explain={$\Yoneda$-def}
    \UU[2](F[u];He)(\beta_e)
    =
    \beta_e\compose(Fu)
  \end{align*}
  On the other hand:
  \begin{align*}
    \alpha_{e, \avec}(u)
    &{}\explain={assumption}
    \parent{(\UU[2](F;-)\hcompose\beta)\vcompose
      (\map F \hcompose G)}_{e,\avec}(u)
    \explain={$(\hcompose$-def$)\vcompose(\map F$-def$)$}
    \UU[2](F;\beta_e)(Fu)
    =
    \beta_e \compose (Fu)
  \end{align*}
  and therefore $\beta$ is uniquely determined by $\alpha$.

  For the existence, set $\beta_e \definedby
  \Yoneda_{\UU,\UU[2](F;He),\seq{Ge}}\alpha_e$ --- this choice
  motivates our notation --- and we'll show:
  \begin{itemize}
  \item $\beta : FG \to H$ is multi-natural; and
  \item $\alpha = (\UU[2](F;-)\hcompose\beta)\vcompose (\map F
    \hcompose G)$
  \end{itemize}
  For multi-naturality, take any $(e_1, \ldots, e_n) \xto{w} e$ in
  $\UU[3]$. We want to prove:
  \[
  Hw \compose (\beta_{e_1}, \ldots, \beta_{e_n}) = \beta_e \compose (FGw)
  \]
  Calculate:
  \begin{align*}
    Hw \compose (\beta_{e_1}, \ldots, \beta_{e_n})
    &{}=
    \UU[2](F;Hw)\seq[i=1][n]{\Yoneda(\alpha_{e_i})}
    \explain={$\Yoneda$-multi-naturality}
    \Yoneda\parent{\UU[2](F;Hw)\compose\seq[i=1][n]{\alpha_{e_i}}}
    \\&{}\explain*={$\alpha : \yoneda(G-) \to \UU[2](F;H-)$\\multi-naturality}
    \Yoneda\parent{\alpha_{e}\compose(\yoneda (Gw))}
    \explain*={$\Yoneda$-naturality}
    \UU[2](F[Gw]; He)(\Yoneda(\alpha_e))
    =
    \beta_e \compose (FGw)
  \end{align*}
  as we wanted.

  To show that we recover $\alpha$ by pasting, take any $e \in
  \UU[3]$, $\avec \in \Obj{\opposite\UU}$, $u : (\avec) \to Ge$. We
  need to show:
  \[
  \alpha_{e, \avec}(u) = \UU[2](F;\beta_e)(\map F u) = \beta_e \compose (F u)
  \]
  Calculate:
  \begin{align*}
    \beta_e \compose (F u)
    &{}=
    \UU[2](F[u]; He)(\Yoneda(\alpha_e))
    \explain={$\Yoneda$-naturality}
    \Yoneda\parent{
      \alpha_e \compose (\yoneda u)
    }
    \explain={$\Yoneda$-multi-naturality}
    \alpha_{e,\avec}(\Yoneda(\yoneda u))
    \explain={Yoneda reduction%
      % vertical alignment hack --- this looks nicer
      \vphantom{y}}
    \alpha_{e,\avec}u
  \end{align*}
  as we wanted.
\end{proof}

\subsection{The presheaf-construction pseudo-functor}
The third axioms concerns internalising the presheaf construction into
a pseudo-functor from the locally-small multi-categories to
multi-categories. Doing so uses an internalisation of the hom
profunctor that extends the postulate structure $\UU[2](F;-)$, first
to an action on $F$, and then to a similar action in the covariant
position.

\paragraph{Internalising the hom profunctor}
The following structure and properties follow from Axiom~1, letting
$\AdMultiCAT$ denote that full sub-$2$-category of $\AdMultiCAT$
consisting of the admissible multi-functors, we have, for every triple
of multi-categories $\UU[1], \UU[2], \UU[3]$, with $\UU[1]$
locally-small, a functor:
\[
\UU[2](-,- -) :
\opposite{\AdMultiCAT(\UU[1], \UU[2])}\times \MultiCAT(\UU[3], \UU[2]) \to
\MultiCAT(\UU[3], \Psh{\UU[1]})
\]
To define it, we first define its partial application with $\UU[3] :=
\UU[2]$, $\Id : \UU[3] \to \UU[2]$, namely:
\[
\UU[2](-,\Id-) :
\opposite{\AdMultiCAT(\UU[1], \UU[2])}\to
\MultiCAT(\UU[2], \Psh{\UU[1]})
\]
given by:
\begin{itemize}
\item sending `objects', i.e., admissible multi-functors $F : \UU[1] \to
  \UU[2]$, to the postulated multi-functor:
  \[
  \UU[2](F,\Id-) \definedby \UU[2](F;-) : \UU[2] \to \Psh{\UU[1]}
  \]
\item sending `morphisms', i.e., multi-natural transformations
  $\quiver{hom-profunctor morphism input type}$, to the unique
  multi-natural transformation:
  \begin{center}
  $\quiver{hom-profunctor morphism output type}$
  \end{center}
  satisfying:
  \begin{center}
    $\quiver{hom-profunctor special case up}$
  \end{center}
  Working this action out explicitly, it's given by the multi-natural
  transformation consisting of natural multi-transformations acting by
  pre-composition with the appropriate component of $\alpha$:
  \[
  \UU[2](\alpha;\Id-)_{x,\avec} :
    \parent{F[\avec] \xto u x} \mapsto
    \parent{G[\avec] \xto{\seq[i]{\alpha_{\avec}}} F[\avec] \xto u}
  \]
\end{itemize}
We then define the more general functor using horizontal composition
as follows. Given:
\begin{center}
  \quiver{hom profunctor general input}
\end{center}
we define:
\begin{center}
  $\quiver{hom profunctor general def via hcompose}$
\end{center}

Explicitly, the functor $\UU[2](-;- -)$ is given by:
\begin{itemize}
\item sending `objects', i.e., admissible multi-functors $F : \UU[1]
  \to \UU[2]$ and arbitrary multifunctor $H : \UU[3] \to \UU[2]$, to
  the composed multi-functor:
  \[
  \UU[2](F,H-) \definedby \UU[2](F;-)\compose H : \UU[3] \to \UU[2]
  \to \Psh{\UU[1]}
  \]
  Its object function is given by:
  \begin{itemize}
  \item for all $x \in \Obj{\UU[2]}$ and $\avec\in \Obj{\opposite\UU}$:
    \[
     \UU[2](F;H-)x\,\avec \definedby \UU[2](F[\avec];Hx)
    \]
  \item for all $x \in \Obj{\UU[2]}$ and $\avec \xfrom{f} \bvec$ in $\UU$:
    \[
    \UU[2](F;H-)x\,f \definedby \UU[2](F[f];Hx) : ((F[\avec]) \xto u Hx) \mapsto
      \parent{(F[\bvec]) \xto{F[f]} (F[\avec]) \xto u Hx}
    \]
  \end{itemize}
  Its hom action is given by:
  \begin{itemize}
\item for all $g : \xvec \to y$ in $\UU[2]$ and $\avec \in \Obj{\opposite\UU}$:
  \[
  \UU[2](F;Hg)_{\avec} \definedby
  \UU[2](F[\avec];Hg) : \seq[i=1][\abs\xvec]{(F[\avec]) \xto{u_i} Hx} \mapsto
  \parent{F[\avec]\xto{Hg\compose \seq[i]{u_i}} Hy}
  \]
  \end{itemize}
\item and its operation on `morphisms':
  \begin{center}
  \quiver{hom profunctor general input-output}
  \end{center}
  is given by the multi-natural transformation whose component at $x$ is the
  natural multi-transformation whose component at $\avec$ is:
  \[
  \UU[2](\alpha;\gamma)_{x, \avec} : \parent{(F[\avec]) \xto{u} Hx} \mapsto
  \parent{(G[\avec]) \xto{\seq[i=1][\abs\avec]{\alpha_{\avec_i}}}
      (F[\avec]) \xto{u} Hx \xto{\gamma_x} Kx}
  \]
\end{itemize}

Street and Walters then use this functor, still only with Axiom~1, to
extend the presheaf construction to a co-lax (2-)functor from the
$2$-category of locally-small multi-categories to multi-categories
$\Psh- : \coop{\LocallySmall\MultiCAT} \to \MultiCAT$, sending:
\begin{itemize}
\item each multi-category $\UU$ to its presheaf construction $\Psh\UU$;
\item each (necessarily admissible) multi-functor $F : \UU \from \UU[2]$
  between locally-small categories to the multi-functor:
  \(
  \Psh F \definedby \Psh{\UU[1]}(\UU[1](\Id;F-), -): \Psh \UU \to \Psh{\UU[2]}
  \) given explicitly by:
  \begin{itemize}
  \item sending each presheaf $P \in \Psh\UU$ to the presheaf $\Psh F
    P := \Psh{\UU[1]}(\UU[1](\Id;F-), P) \in \Psh{\UU[2]}$:
    \begin{itemize}
    \item sending each $\xvec \in \opposite{\UU[2]}$ to the set of
      $\abs\xvec$-ary natural multi-transformations of type $\beta :
      (\UU[1](\Id;F[\xvec])) \to P$, i.e. $\beta : \yoneda_{\UU[1]}(F[\xvec]) \to P$;
    \item sending each $\xvec \xfrom{g} \yvec$ in $\UU[2]$ to the
      natural multi-transformation that acts by pre-composing with the
      multi-transformation that acts by post-composing $Fg$:
      \begin{align*}
      \Psh F g : \parent{\yoneda_{\UU[1]}(F[\xvec]) \xto\beta P}
      &{}\mapsto
      \parent{\yoneda_{\UU[1]}(F[\yvec]) \xto{\yoneda_{\UU[1]}(F[g])}
        \yoneda_{\UU[1]}(F[\xvec])
        \xto\beta P}
      \\&{}=
      \seq[\partit]{
        \seq[\ell=1][\abs\yvec]{(\partit_\ell)\xto{u_\ell} F\yvec_i}
        \mapsto
        \beta_{\partit_{\pair *{\++}}}\compose(Fg_i\compose\seq[j=1][\abs{\yvec_{\pair i*}}]{u_{\pair ij}})
      }
      \end{align*}
    \end{itemize}
  \item sending each natural multi-transformation $\alpha : (\Pvec)
    \to Q$ in $\Psh\UU$ to the natural multi-transformation $\Psh
    F\alpha : (\Psh F [\Pvec]) \to \Psh F Q$ in $\Psh{\UU[2]}$ given
    at partition $\partit$ of $\xvec \in \Obj{\UU[2]}$ by:
    \[
    \seq[i=1][\abs\Pvec]{\yoneda(F[\partit_i]) \xto{\beta_i} P_i}
    \mapsto
    \parent{
      \seq[\ell=1][\abs\xvec]{\yoneda (F\xvec_\ell)}
      \xto{
        \alpha_{F[[\partit]]}\compose\seq[i=1][\abs\Pvec]{\beta_{i, \seq[j=1][1]{\partit_i}}}
      }
      Q\xvec
    }
    \]
  \end{itemize}
\item each multi-natural transformation $\quiver{multi-natural
  transformation coop}$ to the multi-natural transformation whose
  $\quiver{psh 2-cell action type}$ component at $P \in \Psh\UU$ is
  the unary natural multi-transformation whose component at a
  partition $\seq[j=1][1]{\xvec}$ of some $\xvec \in
  \opposite{\UU[2]}$ is:
  \[
  \Psh\alpha_{P,\partit} :
    \seq[j=1][1]{\parent{\yoneda(F[\xvec])} \xto{\beta} P}
    \mapsto
    \parent{\parent{\yoneda(F[\xvec])}
      \xto{\beta \compose \seq[i = 1][\abs\xvec]{\yoneda\alpha_{\xvec_i}}} P}
    \]
\end{itemize}
To exhibit $\Psh- : \coop{\LocallySmall\MultiCAT} \to \MultiCAT$ as a
co-lax $2$-functor, we need two families $2$-cells:
\begin{itemize}
\item For every locally-small multi-category $\UU[1]$, a multi-natural
  transformation:
  \begin{center}
  $\quiver{iota type}$
  \end{center}
  It is given by factorising the identity multi-natural transformation
  between the Yoneda embedding and itself:
  \begin{center}
    $\quiver{internal yoneda lemma}$
  \end{center}
  We denote it by $\Yoneda$ because it is given by the Yoneda
  lemma. We can see that directly by calculation (the factorisation
  maps $\id$ to $(\yoneda \xto u F) \mapsto \Yoneda (\Id u \compose
  \id) = \Yoneda u$), but also by noting:
\begin{itemize}
\item The components $\Yoneda_{F,\avec}$ are natural in $\avec$ and so
  we have a natural multi-transformation $\Yoneda_{F}:
  \Psh\UU(\yoneda; F) \to F$, and these are further multi-natural in
  $F$ and so we have a multi-natural transformation $\Yoneda :
  \Psh\UU(\yoneda;-) \to \Id$.
\item By Yoneda reduction:
  \[
  ((\Yoneda \hcompose \yoneda) \vcompose (\map \yoneda)) (u)
  = \Yoneda(\yoneda u) = u
  \]
  and so $\Yoneda$ is the left extension.
\end{itemize}
\item We will also need, for every triple of locally-small
  multi-catgeories and composable multi-functors $\UU[1] \xfrom F
  \UU[2] \xfrom G \UU[3]$ a $2$-cell:
  \begin{center}
    $\quiver{gamma type}$
  \end{center}

  Street and Walters axiomatise a more general structure with a
  stronger axiom. For a composable pair of multi-functors $\UU[1]
  \xfrom H \UU[2] \xfrom G \UU[3]$ consisting of a (necessarily
  admissible) multi-functor $G$ between locally-small multi-categories
  and admissible functor $K$, a $2$-cell:
  \begin{center}
    $\quiver{theta type}$
  \end{center}
  This $2$-cell specialises to $\gamma_{F,G}$ taking $\UU[1]
  \definedby \Psh\UU$ and
  \[
  H \definedby \Psh\UU(-,(\yoneda \compose F)-) :
  \Psh{\UU[1]} \xfrom{\yoneda} \UU[1] \xfrom F \UU[2]
  \]
  The defining universal property for $\theta_{G, H}$ is that it's
  the unique factorisation:
  \begin{center}
    \quiver{theta UP}
  \end{center}
  It consists, for each $\UU[1]$ object $a$, and
  $\opposite{\UU[3]}$ object $\zvec$, of the functions:
  \[
  \theta_{G, H, a, \seq[j=1][1]{\zvec}} : \seq[j=1][1]{HG[\zvec] \xto{u} a}
    \mapsto
    \seq[\yvec \in \opposite{\UU[2]}]
        { \seq[i=1][\abs\zvec]{\yvec_i \xto{v_i} G[\zvec_i]}
          \mapsto
          \parent{H[\yvec]\xto{u\compose\seq[i=1][\abs\zvec]{Hv_i}} a}
    }
  \]
\end{itemize}

Axiom~3 states that the $2$-cells $\Yoneda$ and $\theta$ are
$2$-isomorphisms, and in particular making $\Psh- :
\coop{\LocallySmall\MultiCAT} \to \MultiCAT$ a pseudo-functor.  Before
we prove it, we'll first characterise the multi-natural isomorphisms,
and in presheaf multi-categories that will require us to characterise
the natural multi-isomorphisms.

\begin{lemma}\emph{Isomorphism lemma.}
  \begin{itemize}
  \item
    Let $\UU$ be a multi-category, $P, Q \in \Psh\UU$ presheaves over
    $\UU$, and $\alpha : (P) \to Q$ a \emph{unary} natural
    multi-transformation. Then $\alpha$ is an isomorphism iff all of
    its components are bijective.
  \item
    A multi-natural transformation $\quiver{multi-natural
      transformation type}$ is an isomorphism (w.r.t. $2$-cell
    vertical composition) iff its components $\alpha_{a} : (Fa) \to
    Ga$ are isomorphisms in $\UU[2]$.
  \end{itemize}
\end{lemma}
\begin{proof}
  The `only if' directions follow since the identities are
  component-wise identity functions/morphisms and the appropriate
  composition is given component-wise (up-to rearrangement of tuples
  via a canonical isomorphism). So the remainder of proof amounts to
  showing that appropriate multi-naturality and naturality conditions
  for the assembled component-wise inverse.

  Starting with a component-wise bijective natural
  multi-transformation $\alpha : (P) \to Q$, define its inverse
  componentwise, at each unary partition $\seq[j=1][1]{\avec}$:
  \[
  (\inv\alpha)_{\seq[j=1][1]{\avec}} : \prod\set{Q\avec} \xto\isomorphic Q\avec
  \xto{\inv{\parent{\alpha_{\seq[j=1][1]{\avec}}}}}
  \prod\set{P\avec}
  \xto\isomorphic
  P\avec
  \]
  A straightforward calculation shows that this multi-transformation
  would be the inverse multi-morphism, provided it is indeed
  natural. Take any $f : \avec \from \bvec$ in $\UU$, and calculate:
  \diagram{34}

  Next, consider a multi-natural transformation $\quiver{multi-natural
    transformation type}$ whose components are
  multi-isomorphisms. Define $(\inv\alpha)_{a} : (Ga) \to Fa$ by:
  \[
  (\inv\alpha)_{a} \definedby \inv{\parent{\alpha_{a}}}
  \]
  and we now need to show $\inv\alpha : G \to F$ is
  multi-natural. Take any $f : (\avec) \to b$ in $\UU[1]$, let $n
  \definedby \abs\avec$, and calculate: \begin{align*}
    Gf \compose
    \seq{(\inv\alpha)_{\avec_i}}
    &{}\explain={identity + $\inv\alpha$-def}
    \id[b]\compose (Gf \compose
    \seq[i=1][n]{\inv{\parent{\alpha_{\avec_i}}}})
    \explain={inverses}
    \parent{\inv{\parent{\alpha_{b}}}
      \compose (\alpha_{b})}
    \compose (Gf \compose
    \seq[i=1][n]{\inv{\parent{\alpha_{\avec_i}}}})
    \\&{}\explain={$\compose$-assoc.}
    \inv{\parent{\alpha_{b}}}
      \compose \parent{\alpha_{b}
    \compose (Gf \compose
    \seq[i=1][n]{\inv{\parent{\alpha_{\avec_i}}}})}
    \explain={multi-naturality}
    \inv{\parent{\alpha_{b}}}
      \compose \parent{Ff \compose
        \seq[i=1][n]{\alpha_{\avec_i}\compose (\inv{\parent{\alpha_{\avec_i}}})}}
    \\&{}\explain={inverses}
    \inv{\parent{\alpha_{b}}}
      \compose \parent{Ff \compose
        (\seq[i=1][n]{\id[\avec_i]})}
    \\&{}\explain={identities}
    \inv{\parent{\alpha_{b}}}
      \compose (Ff)
  \end{align*}
  as we want.  This proof is the usual one presented algebraically
  since we don't have a convenient diagrammatic formalism in the
  multi-categorical setting.
\end{proof}

\begin{proposition}\emph{Axiom 3.}
  \begin{itemize}
  \item
    The multi-natural transformation $\Yoneda : \Psh{\Id[\UU]} \to
      \Id[\Psh\UU]$ is a multi-natural isomorphism.
  \item
    The multi-natural transformations $\theta_{G, H} : \UU[3](H
    \compose G, -) \to \Psh G \compose \UU(H, -)$ are multi-natural
    isomorphisms.
  \end{itemize}
\end{proposition}
\begin{proof}
  By the Yoneda lemma, each component of $\Yoneda$ is a bijection, and
  by the Isomorphism Lemma (1) each component $\Yoneda_P$ is a natural
  multi-isomorphism, and therefore (2) $\Yoneda$ is a multi-natural
  isomorphism.

  Turning to $\theta_{G, H}$ for $\UU[1] \xfrom H \UU[2] \xfrom G
  \UU[3]$, we will use its UP to show that each component:
  \[
  \theta_{G, H, a,
    \seq[k=1][1]\zvec} : \UU[1](HG[\zvec], a) \to
  \Psh{\UU[2]}(\yoneda[G[\zvec]], \UU[1](H, a))
  \]
  is in fact the inverse Yoneda Lemma:
  \[
  \inv\Yoneda_{\UU[1]; \UU(H[-],a), G[\zvec]} : \UU[1](HG[\zvec], a) \to
  \Psh{\UU[2]}(\yoneda[G[\zvec]], \UU[1](H, a))
  \]
  and therefore bijective. Showing the universal property means
  showing that $\inv\Yoneda$ is multi-natural and that:
  \begin{center}
    $\quiver{axiom 3b proof goal}$
  \end{center}
  Multi-naturality follows from the Isomorphism Lemma, and we'll
  establish the equation by direct calculation, taking any $z \in
  \Obj{\UU[3]}, \zvec \in \opposite{\UU[3]}$ and $w : (\zvec) \to z$
  in $\UU[3]$ and showing that:
  \begin{multline}\label{axiom3b RTP}
  \parent{(\Psh G \hcompose (\map H) \hcompose G)\vcompose (\map (\yoneda \compose G))}_{z, \seq[k=1][1]\zvec} w
  \\=
  \parent{
    (\inv\Yoneda \hcompose (G \compose H)) \vcompose
    (\map(H\compose G))}_{z, \seq[k=1][1]\zvec} w
  \mspace{100mu}
  \end{multline}
  Type-checking, this equation is between natural multi-transformations within the multi-hom-set:
  \[
  \Psh{\UU[2]}(\yoneda (G[\zvec]), \UU(H, HGz))
  \]
  and so let $n \definedby \abs\zvec$, and take an $n$-ary partition
  $\partit$ of some $\yvec \in \opposite{\UU[2]}$, and any $n$-tuple
  of multi-morphisms:
  \[
  \seq[i=1][n]{(\partit_i) \xto{v_i} G\zvec_i)} \in \prod_{i=1}^n\yoneda(G\zvec_i)(\partit_i)
  \]
  and apply both sides of \eqref{axiom3b RTP} at the $\partit$-component to this tuple.

  On the LHS:
  \begin{multline*}
    \parent{\parent{(\Psh G \hcompose (\map H) \hcompose G)\vcompose (\map (\yoneda \compose G))}_{z, \seq[k=1][1]\zvec} w}_{\partit}\seq[i=1][n]{v_i}
    \\
    \begin{aligned}
      &{}=
      \parent{(\Psh G (\map H)_{Gz})_{\seq[k=1][1]\zvec} (\yoneda (Gw))}_{\partit}\seq[i=1][n]{v_i}
      \explain={$\Psh G$ action on natural \\multi-transformations}
      (\map H \compose (\yoneda (Gw)))_{\partit}\seq[i=1][n]{v_i}
      \\&{}=
      \map H(Gw \compose \seq[i=1][n]{v_i})
      \explain={$\map H$-def, \\multi-functoriality}
      HGw \compose \seq[i=1][n]{Hv_i}
    \end{aligned}
  \end{multline*}
  On the RHS:
  \begin{multline*}
    \parent{\parent{
    (\inv\Yoneda \hcompose (G \compose H)) \vcompose
        (\map(H\compose G))}_{z, \seq[k=1][1]\zvec} w}_{\partit}\seq[i=1][n]{v_i}
    \\
    \begin{aligned}
      &{}=
      \inv\Yoneda_{;\UU(H;a);\yvec}\parent{\map (H \compose G)w}\seq[i=1][n]{v_i}
      \explain={$\inv\Yoneda$-def,$\map$-def}
      \UU(H[\seq[i=1][n]{v_i}];a)(HGw)
      =
      HGw \compose \seq[i=1][n]{Hv_i}
    \end{aligned}
  \end{multline*}
  and so \eqref{axiom3b RTP} holds, and $\theta_{G,H}$ is a
  multi-natural isomorphism.
\end{proof}

% \subsection{Failure of Axiom $3^*$}
% Street and Walters also suggest a stronger axiom:
% \paragraph{Axiom $3^*$}
% Every multi-natural transformation $\sigma$ as on the left such that
% the diagram on the right is an absolute lifting of $\yoneda$ through
% $H$, must be a multi-natural isomorphism:
% \begin{center}
%   $\quiver{axiom 3* assumption}$
% \end{center}
% \begin{counterexample}\emph{Failure of Axiom $3^*$.}
%   Take:
%   \begin{itemize}
%   \item
%     $\UU$ to be the multi-category with no objects and no
%     multi-morphisms whatsoever;
%   \item $\UU[2]$ to be the multi-category with a single object $x$ and
%     only the identity on $x$;
%   \item $F : \UU \to \UU[2]$ to be the (unique) multi-functor of this type.
%   \end{itemize}
%   The relative Yoneda functor $\UU[2](F,-)$ then sends $x$ to the
%   presheaf constantly taking the empty set.
%   \begin{itemize}
%   \item Take $H : \UU[2] \to \Psh\UU$ to be the presheaf constantly
%     taking the singleton $\terminal \definedby \set{\star}$.
%   \item Take $\sigma : \UU[2](F, -) \to H$ to be the only possible
%     multi-transformation.
%   \end{itemize}
%   Since the sink in $\sigma$'s naturality diagrams is the singleton,
%   it is a natural multi-transformation $\sigma : \UU[2](F, -) \to
%   H$. By the Isomorphism Lemma, since the source and target sets of
%   $\sigma_{x, \seq{}} : \emptyset \to \terminal$ have different
%   cardinalities, $\sigma$ is not a natural multi-isomorphism.

%   However, the appropriate triangle is an absolute lifting. Indeed,
%   take any multifunctors: $W : \UU[3] \to \UU$ and $G : \UU[3] \to
%   \UU[2]$. We will show that there are unique multi-natural
%   transformations $\alpha$ (left) and $\beta$ (right):
%   \begin{center}
%     $\quiver{Axiom 3* counterexample goal}$
%   \end{center}
%   The pasting (on the right) then forms the bijection we want.

%   Since $\UU$ has no objects or morphisms, we must have that $\UU[3]$
%   doesn't have any objects or morphisms either, and so it is identical
%   $\UU$. Since $F$ is the unique multi-functor out of $\UU$, we have
%   that $G$ is the unique multi-functor out of $\UU[3]$, and $\beta$
%   must be the identification of the two. Since $H \compose G$ is
%   componentwise a singleton, each component of $\alpha$ must be the
%   unique map into the singleton, and so $\alpha$ is unique.
% \end{counterexample}

\subsection{Representability}
We can now reap the fruit of our labour:

\begin{proposition}\emph{Characterisation: multi-categorical adjunctions.}
\propositionlabel{representability multi-categorical adjunctions}
Let $\UU[1]$ be a locally-small multi-category. TFAE:
\begin{itemize}
\item An adjunction $\triple{\quiver{adjunction type}}{\Id \xto{\unit}
  UF}{FU \xto\counit \Id}$ in $\MultiCAT$ with $F$ admissible.
\item A multi-natural isomorphism $\quiver{adjunction hom-set bijection}$ with $F$ admissible.
\item A multi-functor $U : \UU[2] \to \UU[1]$, a choice of objects $Fa
  \in \UU[2]$ for all $a \in \UU[1]$ together with a (multi-)morphism
  $\unit_a : (a) \to UFa$ such that:
  \begin{itemize}
  \item All homs $\UU[2](Fa_1, \ldots, Fa_n; x)$ are small; and
  \item The components of the transformation $\unit$ are
    multi-universal: for all $(a_1, \ldots, a_n) \xto f U x$ in
    $\UU[1]$, there exist a unique $(Fa_1, \ldots, Fa_n) \xto{\psi(f)}
    x$ in $\UU[2]$ such that:
    \[
    U (\psi(f))\compose (\unit_{a_1}, \ldots, \unit_{a_n}) = f
    \]
  \end{itemize}
\item
  An admissible multi-functor $F : \UU[1] \to \UU[2]$, a choice of
  objects $Ux \in \UU[1]$ for all $x \in \UU[2]$ together with a
  (multi-)morphism $\counit_x : (FUx) \to x$ that is universal: for
  all $(Fa_1, \ldots, Fa_n) \xto g x$ in $\UU[2]$, there exist a
  unique $(a_1, \ldots, a_n) \xto{\phi g} U x$ in $\UU[1]$ such that:
  \[
  \counit_x\compose (F (\phi g)) = g
  \]
\end{itemize}
\end{proposition}
\begin{proof}
  The first two items are Prop.~8 of \cite{street-walters-yoneda-structures}.
  Given a multi-natural isomorphism $\phi$ with inverse $\psi$, take:
  \[
  \unit_a \definedby \phi\id
  \qquad
  \counit_x \definedby \psi\id
  \]
  Take any $f : (\avec) \to U x$, $n \definedby \abs\avec$, and chase
  $\seq[i=1][n]{\id[a_n]} \in \prod_{i=1}^n\UU[2](Fa_i; Fa_i)$ around
  the multi-naturality condition for $\phi$ when instantiated for
  $\psi f$ in $\UU[2]$.  OTOH:
  \[
  \phi\parent{\UU[2](F;\psi f)\seq[i=1][n]{\id[a_i]}}
  =
  \phi\parent{\psi f \compose \seq[i=1][n]{\id[a_i]}}
  =
  \phi\parent{\psi f}
  =
  f
  \]
  and OTOH:
  \[
  \UU[1](\Id; U(\psi f))\seq[i=1][n]{\phi(\id[a_i])}
  =
  U(\psi f)\compose\seq[i=1][n]{\unit_{a_i}}
  \]
  and we proved the third item.

  Take any $g : (F[\avec]) \to x$ in $\UU[2]$, and chase $\id[x]$
  around the naturality condition for $\psi_x$ when applied to $\phi g
  : \seq[k=1][1]{U x} \from \avec$. OTOH:
  \[
  \psi\parent{\UU(\phi g; Ux)\id}
  =
  \psi\parent{\id \compose (\phi g)}
  = \psi (\phi g) = g
  \]
  and OTOH:
  \[
  \UU[2](F(\phi g); U x)(\psi\id[x])
  =
  \counit_x\compose (F\phi g)
  \]
  as we wanted.

  Therefore, the first two conditions imply the last two conditions.
  We'll now show that the third and fourth conditions imply the third.
  The exists-unique conditions construct the required bijections, and
  we need to equip $F$ or $U$ with a multi-functorial action and show
  that the given bijection is multi-natural in $x$ and natural in
  $\avec$, in each case.

  First, assume we have a choice of multi-universal units. Extend $F$
  to a functor, by sending each $f : (\avec) \to b$ in $\UU[1]$ to:
  \[
  \psi(\unit_b\compose (f))
  \qquad
  \text{i.e., the unique $g : (F[\avec]) \to Fb$ satisfying: }
  \qquad
  Ug \compose \seq[i=1][\abs\avec]{\unit_{\avec_i}} = \unit_b \compose (f)
  \]
  This assignment is functorial:
  \[
  U\id[Fa] \compose (\unit_{a})
  \explain={$U$ multi-functor}
  \unit_{a}
  =
  \unit_{a}\compose(\id[a])
  \qquad\implies\qquad
  F\id[a] = \id[Fa]
  \]
  and for all composable $f : (\bvec) \to c$, $n \definedby \abs\bvec$
  and $\fvec \definedby \seq[i=1][n]{f_i : (\avec_{\pair i*}) \to
    \bvec_i}$, $p \definedby \abs\avec$,$m_i \definedby
  \abs{\avec_{\pair i *}}$, we have:
  \begin{align*}
    U(Ff\compose \seq[i=1][n]{F\fvec_i}) \compose \seq[\ell = 1][p]{\unit_{a_\ell}}
    &{}\explain={$U$ multi-functor}
    (UFf \compose UF\fvec_i)\compose \seq[\ell = 1][p]{\unit_{a_\ell}}
    \explain={$\compose$-assoc.}
    UFf \compose (UF\fvec_i\compose \seq[j = 1][m_i]{\unit_{a_{\pair ij}}})
    \\&{}= UFf \compose (\unit_{\bvec_i}\compose (\fvec_i))
    = (UFf \compose \unit_{\bvec_i})\compose (\fvec)
    = (\unit_c \compose UFf) \compose (\fvec)
    \\&{}= \unit_c \compose (UFf \compose (\fvec))
  \end{align*}
  and so $F(f \compose (\fvec)) = F f \compose (F[\fvec])$.

  Fix $x$. To show that $\psi_x$ is natural, take any $\avec \xfrom \uvec
  \bvec$ in $\UU[1]$, and chasing any $f : (\avec) \to U x$ around the
  naturality condition for $\psi$ gives us the proof obligation:
  \[
  \psi(f \compose (\uvec))
  =
  \psi f \compose (F[\uvec])
  \]
  Calculate:
  \begin{align*}
    U\parent{\psi f \compose (F[\uvec])} \compose
    \seq[\ell = 1][\abs\bvec]{\unit_{\avec_\ell}}
    &{}\explain={$F$-multi-functor\\$\compose$-assoc.}
    U(\psi f) \compose \seq[i=1][\abs\avec]{UF\uvec_i \compose
      \seq[j = 1][\abs{\bvec_{\pair i*}}]{\unit_{\avec_{\pair ij}}}}
    \explain={$F$-def}
    U(\psi f) \compose \seq[i=1][\abs\avec]{\unit_{\bvec}\compose (\uvec_i)}
    \\&\explain={$\compose$-assoc\\$\psi$-def}
    f \compose (\uvec)
  \end{align*}
  and so $\psi_x$ is natural.

  To show $\phi$ is multi-natural, take any $(\xvec)\xto v y$ in
  $\UU[2]$, and chase an arbitrary tuple
  $\seq[i=1][\abs\xvec]{\avec_{\pair i*}\xto{\fvec_i} U\xvec_i}$ in
  the multi-naturality condition to get the proof obligation:
  \[
  \psi(Uv \compose (\fvec)) = v \compose (\psi[\fvec])
  \qquad
  \impliedby
  \qquad
      U\parent{v \compose (\psi[\fvec])} \compose
    \seq[\ell=1][\abs\avec]{\unit_{\avec_\ell}}
    \explain={$U$ multi functor\\$\compose$-assoc.\\$\psi$-def}
    Uv \compose (\seq[i=1][\abs\xvec]{\fvec_i}
  \]
  and so $\psi$ is multi-natural. By the Isomorphism Lemma, it is a
  multi-natural isomorphism, and we proved that the third condition
  implies the second.

  The proof that the final condition implies the second is similar.
  To extend $U$ to a functor, send $g : (\xvec)\to y$ in $\UU[2]$ to:
  \[
  \phi(g\compose\seq[i=1][\abs\xvec]{\counit_{\xvec_i}})
  \qquad
  \text{i.e., the unique $f : (U[\xvec]) \to Uy$ satisfying: }
  \qquad
  \counit_y\compose (Ff) = f\compose\seq[i=1][\abs\xvec]{\counit_{\xvec_i}}
  \]
  This assignment is functorial:
  \[
  \counit \compose (F\id) \explain={$F$ multi-functor}
  \counit \compose (\id) = \id\compose (\counit)
  \qquad\implies\qquad
  U\id = \id
  \]
  and for all composable $(\xvec) \xto{\gvec} (\yvec) \xto g z$:
  \[
  \counit_z \compose(F(Ug \compose (U[\gvec])))
  \explain={$F$ multi-functor\\$\compose$ assoc.}
  (\counit_z\compose FUg)\compose (FU[\gvec])
  =
  g \compose \seq[i]{\counit_{\yvec_i}\compose (FU[\gvec_i])}
  =
  (g \compose (\gvec))\compose \seq[\ell]{\counit_{\xvec_\ell}}
  \]
  and so $U(g \compose (\gvec)) = Ug \compose (U[\gvec])$.

  Fix $x$. To show that $\phi_x$ is natural, take any $\avec
  \xfrom\uvec \bvec$ in $\UU[1]$ and chase any $v : (F[\avec]) \to x$
  through the multi-naturality condition to get the proof obligation:
  \[
  \phi\parent{v \compose F[\uvec]} =
  (\phi v) \compose (\uvec)
  \impliedby
  \counit_x\compose F((\phi v)\compose (\uvec))
  \explain={$F$ multi-functor\\$\compose$ assoc.}
  (\counit_x\compose F\phi v))\compose (F[\uvec])
  \explain={$\phi$-def}
  v \compose (F[\uvec])
  \]

  Finally, to show that $\phi$ is natural, take any $v : (\xvec) \to
  y$ and chase $\fvec : \avec \to U[\xvec]$ in the multi-naturality
  condition to get the proof obligation:
  \[
  \phi\parent{v \compose (\fvec)} = Uv \compose \phi[\fvec]
  \impliedby
  \begin{aligned}[t]
  \counit_y\compose (F\parent{Uv \compose \phi[\fvec]})
  &{}\explain={$F$ multi functor\\$\compose$ assoc.}
  (\counit_y \compose FUv)\compose (F\phi[\fvec])
  \explain={$U$-def}
  v \compose \seq[i]{\counit_{\xvec_i}\compose F\phi[\fvec]}
  \\&{}\explain={$\phi$-def}
  v \compose (\fvec)
  \end{aligned}
  \]
  as we wanted.
\end{proof}

\section{Multi-homomorphic extension}
\label{sec:multihom-extension}

By generalities, the forgetful functor $\carrier - : \Mod(\TT, \Set)
\to \Set$ always has a left adjoint assigning to each set the free
$\TT$-algebra over it. This situation does not extend to
multi-categories: we may not have a left adjoint (in $\MultiCAT$) to
the forgetful multi-functor $\carrier- : \Multi(\TT, \Set) \to
\pair\Set\prod$. However, we will have such an adjoint in many cases
of interest, and our goal is to prove the following characterisation:
\begin{theorem}\theoremlabel{left adjoint iff commutative}
  Let $\TT$ be a theory. The forgetful multi-functor $\carrier- :
  \Multi(\TT, \Set) \to \pair\Set\prod$ has a left adjoint iff $\TT$
  is commutative.
\end{theorem}
The two directions employ quite different techniques, and so we'll
split them into two parts: \propositionref{multi-hom adjunction
  existence} and \propositionref{multi-adjoint forces commutativity}.

In the commutative case, the result follows from, or is already
present in, the work of
\cite{kock:closed-categories-generated-by-commutative-monads} on
commutative monads. Let $T$ be a commutative monad on a finitely
complete cartesian-closed category $\C$. We can internalise the
$T$-homomorphisms, i.e., the hom-set $\C^T(A,B)$, as the equaliser:
\[
\quiver{hom-algebra carrier}
\]
and equip with a $T$-algebra structure defined by this property:
\diagram{35} where the morphism $\lambda : T(Y^X) \to (TY)^X$ is given
by currying the strength and evaluating $T(Y^X)\times X \xto\str
T(Y^X\times X)\xto\eval TY$. This construction extends to a functor
$(\multimap) : \opposite(\C^T) \times \C^T \to \C^T$, as the unique
action that makes $e_{A,B} : \carrier{A \multimap B} \to \carrier
B^{\carrier A}$ a natural transformation. This functor equips the
Eilenberg-Moore category with the structure of a \emph{closed
  category}~\cite{eilenberg-kelly:closed-categories}, which we will
not recount here.

\begin{proposition}
  Assume $\C$ is complete. For every $T$-algebra $A$, the functor $A
  \multimap - : \C^T \to \C^T$ preserves small limits.
\end{proposition}

I didn't find this proposition anywhere so we'll prove it here, but
I'm sure it's well-known.
\begin{proof}
  Consider a small diagram $D : I \to \C^T$ and a limiting cone $\pair
  L\ell$ for it. We need to show that $\pair{A\multimap L}{A\multimap
    \ell}$ is limiting for $A\multimap D-$. We'll do it in 2 stages:
  \begin{itemize}
  \item We'll show that post-composing with the carrier preserves the
    limit: $\pair{\carrier{A\multimap L}}{\carrier{A\multimap \ell}}$
    is limiting for $\carrier{A\multimap D-}$.
  \item We'll use the fact that $\carrier- : \C^T \to \C$ creates
    limits to deduce the result.
  \end{itemize}

  The first step follows because we can partially-evaluate limits in
  any order. Let $D' : I' \to \C$ be the diagram consisting of
  equaliser diagrams:
  \begin{center}
    $\quiver{A multimap limit preservation diagram}$
  \end{center}
  We can calculate this limit in two ways:
  \begin{itemize}
  \item First take the limit w.r.t.~$i$ and, since the exponentials
    $-^{\carrier A}$ and $-^{T\carrier A}$ are right adjoints and so
    preserve limits, get the diagram:
    \begin{center}
      $\quiver{limit Di first}$
    \end{center}
    and then take its limit, which is $\carrier{A \multimap L}$.
  \item First take the limit w.r.t.~the equaliser diagrams to get the
    diagram $\carrier{A\multimap Di}$.
  \end{itemize}
  Therefore $\Lim_{i \in I}\carrier{A\multimap Di} =
  \pair{\carrier{A\multimap L}}{\carrier{A\multimap \ell}}$.

  Since $\carrier- : \C^T \to \C$ creates limits, it in particularly
  reflects limits, and so $\pair{A\multimap L}{A\multimap \ell}$ is
  limiting.
\end{proof}

\begin{proposition}
  The functors $A\multimap- : \Mod(\TT, \Set) \to \Mod(\TT, \Set)$
  have a left adjoint $-\tensor A : \Mod(\TT, \Set) \to \Mod(\TT,
  \Set)$.
\end{proposition}
\begin{proof}
  This proposition follows by generalities.  The category $\Mod(\TT,
  \Set)$ is complete and cocomplete, and the functor $A\multimap-$ is
  continuous and satisfies Freyd's solution-set condition. Therefore,
  by the adjoint functor theorem, the left adjoint exists.
\end{proof}

This tensor product generalises the linear-hom tensor from linear
algebra and its construction in general is non-effective. It equips
$\C^T$ with the structure of a symmetric monoidal closed category.

As any Eilenberg-Moore category for a strong monad over a
cartesian-closed category, $\C^T$ has powers, i.e., for every set $X$
and $T$-algebras $\Gamma$, $A$, we have a representation:
\[
\begin{adjunctions}
  X & \C^T(\Gamma,B) & \text{in $\Set$}
  \\
  \Gamma & B^X & \text{in $\C^T$}
\end{adjunctions}
\]
For $\C = \Set$, $B^X$ is given by the componentwise structure on
$\carrier B^X$.
\begin{proposition}\propositionlabel{powers vs free internal homs}
  Let $T$ be a commutative monad over $\Set$. Then \( B^X \isomorphic
  (FX \multimap B) \) in $\Set^T$.
\end{proposition}
\begin{proof}
  We construct the bijection using naturality:
  \[
  \begin{adjunctions}
    \terminal & \carrier{FX \multimap A} & \text{in $\Set$} \\
    FX & A & \text{in $\Set^T$, by definition} \\
    X & \carrier A & \text{in $\Set$, by definition} \\
    \terminal & \carrier A^X = \carrier{A ^X} & \text{in $\Set$}
  \end{adjunctions}
  \]
  and so we have a bijection of sets $\carrier{FX \multimap A}
  \isomorphic \carrier{A^X}$. Chasing elements, this bijection sends
  $h : FX \to A$ to $h\compose \eta : F \to \carrier A$ one way, and
  $f : X \to \carrier A$ to $\bind f : FX \to A$ the other way.

  We need to show it is a $T$-isomorphism, and it suffices to show one
  of these functions to be a $T$-homomorphism. We'll show that
  \[
  \phi :
  \carrier{FX \multimap A} \xto{e} \carrier A^{TX} \xto{\carrier
    A^{\unit_X}} \carrier A^X
  \] is a $T$-homomorphism:
  \diagram{36}
  as we wanted.
\end{proof}

We can now prove the existence of the adjoint in the commutative case:
\begin{proposition}\propositionlabel{multi-hom adjunction existence}
  Let $\TT$ be a commutative theory. Then there is a left adjoint to
  the multi-functor $\carrier- : \Multi(\TT, \pair\Set\prod) \to
  \pair\Set\prod$.
\end{proposition}
\begin{proof}
  By \propositionref{representability multi-categorical adjunctions}, it
  suffices to show that $FX$, the free $\TT$-algebra on $X$, together with
  the unit $\unit : X \to \carrier {FX}$ for the left adjoint to the
  functor $\carrier- : \Mod(\TT, \Set) \to \Set$, is multi-universal.

  To construct the bijection $\psi : \pair\Set\prod(\Xvec;\carrier A)
  \xto\isomorphic \Multi(\TT,\Set)(F[\Xvec];A)$, we'll note that:
  \begin{itemize}
  \item The homomorphisms $\bigotimes_{i=1}^{\abs\Xvec} F\Xvec_i \to A$
    are precisely the multi-homomorphisms
    $(\seq[i=1][\abs\Xvec]{F\Xvec_i}) \to A$
    (taking $\bigotimes \emptyset := F\terminal$, the unit for $\tensor$).
  \item We have a natural bijection:
    \[
    \begin{adjunctions}
      \bigotimes_{i=1}^{\abs\Xvec} F\Xvec_i & A & \text{in $\Mod(\TT, \Set)$}\\
      \prod_{i=1}^{\abs\Xvec} \Xvec_i & \carrier A & \text{in $\Set$}
    \end{adjunctions}
    \]
    and moreover the bijection, moving from top to bottom, is given by
    $f \mapsto \carrier f \compose (\unit, \ldots, \unit)$.
  \end{itemize}
  We prove the second point by induction on $\abs\Xvec$. For
  $\abs\Xvec = 0$, we have already proved this fact in
  \propositionref{powers vs free internal homs}. Assume precomposition
  by the units provides a bijection for $\Xvec$ and consider $\Xvec'
  := X :: \Xvec$. We then have:
  \[
  \begin{adjunctions}
    FX\tensor \bigotimes_{i=1}^{\abs\Xvec} F\Xvec_i & A & \text{in $\Mod(\TT,\Set)$} \\
    \bigotimes_{i=1}^{\abs\Xvec} F\Xvec_i & FX\multimap A & \text{in $\Mod(\TT,\Set)$} \\
    \bigotimes_{i=1}^{\abs\Xvec} F\Xvec_i & A^X & \text{in
      $\Mod(\TT,\Set)$, by \propositionref{powers vs free internal
        homs}} \\
    \prod_{i=1}^{\abs\Xvec} \Xvec_i & \carrier{A^X} = \carrier A^X & \text{in
      $\Set$, by induction hypothesis}\\
    X \times \prod_{i=1}^{\abs\Xvec} \Xvec_i & \carrier A^X & \text{in $\Set$}
  \end{adjunctions}
  \]
  chasing a generic element $f$ down the chain gives:
  \[
  f \mapsto \fun \xvec\fun x.f(x,\xvec)
    \mapsto \fun \xvec\fun x.f(\unit x,\xvec)
    \mapsto \fun \xvec\fun x.f(\unit x,\prod\unit (\xvec))
    \mapsto f \compose (\unit, \seq[i]{\unit})
  \]
  as we wanted.
\end{proof}

We now turn to showing that if there is a left adjoint, then the
theory must be commutative. First, we show there is at most one
candidate for the left adjoint: the same left adjoint as for
homomorphisms:
\begin{proposition}
  If $\carrier- : \Multi(\TT, \pair\Set\prod) \to \pair\Set\prod$ has
  a left adjoint, then it has a left adjoint given by the free
  $\TT$-algebra and its unit.
\end{proposition}
\begin{proof}
  We use the $2$-functor $\seq- : \MultiCAT \to \CAT$ to transport the
  adjunction, and note that the right adjoint $\carrier- : \Multi(\TT,
  \pair\Set\prod) \to \pair\Set\prod$ maps to the usual $\carrier :
  \Mod(\TT, \Set) \to \Set$. Therefore, the object map of the left
  adjoint is given by the free $\TT$-algebra construction, and the
  unit is the usual unit.
\end{proof}

\begin{proposition}\propositionlabel{multi-adjoint forces commutativity}
  If $\carrier- : \Multi(\TT, \pair\Set\prod) \to
  \pair\Set\prod$ has a left adjoint, $\TT$ is commutative.
\end{proposition}
\begin{proof}
  Take any $t : n$ and $s : m$ in $\TT$. We need to show that:
  \[
  t\seq[i=1][n]{s\seq[i=1][m]{x_{i,j}}} = s\seq[i=1][m]{t\seq[i=1][n]{x_{i,j}}}
  \]
  Consider the multi-function $\eta : [n]\times[m] \to
  T([n]\times[m])$ and extend it to a multi-homomorphism $f : (F[n],
  F[m]) \to F([n]\times[m])$ using the adjunction. Consider the
  (interpretation of the) terms $t\seq[i]{\unit i}$ and
  $s\seq[j]{\unit j}$ in $F[n]$ and $F[m]$, resp., and calculate:
  \[
  f(t\seq[i]{\unit i}, s\seq[j]{\unit j})
  =
  t\seq[i]{f(\unit i, s\seq[j]{\unit j})}
  =
  t\seq[i]{s\seq[j]{f(\unit i, \unit j)}}
  =
  t\seq[i]{s\seq[j]{\eta(i, j)}}
  \]
  and similarly:
  \[
  f(t\seq[i]{\unit i}, s\seq[j]{\unit j})
  =
  s\seq[j]{t\seq[i]{\eta(i, j)}}
  \]
  and so $\TT$ is commutative.
\end{proof}

\section{The distributive combination of theories}
\label{sec:distributivity}

Let $\OO$ be an operad and $\TT$ a theory. A \emph{translation $\Trans
  : \OO \to \TT$} consists of:
\begin{itemize}
\item an assignment, for every $f \in \OO_n$, of a term $t \in \TT_n$;
\end{itemize}
such that:
\begin{itemize}
\item the identity $\id : \OO_1$ is sent to the unique variable
  projection: $\Trans\id = x \in \TT_1$; and
\item the translation respects composition, for all $f \in \OO_n$
  composable with $g_i \in \OO_{m_i}$, namely:
  \[
  \Trans(f\compose(g_1, \ldots,
  g_n)) = \Trans f\coseq[i=1][n]{x_i \mapsto \Trans g_i}
  \]
\end{itemize}
A \emph{renaming} $\ren : n \to m$ is any function from the set of $n$
variables to the set of $m$ variables. Renamings act on terms
covariantly in the evident way. A \emph{surjective} translation
$\Trans : \OO \surtrans \TT$ is one where every term $t \in \TT_n$ can be
expressed as an operator $f \in \OO_m$ followed by a renaming $\ren :
m \to n$: $t = \Trans f[\ren]$. Let $\Trans : \OO \surtrans \MM$ be a surjective
translation. $\MM$ is a \emph{$\Trans$-affine} theory if:
\begin{itemize}
  \item for every operator $g : n$ in $\MM$, there is a chosen operation
  $\tilde g \in \OO_n$ such that
  \[
  n \vdash_\MM \Trans \tilde g = g (x_1, \ldots, x_n)
  \]
  \item the chosen operations $\tilde g$ generate $\OO$ under identities and
  operadic composition
  \item every axiom of the $\MM$ is of the form
  \[
  x_1, \ldots, x_n \vdash_\MM \Trans f_1[\ren_1] = \Trans f_2[\ren_2]
  \]
  for some $f_1 \in \OO_{n_1}$, $f_2 \in \OO_{n_2}$ and injective renamings 
  $\ren_1 : n_1 \hookrightarrow n$, $\ren_2 : n_2 \hookrightarrow n$.
\end{itemize}

\begin{proposition}
  Let $\Trans : \OO \to \TT$ be a translation, and let $\C$ be a category
  with chosen finite products $\prod$.
  There is a carrier-preserving, fully-faithful functor:
    \[
    \Mod(\Trans, \pair\C\prod) : \Mod(\TT, \C) \to \Mod(\OO,
    \pair\C\prod)
    \]
    given by:
    \[
    \Mod(\Trans, \pair\C\prod)A\sem f \definedby A\sem {\Trans f}
    \qquad
    \Mod(\Trans, \pair\C\prod)h \definedby h
    \]
\end{proposition}
\begin{proof}
  The compatibility with composition and identities of the operad
  follows from their preservation by $\Trans$.  Preservation of
  identities and composition, and fully-faithfulness are evident.
\end{proof}

Given two signatures $\Sig[1]$, $\Sig[2]$, their \emph{coproduct}
$\Sig[1] + \Sig[2]$ is the signature consisting of the disjoint union
of operations, but retaining the arities:
\[
\carrier{\Sig[1] + \Sig[2]} := \carrier{\Sig[1]} \amalg \carrier{\Sig[2]}
\qquad
\arity (\injection_i f) \definedby \arity[{\Sig[i]}] f
\]
or succicntly:
\[
\Sig[1] + \Sig[2] := \set{\injection_i f : n \suchthat (f : n) \in \Sig[i], i = 1,2}
\]
Given a $\Sig[i]$-term-in-context $n \types_{\Sig[i]} t$, we denote by
$n \types_{\Sig[1]+\Sig[2]} \injection_i [t]$ the
$\Sig[1]+\Sig[2]$-term-in-context obtained by decorating the
operations with $\injection_i$. Similarly, given a $\Sig[i]$-equation
in context $eq := (n \types t_1 = t_2)$ we denote by $\injection_i[eq]$
the equation $n \types \injection_i[t_1] = \injection_i[t_2]$.

Let $\Trans : \OO \to \MM$ be a translation from an operad $\OO$ to a
theory $\MM$ and $\DD$ be any theory. We define the \emph{distributive
  combination of $\Trans$ over $\DD$} to be the theory $\Trans \dist
\DD$ given by:
\begin{multline*}
\Sig[\Trans\dist \DD] \definedby \Sig[\MM] + \Sig[\DD]
\qquad
\Ax[\Trans\dist\DD] \definedby
\injection_1[\Ax[\MM]] \union
\injection_2[\Ax[\AA]]
\\
\union\set{
  \begin{aligned}
    \seq[\substack{i=1\hfill\\i\neq i_0}][n]{x_i},
      \yvec &{}\types
      \Trans f(x_1,\ldots, x_{i_0-1}, s\yvec, x_{i_0+1}, \ldots, x_n)
      \\&{}=
      s\seq[j=1][\abs\yvec]{\Trans f(x_1,\ldots, x_{i_0-1}, \yvec_j, x_{i_0+1}, \ldots, x_n)}
  \end{aligned}
  \suchthat
  \begin{aligned}
    & (f : n) \in \OO_n
    \\& (s : \abs\yvec) \in \Sig[\DD]
  \end{aligned}
}
\end{multline*}

% \begin{example}\emph{Semiring varities.}
%   In each of the following examples, we take $\OO$ to be the operad
%   freely generated by the signature of each theory $\MM$, and $\Trans$
%   to be the translation mapping each term to itself. All of these
%   translations are surjective:

%   \begin{tabular}{@{}ccl@{}}
%     $\MM$ & $\DD$ & $\Trans \dist \DD$\\
%     \hline
%     $\Theory{Semigroup}$ & $\Theory{CommutativeMonoid}$
%     & non-unital semirings, aka rgs
%     \\
%     $\Theory{Monoid}$ & $\Theory{CommutativeMonoid}$
%     & semirings, aka rigs
%     \\
%     $\Theory{CommutativeMonoid}$ & $\Theory{CommutativeMonoid}$
%     & commutative semirings
%     \\
%     $\Theory{InvMonoid}^{\id}$ & $\Theory{CommutativeMonoid}$
%     & non-reversing involutive semirings
%     \\
%     $\Theory{InvMonoid}^{\rev}$ & $\Theory{CommutativeMonoid}$
%     & reversing involutive semirings.
%     \\
%     $\Theory{Monoid} $ & $\Theory{AbelianGroup}$ &
%     rings
%     \\
%     $\Theory{InvMonoid}^{\id,\rev} $ & $\Theory{AbelianGroup}$
%      & non-reversing/reversing involutive rings
%   \end{tabular}
%   and so on.
% \end{example}

\begin{proposition}\emph{Models of the distributive combination \citep{hyland-power-discrete-lawvere-theories-journal}.}
\propositionlabel{Mod(Trans dist A) is a pullback}
  Let $\Trans : \OO \to \MM$ be a translation from an operad to a
  theory, and $\DD$ be a theory.  The category of models for the
  distributive combination $\Trans \dist \DD$ is given by the
  following pullback:
  \diagram{37}
\end{proposition}
\begin{proof}
  A model $A \models \Trans\dist\DD$ consists of:
  \begin{itemize}
  \item an $\MM$-model $\carrier A_{\MM\dist}$ interpreting the
    operations in $\MM$
  \item an $\DD$-model $\carrier A_{\dist\DD}$ interpreting the
    operations in $\DD$.
  \end{itemize}
  Each of the operation $f \in \OO_n$ distributes over the operations
  in $\DD$, and so we have an $n$-ary multi-homomorphism:
  \[
  \carrier A_{\MM\dist}\sem{\Trans f} :
  (\carrier A_{\dist\DD}, \ldots, \carrier A_{\dist\DD}) \to \carrier A_{\dist\DD}
  \]
  in $\Multi(\DD, \Set)$, and this interpretation is compatible with
  the operadic structure, i.e., we additionally have:
  \begin{itemize}
  \item an $\OO$-model $\carrier A_{\OO\dist\DD}$ in $\Multi(\DD, \Set)$;
  \item whose underlying $\OO$-model in $\pair\Set\prod$ is $\carrier
    A_{\OO\surtrans \MM} := \Mod(\Trans, \Set)\carrier A_{\MM}$.
  \end{itemize}
  I.e., we have an object $\pair{\carrier A_{\MM\dist}}{\carrier
    A_{\OO\dist\DD}}$ in the pullback. Conversely, given such an
  object we can reverse this process and assemble back a corresponding
  $\Trans\dist\DD$-model.
\end{proof}
As Hyland and Power, We can refine this construction into an action
$\pair \OO\DD \mapsto \OO *\DD$ of operads on theories to obtain a
universal characterisation of $\Trans\dist\DD$ as a pushout of
theories:
\diagram{38}
We will not pursue this direction here, but we will use the
construction $(*)$.

\section{Free algebras}
\label{sec:fral}
Consider a surjective translation $\Trans : \OO \surtrans \MM$ from an
operad to a $\Trans$-affine theory, and let $\DD$ be a commutative theory. Let $\TT
\definedby \Trans \dist \DD$ be the distributive combination of
$\Trans$ over $\DD$. Let $X$ be a set, and we'll now
construct $F_{\TT}X$, the free $\TT$-algebra over $X$.

\paragraph{Step 1}%
Let $A_0 \definedby F_{\MM}X \models \MM$ be the free $\MM$-algebra over $X$.

\paragraph{Step 2}%
Applying the functor:
\[
\carrier-_{\Trans} \definedby \Mod(\Trans, \Set) : \Mod(\MM, \Set) \to
\Mod(\OO, \pair\Set\prod)
\]
gives an operadic
algebra $A_1 \definedby \carrier{A_0}_{\Trans} \models \OO$.

\paragraph{Step 3}%
Since $\DD$ is commutative, by \theoremref{left adjoint iff commutative}, 
we have an adjunction 
$F_{\DD} \leftadjointto \carrier- : \Multi(\DD, \Set) \to \pair\Set\prod$ in
$\MultiCAT$. Applying the $2$-functor $\Mod(\OO, -) : \MultiCAT \to
\CAT$ to it, we now have a left adjoint $F_{\OO\to \DD} \definedby
\Mod(\OO, F_{\DD}) : \Mod(\OO, \pair\Set\prod) \to \Mod(\OO,
\Multi(\DD, \Set))$.
Take $A_{\OO\dist\DD} \definedby F_{\OO\to\DD}A_1 \models \OO * \DD$.

% \begin{lemma}\lemmalabel{A_{OO dist DD} normal form}
%   For every $e \in \carrier{A_{\OO \dist \DD}}$ there is some $s \in
%   \DD_k$ and $\evec \in \carrier A_1^k$ such that:
%   \[
%   e = A_{\OO\dist\DD}\sem{s}\seq[i=1][k]{\carrier{\unit^{\OO \to \DD}_{A_1}}\evec_i}
%   \]
% \end{lemma}

% \begin{proof}
%   Let $\carrier B \subset \carrier A_{\OO\dist\DD}$ be the following subset:
%   \[
%   \carrier B \definedby \set{e \suchthat
%     \exists s \in
%   \DD_k, \evec \in \carrier A_1^k:
%   e = A_{\OO\dist\DD}\sem{s}\seq[i=1][k]{\carrier{\unit^{\OO \to \DD}_{A_1}}\evec_i}
%   }
%   \]
%   This subset is closed under both the operations in $\OO$ and $\DD$.

%   \paragraph{Closure under $\DD$}
%   Take any $s \in \DD_n$ and $e_1, \ldots, e_n \in \carrier B$. By
%   assumption, there are $s_i \in \DD_{k_i}$ and $\evec_{\pair i*}$ such that
%   \[
%   e_i = A_{\OO\dist\DD}\sem{s_i}\seq[j=1][k_i]{\carrier{\unit^{\OO \to \DD}_{A_1}}\evec_{\pair ij}}
%   \]
%   and, taking $p \definedby \sum_{i}k_i$ and $s' \definedby
%   s\coseq[i=1][n]{i \mapsto s_i} \in \DD_{p}$, we have:
%   \begin{align*}
%   A_{\OO\dist\DD}\sem{s}\seq[i=1][n]{e_i}
%   &{}=
%   A_{\OO\dist\DD}\sem{s}\seq[i=1][n]{
%     A_{\OO\dist\DD}\sem{s_i}\seq[j=1][k_i]{\carrier{\unit^{\OO \to \DD}_{A_1}}\evec_{\pair ij}}
%   }
%   \\&{}=
%   A_{\OO\dist\DD}\sem{s(s_1, \ldots, s_n)}\seq[\ell=1][p]{
%     \carrier{\unit^{\OO \to \DD}_{A_1}}\evec_{\ell}}
%   \end{align*}
%   and so $A_{\OO\dist\DD}\sem{s}\seq[i=1][n]{e_i} \in \carrier B$.

%   \paragraph{Closure under $\DD$}
%   Take any $f \in \O_n$ and $e_1, \ldots, e_n \in \carrier B$. By
%   assumption, there are $s_i \in \DD_{k_i}$ and $\evec_{\pair i*}$ such that
%   \[
%   e_i = A_{\OO\dist\DD}\sem{s_i}\seq[j=1][k_i]{\carrier{\unit^{\OO \to \DD}_{A_1}}\evec_{\pair ij}}
%   \]
%   Since $A_{\OO\dist\DD}\sem{f}$ is a multi-$\DD$-homomorphism, we have:
%   \begin{align*}
%     A_{\OO\dist\DD}\sem{f}\seq[i=1][n]{e_i}
%     &{}=
%     A_{\OO\dist\DD}\sem{f}\seq[i=1][n]{
%       A_{\OO\dist\DD}\sem{s_i}\seq[j=1][k_i]{
%         \carrier{\unit^{\OO \to \DD}_{A_1}}\evec_{\pair ij}
%       }
%     }
%     \\&{}=
%     A_{\OO\dist\DD}\sem{s_1}\seq[j_1=1][k_1]{\cdots
%       A_{\OO\dist\DD}\sem{s_n}\seq[j_n=1][k_n]{
%           A_{\OO\dist\DD}\sem{f}\seq[i=1][n]{
%             \carrier{\unit^{\OO \to \DD}_{A_1}}\evec_{\pair i{j_i}}
%     }}}
%     \\&{}=
%     A_{\OO\dist\DD}\sem{s_1\seq[j_1=1][k_1]{\cdots s_n\seq[j_n=1][k_n]{x_{j_1, \ldots, j_n}}}}
%     \seq[\ell=1][\prod_{i=1}^nk_i]{
%       A_{\OO\dist\DD}\sem{f}\seq[i=1][n]{
%             \carrier{\unit^{\OO \to \DD}_{A_1}}\evec_{\pair i{\seq[i][-1]\ell}}}
%     }
%   \end{align*}
%   and so $A_{\OO\dist\DD}\sem{f}\seq[i=1][n]{e_i} \in \carrier B$.

%   Therefore, we have a subalgebra $i : B \embed A_{\OO\dist \DD}$, and
%   so $B \models \OO*\DD$.

%   Define an $\OO$-homomorphism $f : A_1 \to \carrier{A_{\OO\dist
%       \DD}}_{\OO\to\DD}$ by noting that the image of the
%   $\OO$-homomorphism $\unit^{\OO\to\DD}_{A_1} : A_1 \to
%   \carrier{A_{\OO\dist\DD}}_{\OO\to\DD}$ lifts along the inclusion
%   homomorphism $i : B \embed A_{\OO\dist \DD}$. We therefore have a
%   unique $\OO*\DD$-homomorphism $g : A_{\OO\dist \DD} \to B$ extending
%   $f$ along the unit $\unit^{\OO\to\DD}$.

%   This homomorphism is given by $\carrier g(e) = e$. Indeed,
%   post-composing with $i : B \embed A_{\OO\dist A}$ we get a
%   homomorphism, and pre-composing it with $\unit^{\OO\dist\DD}$:
%   \[
%   i \compose g \compose \unit^{\OO\dist\DD}
%   \explain={$g$ extends $f$ along $\unit$}
%   i \compose f
%   \explain*={$f$ lifts $\unit$ along $i$}
%   \unit^{\OO\dist\DD}
%   $
%   and therefore $i \compose g = \id$, and so:
%   $
%   g(e) = i(g(e)) = e
%   \]
%   Therefore $\carrier B = \carrier A_{\OO\dist\DD}$.
% \end{proof}

\paragraph{Step 4}%
Let $A_{\OO\dist} \definedby \Mod(\OO, \carrier-)A_{\OO\dist\DD} \models
\OO$ be the underlying operadic algebra. For each operation symbol $g : n$ of $\MM$,
define $A_{\MM} \sem g \definedby A_{\OO\dist} \sem{\tilde g}$.

\begin{lemma}\label{lem:A_M model}
  With the preceding notation:
  \begin{itemize}
    \item For every $f \in \OO_n$, the interpretation $A_\MM \sem{\Trans f}$ is equal
    to $A_{\OO\dist} \sem f$.
    \item The resulting algebra validates $\MM$, i.e.: $A_{\MM} \models \MM$.
  \end{itemize}
\end{lemma}

\begin{proof}
  Let $X \definedby \carrier A_1$ and let $\unit_X : X \to \carrier{F_\DD X}$ be the
  unit of the adjunction $F_\DD \leftadjointto U_\DD$. By construction of 
  $A_{\OO\dist\DD} = F_{\OO\to\DD} A_1$, for every $f \in \OO_n$, the operation 
  $A_{\OO\to\DD} \sem f$ is the unique multi-$\DD$-homomorphism extending
  \[
  X^n \xto{A_1\sem f} X \xto{\unit_X} \carrier{F_\DD x}
  \]
  We first prove that, for every $f \in \OO_n$, the interpretation in $A_\MM$ of the
  term $\Trans f$ is equal to $A_{\OO\dist} \sem f$. Since the chosen operations
  $\tilde g$ generate $\OO$, we may prove this by induction on the operadic expression
  of $f$.
  \begin{itemize}
    \item If $f = \id$, the claim is immediate.
    \item If $f = \tilde g$ is one of the chosen generators, the claim is exactly
    the definition of $A_\MM \sem g$.
    \item If $f = h \compose (f_1, \ldots, f_k)$ then, using the preservation of
    operadic composition by $\Trans$, the induction hypothesis, and the fact that
    $A_{\OO\dist}$ is an $\OO$-model, we get
    \[
    A_\MM \sem{\Trans f} = A_\MM \sem{\Trans h[\Trans f_1, \ldots, \Trans f_k]} =
    A_{\OO\dist} \sem h \compose (A_{\OO\dist} \sem{f_1}, \ldots, A_{\OO\dist} \sem{f_k})
    = A_{\OO\dist} \sem f
    \]
  \end{itemize}
  Now let $x_1, \ldots, x_m \vdash_\MM \Trans f_1 [\ren_1] = \Trans f_2 [\ren_2]$ be an
  axiom of $\MM$. For each $i = 1,2$, define $\Phi_i : X^m \to \carrier{F_\DD X}$ by
  \[
  \Phi_i(e) \definedby \unit_X(A_1 \sem{f_i} (e \compose \ren_i))
  \]
  Since $A_0 \models \MM$, these two maps are equal: $\Phi_1 = \Phi_2$. For each
  $i = 1, 2$, define $\bar \Phi_i : \carrier{F_\DD X}^m \to \carrier{F_\DD X}$ by
  \[
  \bar\Phi_i(e) \definedby A_{\OO\dist\DD} \sem{f_i} (e \compose \ren)
  \]
  Because $\ren_i$ is injective, each coordinate of $e$ is used in at most one argument
  of $A_{\OO\dist\DD}$, and therefore $\bar\Phi_i$ is a multi-$\DD$-homomorphism in each
  variable. By construction, it extends $\Phi_i$. By uniqueness of multi-$\DD$-homomorphic
  extension from generators, we conclude that $\bar\Phi_1 = \bar\Phi_2$. Hence, for
  every environment $e : [m] \to \carrier{A_{\OO\dist\DD}}$,
  \[
  A_{\OO\dist\DD} \sem{f_1} (e \compose \ren_1) = 
  A_{\OO\dist\DD} \sem{f_2} (e \compose \ren_2)
  \]
  Forgetting from $A_{\OO\dist\DD}$ to $A_{\OO\dist}$ preserves this equality, so the two
  sides of the axiom have the same interpretation in $A_\MM$. Hence, all axioms of $\MM$
  hold in $A_\MM$, and therefore $A_\MM \models \MM$.
\end{proof}

\paragraph{Step 5}%
Let $A \models \Trans\dist\DD$ be the pullback of $A_{\OO\dist\DD}$ and $A_\MM$.

\begin{lemma}
\lemmalabel{A is free}
  $A$ is a free $\Trans\dist\DD$-algebra.
\end{lemma}

\begin{proof}
  We already know that $A$ models $\Trans\dist\DD$. Let's prove that it is free. 
  
  Let $B$ by any $\Trans\dist\DD$-model and $f : X \to \carrier B$ be a function.
  Let $B_\OO \definedby \Mod(\OO, \carrier-) \carrier{B}_{\OO\dist\DD}$ and
  $B_\MM \definedby \carrier{B}_{\MM\dist}$.

\begin{itemize}
\item $A_0$ is a free $\MM$-algebra over $X$, therefore we have a unique morphism 
$h_0 : A_0 \to B_\MM$ of $\MM$-algebras such that $h_0 \compose \unit^{A_0} = f$.
\item There is also an homomorphism $h_1 : A_1 \to B_\OO$ of $\OO$-algebras such that 
$h_1 \compose \unit_{A_0}^{A_1} = h_0$, and $h_1 = h_0$ since $\unit_{A_0}^{A_1} = \id$.
\item $B$ has an $\DD$-structure that distributes over $\OO$, therefore thanks to the
adjunction $\Mod(\OO, F_\DD) \leftadjointto \Mod(\OO, \carrier{-})$, we have 
a unique morphism $h_2 : A_{\OO\dist\DD} \to B$ of $\OO \ast \DD$-algebras such that 
$h_2 \compose \unit_{A_1}^{A_{\OO\dist\DD}} = h_1$.
\item Apply the 2-functor $\Mod(\OO, \carrier{-})$ to the morphism 
$h_2 : A_{\OO\dist\DD} \to B$ gives a morphism $h_3 : A_{\OO\dist} \to B_\OO$ of 
$\OO$-algebras such that $h_3 \compose \unit_{A_{\OO\dist\DD}}^{A_{\OO\dist}} = h_2$ and 
$h_3 = h_2$ since $\unit_{A_{\OO\dist\DD}}^{A_{\OO\dist}} = \id$.
\item By the surjectivity of $\Trans$ and
Lemma~\ref{lem:A_M model}, we now have a $\MM$-morphism 
$h_4 = h_3 : A_\MM \to B_\MM$ such that
$h_4 \compose \unit_{A_{\OO\dist}}^{A_\MM} = h_3$.

Let $t \in \mathcal M_n$ and $e : [n] \to \carrier{A_\MM}$. By surjectivity, there 
is some $f \in \mathcal O_m$ and $\theta : m \to n$ such that $t = \Trans f [\theta]$. 
We define $h_4$ as 
$$h_4 \lp A_\MM \sem{t} e \rp = h_3 \lp A_{\OO\dist} \sem{f} e \compose \theta \rp$$ 
This is well-defined by Lemma~\ref{lem:A_M model}. We then check that $h_4$ is a
morphism of $\MM$-algebras. Since $h_3$ is a morphism of $\OO$-algebras, 
\begin{align*}
  h_3 \lp A_{\OO\dist} \sem{f} e \compose \theta \rp &= 
    B_\OO \sem{f} (h_3 \compose e \compose \theta, \dots, h_3 \compose e \compose \theta)\\
  &= B_\MM \sem{f} (h_4 \compose e, \dots, h_4 \compose e)\\
\end{align*}
This proves that $h_4 \lp A_\MM \sem{t} e \rp = B_\MM \sem{t} (h_4 \compose e)$, so $h_4$
is a morphism of $\MM$-algebras.
\end{itemize}
% \[
% \begin{tikzcd}
%   X \arrow[r, "\unit^{A_0}"] \arrow[dr, "f"'] & A_0 \arrow[d, "h_0"]\\
%   & B_\MM
% \end{tikzcd}
% \]

% \[
% \begin{tikzcd}
%   X \arrow[r, "\id \compose \unit^{A_0}"] \arrow[dr, "f"'] & A_1 \arrow[d, "h_1 = h_0"]\\
%   & B_\OO
% \end{tikzcd}
% \]

% \[
% \begin{tikzcd}
%   A_1 \arrow[r, "\unit_{A_1}^{A_{\OO\dist\DD}}"] \arrow[dr, "h_1"'] & A_{\OO\dist\DD} \arrow[d, "h_2"]\\
%   & B_{\OO \ast \DD}
% \end{tikzcd}
% \]

\noindent In the end we have the following diagram:
% https://q.uiver.app/#q=WzAsNyxbMCwwLCJYIl0sWzEsMCwiXFx1bHtBXzB9Il0sWzIsMCwiXFx1bHtBXzF9Il0sWzEsMSwiXFx1bHtCfSJdLFsyLDEsIlxcdWx7QV8yfSJdLFsyLDIsIlxcdWx7QV8zfSJdLFsxLDIsIlxcdWx7QV80fSJdLFswLDMsImYiLDJdLFswLDEsIlxcZXRhXntBXzB9Il0sWzEsMiwiXFxpZCJdLFsxLDMsImhfMCIsMV0sWzIsMywiaF8xIiwxXSxbNCwzLCJoXzIiLDFdLFs1LDMsImhfMyIsMV0sWzIsNCwiXFxldGFfe0FfMX1ee0FfMn0iXSxbNCw1LCJcXGlkIl0sWzYsMywiaF80IiwxXSxbNSw2LCJcXGlkIl1d&macro_url=https%3A%2F%2Fraw.githubusercontent.com%2FS-c-r-a-t-c-h-y%2Fsemiring-solvers%2Frefs%2Fheads%2Fmain%2Fmacros.sty
\[\begin{tikzcd}
	X & {\carrier{A_0}} & {\carrier{A_1}} \\
	& {\carrier{B}} & {\carrier{A_{\OO\dist\DD}}} \\
	& {\carrier{A_\MM}} & {\carrier{A_{\OO\dist}}}
	\arrow["{\unit^{A_0}}", from=1-1, to=1-2]
	\arrow["f"', from=1-1, to=2-2]
	\arrow["\id", from=1-2, to=1-3]
	\arrow["{h_0}"{description}, from=1-2, to=2-2]
	\arrow["{h_1}"{description}, from=1-3, to=2-2]
	\arrow["{\unit_{A_1}^{A_{\OO\dist\DD}}}", from=1-3, to=2-3]
	\arrow["{h_2}"{description}, from=2-3, to=2-2]
	\arrow["\id", from=2-3, to=3-3]
	\arrow["{h_4}"{description}, from=3-2, to=2-2]
	\arrow["{h_3}"{description}, from=3-3, to=2-2]
	\arrow["\id", from=3-3, to=3-2]
\end{tikzcd}\]
We then consider the walking morphism $I : \{0 \xto{u} 1\}$ and the two functors
$\FF_1 : I \to \Mod(\OO, \Multi(\DD, \Set))$ and $\FF_2 : I \to \Mod(\MM, \Set)$ defined by
$\FF_1\, 0 = A_{\OO\dist\DD},\ \FF_1\, 1 = B_{\OO \ast \DD},\ \FF_1\, u = h_2$ and 
$\FF_2\, 0 = A_\MM,\ \FF_2\, 1 = B_\MM,\ \FF_2\, u = h_4$. By the construction, we can then 
use the pullback square to get a unique $\FF[3] : I \to \Mod(\Trans \dist \DD)$
such that the following diagram commutes:
% https://q.uiver.app/#q=WzAsNSxbMSwxLCJcXE1vZChcXG1hdGhmcmFrIFQgXFx0cmlhbmdsZWxlZnRcXG1hdGhjYWwgQSwgXFxTZXQpIl0sWzEsMiwiXFxNb2QoXFxtYXRoY2FsIE0sIFxcU2V0KSJdLFsyLDIsIlxcTW9kKFxcbWF0aGNhbCBPLCBcXGxhIFxcU2V0LCBcXHByb2QgXFxyYSkiXSxbMiwxLCJcXE1vZChcXG1hdGhjYWwgTywgXFxNdWx0aShcXG1hdGhjYWwgQSwgXFxTZXQpKSJdLFswLDAsIkkiXSxbMywyLCJcXE1vZChcXG1hdGhjYWwgTywgXFx1bmRlcmxpbmV7LX0pIl0sWzAsMywiXFx1bmRlcmxpbmV7LX1fe1xcbWF0aGNhbCBPIFxcdHJpYW5nbGVsZWZ0IFxcbWF0aGNhbCBBfSJdLFsxLDIsIlxcTW9kKFxcbWF0aGZyYWsgVCwgXFxTZXQpIiwyXSxbMCwxLCJcXHVuZGVybGluZXstfV97XFxtYXRoY2FsIE19IiwyXSxbNCwzLCJGXzEiLDAseyJjdXJ2ZSI6LTN9XSxbNCwxLCJGXzIiLDIseyJjdXJ2ZSI6M31dLFs0LDAsIkgiLDEseyJzdHlsZSI6eyJib2R5Ijp7Im5hbWUiOiJkYXNoZWQifX19XSxbMCw3LCIiLDEseyJsZXZlbCI6MSwic3R5bGUiOnsibmFtZSI6ImNvcm5lciJ9fV1d&macro_url=https%3A%2F%2Fraw.githubusercontent.com%2FS-c-r-a-t-c-h-y%2Fsemiring-solvers%2Frefs%2Fheads%2Fmain%2Fmacros.sty
\[\begin{tikzcd}
	I && \\
	& {\Mod(\Trans \dist\DD, \Set)} & {\Mod(\OO, \Multi(\DD, \Set))} \\
	& {\Mod(\MM, \Set)} & {\Mod(\OO, \pair\Set\prod)}
	\arrow["{\FF[3]}"{description}, dashed, from=1-1, to=2-2]
	\arrow["{\FF_1}", curve={height=-18pt}, from=1-1, to=2-3]
	\arrow["{\FF_2}"', curve={height=18pt}, from=1-1, to=3-2]
	\arrow["{\carrier{-}_{\OO \dist \DD}}", from=2-2, to=2-3]
	\arrow["{\carrier{-}_{\MM}}"', from=2-2, to=3-2]
	\arrow["{\Mod(\OO, \carrier{-})}", from=2-3, to=3-3]
	\arrow[""{name=0, anchor=center, inner sep=0}, "{\Mod(\Trans, \Set)}"', from=3-2, to=3-3]
	\arrow["\lrcorner"{anchor=center, pos=0.125}, draw=none, from=2-2, to=0]
\end{tikzcd}\]
This gives us a unique morphism $h = \FF[3]\, u : A \to B$ of 
$\Trans \dist \DD$-algebras.

We need to show that $h$ extends $f$ along $\unit$. $h_2 \compose \unit = h_4 \compose \unit = f$,
therefore $h \compose \unit = f$ since $\unit$ comes from the pullback square. 
This morphism is unique for a given $h_2$ and $h_4$, but these are also uniquely
determined by $h_0$ and $h_1$, which are themselves uniquely determined by $f$.
Therefore $h$ is the unique morphism of $\Trans \dist \DD$-algebras such that
the following diagram commutes:
\[\begin{tikzcd}
  X \arrow[r, "\unit"] \arrow[dr, "f"'] & A \arrow[d, "h"]\\
  & B
\end{tikzcd}\]
This proves that $A$ is the free $\Trans \dist \DD$-algebra over $X$.
\end{proof}

\bibliography{ref}
\end{document}

%%% Local Variables:
%%% mode: latex
%%% TeX-master: t
%%% End:
